%********************************************%
%*       Generated from PreTeXt source      *%
%*       on 2026-08-30T20:13:20-05:00       *%
%*   A recent stable commit (2022-07-01):   *%
%* 6c761d3dba23af92cba35001c852aac04ae99a5f *%
%*                                          *%
%*         https://pretextbook.org          *%
%*                                          *%
%********************************************%
\documentclass[oneside,10pt,]{book}
%% Custom Preamble Entries, early (use latex.preamble.early)
%% Default LaTeX packages
%%   1.  always employed (or nearly so) for some purpose, or
%%   2.  a stylewriter may assume their presence
\usepackage{geometry}
%% Some aspects of the preamble are conditional,
%% the LaTeX engine is one such determinant
\usepackage{ifthen}
%% etoolbox has a variety of modern conveniences
\usepackage{etoolbox}
\usepackage{ifxetex,ifluatex}
%% Raster graphics inclusion
\usepackage{graphicx}
%% Color support, xcolor package
%% Always loaded, for: add/delete text, author tools
%% Here, since tcolorbox loads tikz, and tikz loads xcolor
\PassOptionsToPackage{dvipsnames,svgnames,table}{xcolor}
\usepackage{xcolor}
%% begin: defined colors, via xcolor package, for styling
%% end: defined colors, via xcolor package, for styling
%% Colored boxes, and much more, though mostly styling
%% skins library provides "enhanced" skin, employing tikzpicture
%% boxes may be configured as "breakable" or "unbreakable"
%% "raster" controls grids of boxes, aka side-by-side
\usepackage{tcolorbox}
\tcbuselibrary{skins}
\tcbuselibrary{breakable}
\tcbuselibrary{raster}
%% We load some "stock" tcolorbox styles that we use a lot
%% Placement here is provisional, there will be some color work also
%% First, black on white, no border, transparent, but no assumption about titles
\tcbset{ bwminimalstyle/.style={size=minimal, boxrule=-0.3pt, frame empty,
colback=white, colbacktitle=white, coltitle=black, opacityfill=0.0} }
%% Second, bold title, run-in to text/paragraph/heading
%% Space afterwards will be controlled by environment,
%% independent of constructions of the tcb title
%% Places \blocktitlefont onto many block titles
\tcbset{ runintitlestyle/.style={fonttitle=\blocktitlefont\upshape\bfseries, attach title to upper} }
%% Spacing prior to each exercise, anywhere
\tcbset{ exercisespacingstyle/.style={before skip={1.5ex plus 0.5ex}} }
%% Spacing prior to each block
\tcbset{ blockspacingstyle/.style={before skip={2.0ex plus 0.5ex}} }
%% xparse allows the construction of more robust commands,
%% this is a necessity for isolating styling and behavior
%% The tcolorbox library of the same name loads the base library
\tcbuselibrary{xparse}
%% The tcolorbox library loads TikZ, its calc package is generally useful,
%% and is necessary for some smaller documents that use partial tcolor boxes
%% See:  https://github.com/PreTeXtBook/pretext/issues/1624
\usetikzlibrary{calc}
%% We use some more exotic tcolorbox keys to restore indentation to parboxes
\tcbuselibrary{hooks}
%% Save default paragraph indentation and parskip for use later, when adjusting parboxes
\newlength{\normalparindent}
\newlength{\normalparskip}
\AtBeginDocument{\setlength{\normalparindent}{\parindent}}
\AtBeginDocument{\setlength{\normalparskip}{\parskip}}
\newcommand{\setparstyle}{\setlength{\parindent}{\normalparindent}\setlength{\parskip}{\normalparskip}}%% Hyperref should be here, but likes to be loaded late
%%
%% Inline math delimiters, \(, \), need to be robust
%% 2016-01-31:  latexrelease.sty  supersedes  fixltx2e.sty
%% If  latexrelease.sty  exists, bugfix is in kernel
%% If not, bugfix is in  fixltx2e.sty
%% See:  https://tug.org/TUGboat/tb36-3/tb114ltnews22.pdf
%% and read "Fewer fragile commands" in distribution's  latexchanges.pdf
\IfFileExists{latexrelease.sty}{}{\usepackage{fixltx2e}}
%% Footnote counters and part/chapter counters are manipulated
%% April 2018:  chngcntr  commands now integrated into the kernel,
%% but circa 2018/2019 the package would still try to redefine them,
%% so we need to do the work of loading conditionally for old kernels.
%% From version 1.1a,  chngcntr  should detect defintions made by LaTeX kernel.
\ifdefined\counterwithin
\else
    \usepackage{chngcntr}
\fi
%% Text height identically 9 inches, text width varies on point size
%% See Bringhurst 2.1.1 on measure for recommendations
%% 75 characters per line (count spaces, punctuation) is target
%% which is the upper limit of Bringhurst's recommendations
\geometry{letterpaper,total={340pt,9.0in}}
%% Custom Page Layout Adjustments (use publisher page-geometry entry)
\geometry{top=2cm,left=2cm,right=2cm,bottom=2cm}
\usepackage[letter,cam,center]{crop}
%% This LaTeX file may be compiled with pdflatex, xelatex, or lualatex executables
%% LuaTeX is not explicitly supported, but we do accept additions from knowledgeable users
%% The conditional below provides  pdflatex  specific configuration last
%% begin: engine-specific capabilities
\ifthenelse{\boolean{xetex} \or \boolean{luatex}}{%
%% begin: xelatex and lualatex-specific default configuration
\ifxetex\usepackage{xltxtra}\fi
%% realscripts is the only part of xltxtra relevant to lualatex 
\ifluatex\usepackage{realscripts}\fi
%% end:   xelatex and lualatex-specific default configuration
}{
%% begin: pdflatex-specific default configuration
%% We assume a PreTeXt XML source file may have Unicode characters
%% and so we ask LaTeX to parse a UTF-8 encoded file
%% This may work well for accented characters in Western language,
%% but not with Greek, Asian languages, etc.
%% When this is not good enough, switch to the  xelatex  engine
%% where Unicode is better supported (encouraged, even)
\usepackage[utf8]{inputenc}
%% end: pdflatex-specific default configuration
}
%% end:   engine-specific capabilities
%%
%% Fonts.  Conditional on LaTex engine employed.
%% Default Text Font: The Latin Modern fonts are
%% "enhanced versions of the [original TeX] Computer Modern fonts."
%% We use them as the default text font for PreTeXt output.
%% Default Monospace font: Inconsolata (aka zi4)
%% Sponsored by TUG: http://levien.com/type/myfonts/inconsolata.html
%% Loaded for documents with intentional objects requiring monospace
%% See package documentation for excellent instructions
%% fontspec will work universally if we use filename to locate OTF files
%% Loads the "upquote" package as needed, so we don't have to
%% Upright quotes might come from the  textcomp  package, which we also use
%% We employ the shapely \ell to match Google Font version
%% pdflatex: "varl" package option produces shapely \ell
%% pdflatex: "var0" package option produces plain zero (not used)
%% pdflatex: "varqu" package option produces best upright quotes
%% xelatex,lualatex: add OTF StylisticSet 1 for shapely \ell
%% xelatex,lualatex: add OTF StylisticSet 2 for plain zero (not used)
%% xelatex,lualatex: add OTF StylisticSet 3 for upright quotes
%%
%% Automatic Font Control
%% Portions of a document, are, or may, be affected by defined commands
%% These are perhaps more flexible when using  xelatex  rather than  pdflatex
%% The following definitions are meant to be re-defined in a style, using \renewcommand
%% They are scoped when employed (in a TeX group), and so should not be defined with an argument
\newcommand{\divisionfont}{\relax}
\newcommand{\blocktitlefont}{\relax}
\newcommand{\contentsfont}{\relax}
\newcommand{\pagefont}{\relax}
\newcommand{\tabularfont}{\relax}
\newcommand{\xreffont}{\relax}
\newcommand{\titlepagefont}{\relax}
%%
\ifthenelse{\boolean{xetex} \or \boolean{luatex}}{%
%% begin: font setup and configuration for use with xelatex
%% Generally, xelatex is necessary for non-Western fonts
%% fontspec package provides extensive control of system fonts,
%% meaning *.otf (OpenType), and apparently *.ttf (TrueType)
%% that live *outside* your TeX/MF tree, and are controlled by your *system*
%% (it is possible that a TeX distribution will place fonts in a system location)
%%
%% The fontspec package is the best vehicle for using different fonts in  xelatex
%% So we load it always, no matter what a publisher or style might want
%%
\usepackage{fontspec}
%%
%% begin: xelatex main font ("font-xelatex-main" template)
%% Latin Modern Roman is the default font for xelatex and so is loaded with a TU encoding
%% *in the format* so we can't touch it, only perhaps adjust it later
%% in one of two ways (then known by NFSS names such as "lmr")
%% (1) via NFSS with font family names such as "lmr" and "lmss"
%% (2) via fontspec with commands like \setmainfont{Latin Modern Roman}
%% The latter requires the font to be known at the system-level by its font name,
%% but will give access to OTF font features through optional arguments
%% https://tex.stackexchange.com/questions/470008/
%% where-and-how-does-fontspec-sty-specify-the-default-font-latin-modern-roman
%% http://tex.stackexchange.com/questions/115321
%% /how-to-optimize-latin-modern-font-with-xelatex
%%
%% end:   xelatex main font ("font-xelatex-main" template)
%% begin: xelatex math font ("font-xelatex-math" template)
%% end:   xelatex math font ("font-xelatex-math" template)
%% begin: xelatex mono font ("font-xelatex-mono" template)
%% (conditional on non-trivial uses being present in source)
\IfFontExistsTF{Inconsolatazi4-Regular.otf}{}{\GenericError{}{The font "Inconsolatazi4-Regular.otf" requested by PreTeXt output processed by the  xelatex  executable is not available.  Either a file cannot be located in default locations via a filename, or a font is not known by its name as part of your system.}{Consult the PreTeXt Guide for help with LaTeX fonts, or perhaps try using  pdflatex  as a test.}{}}
\IfFontExistsTF{Inconsolatazi4-Bold.otf}{}{\GenericError{}{The font "Inconsolatazi4-Bold.otf" requested by PreTeXt output processed by the  xelatex  executable is not available.  Either a file cannot be located in default locations via a filename, or a font is not known by its name as part of your system.}{Consult the PreTeXt Guide for help with LaTeX fonts, or perhaps try using  pdflatex  as a test.}{}}
\usepackage{zi4}
\setmonofont[BoldFont=Inconsolatazi4-Bold.otf,StylisticSet={1,3}]{Inconsolatazi4-Regular.otf}
%% end:   xelatex mono font ("font-xelatex-mono" template)
%% begin: xelatex font adjustments ("font-xelatex-style" template)
%% end:   xelatex font adjustments ("font-xelatex-style" template)
%%
%% Extensive support for other languages
\usepackage{polyglossia}
%% Set main/default language based on pretext/@xml:lang value
%% document language code is "en-US", US English
%% usmax variant has extra hypenation
\setmainlanguage[variant=usmax]{english}
%% Enable secondary languages based on discovery of @xml:lang values
%% Enable fonts/scripts based on discovery of @xml:lang values
%% Western languages should be ably covered by Latin Modern Roman
%% end:   font setup and configuration for use with xelatex
}{%
%% begin: font setup and configuration for use with pdflatex
%% begin: pdflatex main font ("font-pdflatex-main" template)
\usepackage{lmodern}
\usepackage[T1]{fontenc}
%% end:   pdflatex main font ("font-pdflatex-main" template)
%% begin: pdflatex mono font ("font-pdflatex-mono" template)
%% (conditional on non-trivial uses being present in source)
\usepackage[varqu,varl]{inconsolata}
%% end:   pdflatex mono font ("font-pdflatex-mono" template)
%% begin: pdflatex font adjustments ("font-pdflatex-style" template)
%% end:   pdflatex font adjustments ("font-pdflatex-style" template)
%% end:   font setup and configuration for use with pdflatex
}
%% Micromanage spacing, etc.  The named "microtype-options"
%% template may be employed to fine-tune package behavior
\usepackage{microtype}
%% Symbols, align environment, commutative diagrams, bracket-matrix
\usepackage{amsmath}
\usepackage{amscd}
\usepackage{amssymb}
%% allow page breaks within display mathematics anywhere
%% level 4 is maximally permissive
%% this is exactly the opposite of AMSmath package philosophy
%% there are per-display, and per-equation options to control this
%% split, aligned, gathered, and alignedat are not affected
\allowdisplaybreaks[4]
%% allow more columns to a matrix
%% can make this even bigger by overriding with  latex.preamble.late  processing option
\setcounter{MaxMatrixCols}{30}
%%
%% Symbols for use in text mode (i.e. not in math mode)
\usepackage{pifont}
%% Solid geometric shapes to mark ends of environments
\newcommand{\ptxsquare}{\ding{110}}
\newcommand{\ptxtriangle}{\ding{115}}
\newcommand{\ptxdiamond}{\ding{117}}
%% Markers for unordered list items
\newcommand{\ptxlistdisc}{\raisebox{0.3ex}{\textbullet}}
\newcommand{\ptxlistcircle}{\raisebox{0.3ex}{\textopenbullet}}
\newcommand{\ptxlistsquare}{\raisebox{0.4ex}{\scalebox{0.5}{\ding{110}}}}
%%
%% Division Titles, and Page Headers/Footers
%% titlesec package, loading "titleps" package cooperatively
%% See code comments about the necessity and purpose of "explicit" option.
%% The "newparttoc" option causes a consistent entry for parts in the ToC 
%% file, but it is only effective if there is a \titleformat for \part.
%% "pagestyles" loads the  titleps  package cooperatively.
\usepackage[explicit, newparttoc, pagestyles]{titlesec}
%% The companion titletoc package for the ToC.
\usepackage{titletoc}
%% Fixes a bug with transition from chapters to appendices in a "book"
%% See generating XSL code for more details about necessity
\newtitlemark{\chaptertitlename}
%% begin: customizations of page styles via the modal "titleps-style" template
%% Designed to use commands from the LaTeX "titleps" package
%% Plain pages should have the same font for page numbers
\renewpagestyle{plain}{%
\setfoot{}{\pagefont\thepage}{}%
}%
%% Single pages as in default LaTeX
\renewpagestyle{headings}{%
\sethead{\pagefont\slshape\MakeUppercase{\ifthechapter{\chaptertitlename\space\thechapter.\space}{}\chaptertitle}}{}{\pagefont\thepage}%
}%
\pagestyle{headings}
%% end: customizations of page styles via the modal "titleps-style" template
%%
%% Create globally-available macros to be provided for style writers
%% These are redefined for each occurence of each division
\newcommand{\divisionnameptx}{\relax}%
\newcommand{\titleptx}{\relax}%
\newcommand{\subtitleptx}{\relax}%
\newcommand{\shortitleptx}{\relax}%
\newcommand{\authorsptx}{\relax}%
\newcommand{\epigraphptx}{\relax}%
%% Create environments for possible occurences of each division
%% Environment for a PTX "section" at the level of a LaTeX "section"
\NewDocumentEnvironment{sectionptx}{mmmmmmm}
{%
\renewcommand{\divisionnameptx}{#1}%
\renewcommand{\titleptx}{#2}%
\renewcommand{\subtitleptx}{#3}%
\renewcommand{\shortitleptx}{#4}%
\renewcommand{\authorsptx}{#5}%
\renewcommand{\epigraphptx}{#6}%
\section[{#4}]{#2}%
\label{#7}%
}{}%
%% Environment for a PTX "chapter" at the level of a LaTeX "chapter"
\NewDocumentEnvironment{chapterptx}{mmmmmmm}
{%
\renewcommand{\divisionnameptx}{#1}%
\renewcommand{\titleptx}{#2}%
\renewcommand{\subtitleptx}{#3}%
\renewcommand{\shortitleptx}{#4}%
\renewcommand{\authorsptx}{#5}%
\renewcommand{\epigraphptx}{#6}%
\chapter[{#4}]{#2}%
\label{#7}%
}{}%
%% Environment for a PTX "subsection" at the level of a LaTeX "subsection"
\NewDocumentEnvironment{subsectionptx}{mmmmmmm}
{%
\renewcommand{\divisionnameptx}{#1}%
\renewcommand{\titleptx}{#2}%
\renewcommand{\subtitleptx}{#3}%
\renewcommand{\shortitleptx}{#4}%
\renewcommand{\authorsptx}{#5}%
\renewcommand{\epigraphptx}{#6}%
\subsection[{#4}]{#2}%
\label{#7}%
}{}%
%% Environment for a PTX "references" at the level of a LaTeX "chapter"
\NewDocumentEnvironment{references-chapter}{mmmmmmm}
{%
\renewcommand{\divisionnameptx}{#1}%
\renewcommand{\titleptx}{#2}%
\renewcommand{\subtitleptx}{#3}%
\renewcommand{\shortitleptx}{#4}%
\renewcommand{\authorsptx}{#5}%
\renewcommand{\epigraphptx}{#6}%
\chapter[{#4}]{#2}%
\label{#7}%
}{}%
%% Environment for a PTX "references" at the level of a LaTeX "chapter"
\NewDocumentEnvironment{references-chapter-numberless}{mmmmmmm}
{%
\renewcommand{\divisionnameptx}{#1}%
\renewcommand{\titleptx}{#2}%
\renewcommand{\subtitleptx}{#3}%
\renewcommand{\shortitleptx}{#4}%
\renewcommand{\authorsptx}{#5}%
\renewcommand{\epigraphptx}{#6}%
\chapter*{#2}%
\addcontentsline{toc}{chapter}{#4}
\label{#7}%
}{}%
%%
%% Styles for six traditional LaTeX divisions
\titleformat{\part}[display]
{\divisionfont\Huge\bfseries\centering}{\divisionnameptx\space\thepart}{30pt}{\Huge#1}
[{\Large\centering\authorsptx}]
\titleformat{\chapter}[display]
{\divisionfont\huge\bfseries}{\divisionnameptx\space\thechapter}{20pt}{\Huge#1}
[{\Large\authorsptx}]
\titleformat{name=\chapter,numberless}[display]
{\divisionfont\huge\bfseries}{}{0pt}{#1}
[{\Large\authorsptx}]
\titlespacing*{\chapter}{0pt}{50pt}{40pt}
\titleformat{\section}[hang]
{\divisionfont\Large\bfseries}{\thesection}{1ex}{#1}
[{\large\authorsptx}]
\titleformat{name=\section,numberless}[block]
{\divisionfont\Large\bfseries}{}{0pt}{#1}
[{\large\authorsptx}]
\titlespacing*{\section}{0pt}{3.5ex plus 1ex minus .2ex}{2.3ex plus .2ex}
\titleformat{\subsection}[hang]
{\divisionfont\large\bfseries}{\thesubsection}{1ex}{#1}
[{\normalsize\authorsptx}]
\titleformat{name=\subsection,numberless}[block]
{\divisionfont\large\bfseries}{}{0pt}{#1}
[{\normalsize\authorsptx}]
\titlespacing*{\subsection}{0pt}{3.25ex plus 1ex minus .2ex}{1.5ex plus .2ex}
\titleformat{\subsubsection}[hang]
{\divisionfont\normalsize\bfseries}{\thesubsubsection}{1em}{#1}
[{\small\authorsptx}]
\titleformat{name=\subsubsection,numberless}[block]
{\divisionfont\normalsize\bfseries}{}{0pt}{#1}
[{\normalsize\authorsptx}]
\titlespacing*{\subsubsection}{0pt}{3.25ex plus 1ex minus .2ex}{1.5ex plus .2ex}
\titleformat{\paragraph}[hang]
{\divisionfont\normalsize\bfseries}{\theparagraph}{1em}{#1}
[{\small\authorsptx}]
\titleformat{name=\paragraph,numberless}[block]
{\divisionfont\normalsize\bfseries}{}{0pt}{#1}
[{\normalsize\authorsptx}]
\titlespacing*{\paragraph}{0pt}{3.25ex plus 1ex minus .2ex}{1.5em}
%%
%% Styles for five traditional LaTeX divisions
\titlecontents{part}%
[0pt]{\contentsmargin{0em}\addvspace{1pc}\contentsfont\bfseries}%
{\Large\thecontentslabel\enspace}{\Large}%
{}%
[\addvspace{.5pc}]%
\titlecontents{chapter}%
[0pt]{\contentsmargin{0em}\addvspace{1pc}\contentsfont\bfseries}%
{\large\thecontentslabel\enspace}{\large}%
{\hfill\bfseries\thecontentspage}%
[\addvspace{.5pc}]%
\dottedcontents{section}[3.8em]{\contentsfont}{2.3em}{1pc}%
\dottedcontents{subsection}[6.1em]{\contentsfont}{3.2em}{1pc}%
\dottedcontents{subsubsection}[9.3em]{\contentsfont}{4.3em}{1pc}%
%%
%% Begin: Semantic Macros
%% To preserve meaning in a LaTeX file
%%
%% \mono macro for content of "c", "cd", "tag", etc elements
%% Also used automatically in other constructions
%% Simply an alias for \texttt
%% Always defined, even if there is no need, or if a specific tt font is not loaded
\newcommand{\mono}[1]{\texttt{#1}}
%%
%% \linkhilite macro: wraps hyperlink visible text for highlighting
\newcommand{\linkhilite}[1]{\underline{#1}}
%%
%% Following semantic macros are only defined here if their
%% use is required only in this specific document
%%
%% Used for warnings, typically bold and italic
\newcommand{\alert}[1]{\textbf{\textit{#1}}}
%% Used for inline definitions of terms
\newcommand{\terminology}[1]{\textbf{#1}}
%% End: Semantic Macros
%% Equation Numbering
%% Controlled by  numbering.equations.level  processing parameter
%% No adjustment here implies document-wide numbering
\numberwithin{equation}{section}
%% "tcolorbox" environment for a single image, occupying entire \linewidth
%% arguments are left-margin, width, right-margin, as multiples of
%% \linewidth, and are guaranteed to be positive and sum to 1.0
\tcbset{ imagestyle/.style={bwminimalstyle} }
\NewTColorBox{tcbimage}{mmm}{imagestyle,left skip=#1\linewidth,width=#2\linewidth}
%% Wrapper environment for tcbimage environment with a fourth argument
%% Fourth argument, if nonempty, is a vertical space adjustment
%% and implies image will be preceded by \leavevmode\nopagebreak
%% Intended use is for alignment with a list marker
\NewDocumentEnvironment{image}{mmmm}{\notblank{#4}{\leavevmode\nopagebreak\vspace{#4}}{}\begin{tcbimage}{#1}{#2}{#3}}{\end{tcbimage}%
}%% For improved tables
\usepackage{array}
%% Some extra height on each row is desirable, especially with horizontal rules
%% Increment determined experimentally
\setlength{\extrarowheight}{0.2ex}
%% Define variable thickness horizontal rules, full and partial
%% Thicknesses are 0.03, 0.05, 0.08 in the  booktabs  package
\newcommand{\hrulethin}  {\noalign{\hrule height 0.04em}}
\newcommand{\hrulemedium}{\noalign{\hrule height 0.07em}}
\newcommand{\hrulethick} {\noalign{\hrule height 0.11em}}
%% We preserve a copy of the \setlength package before other
%% packages (extpfeil) get a chance to load packages that redefine it
\let\oldsetlength\setlength
\newlength{\Oldarrayrulewidth}
\newcommand{\crulethin}[1]%
{\noalign{\global\oldsetlength{\Oldarrayrulewidth}{\arrayrulewidth}}%
\noalign{\global\oldsetlength{\arrayrulewidth}{0.04em}}\cline{#1}%
\noalign{\global\oldsetlength{\arrayrulewidth}{\Oldarrayrulewidth}}}%
\newcommand{\crulemedium}[1]%
{\noalign{\global\oldsetlength{\Oldarrayrulewidth}{\arrayrulewidth}}%
\noalign{\global\oldsetlength{\arrayrulewidth}{0.07em}}\cline{#1}%
\noalign{\global\oldsetlength{\arrayrulewidth}{\Oldarrayrulewidth}}}
\newcommand{\crulethick}[1]%
{\noalign{\global\oldsetlength{\Oldarrayrulewidth}{\arrayrulewidth}}%
\noalign{\global\oldsetlength{\arrayrulewidth}{0.11em}}\cline{#1}%
\noalign{\global\oldsetlength{\arrayrulewidth}{\Oldarrayrulewidth}}}
%% Single letter column specifiers defined via array package
\newcolumntype{A}{!{\vrule width 0.04em}}
\newcolumntype{B}{!{\vrule width 0.07em}}
\newcolumntype{C}{!{\vrule width 0.11em}}
%% tcolorbox to place tabular outside of a sidebyside
\tcbset{ tabularboxstyle/.style={bwminimalstyle,} }
\newtcolorbox{tabularbox}[3]{tabularboxstyle, left skip=#1\linewidth, width=#2\linewidth,}
%% Footnote Numbering
%% Specified by numbering.footnotes.level
%% Undo counter reset by chapter for a book
\counterwithout{footnote}{chapter}
\counterwithin*{footnote}{section}
%% Fancy Verbatim for consoles, preformatted, code display, literate programming
\usepackage{fancyvrb}
%% Pre-formatted text, a peer of paragraphs
\DefineVerbatimEnvironment{preformatted}{Verbatim}{}
%% More flexible list management, esp. for references
%% But also for specifying labels (i.e. custom order) on nested lists
\usepackage{enumitem}
%% Lists of references in their own section, maximum depth 1
\newlist{referencelist}{description}{4}
\setlist[referencelist]{leftmargin=!,labelwidth=!,labelsep=0ex,itemsep=1.0ex,topsep=1.0ex,partopsep=0pt,parsep=0pt}
%% hyperref driver does not need to be specified, it will be detected
%% Footnote marks in tcolorbox have broken linking under
%% hyperref, so it is necessary to turn off all linking
%% It *must* be given as a package option, not with \hypersetup
\usepackage[hyperfootnotes=false]{hyperref}
%% configure hyperref's  \href{}{}  and  \nolinkurl  to match listings' inline verbatim
\renewcommand\UrlFont{\small\ttfamily}
%% Hyperlinking active in electronic PDFs, all links without surrounding boxes and blue
\hypersetup{colorlinks=true,linkcolor=blue,citecolor=blue,filecolor=blue,urlcolor=blue}
%% Less-clever names for hyperlinks are more reliable, *especially* for structural parts
%% See comments in the code to learn more about the importance of this setting
\hypersetup{hypertexnames=false}
%%The  hypertexnames  setting then confuses the hyperlinking from the index
%%This patch resolves the incorrect links, see code for StackExchange post.
\makeatletter
\patchcmd\Hy@EveryPageBoxHook{\Hy@EveryPageAnchor}{\Hy@hypertexnamestrue\Hy@EveryPageAnchor}{}{\fail}
\makeatother
\hypersetup{pdftitle={Project Based Engineering Instrumentation}}
%% If you manually remove hyperref, leave in this next command
%% This will allow LaTeX compilation, employing this no-op command
\providecommand\phantomsection{}
%% Division Numbering: Chapters, Sections, Subsections, etc
%% Division numbers may be turned off at some level ("depth")
%% A section *always* has depth 1, contrary to us counting from the document root
%% The latex default is 3.  If a larger number is present here, then
%% removing this command may make some cross-references ambiguous
%% The precursor variable $numbering-maxlevel is checked for consistency in the common XSL file
\setcounter{secnumdepth}{3}
%%
%%
%% A faux tcolorbox whose only purpose is to provide common numbering
%% facilities for most blocks (possibly not projects, 2D displays)
%% Controlled by  numbering.theorems.level  processing parameter
\newtcolorbox[auto counter, number within=section]{block}{}
%%
%% This document is set to number PROJECT-LIKE on a separate numbering scheme
%% So, a faux tcolorbox whose only purpose is to provide this numbering
%% Controlled by  numbering.projects.level  processing parameter
\newtcolorbox[auto counter, number within=section]{project-distinct}{}
%% A faux tcolorbox whose only purpose is to provide common numbering
%% facilities for 2D displays which are subnumbered as part of a "sidebyside"
\makeatletter
\newtcolorbox[auto counter, number within=tcb@cnt@block, number freestyle={\noexpand\thetcb@cnt@block(\noexpand\alph{\tcbcounter})}]{subdisplay}{}
\makeatother
%%
%% tcolorbox, with styles, for FIGURE-LIKE
%%
%% figureptx: 2-D display structure
\tcbset{ figureptxstyle/.style={bwminimalstyle, middle=1ex, blockspacingstyle, fontlower=\blocktitlefont} }
\newtcolorbox[use counter from=block]{figureptx}[4]{lower separated=false, before lower={{\textbf{#1~\thetcbcounter}\space#2}}, phantomlabel={#3}, unbreakable, figureptxstyle, }
%% tableptx: 2-D display structure
\tcbset{ tableptxstyle/.style={bwminimalstyle, middle=1ex, blockspacingstyle, coltitle=black, bottomtitle=2ex, titlerule=-0.3pt, fonttitle=\blocktitlefont, breakable} }
\newtcolorbox[use counter from=block]{tableptx}[4]{title={{\textbf{#1~\thetcbcounter}\space#2}}, phantomlabel={#3}, unbreakable, tableptxstyle, }
%% listptx: 2-D display structure
\tcbset{ listptxstyle/.style={middle=1ex, blockspacingstyle, colback=white, colbacktitle=white, coltitle=black, colframe=black, titlerule=-0.3pt, toprule at break=-0.3pt, bottomrule at break=-0.3pt, sharp corners, fonttitle=\blocktitlefont} }
\newtcolorbox[use counter from=block]{listptx}[4]{title={{\textbf{#1~\thetcbcounter}\space#2}}, phantomlabel={#3}, breakable, listptxstyle, before upper app={\setparstyle}, }
%%
%% xparse environments for introductions and conclusions of divisions
%%
%% introduction: in a structured division
\NewDocumentEnvironment{introduction}{m}
{\notblank{#1}{\noindent\textbf{#1}\space}{}}{\par\medskip}
%% Graphics Preamble Entries
\usepackage{tikz, pgfplots} \usetikzlibrary{positioning,matrix,arrows} \usetikzlibrary{shapes,decorations,shadows,fadings,patterns} \usetikzlibrary{decorations.markings}
%% If tikz has been loaded, replace ampersand with \amp macro
%% Custom Preamble Entries, late (use latex.preamble.late)
%% Begin: Author-provided macros
%% (From  docinfo/macros  element)
%% Plus three from PTX for XML characters
\newcommand{\N}{\mathbb N} \newcommand{\Z}{\mathbb Z} \newcommand{\Q}{\mathbb Q} \newcommand{\R}{\mathbb R}
\newcommand{\lt}{<}
\newcommand{\gt}{>}
\newcommand{\amp}{&}
%% End: Author-provided macros
\begin{document}
%% bottom alignment is explicit, since it normally depends on oneside, twoside
\raggedbottom
\frontmatter
%% begin: half-title
\thispagestyle{empty}
{\titlepagefont\centering
\vspace*{0.28\textheight}
{\Huge Project Based Engineering Instrumentation}\\[2\baselineskip]
{\LARGE High Level Coding and Microcontrollers}\\
}
\clearpage
%% end:   half-title
%% begin: title page
%% Inspired by Peter Wilson's "titleDB" in "titlepages" CTAN package
\thispagestyle{empty}
{\titlepagefont\centering
\vspace*{0.14\textheight}
%% Target for xref to top-level element is ToC
\addtocontents{toc}{\protect\hypertarget{project-based-instrumentation}{}}
{\Huge Project Based Engineering Instrumentation}\\[\baselineskip]
{\LARGE High Level Coding and Microcontrollers}\\[3\baselineskip]
{\Large Carlos Montalvo}\\[0.5\baselineskip]
{\Large University of South Alabama Mobile, AL\footnote{William B. Burnsed Jr. Department of Mechanical, Aerospace and Biomedical Engineering, Associate Professor\label{frontmatter-1-1-2-1-1}}}\\[3\baselineskip]
{\Large Marine Feron Leabeater}\\[0.5\baselineskip]
{\Large Georgia College and State University, Milledgeville, GA\footnote{Lounsbury College of Education, EdD Graduate Student\label{frontmatter-1-2-2-1-1}}}\\[3\baselineskip]
{\Large Lisa Schibelius}\\[0.5\baselineskip]
{\Large Virginia Polytechnic Institute and State University, Blacksburg, VA\footnote{Department of Engineering Education, PhD Graduate Research Assistant\label{frontmatter-1-3-2-1-1}}}\\[3\baselineskip]
{\Large August 30, 2026}\\}
\clearpage
%% end:   title page
%% begin: copyright-page
\thispagestyle{empty}
\hypertarget{front-colophon}{}\vspace*{\stretch{2}}
\noindent{\bfseries Website}: \href{www.aerialsystems.org}{\linkhilite{Facility for Aerospace Systems and Technologies (FAST) Lab}}\par\medskip
\noindent\textcopyright{}2021\textendash{}2026\quad{}Carlos Montalvo\\[0.5\baselineskip]
This work is licensed under the Creative Commons Attribution-ShareAlike 4.0 International License. To view a copy of this license, visit \href{http://creativecommons.org/licenses/by-sa/4.0/}{\linkhilite{CreativeCommons.org}}\par\medskip
\vspace*{\stretch{1}}
\null\clearpage
%% end:   copyright-page
%
%
\typeout{************************************************}
\typeout{Section 1 Current Edition}
\typeout{************************************************}
%
\begin{sectionptx}{Section}{Current Edition}{}{Current Edition}{}{}{current-edition}
This textbook was last updated on August 30, 2026. The latest pdf can be found on \href{https://github.com/cmontalvo251/project-based-instrumentation/blob/main/main.pdf}{\linkhilite{Github}} while the most up to date web based version is on \href{https://cmontalvo251.github.io/project-based-instrumentation/}{\linkhilite{Github Pages}}\hyperlink{Github}{\linkhilite{[{\xreffont 33}]}}.%
\par
A note on AI: Note that many sections have been written with the help of Github’s Copilot \hyperlink{copilot}{\linkhilite{[{\xreffont 1}]}} and even Google’s Gemini \hyperlink{gemini}{\linkhilite{[{\xreffont 2}]}}. Originially their were plans to properly cite each section that was written with the help of AI but the authors lost track given it’s seamless integration with VScode and Android smart phones as well as any browser at this point. As such, it is important to note that after September of 2025 it is safe to assume that most if not all sections were written with the help of AI. The authors are aware of the ethical implications of this issue however are rigorously check all sections and text to ensure that the content is correct and free of errors.%
\end{sectionptx}
%
%
\typeout{************************************************}
\typeout{Section 2 Manuscript Changes}
\typeout{************************************************}
%
\begin{sectionptx}{Section}{Manuscript Changes}{}{Manuscript Changes}{}{}{manuscript-changes}
%
\begin{enumerate}
\item{}July 1st, 2020 - Original tutorials in Google Docs created%
\item{}July 15th, 2020 - Tutorials purchased by \href{https://tangibles-that-teach.gitbook.io/instrumentation-lab-manual/-MbMx70LQzRmEG_hS7Ld/}{\linkhilite{Tangibles that Teach}}%
\item{}December 21st, 2021 - Updated links for manuscript and hardware%
\item{}May 30th, 2022 - Tangibles that teach went out of business and chapter began the move to Github%
\item{}June 28th, 2022 - Work began on a Chromebook. Unfortunately the Figures folder is not back up on Git. As such a main\textunderscore{}latest.pdf has been created that's the latest full version. The main.pdf is the version created by the Chromebook so it has new chapters but none of the older chapters. Figure files are now backup on Git but only figures from the 'Voltage Potentiometer' are currently there.%
\item{}July 2nd, 2022 - All figures backed up and latest manuscript completed%
\item{}October 18th, 2022 - A pedometer lab has been added.%
\item{}November 8th, 2022 - Edited the Servo and feedback control servo lab to be one big lab with 3 parts%
\item{}January 8th, 2023 - Removed all mentions of Tangibles that Teach in the main body of the text.%
\item{}February 6th, 2023 - Changed an assignment description for LEDs and push buttons%
\item{}February 22nd, 2023 - Moved the Bluetooth module to be before the modules lab so this will end up being a few updates to make those projects more uniform. I also created a backup called Bluetooth original in case you wanted to go back to the other version.%
\item{}June 16th, 2023 - Cleaned up the \textasciigrave{}\textasciigrave{}changes needed" list and updated Method 3 with some new software updates.%
\item{}August 9th, 2023 - Edited the preamble for assignments and made a note in the servo lab%
\item{}September 5th, 2023 - Updated hyperlinks to show up on Microsoft Edge using underlines and different color text.%
\item{}September 7th, 2023 - Fixed the TL;DR section%
\item{}September 15th, 2023 - Updated a link in the DAQ lab%
\item{}September 20th, 2023 - Made the modules installation section a standalone tex file%
\item{}October 4th, 2023 - Added a new potentiometer photo%
\item{}November 16th, 2023 - Edited the lists of parts%
\item{}November 22nd, 2023 - Edited acceleration lab%
\item{}December 5th, 2023 - Edited the troubleshooting section%
\item{}January 12th, 2024 - Edited chapters 1 and 2%
\item{}January 22nd, 2024 - Edited the bootloader update notes in the troubleshooting chatper%
\item{}January 25th, 2024 - More edits to the bootloader notes%
\item{}March 20th, 2024 - Added a preamble about where this textbook is located%
\item{}April 27th, 2024 - Previously edited servo lab and pendulum lab just the assignment portion and then added this change log to the manuscript%
\item{}May 7th, 2024 - Title page changed and many sections moved around. Added a new results and discussion section following data analysis from course surveys of this course.%
\item{}August 12th, 2024 - Added a new requirement for the servo lab%
\item{}September 12th, 2024 - A few edits to assignment descriptions%
\item{}October 25th, 2024 - Updated the photocell lab to include the number of data points for the histogram%
\item{}November 11th, 2024 - Updated Method 3 quick list%
\item{}March 26th, 2025 - Updated servo, pitot probe and photocell assignments%
\item{}May 24th, 2025 - Made a few changes to chapters 2 and 19 to transition between the projects and the results and discussion about lab at home frameworks and project based teaching. Also overhauled the feedback control project.%
\item{}June 16th, 2025 - Updated the assignment directions for all projects.%
\item{}June 17th, 2025 - Added more theory to the photocell, thermistor and pitot probe assignment%
\item{}June 20th, 2025 - Added more figures and wiring diagrams to many assignments to help students wire up their circuits.%
\item{}June 21st, 2025 - Added a few problems from a standard instrumentation textbook. Probably need more statistics questions%
\item{}July 2nd, 2025 - Moved changes needed to the preamble. Added learning objectives and parts list for a few projects that had them missing.%
\item{}September 26th, 2025 - Updated the potentiometer assignment to include converting voltage to angle in degrees.%
\item{}October 1st, 2025 - Updated the thermistor modules assignment with the thermistor equations so it's easy to find in two places.%
\item{}November 12th, 2025 - Fixed an equation, link and figure in the feedback control assignment%
\item{}November 14th, 2025 - Fixed a broken link%
\item{}November 20th, 2025 - Added a GPS requirement to the integrating acceleration lab%
\item{}December 5th, 2025 - Added the settling time to the list of parameters for the time constant assignment and edited the overall assignment directions to have a better description of the video requirements. Also edited the servo second order systems equation for an aircraft to be stable since the equations were unstable open loop. I also moved the changes needed to Github. I also edited the purchase assignment to have bold text for some important information.%
\item{}February 2nd, 2026 - Updated the push button and LED assignment to fix a mistake in the voltage source description.%
\item{}February 5th, 2026 - Added an update from Jan 18th that I forgot to push which highlighted the \textasciigrave{}\textasciigrave{}Rubric" for each assignment. Also added a missing grade to the push button and LED assignment.%
\item{}February 9th, 2026 - Moved the future project ideas to Github%
\item{}February 26th, 2026 - Edited the trouble shooting guide for updating the CPB bootloader on Windows%
\item{}March 21st, 2026 - Book officially converted from LaTeX and published to new repo%
\item{}March 22nd, 2026 - Made a few major formatting changes and updated the assignment directions to be better viewable and digestable to students%
\item{}March 23rd, 2026 - Updated the servo lab assignment to fix a broken png and added a link to an equation.%
\item{}March 26th, 2026 - Added an I2C scanner code snippet to the troubleshooting guide as well as many many figures to help with troubleshooting. The IMU wiring photo was also updated to be correct. Also added annotations to the circuitpython kit photo and added the high level sound code to the bluefruit.%
\item{}March 27th, 2026 - Moved the thermistor, accelerometer angles and photocell equations into one place. Also added another simulation to the servo lab.Also added the pitot probe wiring diagram. Also added proper citations for all links. Also started the systems engineering chapter.%
\item{}March 30th, 2026 - Finished the systems engineering chapter and printed and deployed book.%
\item{}March 31st, 2026 - Added a new assignment to the photocell histogram lab and underlined all hyperlinks in LaTeX document%
\item{}April 2nd, 2026 - Edited all the titles to be shorter and made the parts lists\slash{}objectives lists rather than sections and also moved all youtube tutorials into an Extra Help section.%
\item{}April 3rd, 2026 - Updated the pitot wiring photo.%
\item{}April 4th, 2026 - Added more alerts to acceleration lab.%
\item{}April 9th, 2026 - Edited the photocell lab to explain that the resistor can be whatever value the student wants. Also added a CPX wiring diagram for the CPX and also added the voltage divider equations depending on where you measure the voltage.%
\item{}April 10th, 2026 - Added an extra explanation to the servo lab for going into safe mode when drawing too much current.%
\item{}April 13th, 2026 - Updated the servo lab to have annotations in the wiring diagram%
\item{}April 15th, 2026 - Added SFR to the Systems Engineering chapter%
\item{}April 16th, 2026 - Added a new requirement to grading appendices%
\item{}April 17th, 2026 - Fixed an issue with the thermistor equations%
\item{}April 22nd, 2026 - Fixed an equation in the servo lab%
\item{}April 23rd, 2026 - Fixed a typo in the thermistor lab%
\item{}April 26th, 2026 - Changed "plotting in Python" to just "plotting on your computer" to leave room for potentially having students plot everything in Excel\slash{}Sheets. Also added a placeholder for Method 6. Also added an appendix C to the assignment requirements%
\item{}May 4th, 2026 - May the Fourth be with you. Fixed a small error in the servo part 2 lab.%
\item{}May 13th, 2026 - Started adding more information about Arduino%
\item{}June 9th, 2026 - Edited content and title to reflect a more broad ecosystem of Arduino and Circuitpython. Edited men flags to md. Added Google Sheets and Octave to the plotting chapter. Finished the serial Method 6 chapter.%
\item{}June 10th, 2026 - Added two new chapters on GPS and Pressure sensors but they are not done.%
\item{}June 11th, 2026 - Finished the pressure sensor, GPS and load cell chapters.%
\item{}June 12th, 2026 - Added Arduino blinking to the getting started chapter and finished method 5 for both Arduino and CPX\slash{}CPB%
\item{}June 13th, 2026 - Reduced the size of a few jpg files to make the PDF smaller than 50 MB. Added the accelerometer steps for the arduino.%
\item{}June 21st, 2026 - Happy Summer Solstice and Father's Day - I added an Arduino example for the HX711 load cell chapter and started adding the GPS arduino steps%
\item{}June 22nd, 2026 - Finished adding the LPS22 pressure sensor and GPS chapter for Arduino.%
\item{}August 28th, 2026 - Edited the Getting started chapter due to some broken links and confusing assignment requirements.%
\end{enumerate}
%
\end{sectionptx}
%
%
\typeout{************************************************}
\typeout{Section 3 Changes Needed and Future Project Ideas}
\typeout{************************************************}
%
\begin{sectionptx}{Section}{Changes Needed and Future Project Ideas}{}{Changes Needed and Future Project Ideas}{}{}{future-changes}
Needed changes and Future Project Ideas are now tracked on \href{https://github.com/cmontalvo251/project-based-instrumentation/issues}{\linkhilite{Github}}%
\end{sectionptx}
%% begin: table of contents
%% Adjust Table of Contents
\setcounter{tocdepth}{1}
\renewcommand*\contentsname{Contents}
\tableofcontents
%% end:   table of contents
\mainmatter
%
%
\typeout{************************************************}
\typeout{Chapter 1 Introduction}
\typeout{************************************************}
%
\begin{chapterptx}{Chapter}{Introduction}{}{Introduction}{}{}{Preface}
\renewcommand*{\chaptername}{Chapter}
\begin{introduction}{}%
This textbook describes an instrumentation kit used for an Instrumentation and Experimental Methods course at the undergraduate level. This textbook has been designed with the student and faculty member in mind. The kit primarily discusses the CircuitPlayground Express\slash{}Bluefruit running CircuitPython and is used to teach fundamentals of instrumentation and provide a hands-on way of learning. Note that there are some references to Arduino in the text, but the focus is on the CircuitPython implementation. Engineering is usually taught in a traditional lecture format, involving theory in the classroom, homework outside of class, and routine examinations. Progressive forms of learning such as flipped classrooms and project based learning (PBL) have created new and fun ways for professors to interact with students and for students to be more involved in their learning. PBL provides a student-centered method of teaching and learning by posing problems for students to solve with the solution of the project being the primary goal and the theory a secondary goal.%
\par
The course begins with simple plotting and moves into data analysis, calibration and more complex instrumentation techniques such as active filtering and aliasing. This course is designed to get students away from their pen and paper and build something that blinks and moves as well as learn to process real data that they themselves acquire. There is no theory in these projects. It is all applied using the project based learning method. Students will be tasked with downloading code, building circuitry, taking data all from the ground up. By the end of this course students will be well versed in the desktop version of Python while also the variant CircuitPython designed specifically for microelectronics from Adafruit. Again, some chapters have references to the Arduino IDE and thus some students may also gain some knowledge and expertise in the Arduino programming language. After this course students will be able to understand Instrumentation at the fundamental level as well as generate code that can be used in future projects and research to take and analyze data. Python is such a broad and useful language that it will be very beneficial for any undergraduate student to learn this language. To the professors using this textbook, 1 credit hour labs are often hard to work into a curriculum and \textquotedblleft{}live\textquotedblright{} demonstrations in the classroom cost time and money that take away from other faculty duties. I've created this kit and textbook to be completely stand-alone. Students simply need to purchase the required materials and follow along with the lessons. These lessons can be picked apart and taught sequentially or individually on a schedule suited to the learning speed of the course. The authors hope that whomever reads and learns from this textbook will walk away with an excitement to tinker, code and build future projects using microelectronics and programming. The implementation of this kit and overall teaching method has received many positive reviews from students and is reflected in anonymous course evaluations.%
\par
Students can learn in many ways with a variety of different modalities \hyperlink{Gardner}{\linkhilite{[{\xreffont 11}]}}. As such, instructors can choose how to present course material and have students develop content-specific skill sets. College students enter the classroom with existing skills in multi-disciplinary learning that is practiced at the secondary school level \hyperlink{Ralph}{\linkhilite{[{\xreffont 12}]}}. College and university professors have the opportunity to use these existing skills as a foundation for their own instructional practices. Traditional lectures provide students with lectures on theory. The instructor then provides examinations to the students based on that theory. Using this format, an instructor can focus on mathematical principles, the so-called building blocks of engineering or provide specific applications to these models \hyperlink{Learning-Styles}{\linkhilite{[{\xreffont 13}]}}. In a traditional lecture format, the instructor presents theory to the students with academic problems rarely encountered in a students future career, building invaluable theoretical knowledge of concepts that appear disconnected from practical knowledge applied in the field \hyperlink{Yong}{\linkhilite{[{\xreffont 14}]}}. Historically in engineering education, there has been tension between theory and application \hyperlink{PBL1}{\linkhilite{[{\xreffont 15}]}}. Applications include case studies or practical engineering problems that emulate future careers in engineering. This provides students with engineering phenomena that they can see and hear. They can reason through problems intuitively, memorize facts using demos, build mathematical models, or build tangible objects themselves. With these different presentation methods, students are exposed to engineering phenomena in multiple ways rather than just on the whiteboard \hyperlink{Learning-Styles}{\linkhilite{[{\xreffont 13}]}}. Project Based Learning (PBL) provides an alternative for conventional teaching by facilitating problem solving for students in a group setting that requires communication, critical thinking, and creativity. These types of problems have shown that heightened learning can happen when students interact with tangible objects that represent the theory they are presented with in the classroom \hyperlink{PBL-Book}{\linkhilite{[{\xreffont 16}]}}. Rather than learning an equation for heat transfer along a one-dimensional pipe, the students can create a pipe with thermocouples and plot data from multiple sensors along the pipe. This connection to a tangible object reinforces learning and allows students to form bridges in comprehension between their theoretical knowledge and practical application (practical knowledge) of those theories in the field \hyperlink{TANGIBLEINTERFACES}{\linkhilite{[{\xreffont 17}]}}. Research has shown that a classroom that creates an integrated curriculum increases student satisfaction which immediately correlates to satisfactory student performance and increased graduation rates \hyperlink{StudentSatisfaction}{\linkhilite{[{\xreffont 18}]}}. This is an example of the benefits of active student engagement versus passive student engagement. While the traditional lecture format is vital to students understandings, it is understood that application of the concepts covered in lecture are required to reach higher levels of learning \hyperlink{Armstrong}{\linkhilite{[{\xreffont 19}]}}.%
\par
Hands-on projects where students interact with tangible objects is not a new form of teaching. The laboratory environment has been around for decades. However, typically the classroom environment is separate from the laboratory environment. There has been little synergy between the two learning environments. A Mechanical Engineering degree is likely to have around three to four labs in various disciplines. In these courses, students perform an experiment in groups. They take and analyze data as well as create a report documenting their findings. Many of these laboratory experiments of course are choreographed for the student by the instructor. They follow a script and perform the experiment without building anything themselves. The students take no ownership over the experiment and there is no creativity built-in. Over the years, the nature of laboratories has changed including the lack of clear learning objectives \hyperlink{Labs}{\linkhilite{[{\xreffont 20}]}}. Furthermore, this laboratory environment requires a significant amount of instrumentation and hardware to implement. For example, a shock tube or steam pump can cost tens of thousands of dollars. The maintenance and up-keep alone is not practical for smaller institutions. The teaching investment required to prepare the lab every year and the lab itself (often on the order of three hours) can be a time consuming task for a tenure track faculty member who often has a large research load.%
\par
The so-called "Lab at Home" hardware kits are becoming more and more common in the classroom to ease the burden on the instructor and institution itself \hyperlink{LabAtHome}{\linkhilite{[{\xreffont 21}]}}. These kits are small enough to be purchased and shipped to a student for them to perform experiments remotely rather than in the classroom itself \hyperlink{LabAtHome1}{\linkhilite{[{\xreffont 22}]}}. They can also be brought to the classroom to be used as a personal demo aid \hyperlink{LabAtHome2-EE}{\linkhilite{[{\xreffont 23}]}}. In this sense, the take-home lab kit serves as a bridge between theory based lectures and a laboratory setting \hyperlink{LabAtHome3}{\linkhilite{[{\xreffont 24}]}}. This allows both instructors and students to engage with content in a way that promotes enduring understandings and practical application of theoretical knowledge\hyperlink{LabAtHome4-EE}{\linkhilite{[{\xreffont 25}]}}. Since the kits can be used remotely, they can also be used for distance learning courses or other asynchronous activities as well as during the COVID-19 pandemic.%
\par
Two home kits in particular are directly related to instrumentation and circuits. Cyganski and Nicoletti \hyperlink{LabAtHome4-EE}{\linkhilite{[{\xreffont 25}]}} for example created a new curriculum for first year electrical engineering students by creating live demonstrations for the students. Manijikian and Simmons \hyperlink{LabAtHome2-EE}{\linkhilite{[{\xreffont 23}]}} however, combined a popular commercial microcomputer board with their own custom-designed interface board in a kit that students retain for the duration of the course for both in-lab and at-home assignments. The take-home kits however were based on the Motorola 68HC11EVB microcomputer which is programmed in assembly language. Assembly language is a rather difficult language to learn at the undergraduate level. Even with such a difficult learning curve, however, student feedback on their approach was positive \hyperlink{LabAtHome2-EE}{\linkhilite{[{\xreffont 23}]}}.%
\par
It is clear that using a lab at home kit is useful for the instrumentation classroom, however the programming language to be used is a subject of debate in faculty meetings and computing committees. There are multiple languages in use today of varying complexities and use cases. There are scripting languages like Python, Ruby and MATLAB, object oriented languages like Java and C++ and compilation languages like Fortran and C. Note, this is not an exhaustive list. A recent study showed that scripting languages (Python, Ruby, MATLAB) enable writing more concise code while compilation languages (C, Fortran) create smaller executables \hyperlink{codetype}{\linkhilite{[{\xreffont 26}]}}. MATLAB is often considered in engineering given its success in industry and its ability to perform numerical simulations and plotting with ease. Python, however, has become more and more popular with numerical toolboxes like numpy and plotting toolboxes like matplotlib that are free and easy to install in integrated development environments like Thonny or Spyder \hyperlink{IDE}{\linkhilite{[{\xreffont 27}]}}. The Tiobe Index of Programming \hyperlink{Tiobe}{\linkhilite{[{\xreffont 28}]}} has the top 3 programming languages listed as Python, C, and Java. MATLAB is \#20 on the list. In 2004, Hans Fangohr wrote that MATLAB is much better suited than C for engineering computing but the best choice in terms of clarity and functionality is provided by Python \hyperlink{CodeComparison}{\linkhilite{[{\xreffont 29}]}}. It seems then that it would be more practical for educators even in engineering to teach the most popular language. This helps with transferability of skills and has future implications for a students' career. The scripting aspect of Python lowers complexity and allows students to learn the language quickly to apply it as a tool rather than getting stuck memorizing syntax and compilation rules. The language also comes at no cost to the students. This is a plus for students already bearing heavy financial burdens to attend a university which includes tuition and institutional fees.%
\par
Given the success of other kits in many classes and the popularity of Python, the University of South Alabama (USA) has implemented such a kit for Instrumentation \& Experimental Methods (ME 316) using CircuitPython. CircuitPython is a derivative of Python written for embedded systems and designed to simplify experimenting and learning to code \hyperlink{CircuitPython}{\linkhilite{[{\xreffont 30}]}}. As mentioned previously, some experiments have a supplemental Arduino component in the event the students want to use the more robust yet more complex Arduino programming language\hyperlink{Arduino}{\linkhilite{[{\xreffont 35}]}}. The learning objectives for ME316 include: statistics, dynamic response of measurement systems, operational amplifiers, signal conditioners and fundamentals of microprocessors. This course could be taught with theory as the main focus of the lecture. However, this course includes the use of an instrumentation kit and Project Based Learning methods in addition to theory in the classroom. Each student taking the course purchases the kit and downloads a free accompanying project list \hyperlink{Gitbook}{\linkhilite{[{\xreffont 31}]}}. Every Friday, the students complete one of the multiple projects. The following week they submit a report which includes a demo of the working project in the form of photos and videos, any plots generated if they are required to take data, all code used and a small write-up explaining their findings. The sections that follow describe the kit in more detail as well as student evaluations who have taken the course.%
\par
Note that, the Adafruit Learn page contains many tutorials on how to operate the CircuitPlayground \hyperlink{Adafruit}{\linkhilite{[{\xreffont 32}]}}. However, Adafruit sells much more than just the CPX and thus it is often difficult for anyone to find the correct tutorial needed for the CPX. A simple search on the Adafruit Learn System yields 21 results for "servo" just on the first page with over 48 pages of results. The tutorial needed to run a servo with CircuitPython is the 10th result and it's for a different development board. Although the software works with the CPX, it does not explicitly say so. Due to the complexity of some of these systems, documentation was made and freely given to the community \hyperlink{Gitbook}{\linkhilite{[{\xreffont 31}]}}. This documentation was custom built using a combination of multiple sources across the internet. All software for the kit is also on \href{www.github.com/cmontalvo251/Microcontrollers.git}{\linkhilite{Github}} in a separate repository. Each chapter in the book contains one project for the students to complete on their own. A list of chapters and a brief description of the assignment is shown in the table of contents. Currently there are about 20 projects with multiple subsections for the students to work on during the \textquotedblleft{}Funday Friday\textquotedblright{} lab days.%
\end{introduction}%
\end{chapterptx}
%
%
\typeout{************************************************}
\typeout{Chapter 2 Course Description}
\typeout{************************************************}
%
\begin{chapterptx}{Chapter}{Course Description}{}{Course Description}{}{}{course-descripion}
\renewcommand*{\chaptername}{Chapter}
\begin{introduction}{}%
The course is designed to coincide with lecture content in a standard engineering instrumentation course. The students are guided through numerous projects detailed in the subsequent chapters (See \hyperref[purchase-equipment]{\linkhilite{Chapter~{\xreffont\ref{purchase-equipment}}}}-\hyperref[second-order]{\linkhilite{Chapter~{\xreffont\ref{second-order}}}}). During the semester and in addition to the weekly projects, there is a large project which involves some systems engineering techniques (See \hyperref[systems-engineering]{\linkhilite{Chapter~{\xreffont\ref{systems-engineering}}}}). Finally, a results and discussion chapter details the over analysis of the course as experienced by anonymous survey results and quotes from actual students (See \hyperref[results-discussion]{\linkhilite{Chapter~{\xreffont\ref{results-discussion}}}}).%
\par
The specific course learning objectives, which include statistics, dynamic responses, operational amplifiers, and electronic circuits both analog and digital are shown below:%
\begin{enumerate}
\item{}Apply statistical concepts of error and uncertainty analysis using normal,\textquotedblleft{}t\textquotedblright{} distribution and the \(\chi^2\) distribution%
\item{}Use propagation of error methods to determine the uncertainty of calculated quantities%
\item{}Apply the concepts of harmonic response to predict the response of measuring systems to input signals%
\item{}Apply the fundamentals of operational amplifiers and electronic circuits to design signal conditioning circuits including%
\item{}Apply microprocessor fundamentals to gain an understanding of digital circuits used for digital to analog and analog to digital conversion%
\end{enumerate}
%
\par
This course is not designed to be a traditional lecture format, rather hands-on projects will be assigned every Friday using the lab at home kit. The projects every Friday are created to enhance the course learning objectives as dictated above.%
\par
The course is taught on a 2+1 style where the course meets three days a week for 50 minutes each on Monday, Wednesday and Friday. On Monday and Wednesday the class is similar to what would typically be found in a standard lecture format. The instructor goes through theory and examples on a white board. On Fridays, the students come to class with their hardware kits and laptops to perform one of 20+ projects defined in the tutorials in this textbook. Fridays are nicknamed "Funday Fridays" to differentiate the format of the class and to get students excited about the projects. On Fridays, the instructor can walk around the classroom and assist students on their projects as well as highlight fundamental concepts. This is done by showing them what they learned on a whiteboard with their lab at home kit. During the COVID-19 pandemic when Universities were meeting via Zoom, class time on Friday was used to assist students on their projects while they showed their circuits via webcam in a virtual space. Since all students purchased their own kit, there was no issue with students getting hands-on experience even during COVID-19.%
\par
It is important to note here that Project Based Learning (PBL) in this course did not replace existing teaching frameworks; rather, PBL supplemented existing standard teaching measures by allowing students to apply theory to practice. The role of connectivism comes from the online networking tools that were added to the curriculum. The tutorials and assignments as well as example code are all hosted online on Github \hyperlink{Github}{\linkhilite{[{\xreffont 33}]}}. The availability of these tutorials and software opens the students to an ecosystem of free and open source software. Furthermore, the Adafruit Learning System is also used occasionally which opens them to an ecosystem of hardware and software tutorials as well as Adafruit IO (an internet of things site) and the Adafruit Forums \hyperlink{Adafruit}{\linkhilite{[{\xreffont 32}]}}. The Adafruit Forums allow the students to comment and communicate with other groups outside the university that utilize the same hardware kits. If the instructor cannot find a solution online in a Wiki or Tutorial page, a student or the instructor can post on the forum and wait for a response from the larger Adafruit community.%
\par
Students in this course are also required to do a final project where they must create something new by using the microcontroller in the kit plus three other electrical components. One of the three components must be something not originally included in the kit. The students may also work in a group. PBL typically consists of ownership, creativity, collaboration and critical thinking. These four aspects are clearly a big part of this lab at home kit and the final project. They can work in groups: collaboration. They must make something new: creativity. They must build the item themselves: ownership. They must apply fundamental principles of the course: critical thinking. In addition to PBL aspects, the students are exposed to the larger network of learning through Github, the Adafruit Learning system and Forums as well as general Wikis online. This network of learning is seen in the connectivism framework. The chapters that follow detail essential systems engineering techniques for the final project as well as the individual weekly projects. The ending chapters explain results from anonymous course surverys from students taking the class as well as conclusions drawn from the authors and faculty who have taught the course.%
\end{introduction}%
\end{chapterptx}
%
%
\typeout{************************************************}
\typeout{Chapter 3 Project Based Engineering}
\typeout{************************************************}
%
\begin{chapterptx}{Chapter}{Project Based Engineering}{}{Project Based Engineering}{}{}{systems-engineering}
\renewcommand*{\chaptername}{Chapter}
\begin{introduction}{}%
Systems engineering is a branch of engineering used to help in the design of complex systems. The general idea is to break down a complex system into smaller more manageable pieces. This is done by following a life cycle which includes the following steps: requirements analysis, functional analysis and allocation, design synthesis, and system analysis and control. A diagram for a systems engineering life cycle can be seen below. \begin{figureptx}{Figure}{The Systems Engineering Life Cycle: Example Verification and Validation (V\&V) Process\hyperlink{vandv}{\linkhilite{[{\xreffont 3}]}}}{v-and-v}{}%
\begin{image}{0.05}{0.9}{0.05}{}%
\includegraphics[width=\linewidth]{external/VandV.jpg}
\end{image}%
\tcblower
\end{figureptx}%
 The left side of the V\&V diagram highlights all of the required design while the right side details the build. For undergraduate engineering courses the Design, Build, Test (DBT) framework is often used and has been used at the University of South Alabama for many years with much success. The systems engineering life cycle is a more detailed version of the DBT framework. The systems engineering life cycle is used in many industries including aerospace, automotive, and software development. The sections that follow explain how this instrumentation course utilizes the DBT framework and the systems engineering life cylce.%
\end{introduction}%
%
%
\typeout{************************************************}
\typeout{Section 3.1 Project Schedule}
\typeout{************************************************}
%
\begin{sectionptx}{Section}{Project Schedule}{}{Project Schedule}{}{}{project-overview}
A standard college semester is 16 weeks with the 17th week typically reserved for finals week. In order to help the students with engineering life cycle, the requirements are given to the students and only a prototype is requested which removes the very top of both the left and right side of the V\&V diagram. At the undergraduate level the following steps are used for a standard 16 week curriculum.%
\begin{enumerate}
\item{}\terminology{Week 1:} \emph{Title Page and Introduction} - Students are given an explanation of the project scope, rubric and are asked to split into teams to create a title page, team name, a simple hand sketch and an introduction of what they intend to build for the semester. In this course, the project is to design, build and test a data acquisition system with at least 3 electrical components not counting the CPU, flash memory, or power supply. This also includes both mechanical, electrical and software ubsystems so that students practice both hardware and software integration with mechanical components. The teams are typically 1-3 students each depending on class size.%
\item{}\terminology{Week 2:} \emph{Conceptual Design and System Functional Review (SFR)} - Students are given an explanation of the major project requirements and asked to create a simple mission goal for their project as well as include a list of specific requirements which is based on the project scope given by the instructor as well as any additional requirements the students feel are necessary. Course specific requirements are explained in \hyperref[requirements]{\linkhilite{Section~{\xreffont\ref{requirements}}}}. This design phase also includes a list of subsystems and components that will be needed to complete the project which requires the students to create a Block Definition Diagram (BDD) as well as an Internal Block Diagram (IBD) detailed BDDs and IBDs for the mechanical, electrical and software subsystems. The instructor of the course must then review the requirements list and provide feedback to the students. This consitutes and satisfies the System Functional Review (SFR) so that the students do not begin designing a project that is innappropriate, out of scope or too difficult for a 16 week semester.%
\item{}\terminology{Week 3:}\emph{Preliminary Design} - For the preliminary design phase the task more detailed diagrams for every subsystem. This means the mechanical lead must create a general top down layout of the system with dimensions, the electrical lead must figure out what components they plan on using and what they will connect to. The software lead must create a psuedo code flow chart to determine the flow of their main software routine.%
\item{}\terminology{Weeks 4-6:}\emph{Detailed Design} - The detailed design phase is broken into 3 separate weeks. For this course the students need to design the mechanical, electrical and software subsystems creating schematics, wiring diagrams and a list with links of all software needed to run their system. The students also need to create a detailed bill of materials with prices, weights and links. At this point the system is ready to be built.%
\item{}\terminology{Week 7:}\emph{Critical Design Review (CDR)} - The critical design review is a presentation to the instructor and peers where the students present their design and receive feedback. This is an important step in the design process as it allows the students to get feedback on their design before they start building. This can help them avoid costly mistakes and ensure that they are on the right track. The students are expected to present their BDDs, IBDs, detailed diagrams, bill of materials and any other relevant information about their design. The students should also be prepared to answer questions from the instructor and peers about their design.%
\item{}\terminology{Week 8:}\emph{Build Plan, Purchase materials + Break} - In the Fall semester, Fall break typically lands on Week 8 and in the spring Week 8 is typically spring break. As such Students are asked to create a build plan which is not something to be turned in but just something helpful to keep students on track as well as purchase all parts necessary while they are on holiday. If parts are purchased the Friday before break that typically gives the retail stores over 5 business days or 9 calendar days to ship the parts to the students before they return to class. This way students can immediately start building and testing during Week 9 when they return. Of course, in practice I have found many procrastinating students who don't order parts until well after they return and end up rushing to get their parts delivered and operational in time. However, it is not the job of the instructor to micromanage or police students. This is where time management skills are important.%
\item{}\terminology{Weeks 9-11:}\emph{Build and Subsystem Testing} - The build process includes building everything component by component and testing each subsystem as it is built. This is important to ensure that each subsystem works properly independently before integrating everything together. For Instrumentation and Experimental Methods the students will test each electrical component one by one and test the software for that specific component. The mechanical system is something to be built piece by piece and integrated in Week 12. Again, time management is key.%
\item{}\terminology{Weeks 12-13:}\emph{System Integration and Test Readiness Review (TRR)} - The integration process includes integrating all subsystems together and testing the entire system as a whole. For Instrumentation and Experimental Methods, the students will integrate all electrical and software subsystems (all 3 components) and perform a test without the mechanical subsystem in the way initially. The mechanical subsystem must also be fully completed at this time.%
\item{}\terminology{Weeks 14-16:}\emph{Final Testing and Validation} - The final testing and validation phase includes performing the final test of the entire system. For Instrumentation and Experimental Methods the students will fully integrate the mechanical subsystem and perform a final test of the entire system. Note that 3 weeks are alloted for this phase to account for weather delays, part delays and any other unforeseen issues that may arise. Many times the system fails sometimes catastrophically as in a crash or an overheated component and the students need time to rebuild and retest before the final presentation.%
\item{}\terminology{Week 17:}\emph{Final Report and Systems Verification Review (SVR)} - Although a semester is exactly 16 weeks, the final report and presentation is typically due the first week of finals which ends up being Week 17. The final report includes a summary of the entire project including all diagrams, analyses, test results and conclusions. The final presentation is a 10-15 minute presentation to the class and any invited guests which summarizes the entire project. Note that the final report and presentation is typically worth 30-40\% of the overall grade so it is important that students take this seriously and put in their best effort. Furthermore, there are no exams or quizzes in these project based courses so the final report and presentation is the only way to demonstrate what the students have learned throughout the semester. This gives the students ownership of the  project and motivates them to do their best work.%
\end{enumerate}
%
\end{sectionptx}
%
%
\typeout{************************************************}
\typeout{Section 3.2 Specific Final Project Requirements}
\typeout{************************************************}
%
\begin{sectionptx}{Section}{Specific Final Project Requirements}{}{Specific Final Project Requirements}{}{}{requirements}
The final project will utilize the Circuit Playground (or any other comparable or better microcontroller such as the Arduino or Raspberry Pi) and at least 3 other electrical components to move a mechanical system that helps people in some way. The specific requirements are shown below. Note that every group will have a few additional requirements specific to their project based on how creative they decide to be.%
\begin{enumerate}
\item{}System must incorporate 3 electrical components (accelerometer, lights, sound and microphone to name a few). Your power source, microcontroller, speaker amplifier or the Crikit does not count as a component. An interface is not considered a component%
\item{}System must be unique as well as built and designed by your team (Tutorials may be used but you must make the system unique in some way).%
\item{}System must have one of the three components as an external part to the main microcontroller and not be something already taught in a weekly project (buttons, servos or potentiometers to name a few).%
\item{}System must take some sort of measurement to be used in the code.%
\item{}System must utilize a data acquisition method so that data can be plotted on your computer as a function of time.%
\item{}System must help someone in some way.%
\item{}System must have some sort of mechanical component to it and cannot be purely software and electrical components.%
\end{enumerate}
%
\par
\terminology{Microcontroller Choice:} Note, you have the choice to use any type of microcontroller (MCU) for your final project. In my experience teaching this course, many groups choose to use an Arduino UNO\slash{}Mega\slash{}Nano or ESP32 for their final project. Although the CPX\slash{}CPB is a cheap MCU that has many built-in sensors, it is severely lacking in stability, memory and power capabilities. The course is designed around CircuitPython and the CPX\slash{}CPB due to the simplicity of the programming language and versatility of the MCU. However, for your final project, using an Arduino will open your project up to a larger world of possibilities and your CPX\slash{}CPB programming\slash{}hardware skills will transfer very easily to Arduino. Combine that with GenAI and you’ll have an easy time getting code to work on the Arduino. I would say about half the class ends up using an Arduino so although you can absolutely make some very cool projects with the CPX\slash{}CPB, my experience is that about 50\% of the class would benefit from using an Arduino. Why isn’t the entire class just built around Arduino? Because for \textdollar{}20 you get two push buttons, a slide switch, an accelerometer, a temperature sensor, 8 neopixels, PWM support, a light sensor, multiple 16 bit ADCs, a microphone, a speaker and a much simpler programming language. The CPX also has an IR sensor and the CPB supports bluetooth communication. I think students get more out of learning on a CPX\slash{}CPB and then they can take those skills to something more reliable like the Arduino MCU or even something more complex like a Raspberry Pi.%
\end{sectionptx}
%
%
\typeout{************************************************}
\typeout{Section 3.3 Project Examples}
\typeout{************************************************}
%
\begin{sectionptx}{Section}{Project Examples}{}{Project Examples}{}{}{project-examples}
The project schedule and requirements previously shared is a general framework that can be applied to any project and has been for many many semesters. When first formulating this course I came up with a few project ideas for students to use as inspiration but of course the \href{https://learn.adafruit.com/}{\linkhilite{Adafruit Learning System}} provides a wealth of resources for project development\hyperlink{Adafruit}{\linkhilite{[{\xreffont 32}]}}. Still, here are a few examples that are shown below that can still be used as inspiration for future students. I've also included a few examples of projects that students have built in the past to give you an idea of what is possible with the CPX and a little creativity.%
\begin{enumerate}
\item{}A light sensor to turn a servo and then external leds%
\item{}Pressure plate outside that sends a text when package is stolen%
\item{}Log weather data on your wifi network with a raspberry pi and Adafruit IO\hyperlink{Adafruit}{\linkhilite{[{\xreffont 32}]}}%
\item{}Use a computer fan you own to create a solder station filter%
\item{}\alert{The examples below are from peers that have taken the class in the past}%
\item{}Little Robot with Adafruit Crickit%
\item{}Automatic trash can - Uses an ultrasonic sensor and a servo to open a trash can%
\item{}Light saber - an LED strip is used%
\item{}Automatic shades - a light sensor is used to turn a relay on which runs a dc motor to open your blinds%
\item{}Auto plant watering device - uses a waterproof pump and a humidity sensor to water a plant%
\item{}Auto dog bowl water feeder - an ultrasonic sensor is used to detect water level and open a solenoid valve to let water flow into a bowl%
\item{}Mobile for a child - a rotating platform with toys is attached to a continuous servo%
\item{}Laser Tag - Two CPXes are used to play laser tag by sending and receiving IR signals%
\item{}Mailbox - When mail is delivered light enters from the box from opening the mail box. This causes the CPX to send an IR signal to another CPX with says “You’ve got mail” with the built in speaker%
\item{}Door Opener with Dog Bark - using a 5V relay and a large motor a door open when a sound byte over a certain value is read using a microphone%
\item{}Shade Sense - Open blinds at a particular time or when sun light is bright enough using a light sensor and a servo%
\item{}Security Box - The CPX is used to rotate a “lock” using a servo by pressing buttons in a predetermined sequence%
\item{}Underglow that Bounces to Music - An LED strip is used that changes color by using the microphone to detect sound levels%
\item{}Auto-soap dispenser - A servo presses down on soap when a hand is placed under an ultrasonic sensor.%
\item{}Remote control car - A Circuit Playground Bluefruit (CPB) and Cricket are used to control a remote control car. Continuous servos are used to control the rear wheels and a front servo is used to control the front linkage. The CPB has bluetooth to allow control of the car using Bluetooth and an app on your phone.%
\item{}Fortune teller - Student bought a coin machine so you could put a quarter in and it would then tell a fortune by rotating one of 4 servos.%
\item{}Putt putt \slash{} skeeball \slash{} miniature bowling - 2 sonar sensors can be used to get speed of the ball%
\item{}Tattle tale robot - Students used a motion sensor to detect when someone walked in the room%
\end{enumerate}
%
\end{sectionptx}
%
%
\typeout{************************************************}
\typeout{Section 3.4 Pitot Probe Example}
\typeout{************************************************}
%
\begin{sectionptx}{Section}{Pitot Probe Example}{}{Pitot Probe Example}{}{}{pitot-probe-example}
\begin{introduction}{}%
One example project that I will use to explain future diagrams is from a project I built a few years ago. The overall mission goal was to \textquotedblleft{}use a pitot probe to measure the velocity of my bicycle while riding around town\textquotedblright{}. My initial hand sketch showed the pitot probe on the front of my bike which would use the pressure difference between the static and dynamic ports to calculate velocity (See \hyperref[pitot]{\linkhilite{Chapter~{\xreffont\ref{pitot}}}}). The project was successful and I was able to get some good data on my bike rides around town. The project also included a mechanical subsystem which was a simple plastic case for the electronics and zip ties to mount the pitot probe as well as a data acquisition system using the CPB and a Bluetooth module to send the data to my phone in real time. The figure below showed my initial sketch for the project. \begin{figureptx}{Figure}{Initial Sketch of Pitot Probe Project}{pitot-sketch}{}%
\begin{image}{0.05}{0.9}{0.05}{}%
\includegraphics[width=\linewidth]{external/pitot_sketch.png}
\end{image}%
\tcblower
\end{figureptx}%
 After the sketch it's important to be sure that all requirements set by the instructor are satisfied. In this case I needed 3 components so I decided to use the bluetooth transmitter, pitot probe, neopixels and button on the CPB which ended up being 4 components but I thought it would be nice to have visual feedback so I added the neopixels. The pitot probe would be the external component and the button would be used to calibrate the pitot probe. Finally, the bluetooth would be used to send the data to my phone. Although the course has a set of general requirements defined in \hyperref[requirements]{\linkhilite{Section~{\xreffont\ref{requirements}}}}, I still needed to add a few specific requirements as shown below.%
\begin{enumerate}
\item{}The system shall measure the windspeed of the vehicle.%
\item{}The system shall be calibrated on startup or during operation by pressing a button.%
\item{}The system will be powered during operation of the vehicle.%
\item{}The system shall not adversely affect the vehicle performance.%
\item{}The system shall send windspeed data to the operators smart phone.%
\item{}The system shall light up during calibration to indicate to the user that the system shall not be touched.%
\end{enumerate}
Note however that I failed to satisfy the requirement of \emph{System must have one of the three components as an external part to the main microcontroller and not be something already taught in a weekly project}. This is because the pitot probe project is a weekly project (See \hyperref[pitot]{\linkhilite{Chapter~{\xreffont\ref{pitot}}}}). However, I still think this is a good example of how to use the systems engineering life cycle to design an electrial subsystem so I will continue with this example.%
\end{introduction}%
%
%
\typeout{************************************************}
\typeout{Subsection 3.4.1 Diagrams}
\typeout{************************************************}
%
\begin{subsectionptx}{Subsection}{Diagrams}{}{Diagrams}{}{}{diagrams}
Upon completion of the initial sketch and the requirements it is now time to move away from a mission goal with requirements and begin the conceptual design process. This requires using the requirements to figure out what components and modules are necessary to accomplish the system you are trying to build. In order to do this, many diagrams are going to be created namely Block Definition Diagrams (BDDs) and Internal Block Diagrams (IBDs). I am going to be using \href{https://www.drawio.com/}{\linkhilite{DrawIO}} to create these diagrams\hyperlink{drawio}{\linkhilite{[{\xreffont 4}]}}.%
\end{subsectionptx}
%
%
\typeout{************************************************}
\typeout{Subsection 3.4.2 Electrical Subsystem}
\typeout{************************************************}
%
\begin{subsectionptx}{Subsection}{Electrical Subsystem}{}{Electrical Subsystem}{}{}{pitot-probe-example-4}
Block Definition Diagrams (BDD) are used to show the different components and modules of a system and how they are related to each other. For the pitot probe project, the electrical BDD is the easiest and most intuitive to complete first and is shown below. \begin{figureptx}{Figure}{BDD of Electrical Components}{bdd-electrical}{}%
\begin{image}{0.05}{0.9}{0.05}{}%
\includegraphics[width=\linewidth]{external/BDD_Electrical.jpg}
\end{image}%
\tcblower
\end{figureptx}%
 The BDD shows the different components of the system which includes general names for all subsystems included. Notice that instead of writing pitot probe specifically I wrote \textquotedblleft{}Windspeed Sensor\textquotedblright{}. This was I can elect to change to something different later on in the design process if I decide to do a last minute design change. Once you have the BDD completed it is time to move on to the Internal Block Diagram (IBD) which shows how the different components are connected to each other. The IBD for the electrical subsystem is shown below. \begin{figureptx}{Figure}{IBD of Electrical Components}{ibd-electrical}{}%
\begin{image}{0.05}{0.9}{0.05}{}%
\includegraphics[width=\linewidth]{external/IBD_Electrical.jpg}
\end{image}%
\tcblower
\end{figureptx}%
 The IBD shows how the different components are connected to each other. In this case, the Windspeed Sensor is connected to the CPU, the Notification UI(most likely neopixels) is connected to the CPU as well as the Calibration UI (most likely a button) and the CPU will send data to the Wireless Communications module (most likely bluetooth). The IBD also shows the flow of data between the different components and using annotations via text boxes also shows what information is passed from one block to the next. Notice also that every block shown in the BDD is in the IBD except for the central block which is merely the name of the system. Finally, I want to draw attention to the fact that I do not have a data storage module because all data will be sent to my smart phone and logged there which satisfies one of my major requirements. The IBD is an important step in the design process as it helps to ensure that all components are properly connected and that the flow of data is correct. It also helps to identify any potential issues with the design before moving on to the detailed design phase.%
\end{subsectionptx}
%
%
\typeout{************************************************}
\typeout{Subsection 3.4.3 Software Subsystem}
\typeout{************************************************}
%
\begin{subsectionptx}{Subsection}{Software Subsystem}{}{Software Subsystem}{}{}{subsec-Software-Subsystem}
Once you have completed the electrical BDD and IBD the software IBD is almost identical. The only difference is you need to first identify which components require software. For example, the power supply does not require any software but the wireless communication module does. The software BDD is shown below.%
\begin{figureptx}{Figure}{BDD of Software Components}{bdd-software}{}%
\begin{image}{0.05}{0.9}{0.05}{}%
\includegraphics[width=\linewidth]{external/BDD_Software.jpg}
\end{image}%
\tcblower
\end{figureptx}%
Notice that every software name contains () to indicate that it will eventually be a function. This is useful at this stage to help later on in the preliminary design phase. The CPU requires much software to run the system but since we're programming everything in a very high level programming language like Python, CircuitPython or Arduino, we're assuming that the only software required by the user is just a main loop. In Python\slash{}CircuitPython the main loop is typically just a \emph{while True:} loop and in Arduino it is the \emph{loop()} function. Once the BDD is completed, the blocks in the BDD can be used to create the IBD which is shown below.%
\begin{figureptx}{Figure}{IBD of Software Components}{ibd-software}{}%
\begin{image}{0.05}{0.9}{0.05}{}%
\includegraphics[width=\linewidth]{external/IBD_Software.jpg}
\end{image}%
\tcblower
\end{figureptx}%
Again, The flow of data is similar to the electrical IBD but now we can see what information is being passed from one function to the next which again will be helpful when we get to the preliminary design phase.%
\end{subsectionptx}
%
%
\typeout{************************************************}
\typeout{Subsection 3.4.4 Mechanical Subsystem}
\typeout{************************************************}
%
\begin{subsectionptx}{Subsection}{Mechanical Subsystem}{}{Mechanical Subsystem}{}{}{subsec-Mechanical-Subsystem}
The mechanical subsystem is a bit more difficult to create BDDs and IBDs for mostly because of how trivial and simple the systems are. Typically these diagrams are so simple that students often get confused and over complicate these diagrams. However, if you look at the diagrams below it's easy to see that mechanical system is important but much smaller and less complex than the electrical and software subsystems. The BDD for the mechanical subsystem is shown below. \begin{figureptx}{Figure}{BDD of Mechanical Components}{bdd-mechanical}{}%
\begin{image}{0.05}{0.9}{0.05}{}%
\includegraphics[width=\linewidth]{external/BDD_Mechanical.jpg}
\end{image}%
\tcblower
\end{figureptx}%
 In this case my plan was to attach the windspeed sensor to my handelbars somehow and then put all the other electrical components into a main houseing. Then all I would have to do is figure out how to run the wires from the main housing to the windspeed sensor. So the BDD only has 3 components. With the BDD completed, I moved on to the mechanical IBD which is shown below. \begin{figureptx}{Figure}{IBD of Mechanical Components}{ibd-mechanical}{}%
\begin{image}{0.05}{0.9}{0.05}{}%
\includegraphics[width=\linewidth]{external/IBD_Mechanical.jpg}
\end{image}%
\tcblower
\end{figureptx}%
 In this case I elected to add the bicycle to the IBD because even though the bicycle is an interface I still wanted to show how the system would be attached to the bicycle and how the different components would be connected to it. The IBD shows that the main housing is attached to the bicycle via velcro strap and the windspeed sensor is attached to the bicycle via zip ties.%
\end{subsectionptx}
%
%
\typeout{************************************************}
\typeout{Subsection 3.4.5 Preliminary Diagrams}
\typeout{************************************************}
%
\begin{subsectionptx}{Subsection}{Preliminary Diagrams}{}{Preliminary Diagrams}{}{}{detailed-diagrams}
Following the conceptual design phase where missions goals, requirements, hand sketches, BDDs and IBDs are created, it is time to work out component selection. For example, instead of calling the module \textquotedblleft{}Windspeed Sensor\textquotedblright{} it is time to select the component and call it a pitot probe. Similarly, buttons, neopixels and bluetooth must be selected. In the preliminary design phase the mechanical lead must create a simple layout of the system. This requires a simple top down layout of the system with dimensions to give the designer an idea of how that's going to go together. For this, a simple sketchup tool can be used like \href{inkscape.org}{\linkhilite{Inkscape}}\hyperlink{inkscape}{\linkhilite{[{\xreffont 10}]}}. Because my project was so simple I did not create a simple layout for this system. For the electrical lead, a connections diagram is necessary to show how the different components are connected to each other. This would be similar to the electrical IBD except this would be actual components with references to actual parts. For this project the connections diagram is super simple because the CPU, clock, bluetooth module, buttons, and the neopixels are all built into the CPB so the only external component is the pitot probe which is connected to the CPB. So my connections diagram just has 3 components: the CPB, pitot probe and power supply. \begin{figureptx}{Figure}{Connections Diagram for Bike Project}{connections-diagram}{}%
\begin{image}{0.05}{0.9}{0.05}{}%
\includegraphics[width=\linewidth]{external/Connections_Diagram.jpg}
\end{image}%
\tcblower
\end{figureptx}%
 For the software lead it is important to create psuedo code in the form of a software flow chart. These charts are something that is often overlooked but can be very helpful in the design process. The flowchart created for this pitot probe bike project is shown below. \begin{figureptx}{Figure}{Psuedo Code Flow Chart for Bike Project}{pitot-psuedo-code}{}%
\begin{image}{0.05}{0.9}{0.05}{}%
\includegraphics[width=\linewidth]{external/Psuedo_Code_FlowChart.jpg}
\end{image}%
\tcblower
\end{figureptx}%
 In this flow chart the code starts by initializing all the different modules, then turns on the neopixels. After that it checks if the bluetooth is connected. If it is not connected, then the neopixels turn blue and the system waits until a connection is made. Once a connection is made, the system checks if the button is pressed. If the button is pressed, then the system goes into calibration mode where it takes 100 samples of the voltage and averages them to get a bias voltage which is used to calibrate the system. During calibration, the neopixels turn green to indicate to the user that they should not touch the system. If the button is not pressed, then the system takes 10 samples of the voltage and averages them to get an average voltage which is then converted to windspeed using a conversion function. Finally, the data is sent to the user's phone via bluetooth and also printed to serial for debugging purposes.%
\end{subsectionptx}
%
%
\typeout{************************************************}
\typeout{Subsection 3.4.6 Detailed Design}
\typeout{************************************************}
%
\begin{subsectionptx}{Subsection}{Detailed Design}{}{Detailed Design}{}{}{detailed-design}
Once the conceptual and preliminary design phase is over it is time to create the detailed design. The detailed design needs to be so detailed that anyone following the design on paper to turn the system into a reality just by reading all of the schematics. Thus, for the mechanical lead, some sort of CAD (Computer Aided Design) tool is necessary for dimensioned drawings and detailed schematics. In this case I recommend using your favourite CAD program such as \href{https://www.blender.org/}{\linkhilite{Blender}}, \href{https://www.solidworks.com/}{\linkhilite{SolidWorks}}, \href{https://www.autodesk.com/products/fusion-360}{\linkhilite{Fusion 360}} or anything you're comfortable with \hyperlink{blender}{\linkhilite{[{\xreffont 5}]}}\hyperlink{solidworks}{\linkhilite{[{\xreffont 7}]}}\hyperlink{autodesk}{\linkhilite{[{\xreffont 8}]}}. For the electrical lead on the team, wiring diagrams are necessary and I recommend using either \href{https://app.cirkitdesigner.com/}{\linkhilite{CirKit Designer}} or \href{https://fritzing.org/}{\linkhilite{Fritzing}}\hyperlink{cirkit}{\linkhilite{[{\xreffont 9}]}}\hyperlink{fritzing}{\linkhilite{[{\xreffont 6}]}}. CAD and wiring diagram examples will not be shown in this chapter simply because numerous examples of wiring diagrams are shown in virtually every weekly project and 3D CAD is the subject of an entire course in and of itself.%
\par
For the software lead it is time to figure out what software is required to run each component. For example, the pitot probe as explained in \hyperref[pitot]{\linkhilite{Chapter~{\xreffont\ref{pitot}}}} is an analog sensor which means you need the \emph{analogio} module. The neopixels need the \emph{neopixel}. This process must be repeated for every component to be sure that you have all the modules necessary to run the program. Although this is a lot of work it helps tremendously later on in the subsystem testing phase. In this case the list below contains all the modules I need to run this system\hyperlink{CircuitPython}{\linkhilite{[{\xreffont 30}]}}.%
\begin{enumerate}
\item{}\href{https://docs.circuitpython.org/en/latest/shared-bindings/time/}{\linkhilite{time}} - for the clock%
\item{}\href{https://docs.circuitpython.org/en/latest/shared-bindings/analogio/}{\linkhilite{analogio}} - for the pitot probe%
\item{}\href{https://docs.circuitpython.org/projects/ble/en/latest/}{\linkhilite{adafruit\textunderscore{}ble}} - for the bluetooth connection%
\item{}\href{https://docs.circuitpython.org/en/latest/shared-bindings/digitalio/}{\linkhilite{digitalio}} - for the button%
\item{}\href{https://docs.circuitpython.org/projects/neopixel/en/latest/}{\linkhilite{neopixel}} - for the neopixels%
\item{}\href{https://circuitpython.org/}{\linkhilite{CircuitPython}} - for the main programming language%
\end{enumerate}
Notice that I didn't need anything for the battery since that's just hardware. An interesting corollary is that it's actually possible to create a software flow chart that is so good that at the time of this writing if you upload the screenshot of the psuedo code flowchart to Gemini it will give you the actual code for it in any language you like\hyperlink{gemini}{\linkhilite{[{\xreffont 2}]}}. In the youtube tutorial discussed in \hyperref[sys-eng-extra-help]{\linkhilite{Section~{\xreffont\ref{sys-eng-extra-help}}}}, I show how to do this but for the sake of this book I will show the code that was generated from the flow chart.\begin{preformatted}
import math
import time
import board
import analogio
import digitalio
import adafruit_ble
from adafruit_ble.advertising.standard import SolitaryAdvertisement
from adafruit_led_animation.animation.solid import Solid
import neopixel

# --- Configuration & Initialization ---
rho = 1.225
v_bias = 0.0  # Initial calibration value

# Setup Analog Input (Change 'A1' to your sensor pin)
airspeed_sensor = analogio.AnalogIn(board.A1)

# Setup Button A
button = digitalio.DigitalInOut(board.BUTTON_A)
button.direction = digitalio.Direction.INPUT
button.pull = digitalio.Pull.DOWN # Adjust to Pull.UP if your button logic is inverted

# Setup NeoPixels
pixels = neopixel.NeoPixel(board.NEOPIXEL, 10, brightness=0.2)
RED = (255, 0, 0)
BLUE = (0, 0, 255)
GREEN = (0, 255, 0)
OFF = (0, 0, 0)

# Setup Bluetooth
ble = adafruit_ble.BLERadio()
advertisement = SolitaryAdvertisement()

def get_voltage():
    # Helper to convert 16-bit ADC value to Voltage (0-3.3V)
    return (airspeed_sensor.value * 3.3) / 65535

def get_wind_data(samples):
    """Measures analog value, converts to voltage, and averages."""
    total_voltage = 0
    for _ in range(samples):
        total_voltage += get_voltage()
    return total_voltage / samples

def convert_2_windspeed(avg_voltage):
    """Converts voltage to windspeed using dP = 1000 * dV."""
    dv = avg_voltage - v_bias
    dp = 1000 * dv
    # Ensure dp is not negative before square root
    if dp < 0: dp = 0 
    windspeed = math.sqrt((2 * dp) / rho)
    return windspeed

def neo_pixel_run(state, color=OFF):
    """Sets the NeoPixel state and color."""
    if state:
        pixels.fill(color)
    else:
        pixels.fill(OFF)
    pixels.show()

# --- Main Loop ---
print("System Initialized...")

while True:
    # Bluetooth Connection Check
    if not ble.connected:
        neo_pixel_run(True, BLUE)
        if not ble.advertising:
            ble.start_advertising(advertisement)
        # Flowchart shows loop back to start/check if disconnected
        continue 
    
    # Check Button Routine
    if button.value:
        # Calibration Branch
        neo_pixel_run(True, GREEN)
        avg_v = get_wind_data(100)
        v_bias = avg_v  # Calibrate() routine
        print(f"Calibrated! New V_Bias: {v_bias}")
    else:
        # Standard Measurement Branch
        neo_pixel_run(False)
        avg_v = get_wind_data(10)
        speed = convert_2_windspeed(avg_v)
        
        current_time = time.monotonic() # GetTime()
        
        # SendData() routine (Printing to Serial/BLE)
        data_string = f"Time: {current_time:.2f}s | Windspeed: {speed:.2f} m/s"
        print(data_string)
        # If you have a specific BLE UART service, you'd send data_string there.
\end{preformatted}
You can see from the code generated that with the help of creating a psuedo code flowchart and the power of AI you can create code from scratch that actually flows the way you want it to. Note that the code above was not tested and may not actually work the way you intended. However, just by glancing at it, it is a great start and shows the power of creating this detailed flow chart. Just to be thorough, \href{https://github.com/cmontalvo251/Microcontrollers/blob/master/Circuit_Playground/CircuitPython/CPB_DataLogger/pitot_datalogger_bluetooth.py}{\linkhilite{the code I ended up using}} is on Github and you can see how similar the code is between the autogenerated code and my code that I built from scratch before AI was widely adopted \hyperlink{Github}{\linkhilite{[{\xreffont 33}]}}.%
\end{subsectionptx}
%
%
\typeout{************************************************}
\typeout{Subsection 3.4.7 Critical Design Review (CDR)}
\typeout{************************************************}
%
\begin{subsectionptx}{Subsection}{Critical Design Review (CDR)}{}{Critical Design Review (CDR)}{}{}{subsec-final}
After the detailed design, the Critical Design Review (CDR) is conducted to evaluate the design against its requirements. This involves a live presentation to stakeholders and peers. In this presentation the students are required to explain their current design and respond to any quesitons from the audience. This is a good way to receive feedback about the project before moving into subsystem testing and final verification and validation. Of course, for my project I did not do a CDR and moved directly into subsystem testing.%
\end{subsectionptx}
%
%
\typeout{************************************************}
\typeout{Subsection 3.4.8 Subsystem Testing}
\typeout{************************************************}
%
\begin{subsectionptx}{Subsection}{Subsystem Testing}{}{Subsystem Testing}{}{}{subsec-subsystem-testing}
For my first test, I elected to deal with the pitot probe and used Method 1 (\hyperref[daq]{\linkhilite{Chapter~{\xreffont\ref{daq}}}}) to record data with my laptop. Then I had my parents drive down the highway while I hung the pitot probe out the window. I was able to get a nice plot of windspeed. One issue I ran into was that when the car came to a stop the speed did not. I realized I needed to go ahead and add the button press so that I could calibrate often. Hence why having the button press was so crucial. Once I was sure the pitot probe worked and my button was able to calibrate I added the bluetooth module and repeated the same experiment only this time I plotted the data on my phone as shown below. \begin{figureptx}{Figure}{Plot of Windspeed Data During Calibration Test}{calibration-plot}{}%
\begin{image}{0.25}{0.5}{0.25}{}%
\includegraphics[width=\linewidth]{external/Calibration_Plot.png}
\end{image}%
\tcblower
\end{figureptx}%
 Once I knew the button, pitot probe and bluetooth were working I added the neopixels and tested that as well although I do not have a photo of that particular test.%
\end{subsectionptx}
%
%
\typeout{************************************************}
\typeout{Subsection 3.4.9 Test Readiness Review (TRR) and Systems Verification Review (SVR)}
\typeout{************************************************}
%
\begin{subsectionptx}{Subsection}{Test Readiness Review (TRR) and Systems Verification Review (SVR)}{}{Test Readiness Review (TRR) and Systems Verification Review (SVR)}{}{}{subsec-final-verification}
Once I tested all the components it was time to integrate everything together by placing everything into a small case on my bike. I then took the system on a bike ride around town and was able to get some good data. The final product is shown below. \begin{figureptx}{Figure}{Main Housing of Pitot Probe Bike Project}{final-product}{}%
\begin{image}{0.15}{0.7}{0.15}{}%
\includegraphics[width=\linewidth]{external/pitot_box.jpg}
\end{image}%
\tcblower
\end{figureptx}%
 Overall the prototype worked as intended. It recorded wind speed and had a button to re-calibrate. While the system is recalibrating the neopixels light up to notify the user not to start biking yet. The housing mounted to the top tube of my bike as shown in the figure above and I ended up using zip ties to secure the pitot tube to the front of my bike like in the photo below. \begin{figureptx}{Figure}{Final Product with Pitot Probe Mounted to Front of Bike}{final-product-2}{}%
\begin{image}{0.15}{0.7}{0.15}{}%
\includegraphics[width=\linewidth]{external/pitot_zip_tie.jpg}
\end{image}%
\tcblower
\end{figureptx}%
 I have future ideas where I could upgrade the power supply and even install a waterproof toggle switch so that I can turn the device on and off without having to remove the entire housing. It would also be nice to have a button on the outside to recalibrate without having to remove the entire housing. Still, all requirements were satisfied and the project was a success as shown in the table below. Note that satisfied requirements are in plain text, partially satisfied comments are in italics and bold text are requirements that were not satisfied. In this case, all requirements were satisfied except for 1 which was partially satisfied.%
\begin{tableptx}{Table}{\textbf{System Requirements Verification Matrix}}{table-system-verification}{}%
\begin{tabularbox}{0}{1}{0}%
\resizebox{\ifdim\width > \linewidth\linewidth\else\width\fi}{!}{%
{\tabularfont%
\begin{tabular}{AlAlA}\hrulethin
{\bfseries{}Requirement}&{\bfseries{}Verification \slash{} Status}\tabularnewline\hrulethin
\multicolumn{1}{Am{0.5\linewidth}A}{\raggedright%
The system shall measure the windspeed of the vehicle%
}&\multicolumn{1}{m{0.5\linewidth}A}{\raggedright%
Yes, with a pitot probe%
}\tabularnewline\hrulethin
\multicolumn{1}{Am{0.5\linewidth}A}{\raggedright%
The system shall be calibrated on startup or during operation by pressing a button%
}&\multicolumn{1}{m{0.5\linewidth}A}{\raggedright%
Yes, on start up and by pressing button A%
}\tabularnewline\hrulethin
\multicolumn{1}{Am{0.5\linewidth}A}{\raggedright%
The system will be powered during operation of the vehicle%
}&\multicolumn{1}{m{0.5\linewidth}A}{\raggedright%
Yes, with a portable battery pack%
}\tabularnewline\hrulethin
\multicolumn{1}{Am{0.5\linewidth}A}{\raggedright%
The system shall not adversely affect the vehicle performance.%
}&\multicolumn{1}{m{0.5\linewidth}A}{\raggedright%
Yes, the user did not report any issues while biking.%
}\tabularnewline\hrulethin
\multicolumn{1}{Am{0.5\linewidth}A}{\raggedright%
The system shall send windspeed data to the operators smart phone.%
}&\multicolumn{1}{m{0.5\linewidth}A}{\raggedright%
Yes, with bluetooth%
}\tabularnewline\hrulethin
\multicolumn{1}{Am{0.5\linewidth}A}{\raggedright%
The system shall light up during calibration to indicate to the user that the system shall not be touched.%
}&\multicolumn{1}{m{0.5\linewidth}A}{\raggedright%
Yes, it lights up green%
}\tabularnewline\hrulethin
\multicolumn{1}{Am{0.5\linewidth}A}{\raggedright%
System must incorporate 3 electrical components%
}&\multicolumn{1}{m{0.5\linewidth}A}{\raggedright%
System had four: pitot probe, neopixels, button, and bluetooth.%
}\tabularnewline\hrulethin
\multicolumn{1}{Am{0.5\linewidth}A}{\raggedright%
System must be unique as well as built and designed by your team%
}&\multicolumn{1}{m{0.5\linewidth}A}{\raggedright%
Yes, I created the entire project from scratch and only grabbed inspiration from other projects.%
}\tabularnewline\hrulethin
\multicolumn{1}{Am{0.5\linewidth}A}{\raggedright%
\emph{System must have one of the three components as an external part to the main microcontroller and not be something already taught in a weekly project.}%
}&\multicolumn{1}{m{0.5\linewidth}A}{\raggedright%
\emph{Partially, the pitot probe was external but already a part of the weekly projects.}%
}\tabularnewline\hrulethin
\multicolumn{1}{Am{0.5\linewidth}A}{\raggedright%
System must take some sort of measurement to be used in the code%
}&\multicolumn{1}{m{0.5\linewidth}A}{\raggedright%
Yes, it took time and windspeed data.%
}\tabularnewline\hrulethin
\multicolumn{1}{Am{0.5\linewidth}A}{\raggedright%
System must utilize a data acquisition method so that data can be plotted on your computer as a function of time%
}&\multicolumn{1}{m{0.5\linewidth}A}{\raggedright%
Yes, Method 5 was utilized%
}\tabularnewline\hrulethin
\multicolumn{1}{Am{0.5\linewidth}A}{\raggedright%
System must help someone in some way%
}&\multicolumn{1}{m{0.5\linewidth}A}{\raggedright%
Yes, it allowed the user to ascertain its wind windspeed.%
}\tabularnewline\hrulethin
\multicolumn{1}{Am{0.5\linewidth}A}{\raggedright%
System must have some sort of mechanical component to it and cannot be purely software and electrical components.%
}&\multicolumn{1}{m{0.5\linewidth}A}{\raggedright%
Yes, it had a main housing and pitot probe zip tied to the handlebars.%
}\tabularnewline\hrulethin
\end{tabular}
}%
}%
\end{tabularbox}%
\end{tableptx}%
This is just one example of how to use the systems engineering life cycle to design and build a project. The same process can be applied to any project and can be used to help ensure that the project is successful. The key is to follow the process and not skip any steps. By following the process, you can ensure that your project is well designed, well built and meets all of its requirements. Hopefully this example gets you excited about the CPX\slash{}CPB and the entire Adafruit ecosystem as well as a quick preview of what they can do. Hopefully you enjoy the following chapters which details individual weekly projects geared to help you learn simple sensors and electrical components that you can piece together and build your own DIY electro-mechanical system.%
\end{subsectionptx}
\end{sectionptx}
%
%
\typeout{************************************************}
\typeout{Section 3.5 Extra Help}
\typeout{************************************************}
%
\begin{sectionptx}{Section}{Extra Help}{}{Extra Help}{}{}{sys-eng-extra-help}
I have a youtube tutorial on \href{https://www.youtube.com/watch?v=l4JhueUcw4U}{\linkhilite{how to create an airplane in Blender}} which you could easily adapt to anything else if you're having trouble making 3D objects. I similarly, have a tutorial on \href{https://youtu.be/G-vQ3v1j7ZU}{\linkhilite{creating BDDs and IBDs in DrawIO}} I also have a youtube tutorial where I \href{https://www.youtube.com/watch?v=FSF-TeceGQ8}{\linkhilite{create a psuedo code flow chart}} just like the chart in this chapter. Finally, I have another youtube tutorial on \href{https://www.youtube.com/watch?v=xZ-j7GbpLV0}{\linkhilite{how to make a simple circuit in Cirkit Designer}}\hyperlink{youtube}{\linkhilite{[{\xreffont 50}]}}.%
\end{sectionptx}
\end{chapterptx}
%
%
\typeout{************************************************}
\typeout{Chapter 4 Purchase Equipment}
\typeout{************************************************}
%
\begin{chapterptx}{Chapter}{Purchase Equipment}{}{Purchase Equipment}{}{}{purchase-equipment}
\renewcommand*{\chaptername}{Chapter}
\begin{introduction}{}%
\begin{listptx}{List}{\textbf{Parts List}}{purchase-equipment-2-1}{}%
%
\begin{enumerate}
\item{}Laptop with internet connectivity%
\item{}Credit\slash{}Debit Card to Purchase items%
\end{enumerate}
\end{listptx}%
\begin{listptx}{List}{\textbf{Learning Objectives}}{purchase-equipment-2-2}{}%
%
\begin{enumerate}
\item{}Understand the world of microcontrollers%
\item{}Learn about the Adafruit eco-system%
\item{}Learn about the Circuit Playground Express\slash{}Bluefruit%
\item{}Learn about the main alternative: the Arduino%
\end{enumerate}
\end{listptx}%
In this class you're going to build some circuits that will enhance your learning experience. Rather than just solving problems by hand you're going to take and analyze data. Over the summer of 2020, I began to work with Tangibles that Teach and at the time they bundled all components together. Unfortunately, that company has gone out of business and as such for the time being you must build the kit yourself. This project is all about what parts you need to complete the subsequent chapters and what to purchase to get started.%
\end{introduction}%
%
%
\typeout{************************************************}
\typeout{Section 4.1 CircuitPython Kit}
\typeout{************************************************}
%
\begin{sectionptx}{Section}{CircuitPython Kit}{}{CircuitPython Kit}{}{}{purchase-equipment-3}
Here is a list of all the components in the kit \alert{Note that every semester is different and every instructor is different. Some instructors may choose to add or remove components from the kit. The list below is just a recommendation and is based on the components used in the labs in this book. Given a standard 16 week semester, I have personally found it impossible to do more than 12 or 13 labs which means we don't use everything in this list. Be sure to check with your instructor before buying everything on this list.}%
\begin{enumerate}
\item{}Circuit Playground Bluefruit or Express + Included USB Cable%
\item{}Micro Servo%
\item{}Potentiometer%
\item{}Photocell%
\item{}Two Resistors (330 Ohm and 1K Ohm)%
\item{}Alligator Clips x3%
\item{}External Battery Pack (optional)%
\item{}AAA Batteries x3 (optional)%
\item{}Breadboard%
\item{}Push Button%
\item{}LEDs x2%
\item{}LSM6DS33 + LIS3MDL - 9 DoF IMU (\hyperref[imu]{\linkhilite{Chapter~{\xreffont\ref{imu}}}})%
\item{}Pitot Probe (not pictured below)(\hyperref[pitot]{\linkhilite{Chapter~{\xreffont\ref{pitot}}}})%
\item{}LPS22 Pressure Sensor (not pictured below)(\hyperref[pressure]{\linkhilite{Chapter~{\xreffont\ref{pressure}}}})%
\item{}GPS Sensor (not pictured below)(\hyperref[gps]{\linkhilite{Chapter~{\xreffont\ref{gps}}}})%
\end{enumerate}
%
\par
Below is also a photo of all the components listed above. When you get your kit familiarize yourself with all of the components. The hardware kit is a collection of multiple items. At the time of this writing the entire kit cannot be purchased however the CircuitPlayground Express (CPX) can be purchased on Adafruit \hyperlink{AdafruitMAIN}{\linkhilite{[{\xreffont 34}]}}. A photo of the kit is shown in the figure below. The photo shows the CPX on the left inside it's protective zip lock bag.%
\begin{figureptx}{Figure}{Photo of CircuitPython Kit Components%
}{purchase-equipment-3-4}{}%
\begin{image}{0.1}{0.8}{0.1}{}%
\includegraphics[width=\linewidth]{external/components_w_annotations.png}
\end{image}%
\tcblower
\end{figureptx}%
The CPX is a relatively new development board that operates in the same style as Arduino\hyperlink{Arduino}{\linkhilite{[{\xreffont 35}]}}. That is, the student can create a relatively simple script and then flash the software to the CPX via USB. The Arduino operates in the same way. The benefit of the CPX vs the Arduino is that the software is CircuitPython which is indistinguishable from Python. As mentioned earlier, Python is in the top 3 programming languages as stated by the Tiobe Index of Programming\hyperlink{Tiobe}{\linkhilite{[{\xreffont 28}]}}. Other benefits include the components built in to the CPX. The CPX itself contains a built in microphone, accelerometer, light sensor, speaker, three push buttons, slide switch and infrared sensor. The CircuitPlayground Express (CPX) also contains 10 smart NeoPixel LEDs and a ring of copper plated pads that can be attached to the included alligator clips and breadboard. Students can also elect to purchase the CircuitPlayground Bluefruit (CPB) which does not have an infrared sensor but does contain bluetooth functionality to send data wirelessly to the Adafruit Connect App which runs on Apple\slash{}Android \hyperlink{AdafruitBLE}{\linkhilite{[{\xreffont 36}]}}.%
\end{sectionptx}
%
%
\typeout{************************************************}
\typeout{Section 4.2 Arduino Kit}
\typeout{************************************************}
%
\begin{sectionptx}{Section}{Arduino Kit}{}{Arduino Kit}{}{}{purchase-equipment-4}
Note, that the instructor and student has the choice to use any type of microcontroller (MCU). In reality, every project detailed in this book can be done with an Arduino with the notable exception being Method 3 Data Acquisition lab. In that lab, data is stored on the MCU itself since the CPX has built-in storage. However, Arduino does not have built-in storage which is one of the limitations with Arduino. However, it the experience of the primary author teaching this course that many groups choose to use an Arduino UNO\slash{}Mega\slash{}Nano or ESP32 for their final project. This is because although the CPX\slash{}CPB is a cheap MCU that has many built-in sensors, it is severely lacking in stability, memory and power capabilities. Using an Arduino will allow a larger world of possibilities and thankfully CPX\slash{}CPB programming\slash{}hardware skills will transfer very easily to Arduino. As such, many labs will have an Arduino alternative in the event students have found the Arduino to be more useful for that specific lab. So the question of course is why isn’t the entire class just built around Arduino? The answer is because for \textdollar{}20 the CPX\slash{}CPB comes with two push buttons, a slide switch, an accelerometer, a temperature sensor, 8 neopixels, a light sensor, a microphone, a speaker and a much simpler programming language. The CPX also has an IR sensor and the CPB supports bluetooth communication. The Arduino does not have any of those extra capabilities. This means in order to perform many of the experiments in this book, the students will need to buy extra supplies if you use an Arduino. For example, in order to store data on a standalone project using an Arduino requires the use of an extra micro SD card breakout board and a \href{https://www.adafruit.com/product/4682?gad_source=1\&gad_campaignid=23438252138\&gbraid=0AAAAADx9JvSwU8AgEo6zbpyqSObyRYEag\&gclid=CjwKCAjwwpDQBhAuEiwAa-4Wo0MRxAOx4gdkI3ZPRh7lYKrN5u-dFIkiIfKsLBkJC6I18aCrNlCvEBoC0rgQAvD_BwE}{\linkhilite{micro SD card}}. Furthermore, for the accelerometer lab, the students would need to purchase an extra \href{https://www.adafruit.com/product/2019}{\linkhilite{external accelerometer}}. In this case, at the time of this writing, that represents an extra \textdollar{}10 on top of the microcontroller itself which is \textdollar{}27.60 on the \href{https://store-usa.arduino.cc/products/arduino-uno-rev3?utm_source=google\&utm_medium=cpc\&utm_campaign=US-Pmax\&gad_source=1\&gad_campaignid=21317508903\&gbraid=0AAAAACbEa86LJ3YDsL-3aDjMOhB_LzEaL\&gclid=CjwKCAjwwpDQBhAuEiwAa-4Wo3mpltsHJoDltwWD1lLRXMur3mj01diTFH-Fb50kmQlrHxQXUywfVRoCLKsQAvD_BwE}{\linkhilite{Official Arduino Store}}. The bottom line is that the authors feel that students get more out of learning on a CPX\slash{}CPB given the seamless integration of sensors and the much easier programming language. Once they gain mastery over the CPX\slash{}CPB they can then take those skills to something more reliable like the Arduino MCU or even something more complex like a Raspberry Pi. Still, in the event the student would like to purchase an Arduino instead of the CPX\slash{}CPB, the recommended Arduino is the Arduino UNO R3 which can be purchased on the official Arduino store \hyperlink{Arduino}{\linkhilite{[{\xreffont 35}]}}. The UNO is a great starting point for learning about microcontrollers and has a large online community for support. However, if you want to do more complex projects in the future, you may want to consider purchasing an Arduino Mega or an ESP32 which have more memory and capabilities than the UNO. However, instead of purchasing just the Arduino it's recommended to buy a kit that comes with a lot of extra components. For example, the \href{https://www.amazon.com/-/es/inicial-completo-Arduino-tutorial-art\%C3\%ADculos/dp/B01CZTLHGE/ref=pd_sbs_d_sccl_1_1/147-7312721-2903556?pd_rd_w=BdloZ\&content-id=amzn1.sym.aa738fbd-ad05-4d11-aae2-04b598db6305\&pf_rd_p=aa738fbd-ad05-4d11-aae2-04b598db6305\&pf_rd_r=XKB0NC1YZ871551J9WT5\&pd_rd_wg=SaZrU\&pd_rd_r=2b9dc346-4385-4898-9644-aacc74dfa69c\&pd_rd_i=B01CZTLHGE\&th=1}{\linkhilite{Elegoo UNO R3 Kit}} \hyperlink{Elegoo}{\linkhilite{[{\xreffont 67}]}} comes with a lot of extra components that will allow you to do many of the projects in this book without having to purchase extra components. At the time of this writing the kit is \textdollar{}47.99 and all the components is shown below. \begin{figureptx}{Figure}{Photo of Elegoo UNO R3 Kit Components%
}{purchase-equipment-4-2-7}{}%
\begin{image}{0.1}{0.8}{0.1}{}%
\includegraphics[width=\linewidth]{external/ArduinoUNO_Kit.jpg}
\end{image}%
\tcblower
\end{figureptx}%
%
\end{sectionptx}
%
%
\typeout{************************************************}
\typeout{Section 4.3 Extra Help}
\typeout{************************************************}
%
\begin{sectionptx}{Section}{Extra Help}{}{Extra Help}{}{}{purchase-equipment-5}
I created an \href{https://youtu.be/6sNNQrhnzLE}{\linkhilite{unboxing video on Youtube}} for you to take a look and see the actual hardware as I talk about each piece of equipment\hyperlink{youtube}{\linkhilite{[{\xreffont 50}]}}.%
\end{sectionptx}
%
%
\typeout{************************************************}
\typeout{Section 4.4 Assignment}
\typeout{************************************************}
%
\begin{sectionptx}{Section}{Assignment}{}{Assignment}{}{}{purchase-equipment-6}
\begin{introduction}{}%
Your assignment for this module is to purchase the required equipment for your class. \alert{Note that the specific equipment will be instructor dependent so be sure you discuss with the instructor of your course before buying everything.} At a minimum you must purchase a CPX or CPB but the rest of the items in your kit will depend on which modules your instructor wants you to do during the semester. \alert{Also, the CPX\slash{}CPB requires a microUSB cable and that cable must have a data line. Many USB cables are just power and ground.}%
%
\begin{enumerate}
\item{}Title page with name,title and date - 5\%%
\item{}Introduction - 5\%%
\item{}All figures must have appropriate annotations, figure captions, figure labels and a paragraph of supporting text - 10\%%
\item{}A figure of your receipt of your purchases this semester (If you already own the MCU you may just take a photo of you holding the device). - 60\%%
\item{}Using the number of students in your class, compute the probability that at least one MCU will be shipped DOA (dead on arrival) assuming a low (p=1\%) failure rate. - 20\%%
\item{}Appendix A - Video Details - Use the webcam on your computer to film yourself introducing yourself to me and either show me your purchased MCU or film your screen and show the receipt that shows you've purchased the MCU. - Pass\slash{}Fail%
\item{}Appendix B - Code used to take data and code used to plot in plain text (no screenshots) - Pass\slash{}Fail%
\item{}Appendix C - Screenshot of Excel or Sheets if applicable - Pass\slash{}Fail%
\item{}\alert{If you are asked to fix something in the Appendices such as missing content in your video or screenshots of your code, you will lose 10\% of your grade everytime you submit an update after the deadline.}%
\end{enumerate}
\end{introduction}%
\end{sectionptx}
\end{chapterptx}
%
%
\typeout{************************************************}
\typeout{Chapter 5 Plotting}
\typeout{************************************************}
%
\begin{chapterptx}{Chapter}{Plotting}{}{Plotting}{}{}{plotting}
\renewcommand*{\chaptername}{Chapter}
\begin{introduction}{}%
\begin{listptx}{List}{\textbf{Parts List}}{plotting-2-1}{}%
%
\begin{enumerate}
\item{}Laptop%
\item{}\href{https://www.python.org/}{\linkhilite{Python}} \hyperlink{python}{\linkhilite{[{\xreffont 51}]}} Integrated Development Environment (IDE) (\href{https://thonny.org/}{\linkhilite{Thonny}} \hyperlink{thonny}{\linkhilite{[{\xreffont 52}]}} or \href{https://www.spyder-ide.org/}{\linkhilite{Spyder}} \hyperlink{spyder}{\linkhilite{[{\xreffont 53}]}} - Do not use Python IDLE)%
\item{}Optional: \href{https://colab.research.google.com/}{\linkhilite{Google Colab}} \hyperlink{colab}{\linkhilite{[{\xreffont 54}]}}%
\item{}Optional: \href{https://sheets.google.com/}{\linkhilite{Google Sheets}} \hyperlink{sheets}{\linkhilite{[{\xreffont 69}]}}%
\item{}Optional: \href{https://www.gnu.org/software/octave/}{\linkhilite{Octave}} \hyperlink{octave}{\linkhilite{[{\xreffont 68}]}}%
\end{enumerate}
\end{listptx}%
\begin{listptx}{List}{\textbf{Learning Objectives}}{plotting-2-2}{}%
%
\begin{enumerate}
\item{}How to download and use a Python IDE (Thonny or Spyder is suggested. Do not use Python IDLE)%
\item{}How to install Python modules such as numpy and matplotlib%
\item{}Learn how to use Google Colab, Google Sheets and Octave if you choose to use those platforms%
\item{}Plot a simple function in Python, Octave or Sheets%
\end{enumerate}
\end{listptx}%
As you learn Instrumentation throughout the semester, you will be tasked with creating computer programs on the Circuit Playground Express (CPX) or Arduino. The CPX itself has it's own RAM, CPU, HDD and many sensors while the Arduino is a simple microcontroller with no sensors, or harddrive. The CPX though is kind of like a mini computer! You can plug the CPX into your computer via USB and access the hard drive (HDD) from your own computer. When you program on the CPX you need to write programs on the CPX itself so that the mini computer can run the program you wrote. The CPX knows how to read multiple different languages but in this class the main languate will be the \href{https://www.python.org/}{\linkhilite{Python}} language which has been ported to the CPX and called \href{https://circuitpython.org/}{\linkhilite{CircuitPython}} \hyperlink{python}{\linkhilite{[{\xreffont 51}]}}\hyperlink{CircuitPython}{\linkhilite{[{\xreffont 30}]}}. Since many programs will be written in CircuitPython it is import to first learn how to program some things in Python. Python can be easily downloaded by itself but it's nice to get what's called an Integrated Development Environment (IDE). This way programs can be written on the computer for practice and better debugging. This chapter is all about getting Python and an IDE downloaded on the computer and learning how to use it. Extra sections have been written to briefly discuss Octave, Google Sheets and Google Colab as well. You can use any of those platforms to complete the assignment.%
\end{introduction}%
%
%
\typeout{************************************************}
\typeout{Section 5.1 Python IDEs}
\typeout{************************************************}
%
\begin{sectionptx}{Section}{Python IDEs}{}{Python IDEs}{}{}{plotting-3}
\begin{introduction}{}%
So which IDE is the best and which is recommended? Two IDEs are recommended for Python. They are explained in the next two sections. Either one is totally fine. If you just Google \textquotedblleft{}Python download\textquotedblright{} you will find a humongous list of editors (Scratch, Anaconda, Canopy, Eclipse, PyDev, etc). It's easy to get lost when searching for something so broad. Again, just using Spyder or Thonny.%
\end{introduction}%
%
%
\typeout{************************************************}
\typeout{Subsection 5.1.1 Thonny}
\typeout{************************************************}
%
\begin{subsectionptx}{Subsection}{Thonny}{}{Thonny}{}{}{plotting-3-3}
Thonny - \href{https://thonny.org/}{\linkhilite{https:\slash{}\slash{}thonny.org}}\hyperlink{thonny}{\linkhilite{[{\xreffont 52}]}} is a great, easy, simple and lightweight IDE that has been around for years. Note that there is a 32 bit and 64 bit version. Which one is necessary? Well that depends on how many bits the computer has. This is a great thing to Google. Type the following: "do I have a 32 bit or 64 bit computer" into Google. It's more than likely that everyone reading this has a 64 bit computer but it's always a good idea to check. Learning the difference between 32 and 64 bit computers is the topic of another chapter (\hyperref[voltage]{\linkhilite{Chapter~{\xreffont\ref{voltage}}}}). The figure below is a screenshot of what Thonny looks like with a sine graph plotted.%
\begin{figureptx}{Figure}{Example Sine Graph in Thonny}{plotting-3-3-3}{}%
\begin{image}{0.135}{0.73}{0.135}{}%
\includegraphics[width=\linewidth]{external/Thonny.png}
\end{image}%
\tcblower
\end{figureptx}%
\end{subsectionptx}
%
%
\typeout{************************************************}
\typeout{Subsection 5.1.2 Spyder}
\typeout{************************************************}
%
\begin{subsectionptx}{Subsection}{Spyder}{}{Spyder}{}{}{plotting-3-4}
Spyder - \href{https://www.spyder-ide.org/}{\linkhilite{https:\slash{}\slash{}www.spyder-ide.org\slash{}}}\hyperlink{spyder}{\linkhilite{[{\xreffont 53}]}} is a powerful IDE that is widely used in the scientific Python community. The figure below is a screenshot of what Spyder looks like with a sine graph plotted.%
\begin{figureptx}{Figure}{Example Sine Graph in Spyder}{plotting-3-4-3}{}%
\begin{image}{0.135}{0.73}{0.135}{}%
\includegraphics[width=\linewidth]{external/Spyder.png}
\end{image}%
\tcblower
\end{figureptx}%
\end{subsectionptx}
%
%
\typeout{************************************************}
\typeout{Subsection 5.1.3 Setting up the IDE}
\typeout{************************************************}
%
\begin{subsectionptx}{Subsection}{Setting up the IDE}{}{Setting up the IDE}{}{}{plotting-3-5}
Once Thonny or Spyder is installed, numpy and matplotlib are necessary modules within Python that allow Python to do some extra things like numerical computation with Python (numpy) and Matlab style Plotting libraries (matplotlib). How to install modules in Youtube is in the Extra Help section \hyperref[extra-help]{\linkhilite{Section~{\xreffont\ref{extra-help}}}}; however, the simple way is to head over to Tools\textgreater{}Manage Packages in Thonny. The image below shows that version 3.1.2 is already installed but can be upgraded to 3.2.2%
\begin{figureptx}{Figure}{Installing Matplotlib in Thonny}{plotting-3-5-3}{}%
\begin{image}{0.135}{0.73}{0.135}{}%
\includegraphics[width=\linewidth]{external/IDE_upgrades.png}
\end{image}%
\tcblower
\end{figureptx}%
If numpy or matplotlib is not already included in Spyder then type the following into the Python Console in the lower right hand corner of Spyder which is called the IPython console.%
\begin{preformatted}
!pip install matplotlib
\end{preformatted}
If that doesn't work try%
\begin{preformatted}
!pip3 install matplotlib
\end{preformatted}
The output example below shows that matplotlib is already installed installed as it says \textquotedblleft{}requirement already satisfied\textquotedblright{}. Assuming there is a valid internet connection it will install the necessary module.%
\begin{figureptx}{Figure}{Installing Python Modules via Thonny Console}{plotting-3-5-9}{}%
\begin{image}{0.25}{0.5}{0.25}{}%
\includegraphics[width=\linewidth]{external/console_thonny.png}
\end{image}%
\tcblower
\end{figureptx}%
\end{subsectionptx}
%
%
\typeout{************************************************}
\typeout{Subsection 5.1.4 Scripting}
\typeout{************************************************}
%
\begin{subsectionptx}{Subsection}{Scripting}{}{Scripting}{}{}{plotting-3-6}
Once numpy and matplotlib is installed it's time to make a plot. Here is a simple example.%
\begin{figureptx}{Figure}{Close up text of plotting in Python}{plotting-3-6-3}{}%
\begin{image}{0.25}{0.5}{0.25}{}%
\includegraphics[width=\linewidth]{external/matplotlib_example.png}
\end{image}%
\tcblower
\end{figureptx}%
The code above will plot a sine wave from 0 to 2pi. The two lines at the top are importing the numpy and matplotlib modules you installed earlier. When they are imported we give them shorter names so it's easier to reference them so numpy will now be called np and matplotlib.pyplot will be called plt. The next two lines then create a vector \textquotedblleft{}x\textquotedblright{} from 0 to 2pi using 1000 data points. The next line then uses the sine function to create the vector \textquotedblleft{}y\textquotedblright{}. Finally \textquotedblleft{}x\textquotedblright{} and \textquotedblleft{}y\textquotedblright{} are plotted and the figure is instructed to pop up on your screen using the show() function.%
\end{subsectionptx}
\end{sectionptx}
%
%
\typeout{************************************************}
\typeout{Section 5.2 Other Plotting Tools}
\typeout{************************************************}
%
\begin{sectionptx}{Section}{Other Plotting Tools}{}{Other Plotting Tools}{}{}{plotting-4}
%
%
\typeout{************************************************}
\typeout{Subsection 5.2.1 Google Colab}
\typeout{************************************************}
%
\begin{subsectionptx}{Subsection}{Google Colab}{}{Google Colab}{}{}{plotting-4-2}
It is possible to use \href{https://colab.research.google.com}{\linkhilite{Google Colab}}\hyperlink{colab}{\linkhilite{[{\xreffont 54}]}} to collaborate on Python projects, run in the cloud\slash{}browser. In order to get Colab to work a Google account is necessary. Colab is a great tool for collaborating on Python projects and sharing code with others. It also has the advantage of running in the cloud so it can be accessed from any device with an internet connection. Although it's possible to head directly to the Colab website, it is recommended to open Colab through Google Drive. To do that simply open Google Drive\textgreater{}+New\textgreater{}More\textgreater{}Connect more apps and then install Google Colab as shown in the figure below.%
\begin{figureptx}{Figure}{Installing Google Colab through Google Drive}{plotting-4-2-3}{}%
\begin{image}{0.135}{0.73}{0.135}{}%
\includegraphics[width=\linewidth]{external/Google_Colab_Install.png}
\end{image}%
\tcblower
\end{figureptx}%
Once that's done it will be possible to open a new Google Colab notebook. The figure below is a screenshot of what Colab looks like with a sine graph plotted. The other benefit to using Google Colab is that Gemini is imbedded into Google Colab for easy AI integration.%
\begin{figureptx}{Figure}{Example Sine Graph in Google Colab}{plotting-4-2-5}{}%
\begin{image}{0.135}{0.73}{0.135}{}%
\includegraphics[width=\linewidth]{external/Colab.png}
\end{image}%
\tcblower
\end{figureptx}%
\end{subsectionptx}
%
%
\typeout{************************************************}
\typeout{Subsection 5.2.2 Octave}
\typeout{************************************************}
%
\begin{subsectionptx}{Subsection}{Octave}{}{Octave}{}{}{plotting-4-3}
Octave is very similar to MATLAB but free. The only major downside is that Octave is much slower and doesn't have a lot of extra utilities that MATLAB does. However, if it's not possible to get MATLAB, Octave is a great alternative. The figure below is a screenshot of what Octave looks like with a sine graph plotted.%
\begin{figureptx}{Figure}{Example Sine Graph in Octave}{plotting-4-3-3}{}%
\begin{image}{0.135}{0.73}{0.135}{}%
\includegraphics[width=\linewidth]{external/Octave.png}
\end{image}%
\tcblower
\end{figureptx}%
Notice that the syntax is very similar to Python but slightly different. Deep diving into programming is not the focus of this class but anyone wanting to learn more about programming in Python or Octave can find plenty of resources online. The best way to learn is to just start writing code and then Google any errors that come up. The error messages are usually pretty descriptive and will lead to the answer. Furthermore, with the spawn of Generative AI many students often code better than ever before with a few prompts into any AI tool like Gemini, Copilot, Grok or ChatGPT. Also, the extra help section (xref ref="extra-help"\slash{}\textgreater{}) has many youtube videos on how to program if videos are preferred to AI or text.%
\end{subsectionptx}
%
%
\typeout{************************************************}
\typeout{Subsection 5.2.3 Google Sheets}
\typeout{************************************************}
%
\begin{subsectionptx}{Subsection}{Google Sheets}{}{Google Sheets}{}{}{plotting-4-4}
Google Sheets is typically already installed in all browsers again provided a Google account is created first. However, Microsoft Excel can also be used if a Google account is not available. The figure below is a screenshot of what Google Sheets looks like with a sine graph plotted. Plotting in spreadsheet programs can be cumbersome but it is possible to do and is a good option for those who are not interested in learning Python or Octave or programming in general. Futhermore, there is a tongue-in-cheek comment from recently graduated engineers that no one knows how to code and everyone just uses Excel. So if that's the future of engineering, Google Sheets or Excel is the way to go.%
\begin{figureptx}{Figure}{Example Sine Graph in Google Sheets}{plotting-4-4-3}{}%
\begin{image}{0.135}{0.73}{0.135}{}%
\includegraphics[width=\linewidth]{external/Sheets.png}
\end{image}%
\tcblower
\end{figureptx}%
\end{subsectionptx}
\end{sectionptx}
%
%
\typeout{************************************************}
\typeout{Section 5.3 Extra Help}
\typeout{************************************************}
%
\begin{sectionptx}{Section}{Extra Help}{}{Extra Help}{}{}{extra-help}
\begin{introduction}{}%
If extra help is needed with this assignment there is a \href{https://www.youtube.com/watch?v=qaxukpYRqfA\&list=PL_D7_GvGz-v1RsDs_OdNW65qRjEjmpfQx\&index=14}{\linkhilite{youtube video on how to install Thonny}} as well as a \href{https://www.youtube.com/watch?v=OjYwET-6QtE\&list=PL_D7_GvGz-v1RsDs_OdNW65qRjEjmpfQx\&index=35}{\linkhilite{youtube video on how to install Spyder}}. There is also a pretty \href{https://www.youtube.com/watch?v=7nAzHPURYW0\&list=PL_D7_GvGz-v1RsDs_OdNW65qRjEjmpfQx\&index=8}{\linkhilite{comprehensive youtube video on how to plot in matplotlib}} \hyperlink{youtube}{\linkhilite{[{\xreffont 50}]}}.%
\end{introduction}%
%
%
\typeout{************************************************}
\typeout{Subsection 5.3.1 Built-In Help Function and dir()}
\typeout{************************************************}
%
\begin{subsectionptx}{Subsection}{Built-In Help Function and dir()}{}{Built-In Help Function and dir()}{}{}{extra-help-3}
Running code will always create syntax errors. Typing your syntax error into Google will yield so many results you might get lost. Sometimes it helps to know how to learn things just from your computer. For example, type in the commands below in the IPython console or the Shell.%
\par
\emph{import numpy as np}%
\par
\emph{dir(np)}%
\begin{figureptx}{Figure}{Output of dir(numpy)}{extra-help-3-5}{}%
\begin{image}{0.135}{0.73}{0.135}{}%
\includegraphics[width=\linewidth]{external/dir.png}
\end{image}%
\tcblower
\end{figureptx}%
I included a photo of the output from the dir function. You'll notice there are a ton of functions in numpy. Every function in Python has a%
\begin{preformatted}
__doc__
\end{preformatted}
function. That's two underscores followed by \textquotedblleft{}doc\textquotedblright{} and then another two underscores. If you're ever curious about what a particular function does you can just run the command below again in the IPython console or Shell. In this example I'm looking at what \emph{arctan2} does.%
\begin{preformatted}
print(np.arctan2.__doc__)
\end{preformatted}
\begin{figureptx}{Figure}{Output of print(np.arctan2.\textunderscore{}\textunderscore{}doc\textunderscore{}\textunderscore{})}{extra-help-3-10}{}%
\begin{image}{0.135}{0.73}{0.135}{}%
\includegraphics[width=\linewidth]{external/arctan2_doc.png}
\end{image}%
\tcblower
\end{figureptx}%
You'll see that arctan2 takes 2 input arguments \textquotedblleft{}x1\textquotedblright{} and \textquotedblleft{}x2\textquotedblright{}. I didn't include the entire output but if you continue to scroll through the output it will even include examples on how to use the function.%
\par
Another way to learn certain functions is by visiting the appropriate documentation. The \href{https://numpy.org/doc/}{\linkhilite{Numpy Docs}} website for example has all the documentation you need for Numpy \hyperlink{python}{\linkhilite{[{\xreffont 51}]}}. Navigating that website you can find the same \href{https://numpy.org/doc/1.19/reference/generated/numpy.arctan2.html?highlight=arctan2\#numpy.arctan2}{\linkhilite{documentation for arctan2}}.%
\par
As a last resort you can always Google \textquotedblleft{}how to compute the inverse tangent 2 function in Python\textquotedblright{}. Note though that there is so much content out there on Google that you could easily get lost. Still, there's also so much information that the answers are out there for just about anything.%
\par
So you have three methods for finding out how to program in python. The dir and \emph{\textunderscore{}\textunderscore{}doc\textunderscore{}\textunderscore{}} functions in Python, using the appropriate documentation online and of course Google. I'm lumping Youtube in with Google which is also another way to learn information although when I want to find information quickly I just use the documentation. It's the best in my opinion.%
\end{subsectionptx}
\end{sectionptx}
%
%
\typeout{************************************************}
\typeout{Section 5.4 Assignment}
\typeout{************************************************}
%
\begin{sectionptx}{Section}{Assignment}{}{Assignment}{}{}{plotting-6}
\begin{introduction}{}%
The assignment for this project is to plot the equation below from 0 to 10 seconds and include that Figure in a report. The Figure must include a grid and an x-axis label (‘Time (sec)') and y-axis label (‘Temperature (F)'). Any plotting tool Sheets, Excel, Colab, Python, Octave or similar is acceptable provided the figure conforms to the requirements above and the specific rubric below.%
\begin{equation}
T(t) = 60(1-e^{-5t})+30\label{plotting-6-2-1-1}
\end{equation}
%
\begin{enumerate}
\item{}Title page with name,title and date - 5\%%
\item{}Introduction - 5\%%
\item{}All figures must have appropriate annotations, figure captions, figure labels and a paragraph of supporting text - 10\%%
\item{}Include a screenshot (entire screen) of your plotting environment with your code pasted into the workspace - 40\%%
\item{}Include the plot of temperature vs time being sure to save the figure so it is in high resolution - 40\%%
\item{}Appendix A - Video Details - Film yourself and your entire screen running the example code and showing the result of \emph{plt.show()} - Pass\slash{}Fail%
\item{}Appendix B - Code used to take data and code used to plot in plain text (no screenshots) - Pass\slash{}Fail%
\item{}Appendix C - Screenshot of Excel or Sheets if applicable - Pass\slash{}Fail%
\item{}\alert{If you are asked to fix something in the Appendices such as missing content in your video or screenshots of your code, you will lose 10\% of your grade everytime you submit an update after the deadline.}%
\end{enumerate}
%
\end{introduction}%
\end{sectionptx}
\end{chapterptx}
%
%
\typeout{************************************************}
\typeout{Chapter 6 Setting up your Microcontroller}
\typeout{************************************************}
%
\begin{chapterptx}{Chapter}{Setting up your Microcontroller}{}{Setting up your Microcontroller}{}{}{getting-started}
\renewcommand*{\chaptername}{Chapter}
\begin{introduction}{}%
\begin{listptx}{List}{\textbf{Parts List}}{getting-started-2-1}{}%
%
\begin{enumerate}
\item{}Laptop%
\item{}\href{https://codewith.mu/en/download}{\linkhilite{Mu}} \hyperlink{mu}{\linkhilite{[{\xreffont 57}]}} or \href{https://www.arduino.cc/en/software/}{\linkhilite{Arduino IDE}} \hyperlink{Arduino}{\linkhilite{[{\xreffont 35}]}}%
\item{}CPX\slash{}CPB \href{https://circuitpython.org}{\linkhilite{UF2 File}} \hyperlink{CircuitPython}{\linkhilite{[{\xreffont 30}]}}%
\item{}CPX\slash{}CPB or Arduino Family (UNO, MEGA, or similar)%
\item{}USB Cable (with a data line. Not all USB cables have data lines)%
\end{enumerate}
\end{listptx}%
\begin{listptx}{List}{\textbf{Learning Objectives}}{getting-started-2-2}{}%
%
\begin{enumerate}
\item{}Learn how to setup your Microcontroller%
\item{}Learn how to install Mu or Arduino IDE%
\item{}Learn how to code in CircuitPython or Arduino%
\item{}Learn about hardware and software debugging%
\item{}Learn the difference between Python and CircuitPython and potentially Arduino%
\end{enumerate}
\end{listptx}%
By now you hopefully have your Circuit Playground Express (CPX), Circuit Playground Bluefruit (CPB) or an Arduino and it's time to get your microcontroller up and running. A very in depth and detailed tutorial can be found on the \href{https://learn.adafruit.com/circuitpython-made-easy-on-circuit-playground-express/first-things-first}{\linkhilite{Adafruit Learn site}}\hyperlink{Adafruit}{\linkhilite{[{\xreffont 32}]}} for the CPX and similar tutorials exists for the Arduino. This chapter is a summary of what you need to do to get the CPX or Arduino up and running as well as writing a simple blink code.%
\end{introduction}%
%
%
\typeout{************************************************}
\typeout{Section 6.1 Setting up the CPX\slash{}CPB}
\typeout{************************************************}
%
\begin{sectionptx}{Section}{Setting up the CPX\slash{}CPB}{}{Setting up the CPX\slash{}CPB}{}{}{getting-started-setup}
When you get your CPX and plug it into the computer via USB it actually won't run CircuitPython just yet. First you need to double click the reset button (the button in the center. It says RESET above the button) and put it into boot mode. All the neopixels (the ring lights on the CPX) will light up green.%
\begin{figureptx}{Figure}{The Circuit Playground Express in Boot mode}{getting-started-setup-3}{}%
\begin{image}{0.135}{0.73}{0.135}{}%
\includegraphics[width=\linewidth]{external/CPX_boot.jpeg}
\end{image}%
\tcblower
\end{figureptx}%
Something called CPLAYBOOT will pop up on your computer just like a USB stick or external harddrive. A couple files with be in there but it doesn't matter what they say right now.%
\begin{figureptx}{Figure}{Contents of the CPLAYBOOT Folder}{getting-started-setup-5}{}%
\begin{image}{0.135}{0.73}{0.135}{}%
\includegraphics[width=\linewidth]{external/CPLAYBOOT.png}
\end{image}%
\tcblower
\end{figureptx}%
You then need to download what's called a \href{https://circuitpython.org/board/circuitplayground_express/}{\linkhilite{UF2 file}} and transfer it onto the CPX. Note if you purchased a kit with the Bluefruit you need to download a \href{https://circuitpython.org/board/circuitplayground_bluefruit/}{\linkhilite{different UF2}}\hyperlink{CircuitPython}{\linkhilite{[{\xreffont 30}]}}. Make sure you get the right one. Once the UF2 is downloaded you need to drag the UF2 over to the CPLAYBOOT drive on your computer. After a bit of time a USB drive called CIRCUITPY will pop up as a flash drive on your computer. The CPX is now like a USB stick with 2MB of storage. \alert{Note that if you have an old bootloader you must update the \href{https://learn.adafruit.com/adafruit-circuit-playground-bluefruit/update-bootloader-use-command-line}{\linkhilite{bootloader}} \hyperlink{Adafruit}{\linkhilite{[{\xreffont 32}]}}.}%
\begin{figureptx}{Figure}{Contents of the CIRCUITPY Folder}{getting-started-setup-7}{}%
\begin{image}{0.135}{0.73}{0.135}{}%
\includegraphics[width=\linewidth]{external/CIRCUITPY.png}
\end{image}%
\tcblower
\end{figureptx}%
At this point the CPX is like a mini computer. If you put CircuitPython code on the \textquotedblleft{}flash drive\textquotedblright{} it will run the code. Since I've done this before, there are a number of files already on my CPX. You may have some or none of these. The \textquotedblleft{}boot\textunderscore{}out.txt\textquotedblright{} file will tell you the version of CircuitPython you have on the CPX. Mine says this:%
\begin{preformatted}
Adafruit CircuitPython 5.3.0 on 2020-04-29; Adafruit CircuitPlayground Express with samd21g18
\end{preformatted}
Which means I have CircuitPython version 5.3.0 last updated on April 29th, 2020. The CPX itself is using the \href{https://www.microchip.com/wwwproducts/en/ATsamd21g18}{\linkhilite{ATSAMD21G18}} microprocessor \hyperlink{CPXchip}{\linkhilite{[{\xreffont 58}]}}. The folder lib is a folder with extra libraries that you may need to install later. The file main.py is the CircuitPython file that the CPX is currently running. \alert{The CPX can store as many CircuitPython files as possible for a 2MB flash drive but it will only run ONE CircuitPython script at a time and that file must be named main.py} The folder \emph{System\textunderscore{}Volume\textunderscore{}Information} is a file management folder that we will never use.%
\par
If you want the CPX to run code you simply need to edit the file \emph{main.py} or if that file does not exist you just need to create it. You could just open Notepad or any other text app (Sublime,TexWrangler, Emacs, Vi, Nano, Gedit, Notepad++, Wordpad, VSCode, etc) but the CPX has alot of debugging options and it is recommended to use a program called \textquotedblleft{}\href{https://codewith.mu/en/download}{\linkhilite{Mu}}\textquotedblright{}\hyperlink{mu}{\linkhilite{[{\xreffont 57}]}}. Mu is a good way to write and debug code on the CPX. Note that Mu is only used to program the CPX. If you want to run Python code on your laptop you need to use Thonny or Spyder (or whatever other IDE you downloaded). If you want to run CircuitPython code on your CPX you must use Mu. Once you've downloaded and installed Mu and open it up it will look like this (Note it's possible the the software gets updated from the time this book is published. As such be sure to select the CircuitPython Mode for board development).%
\begin{figureptx}{Figure}{Opening Screen for Mu}{getting-started-setup-12}{}%
\begin{image}{0.135}{0.73}{0.135}{}%
\includegraphics[width=\linewidth]{external/Mu.png}
\end{image}%
\tcblower
\end{figureptx}%
Make sure to select the \textquotedblleft{}Adafruit CircuitPython\textquotedblright{} option. If the software has been updated and that option no longer exists, be sure to select the option that says CircuitPython for board development.%
\par
Alright, let's start writing code! If you have a file called \emph{main.py} on your CPX click the Load button and load \emph{main.py} (make sure to load the main.py that is stored on your CPX and not somewhere else on your computer.) If a file \emph{main.py} does not exist on your CPX simply click the New button and then Save the file as \emph{main.py} (again make sure you save it to the CPX and not to your computer)%
\begin{figureptx}{Figure}{Opening main.py in Mu}{getting-started-setup-15}{}%
\begin{image}{0.135}{0.73}{0.135}{}%
\includegraphics[width=\linewidth]{external/Mu2.png}
\end{image}%
\tcblower
\end{figureptx}%
You'll see that I am accessing the CIRCUITPY drive and saving the file as main.py The file itself is empty and just has a comment using the \# symbol. At this point since the file is blank the CPX won't do anything.%
\par
\alert{Now this is really important. Your CPX is a USB stick that can hold as many CircuitPython files as 2MB will allow but it can only run or execute one CircuitPython script at a time. Furthermore it will only run two types of files. It will run code.py if it exists and if it can't find code.py it will run main.py If the CPX can't find main.py or code.y it will just not do anything. If you have two versions of main.py or a combination of main.py or code.py it will run one of them and not the other. Make sure you only have one version of main.py or code.by but not both! Some common things to check if your CPX isn't working.}%
\begin{enumerate}
\item{}Make sure you're using Mu in the right Mode%
\item{}Make sure main.py or code.py is on the CIRCUITPY drive and not somewhere on your computer.%
\item{}Make sure you are editing the right file in Mu. Do you have two versions of main.py?%
\item{}Are you editing using Thonny or Spyder?%
\item{}Are you editing a file on your computer? Make sure you are writing to the CIRCUITPY drive.%
\item{}Unplug the CPX, close Mu and try again.%
\end{enumerate}
%
\par
So let's get the CPX to do something simple like blink an LED. I have an entire \href{https://github.com/cmontalvo251/Microcontrollers}{\linkhilite{Github}} devoted to Microelectronics. Specifically I have a \href{https://github.com/cmontalvo251/Microcontrollers/tree/master/Circuit_Playground/CircuitPython}{\linkhilite{folder with all of my Circuit Playground files}}. The easiest program to run is the \href{https://github.com/cmontalvo251/Microcontrollers/blob/master/Circuit_Playground/CircuitPython/blink/blink.py}{\linkhilite{blink.py}} script. I've attached a screenshot of the script below\hyperlink{Github}{\linkhilite{[{\xreffont 33}]}}.%
\begin{figureptx}{Figure}{CircuitPython Blink Code for the Circuit Playground Express}{getting-started-setup-19}{}%
\begin{image}{0.135}{0.73}{0.135}{}%
\includegraphics[width=\linewidth]{external/blink.png}
\end{image}%
\tcblower
\end{figureptx}%
We will talk about what this code is doing later on. For now copy and paste these 17 or so lines of code and paste them into Mu specifically the \emph{main.py} script. It will hopefully look like this.%
\begin{figureptx}{Figure}{Blink Code inside Mu}{getting-started-setup-21}{}%
\begin{image}{0.135}{0.73}{0.135}{}%
\includegraphics[width=\linewidth]{external/blinkmu.png}
\end{image}%
\tcblower
\end{figureptx}%
Make sure to save. You can click \emph{Zoom-In} and \emph{Zoom-out} to zoom in and out to change the font size and see more of the output. If the blink code is working you will see a red led labeled D13 on the CPX blink back and forth. D13 stands for digital pin 13. You'll notice there are analog pins labeled A5 and A6 among other numbers. Let's talk about the code a bit more and explain why it's doing what it's doing. The three lines at the top of the code are \emph{import} commands to import different modules just like we did for \emph{numpy} and \emph{matplotlib}. In this case the modules being imported are \emph{board}, \emph{digitalio} and \emph{time}. The module \emph{board} is used to import the layout of the CPX so we can access different pins on the CPX. The module \emph{digitalio} stands for digital input output which means we are inputting and outputting digital signals. Since we can combine this module with the \emph{board} module we will be able to output digital signals to different pins on the CPX board. Remember that PCB stands for printed circuit board so \emph{board} implies we are accessing pins on the CPX PCB. Hopefully that makes sense. The final module we are importing is \emph{time} which acts just like the \emph{time} module on your desktop computer. It will let us access the CPX's internal clock.%
\par
Moving along, lines 7 and 8 create an LED object using the \emph{board} module and \emph{digitalio}. It's a long line of code that basically says, create a variable called \emph{led} that lets us output a digital signal to pin D13. We also set the direction of the LED to output since we only want to write to the LED.%
\par
Lines 10 through 19 kick off an infinite loop that never ends. The line that says \emph{while True:} means loop while \emph{True}. Well \emph{True} is always true which means it will loop forever. The colon at the end of the line tells CircuitPython that the loop condition statement ends and to begin looping from 11 through 19.%
\par
Line 11 specifically says \emph{print(time.monotonic())}. First the \emph{print()} function is used to print things so that you and I can see it. Rather than just seeing a blinking LED we want to see the time printed. The \emph{time.monotonic()} is using the module \emph{time} which we imported and using a function from that module called \emph{monotonic()} which calls the internal clock of the CPX. So how do you see the output from the print statement? Hit the \emph{Serial} button on Mu and you will hopefully see some output. Here's what it looks like on my machine.%
\begin{figureptx}{Figure}{Serial Monitor Open in Mu}{getting-started-setup-26}{}%
\begin{image}{0.135}{0.73}{0.135}{}%
\includegraphics[width=\linewidth]{external/blinkmu2.png}
\end{image}%
\tcblower
\end{figureptx}%
In this case you can see the time printed every time it goes through the loop. You may even see an error. This \emph{Serial} button is great for debugging because it will tell you the error in your code. For example, in the picture below I have an error in my code and the \emph{Serial} output is letting me know.%
\begin{figureptx}{Figure}{Example Error Shown in Serial Monitor}{getting-started-setup-28}{}%
\begin{image}{0.135}{0.73}{0.135}{}%
\includegraphics[width=\linewidth]{external/blinkerror.png}
\end{image}%
\tcblower
\end{figureptx}%
In this case I have an error on line 5. It's saying there is no module named \emph{timed}. The reason that module doesn't exist is because the module is actually \emph{time} not \emph{timed}. You can use the \emph{Serial} monitor to check on your program and see what errors you may have. Ok so there are two more lines of code to discuss. They are \emph{led.value = True} and \emph{time.sleep(0.5)}.%
\par
These lines of code are repeated throughout the while loop and do two things. First, the \emph{led.value} either sets the value of the LED to True which turns the light on or False which turns the light off. The LED is digital which means the signal can either be on or off. There's no in between for digital signals. The \emph{time.sleep()} function tells the CPX to pause for half a second. You can change the number in the parentheses if you want to change length of time the code pauses. Note that the CPX completely pauses. That is no code runs during a sleep.%
\par
If you've gotten the LED to blink you're all set for this lab. However, I'd like to you learn a few more things about documentation. Just like Python on your desktop you can lookup the documentation on the CPX itself. For example, I've added a print statement to print the directory of time.%
\begin{figureptx}{Figure}{Showing Directory of the time module in Mu}{getting-started-setup-32}{}%
\begin{image}{0.135}{0.73}{0.135}{}%
\includegraphics[width=\linewidth]{external/printtimeMu.png}
\end{image}%
\tcblower
\end{figureptx}%
In this case I've added \emph{print(dir(time))} to line 7. The output shows that the \emph{time} module has 9 functions including \emph{monotonic()}. Unfortunately CircuitPython does not have \emph{\textunderscore{}\textunderscore{}doc\textunderscore{}\textunderscore{}} functions built in which means if you want to learn about a specific function, you need to visit the \href{https://circuitpython.readthedocs.io/en/5.3.x/README.html}{\linkhilite{documentation for CircuitPython}}. Here's the \href{https://circuitpython.readthedocs.io/en/latest/shared-bindings/time/index.html}{\linkhilite{specific documentation for the}} \href{https://circuitpython.readthedocs.io/en/latest/shared-bindings/time/index.html}{\linkhilite{time}} \href{https://circuitpython.readthedocs.io/en/latest/shared-bindings/time/index.html}{\linkhilite{module}}. A lot of \href{https://github.com/adafruit/Adafruit_Learning_System_Guides/tree/master/CircuitPython_Essentials}{\linkhilite{example code from the Adafruit Learning System is also on Github}}\hyperlink{CircuitPython}{\linkhilite{[{\xreffont 30}]}}.%
\par
Finally, in order to keep practicing using Python on your desktop I want you to modify the code above to run on your desktop computer. You're going to have to modify a few things. First, make sure to open Thonny or Spyder depending on which version of Python IDE you downloaded. Then, only import the \emph{time} module. All the other modules don't exist on your desktop. Also, get rid of all the lines of code that blink the LED. We just want to print time in a while loop. Finally, the time module on your desktop uses a function called time.time() (unless you have Python3 installed) so when you print time make sure to use that module instead. Again visit the \href{https://docs.python.org/3/library/time.html\#functions}{\linkhilite{documentation for time for Python if you want to learn more}}\hyperlink{python}{\linkhilite{[{\xreffont 51}]}}. Thonny by default will load Python version 3 but it's possible you may have Python version 2 so make sure you look up the documentation for the appropriate version of Python. After searching through the documentation you can use a function called asctime() on your desktop. This is the output I get in Thonny when I add a sleep of 1 second. The exercise for you is to do something similar.%
\begin{figureptx}{Figure}{Output of datetime in Thonny}{getting-started-setup-35}{}%
\begin{image}{0.135}{0.73}{0.135}{}%
\includegraphics[width=\linewidth]{external/blinkcomputer.png}
\end{image}%
\tcblower
\end{figureptx}%
\end{sectionptx}
%
%
\typeout{************************************************}
\typeout{Section 6.2 TL;DR \slash{} Extra Help for the CPX\slash{}CPB}
\typeout{************************************************}
%
\begin{sectionptx}{Section}{TL;DR \slash{} Extra Help for the CPX\slash{}CPB}{}{TL;DR \slash{} Extra Help for the CPX\slash{}CPB}{}{}{getting-started-4}
%
\begin{enumerate}
\item{}Plug in your CPX and double tap the reset button to go into reset mode. CPLAYBOOT will mount to your computer. (Some CPXes you need to hold the button down instead of double tap).%
\item{}Download the \href{https://circuitpython.org/downloads}{\linkhilite{UF2 file}}\hyperlink{CircuitPython}{\linkhilite{[{\xreffont 30}]}}.%
\item{}Drag the UF2 file to your CPLAYBOOT drive. After a few seconds CIRCUITPY will mount.%
\item{}You need to then download \href{https://codewith.mu/en/download}{\linkhilite{Mu}}\hyperlink{mu}{\linkhilite{[{\xreffont 57}]}}.%
\item{}Open Mu and make sure to select the mode Adafruit CircuitPython (or just CircuitPython if you have a new version)%
\item{}Open main.py (or code.py) in Mu from the CIRCUITPY drive in Mu. (It's possible you have code.py on your drive instead of main.py which is fine. Just open one or the other)%
\item{}Copy the \href{https://github.com/cmontalvo251/Microcontrollers/blob/master/Circuit_Playground/CircuitPython/blink/blink.py}{\linkhilite{blink.py}}\hyperlink{Github}{\linkhilite{[{\xreffont 33}]}} script into main.py (or code.py)%
\end{enumerate}
There is an accompanying \href{https://www.youtube.com/watch?v=XFvLn6rwm3I}{\linkhilite{youtube video to help you see me perform the 8 steps above}}\hyperlink{youtube}{\linkhilite{[{\xreffont 50}]}}.%
\end{sectionptx}
%
%
\typeout{************************************************}
\typeout{Section 6.3 Setting up an Arduino}
\typeout{************************************************}
%
\begin{sectionptx}{Section}{Setting up an Arduino}{}{Setting up an Arduino}{}{}{getting-started-5}
Setting up the Arduino is a bit different than the CPX\slash{}CPB family of microcontrollers. The main difference is that the code is done on your computer in the Arduino IDE and then sent or compiled to the Arduino. This means that the code from the Arduino is unrecoverable if you don't have the code on your computer. It's also a lot harder to debug code on the Arduino because it isn't like CircuitPython where the language will tell you which line of code is incorrect. My strategy with coding on the Arduino is pretty archaic in that I like to add one line of code at a time and have a lot of print statements to see where the code fails. This won't solve memory leaks or hardware issues but it does help a bit. In this section we will dive into getting the blink code to work. The pro of using an Arduino is highlighted here in that the only setup is in the IDE itself. To start you just plug in your Arduino to your laptop and then open the Arduino IDE.%
\begin{figureptx}{Figure}{Arduino IDE Setup}{getting-started-5-3}{}%
\begin{image}{0.135}{0.73}{0.135}{}%
\includegraphics[width=\linewidth]{external/ArduinoIDE_Setup.png}
\end{image}%
\tcblower
\end{figureptx}%
When you open the Arduino IDE for the first time the window will open up with the \emph{void setup()} routine and the \emph{void loop()} routine. The \emph{setup} routine runs once and is where all the one time setup code goes. The \emph{loop} routine runs over and over again and never stops and unlike the CPX\slash{}CPB you can never stop it unless you have a line of code that stops the code automatically. To start programming on your microcontroller you first need to select the Tools\textgreater{}\emph{Board, Processor} and \emph{Port}. The \emph{Board} depends on the kind of board you purchased. For example, I purchased the Arduino MEGA so I selected Mega for the Board and Processor but you might have to select UNO if you're using an Arduino UNO. The \emph{Port} is the serial port that the Arduino is connected to. Since I'm on Linux the port the Arduino is connected to is \emph{\slash{}dev\slash{}ttyACMO} while it would be \emph{COM3} or something similar on Windows or \emph{\slash{}dev\slash{}tty.usbmodem14101} on Mac. This is shown in the Figure above. When you're ready to code the blink example the easy part is that the code for that is already included in the examples in the Arduino IDE. To do that click \emph{File\textgreater{}Examples\textgreater{}01.Basics\textgreater{}Blink}%
\begin{figureptx}{Figure}{Arduino Blink Example Using Pin 13}{getting-started-5-5}{}%
\begin{image}{0.135}{0.73}{0.135}{}%
\includegraphics[width=\linewidth]{external/ArduinoBlink_13.png}
\end{image}%
\tcblower
\end{figureptx}%
The figure above shows the code for blinking the LED connected to pin 13. The Arduino UNO\slash{}MEGA and also the CPX\slash{}CPB all share the same default LED on pin 13. The default example program will say \emph{LED\textunderscore{}BUILTIN} but I changed the pin to pin 13 to be more specific. The code is pretty straightforward where the setup routine makes pin 13 an output pin using the \emph{pinMode()} function and then the loop routine sets pin 13 to HIGH and LOW using the \emph{digitalWrite()} function with a \emph{delay} function used to pause for 1000 ms or 1 second. In this case HIGH is on and LOW is off. If you click the \emph{Upload} button which looks like a right arrow the code will be compiled to binary and then uploaded to the Arduino. If it's working properly the LED on pin 13 will blink.%
\begin{figureptx}{Figure}{Arduino MEGA with Pin 13 Illuminated}{getting-started-5-7}{}%
\begin{image}{0.135}{0.73}{0.135}{}%
\includegraphics[width=\linewidth]{external/Arduino_MEGA_Blink.png}
\end{image}%
\tcblower
\end{figureptx}%
Again there isn't a lot of ability to debug if it isn't working right. Note that if you want to save this code you must click Save and then save the \emph{.ino} file to somewhere on your computer. If you want to share code you can't just give someone your Arduino, you have to send them the code because the code is unrecoverable from the binary format on the Arduino itself. As for the code itself, the \slash{}\slash{} is used for comments and semicolons ; are used to end every single line. Unlike python you don't need to worry about tabs to be setup properly. Loops and if statements are separated by brackets \textbraceleft{} and \textbraceright{}.%
\begin{figureptx}{Figure}{Arduino Blink Example with Serial Monitor}{getting-started-5-9}{}%
\begin{image}{0.135}{0.73}{0.135}{}%
\includegraphics[width=\linewidth]{external/Arduino_Blink_Serial.png}
\end{image}%
\tcblower
\end{figureptx}%
As I mentioned however, it's possible to add many print statements to the code and open up the Serial Monitor. The figure above shows the edited code with \emph{Serial.begin(9600)} added to the setup routine which opens up serial communication between your computer and the Arduino at a baud rate of 9600 bits per second. In the loop routine I have added many \emph{Serial.print("");} commands to print to the Serial monitor which can be opened by clicking the button that looks like a magnifying glass. Make sure the baud rate in the serial monitor is set to the same 9600 value or the values will be gibberish. Also \emph{millis()} is the internal timer on all Arduinos that output time in milliseconds. Note because this program is so simple, the code is left to the student to recreate.%
\end{sectionptx}
%
%
\typeout{************************************************}
\typeout{Section 6.4 Assignment}
\typeout{************************************************}
%
\begin{sectionptx}{Section}{Assignment}{}{Assignment}{}{}{getting-started-6}
\begin{introduction}{}%
For this assignment you must appropriately set up your CPX\slash{}CPB, get the blink code to run and include some screenshots (your entire computer screen) of relevant steps. \alert{Note for the blink code, modify blink pattern to be different from the default on Github}. The specific items requested are shown below.%
%
\begin{enumerate}
\item{}Title page with name,title and date - 5\%%
\item{}Introduction - 5\%%
\item{}All figures must have appropriate annotations, figure captions, figure labels and a paragraph of supporting text - 10\%%
\item{}Include a screenshot (entire computer screen) of your computer showing the CPLAYBOOT drive on your computer - 20\%%
\item{}Include a screenshot (entire computer screen) of your computer showing the CIRCUITPY drive on your computer - 20\%%
\item{}Include a screenshot (entire computer screen) of Mu with the Serial monitor open showing the output of the blink code. Remember to modify the blink pattern - 20\%%
\item{}Include a photo of your CPX\slash{}CPB with the red led on - 20\%%
\item{}Appendix A - Video Details - Film yourself showing the CPX\slash{}CPB blinking and your entire screen with your blink code running with the Serial monitor open - Pass\slash{}Fail%
\item{}Appendix B - Code used on your CPX\slash{}CPB and code used to plot in plain text (no screenshots) - Pass\slash{}Fail%
\item{}Appendix C - Screenshot of Excel or Sheets if applicable - Pass\slash{}Fail%
\item{}\alert{If you are asked to fix something in the Appendices such as missing content in your video or screenshots of your code, you will lose 10\% of your grade everytime you submit an update after the deadline.}%
\end{enumerate}
\end{introduction}%
\end{sectionptx}
\end{chapterptx}
%
%
\typeout{************************************************}
\typeout{Chapter 7 Help}
\typeout{************************************************}
%
\begin{chapterptx}{Chapter}{Help}{}{Help}{}{}{troubleshooting}
\renewcommand*{\chaptername}{Chapter}
\begin{introduction}{}%
Are you getting tons of errors because a certain moduble is not installed? Is your CPX\slash{}CPB broken or not running code? This chapter walks through a few common mishaps that will happen as you work through these code examples.%
\end{introduction}%
%
%
\typeout{************************************************}
\typeout{Section 7.1 Updating your Bootloader}
\typeout{************************************************}
%
\begin{sectionptx}{Section}{Updating your Bootloader}{}{Updating your Bootloader}{}{}{troubleshooting-3}
Some of the newer CircuitPython UF2s do not install correctly on old boot loaders. If you put the CPX\slash{}CPB into CPLAYBOOT mode and open the \emph{INFO\textunderscore{}UF2.TXT} file it will tell you the version of the bootloader that you have just like the figure below. \begin{figureptx}{Figure}{Old Bootloader Version}{old-bootloader}{}%
\begin{image}{0.1}{0.8}{0.1}{}%
\includegraphics[width=\linewidth]{external/old_bootloader.png}
\end{image}%
\tcblower
\end{figureptx}%
 In this case the version of the bootloader is 0.2.11 which is older than most new UF2s require. If you're having an issue with an old bootloader be sure to update the \href{https://learn.adafruit.com/adafruit-circuit-playground-bluefruit/update-bootloader-use-command-line}{\linkhilite{bootloader for the CPB}} or alternatively the \href{https://learn.adafruit.com/adafruit-circuit-playground-express/updating-the-bootloader}{\linkhilite{bootloader for the CPX}}. At the time of this writing here is the command you need to upload the new bootloader the CPB. If you have version 0.4.0 or later you can use the \href{https://learn.adafruit.com/adafruit-circuit-playground-bluefruit/update-bootloader-use-uf2}{\linkhilite{UF2 method}} \hyperlink{Adafruit}{\linkhilite{[{\xreffont 32}]}}. If your bootloader is too old you need \emph{adafruit-nrfutil}. Here is the command you need to upload the new bootloader to the CPB.\begin{preformatted}
adafruit-nrfutil --verbose dfu serial --package 
circuitplayground_nrf52840_bootloader-0.8.2_s140_6.1.1.zip -p /dev/ttyACM0 
-b 115200 --singlebank --touch 1200
\end{preformatted}
\terminology{For windows users} you use the \emph{adafruit-nrfutil.exe} file from the Adafruit bootloader website. Then, usingcommand prompt to run the command above you need to change the tty port to the correct COM port which can be found using \emph{Device Manager}.%
\par
\terminology{For linux users} you will need a few things%
\begin{enumerate}
\item{}You need pip3 \(->\) \textdollar{} sudo apt install python3-pip%
\item{}You need adafruit-nrfutil \(->\) \textdollar{} pip3 install -{}-{}user adafruit-nrfutil%
\item{}adafruit-nrfutil gets installed into \textbackslash{}textasciitilde\slash{}.local\slash{}bin%
\item{}You need to add your user to dialout \(->\) \textdollar{} sudo usermod -a -G your\textunderscore{}username%
\item{}Restart your computer%
\end{enumerate}
This is what it looks like when running the \emph{adafruit-nrfutil} command on linux. \begin{figureptx}{Figure}{Output from running the bootloader update command on linux}{troubleshooting-3-3-4}{}%
\begin{image}{0.1}{0.8}{0.1}{}%
\includegraphics[width=\linewidth]{external/bootloader_linux_command.png}
\end{image}%
\tcblower
\end{figureptx}%
 For the CPX there is a UF2 that you transfer over to the CPLAYBOOT drive. At the time of this writing the version you're looking for is \emph{update-bootloader-circuitplay\textunderscore{}m0-v3.15.0.uf2}. Once you update the bootloader you can check the version by opening the \emph{INFO\textunderscore{}UF2.TXT} file again and it should show the new version. The figure below shows an example of the new bootloader version after updating. In this case the version is 0.10.0 which is much newer than the old 0.2.11 version. \begin{figureptx}{Figure}{Updated Bootloader Version}{troubleshooting-3-3-7}{}%
\begin{image}{0.1}{0.8}{0.1}{}%
\includegraphics[width=\linewidth]{external/updated_bootloader.png}
\end{image}%
\tcblower
\end{figureptx}%
%
\end{sectionptx}
%
%
\typeout{************************************************}
\typeout{Section 7.2 Installing Modules}
\typeout{************************************************}
%
\begin{sectionptx}{Section}{Installing Modules}{}{Installing Modules}{}{}{installing-modules}
\alert{One issue you're going to run into when you run certain codes is that you won't have some of the modules on your CPB\slash{}CPX.} To fix this you need to download the \href{https://circuitpython.org/libraries}{\linkhilite{CircuitPython Libraries}}\hyperlink{CircuitPython}{\linkhilite{[{\xreffont 30}]}}. You need to download the appropriate version: 6.x, 7.x or 8.x. The figure below is a screenshot of the libraries homepage on the CircuitPython website. \begin{figureptx}{Figure}{CircuitPython Libraries homepage}{installing-modules-2-4}{}%
\begin{image}{0.1}{0.8}{0.1}{}%
\includegraphics[width=\linewidth]{external/libraries.png}
\end{image}%
\tcblower
\end{figureptx}%
 How do you know what version of CircuitPython you have? Well head over to your CIRCUITPY drive and open the \emph{boot\textunderscore{}out.txt} file and it will tell you the version. The Figure below shows an example \emph{boot\textunderscore{}out.txt} file. In this case the version is 10.1.4 so, you would download the 10.x version of the libraries. \begin{figureptx}{Figure}{Example boot\textunderscore{}out.txt file showing the version of CircuitPython}{installing-modules-2-7}{}%
\begin{image}{0.1}{0.8}{0.1}{}%
\includegraphics[width=\linewidth]{external/example_boot_out.png}
\end{image}%
\tcblower
\end{figureptx}%
 Note that this is the same version as the .UF2 file installed back in the Getting Started labs (See \hyperref[getting-started]{\linkhilite{Chapter~{\xreffont\ref{getting-started}}}}). When you download the modules it will download a .zip file. Extract the .zip file on your desktop computer and then open the \emph{lib} folder on your desktop and your CIRCUITPY. You then need to transfer the modules (ONLY THE ONES YOU NEED) from your desktop to your CPX\slash{}CPB \emph{lib} folder. The reason why you can't copy the entire folder is because the CPB\slash{}CPX only has 2MB of flash and the CircuitPython download is 4.1 MB at the time of this writing. If your CIRCUITPY drive does not have a \emph{lib} folder you can easily create one as shown in the Figure below. \begin{figureptx}{Figure}{CIRCUITPY Drive showing lib folder}{installing-modules-2-12}{}%
\begin{image}{0.1}{0.8}{0.1}{}%
\includegraphics[width=\linewidth]{external/CIRCUITPY.png}
\end{image}%
\tcblower
\end{figureptx}%
%
\par
After extracting the zip folder to your computer you can begin dragging files over to your CPX\slash{}CPB. The Figure below shows an example lib folder full of modules that have been dragged from the desktop to the CIRCUITPY drive. \begin{figureptx}{Figure}{Example lib folder on CIRCUITPY drive with modules dragged over from desktop}{installing-modules-3-1}{}%
\begin{image}{0.1}{0.8}{0.1}{}%
\includegraphics[width=\linewidth]{external/lib_folder_full.png}
\end{image}%
\tcblower
\end{figureptx}%
 The follow up question of course is \emph{Which modules do I need?} The answer to that question is in the code. When you run a code that needs a specific module, the Serial monitor will say \emph{ImportError: No module named 'x'}. In this case you will then know to install the module 'x'. You can also look at the import statements at the top of the code to see which modules are being imported and then make sure you have those modules installed. The figure below is an example of that error where the code is specifically looking for the module named \emph{adafruit\textunderscore{}circuitplayground}. In this case, the user would need to copy that specific module from the desktop to the CIRCUITPY drive \emph{lib} folder. \begin{figureptx}{Figure}{Example of a missing module error}{moduleerror}{}%
\begin{image}{0.1}{0.8}{0.1}{}%
\includegraphics[width=\linewidth]{external/missing_module.png}
\end{image}%
\tcblower
\end{figureptx}%
 Note that it's possible to import too many modules on the CPX\slash{}CPB and run out of memory. It turns out that the CPB has much more memory then the CPX. Run the commands below to see how much memory your CPX\slash{}CPX has left.\begin{preformatted}
import gc
gc.collect()
print(gc.mem_free())
\end{preformatted}
At the time of this writing, the CPX has about 17 KB of RAM and the CPB has around 140 KB of RAM. Note that RAM is different than flash. The CPX\slash{}CPB can hold about 2MB of data but only operate with 17\slash{}140 KB of working memory.%
\end{sectionptx}
%
%
\typeout{************************************************}
\typeout{Section 7.3 Reformatting your CPX\slash{}CPB}
\typeout{************************************************}
%
\begin{sectionptx}{Section}{Reformatting your CPX\slash{}CPB}{}{Reformatting your CPX\slash{}CPB}{}{}{troubleshooting-5}
If your CPX\slash{}CPX is completely bricked you can completely reformat it. In order to do that you simply put the CPX\slash{}CPB into CPLAYBOOT mode. Then drag either the \href{https://cdn-learn.adafruit.com/assets/assets/000/048/745/original/flash_erase_express.ino.circuitplay.uf2?1512152080}{\linkhilite{CPX eraser tool}} or the \href{https://cdn-learn.adafruit.com/assets/assets/000/082/950/original/CP_Bluefruit_QSPI_Erase.UF2?1572026649}{\linkhilite{CPB eraser tool}}. After that, you can reinstall the CircuitPython UF2 and start fresh (See \hyperref[getting-started]{\linkhilite{Chapter~{\xreffont\ref{getting-started}}}}) \hyperlink{Adafruit}{\linkhilite{[{\xreffont 32}]}}. If you have a linux computer you can also access the REPL of your CPX\slash{}CPB using the commands below\begin{preformatted}
sudo screen /dev/ttyACM0
import storage
storage.erase_filesystem()
\end{preformatted}
Then you would proceed with installing the CircuitPython UF2 again and starting fresh. Note that this will erase everything on your CPX\slash{}CPB so make sure you have a backup of your code and files before doing this. If you feel like you've tried everything this \href{https://learn.adafruit.com/adafruit-circuit-playground-express/troubleshooting}{\linkhilite{troubleshooting guide}} may have some other steps outlined to help you. However, I have had great experience with the eraser tool above \hyperlink{Adafruit}{\linkhilite{[{\xreffont 32}]}}.%
\end{sectionptx}
%
%
\typeout{************************************************}
\typeout{Section 7.4 Using an I2C Scanner}
\typeout{************************************************}
%
\begin{sectionptx}{Section}{Using an I2C Scanner}{}{Using an I2C Scanner}{}{}{I2C_scanner}
If you're having trouble getting an I2C device to work, it's a good idea to run an I2C scanner to make sure you can see the device and that you have the correct address. Here is a simple I2C scanner code that you can run on your CPX\slash{}CPB.%
\begin{preformatted}
import time
import board
import busio

i2c = busio.I2C(board.SCL, board.SDA)
while not i2c.try_lock():
    pass

try:
  while True:
      print("I2C addresses found:", [hex(device_address) for device_address in i2c.scan()])
  time.sleep(2)
\end{preformatted}
\end{sectionptx}
%
%
\typeout{************************************************}
\typeout{Section 7.5 Common Mishaps}
\typeout{************************************************}
%
\begin{sectionptx}{Section}{Common Mishaps}{}{Common Mishaps}{}{}{troubleshooting-7}
%
\begin{enumerate}
\item{}Your CPX is a USB stick that can hold as many CircuitPython files as 2MB will allow but it can only run or execute one CircuitPython script at a time. It will look for code.py and then main.py and that's it.%
\item{}Make sure you're using Mu in the right Mode (CircuitPython)%
\item{}Make sure main.py or code.py is on the CIRCUITPY drive and not somewhere on your computer.%
\item{}Make sure you are editing the right file in Mu. Do you have two versions of main.py or perhaps main.py and code.py?%
\item{}Do you have boot.py on there when you don't need it?%
\item{}Are you editing using Thonny or Spyder? You're supposed to use Mu.%
\item{}Are you editing a file on your computer? Make sure you are writing to the CIRCUITPY drive.%
\end{enumerate}
\end{sectionptx}
%
%
\typeout{************************************************}
\typeout{Section 7.6 Extra Help}
\typeout{************************************************}
%
\begin{sectionptx}{Section}{Extra Help}{}{Extra Help}{}{}{troubleshooting-8}
I have a youtube video where I \href{https://www.youtube.com/watch?v=xVMPgU5z9gE}{\linkhilite{update the bootloader on a CPB}}\hyperlink{youtube}{\linkhilite{[{\xreffont 50}]}} That may help you if you're having trouble following the bootloader tutorial above.%
\end{sectionptx}
\end{chapterptx}
%
%
\typeout{************************************************}
\typeout{Chapter 8 External LEDs and Push Buttons}
\typeout{************************************************}
%
\begin{chapterptx}{Chapter}{External LEDs and Push Buttons}{}{External LEDs and Push Buttons}{}{}{led-buttons}
\renewcommand*{\chaptername}{Chapter}
\begin{introduction}{}%
\begin{listptx}{List}{\textbf{Parts List}}{led-buttons-2-1}{}%
%
\begin{enumerate}
\item{}Laptop%
\item{}CPX + USB Cable%
\item{}2 Alligator clips%
\item{}Push Button%
\item{}Breadboard%
\item{}LED (Light Emitting Diode) (x3 in case you fry some)%
\item{}Resistor (300 to 1000 Ohms)%
\end{enumerate}
\end{listptx}%
\begin{listptx}{List}{\textbf{Learning Objectives}}{led-buttons-2-2}{}%
%
\begin{enumerate}
\item{}The VOUT and 3.3V pins are always "ON" even when code is not running on the CPX. So long as your CPX is plugged in via USB or a Lipo Battery%
\item{}LEDs are Light Emitting \emph{Diodes} which means current only flows in one direction%
\item{}LEDs need resistors in series otherwise they will get too hot and burn up%
\item{}Breadboard pinout diagrams%
\item{}Analog pins can be controlled by simply using the digitalio module%
\item{}LEDs can be hooked up to analog pins and set to blink by changing the board pin%
\end{enumerate}
\end{listptx}%
In this project we're going to use the same blink code as Chapter \hyperref[getting-started]{\linkhilite{Chapter~{\xreffont\ref{getting-started}}}} but modify it to blink an external LED. Another goal of this lab is to familiarize yourself with some of the pins on the CPX and create a simple circuit using the 5V (VOUT) pin on the CPX and one of the Analog pins.%
\par
Your laptop has a battery with something between 10 to 20V. There are DC to DC converters in your laptop that provide 5V to your USB ports. These USB ports can be used to power your CPX as you have done in the past few labs.%
\par
If you purchased the optional \href{https://www.adafruit.com/product/3287}{\linkhilite{battery pack}} you can also power the CPX using 3 AA batteries \hyperlink{AdafruitMAIN}{\linkhilite{[{\xreffont 34}]}}. These batteries nominally have 1.5 V but fully charged it's actually something like 1.8 V. So 1.8 times 3 is 5.4V which is enough to power the CPX. If you have the battery pack and some AA batteries, give it a try. If you still have the blink code from the last project on board you'll see the D13 LED blink as before. You won't be able to see the serial print() output as before but that code will be running which is why D13 is blinking. \alert{I have noticed that some of the battery packs have power and ground wires swapped. If the battery pack doesn't work it may be because those two wires are backwards.}%
\par
The CPX itself uses 3.3V logic which means when it converts numbers to binary a 0 (False) represents 0 volts and a 1 (True) represents 3.3V. The CPX has ports that are labeled various things. GND stands for ground and you need to hook the negative end of your circuit to this and it also has VOUT which supplies 5V to any circuit you build. Hook the positive end of your circuit to the VOUT pin. There is also a port labeled 3.3V and obviously that outputs 3.3V. The Figure below shows an \href{https://learn.adafruit.com/circuitpython-digital-inputs-and-outputs/digital-outputs }{\linkhilite{Arduino Metro}} hooked up to an LED on a breadboard (Courtesy of \href{https://learn.adafruit.com/u/tdicola}{\linkhilite{Tony Dicola}} \hyperlink{Metro}{\linkhilite{[{\xreffont 37}]}}).%
\begin{figureptx}{Figure}{Fritzing diagram of an Arduino Metro hooked up to an LED \hyperlink{Metro}{\linkhilite{[{\xreffont 37}]}}}{led-buttons-2-7}{}%
\begin{image}{0.135}{0.73}{0.135}{}%
\includegraphics[width=\linewidth]{external/circuitpython_01_digital_io_figure_1.png}
\end{image}%
\tcblower
\end{figureptx}%
Your lab today specifically involves an external LED as shown above. You can read about \href{https://learn.sparkfun.com/tutorials/light-emitting-diodes-leds/}{\linkhilite{LEDs}} more online if you wish\hyperlink{sparkfun}{\linkhilite{[{\xreffont 59}]}}. \alert{Remember that the long leg of the LED is the positive end and the short leg is the negative end.} The task today is to wire an LED up to the CPX in multiple ways. The figure above shows the LED connected to the A1 pin but we will be wiring up the LED to multiple pins on the CPX\slash{}CPB to gain understanding in how the circuit works. \alert{Note, whenever you modify a circuit on the breadboard, always be sure to remove power from the CPX. You can damage multiple components if you're not careful.}%
\end{introduction}%
%
%
\typeout{************************************************}
\typeout{Section 8.1 LED with no Code}
\typeout{************************************************}
%
\begin{sectionptx}{Section}{LED with no Code}{}{LED with no Code}{}{}{led-buttons-3}
For this part we are going to light up the LED without the use of any code on the CPX. First, wire up the circuit with the positive end connected to 5V. This is how my circuit looks. Make sure to use a resistor between 300 and 1000 Ohms. An LED does not have that much internal resistance so you need a resistor in series with an LED to reduce the amount of current flowing through the LED or the entire LED will fry. If you use a resistor that has too much impedance the LED just won't turn on because the voltage\slash{}current through the LED will be below the activation voltage of the circuit.%
\begin{figureptx}{Figure}{CircuitPlayground Express connected to an LED on a breadboard}{led-buttons-3-3}{}%
\begin{image}{0.135}{0.73}{0.135}{}%
\includegraphics[width=\linewidth]{external/LED1.jpeg}
\end{image}%
\tcblower
\end{figureptx}%
Once you have that circuit working, wire up the circuit again with the 3.3V output. Do you notice anything different when you hook up the circuit with different pins? Here's my circuit. Do you notice something different about the intensity of the LED? Why is it different?%
\begin{figureptx}{Figure}{CircuitPlayground Express connected to an LED using 3.3V}{led-buttons-3-5}{}%
\begin{image}{0.135}{0.73}{0.135}{}%
\includegraphics[width=\linewidth]{external/LED2.jpeg}
\end{image}%
\tcblower
\end{figureptx}%
\end{sectionptx}
%
%
\typeout{************************************************}
\typeout{Section 8.2 LED with a push button}
\typeout{************************************************}
%
\begin{sectionptx}{Section}{LED with a push button}{}{LED with a push button}{}{}{led-buttons-4}
Now that we understand breadboards a bit, we're now going to manually blink the LED using a push button placed onto the breadboard and have it act like a switch. Therefore, when the button is pressed, the LED will turn on and when the button is released the LED will turn off. The button just acts like a wire so you can plug in the button anywhere in the circuit.%
\begin{figureptx}{Figure}{CircuitPlayground Express turning on an LED with an external push button}{led-buttons-4-3}{}%
\begin{image}{0.135}{0.73}{0.135}{}%
\includegraphics[width=\linewidth]{external/LED3.jpeg}
\end{image}%
\tcblower
\end{figureptx}%
\end{sectionptx}
%
%
\typeout{************************************************}
\typeout{Section 8.3 LED with code}
\typeout{************************************************}
%
\begin{sectionptx}{Section}{LED with code}{}{LED with code}{}{}{led-buttons-5}
Next I want you to remove the button from the circuit and wire up the LED like you had it when the positive end was connected to VOUT or 3.3V. Except this time I want you to hook the positive end of the circuit to pin A2. Then edit your blink code to blink pin A2. Take a look at the blink code. Right now the code is blinking pin D13. How do you think you need to change the code to blink pin A2? Here's what my circuit looks like for this one. I won't include code for this one since you just need to change one line of code.%
\begin{figureptx}{Figure}{CircuitPlayground Express connected to an LED with an analog pin}{led-buttons-5-3}{}%
\begin{image}{0.135}{0.73}{0.135}{}%
\includegraphics[width=\linewidth]{external/LED4.jpeg}
\end{image}%
\tcblower
\end{figureptx}%
\end{sectionptx}
%
%
\typeout{************************************************}
\typeout{Section 8.4 LED with CPX button}
\typeout{************************************************}
%
\begin{sectionptx}{Section}{LED with CPX button}{}{LED with CPX button}{}{}{led-buttons-6}
Finally, I want to use one of the buttons on the CPX to blink the LED hooked up to pin A2. For this code to work you first need to detect a button press and then tell the program to change the light from True to False depending on what it's current status is. This one is a bit more difficult so I'll include the code here and discuss the code itself. Here is my circuit (identical) to the previous one with Button A on the CPX pressed down.%
\begin{figureptx}{Figure}{CircuitPlayground Express powering an LED with an on-board button}{led-buttons-6-3}{}%
\begin{image}{0.135}{0.73}{0.135}{}%
\includegraphics[width=\linewidth]{external/LED5.jpeg}
\end{image}%
\tcblower
\end{figureptx}%
Alright so how do we detect a button press? Well the documentation on this is not so straight forward. What we want to do is detect the INPUT of a digital signal and then do something if we detect that signal. Here's the code I created to get it to work.%
\begin{figureptx}{Figure}{CircuitPython code to power an LED with an on-board button}{led-buttons-6-5}{}%
\begin{image}{0.135}{0.73}{0.135}{}%
\includegraphics[width=\linewidth]{external/button_blink.png}
\end{image}%
\tcblower
\end{figureptx}%
The code above is an image. You could type exactly what you see and hope you don't type any errors (which is highly unlikely) or you could look for my code on my \href{https://github.com/cmontalvo251/Microcontrollers/tree/master/Circuit_Playground/CircuitPython}{\linkhilite{Github}}\hyperlink{Github}{\linkhilite{[{\xreffont 33}]}}. Have you bookmarked this link yet? I recommend you do so!!! So you're making a code about Buttons....hmmm. Click that Github link. Where do you think the code about Buttons is located? Anywho, let's talk about the code.%
\par
The first 3 lines are exactly the same as before. Line 5 through 7 are similar to creating the \textquotedblleft{}led\textquotedblright{} variable except we're using \emph{BUTTON\textunderscore{}A} as the board value and setting the direction to \emph{INPUT}. Finally we're setting the pull direction to \emph{DOWN}. This means the button acts like a pull down resistor and when it's pressed the value of the button goes \emph{HIGH}.%
\par
Lines 9-11 are the same as before. We create an led variable and tell the CPX that the led is hooked up to pin A2. We then start the while loop on line 13. First if you click the Serial button you'll see the text \emph{Button value is False}. It's False because \emph{buttonA.value} is not being pressed. On line 15 the led is set to False (turned off). On line 16 the button value is checked using an if statement as to whether or not the button is pressed (True). If the button is pressed, the user will be notified that the button is pressed on line 17 and the led will be set to True (on) in line 18. Lines 19-22 are while loop that will notify the Serial monitor that you must let go of the button before the code can continue to the main while loop. The \emph{time.sleep} functions are there to make sure a human can operate the button without code running faster than a human can press a button. When I press the button down here is the output I get from the \emph{Serial} monitor.%
\begin{figureptx}{Figure}{CircuitPython code to power an LED with an on-board button in Mu}{led-buttons-6-9}{}%
\begin{image}{0.135}{0.73}{0.135}{}%
\includegraphics[width=\linewidth]{external/MuButton.png}
\end{image}%
\tcblower
\end{figureptx}%
Here you'll see 4 lines that say \textquotedblleft{}Button value is False\textquotedblright{} and then two lines that say \textquotedblleft{}Button value is True\textquotedblright{} followed by 5 lines that say \textquotedblleft{}Waiting for you to let go...\textquotedblright{}. See if you can get this code to work and play with it and modify it as you see fit. By the way, the LED connected to pin D13 has this exact same circuitry, an LED a resistor, it's all just soldered to the PCB so you don't have to build it using a breadboard. Hopefully now you have some appreciation for buttons and LEDs!!%
\end{sectionptx}
%
%
\typeout{************************************************}
\typeout{Section 8.5 Extra Help}
\typeout{************************************************}
%
\begin{sectionptx}{Section}{Extra Help}{}{Extra Help}{}{}{led-buttons-7}
If you're not familiar with how breadboards work I would recommend watching this \href{https://www.youtube.com/watch?v=mE33WpRWrXs}{\linkhilite{video on how breadboards work}}\hyperlink{youtube}{\linkhilite{[{\xreffont 50}]}}.%
\end{sectionptx}
%
%
\typeout{************************************************}
\typeout{Section 8.6 Assignment}
\typeout{************************************************}
%
\begin{sectionptx}{Section}{Assignment}{}{Assignment}{}{}{led-buttons-8}
\begin{introduction}{}%
For this specific project you must get your LED to turn on and off using buttons as well as software. The different scenarios are shown below and the grading rubric is presented after this list.%
\begin{enumerate}
\item{}Get your LED to turn on with the VOUT pin.%
\item{}Get your LED to turn on with the 3.3V pin.%
\item{}Get your LED to blink automatically using an analog pin on the CPX\slash{}CPB.%
\item{}Get your LED to turn on using a button on the breadboard.%
\item{}Get your LED to blink automatically by pressing a button on the CPX\slash{}CPB.%
\end{enumerate}
%
%
\begin{enumerate}
\item{}Title page with name,title and date - 5\%%
\item{}Introduction with wiring diagram of LED connected to VOUT - 5\%%
\item{}All figures must have appropriate annotations, figure captions, figure labels and a paragraph of supporting text - 10\%%
\item{}Include a photo of your circuit with your LED turned on using VOUT. Based on the forward voltage in your LED and the resistance in your circuit, what is the voltage drop across the resistor and current through the LED? Assuming your LED acted like a resistor, what would the resistance of your LED be? - 10\%%
\item{}Include a photo of your circuit with your LED turned on using 3.3V. Based on the forward voltage in your LED and the resistance in your circuit, what is the voltage drop across the resistor and current through the LED? Assuming your LED acted like a resistor, what would the resistance of your LED be? - 10\%%
\item{}The USB port from your computer is 5V. The DC to DC converter on the CPX\slash{}CPB converts that 5V to 3.3V. Compute the Gain (\(G\)) for this converter both in magnitude and in \(dB\). - 10\%%
\item{}Include a photo of your circuit with your LED turned on using a button on your breadboard - 10\%%
\item{}Include a photo of your circuit with your LED turned on using an analog pin and then also include a screenshot (entire computer screen) of your Mu code showing the same analog pin used to blink your LED - 20\%%
\item{}Include a photo of your circuit with your LED turned on by pressing a button on the CPX\slash{}CPB and then also include a screenshot (entire computer screen) of your Mu code showing the line of code used to read the button press and blink the LED. - 20\%%
\item{}Appendix A - Video Details - Film yourself and your entire screen with the Serial Monitor open in Mu while you press and hold a button on the CPX\slash{}CPB and show the LED on your breadboard blink while you hold the button down%
\item{}Appendix B - Code used to take data and code used to plot in plain text (no screenshots) - Pass\slash{}Fail%
\item{}Appendix C - Screenshot of Excel or Sheets if applicable - Pass\slash{}Fail%
\item{}\alert{If you are asked to fix something in the Appendices such as missing content in your video or screenshots of your code, you will lose 10\% of your grade everytime you submit an update after the deadline.}%
\end{enumerate}
\end{introduction}%
\end{sectionptx}
\end{chapterptx}
%
%
\typeout{************************************************}
\typeout{Chapter 9 Logging Data}
\typeout{************************************************}
%
\begin{chapterptx}{Chapter}{Logging Data}{}{Logging Data}{}{}{daq}
\renewcommand*{\chaptername}{Chapter}
\begin{introduction}{}%
\begin{listptx}{List}{\textbf{Parts List}}{daq-2-1}{}%
%
\begin{enumerate}
\item{}Laptop%
\item{}CPX\slash{}CPB%
\item{}USB Cable%
\end{enumerate}
\end{listptx}%
\begin{listptx}{List}{\textbf{Learning Objectives}}{daq-2-2}{}%
%
\begin{enumerate}
\item{}Record real time measurements from the CircuitPlayground using the Serial monitor%
\item{}Learn how to use the typing module to type data directly into a spreadsheet or text file%
\item{}Learn how to log data on the on board memory of the CircuitPlayground%
\end{enumerate}
\end{listptx}%
Taking data is the core of any instrumentation project. Data Acquisition Systems (DAQ) come in all shapes and sizes. Believe it or not the CPX can be used as a crude and cheap DAQ. The CPX can easily take temperature data and monitor the temperature in a greenhouse or take humidity readings of a plant to monitor soil content. Before we learn about the different sensors on board the CPX, we want to make sure we can store that data later rather than just having it spit data out via the serial monitor. In this chapter you're going to make a simple script that prints button presses and then log that in multiple different ways. The button presses are just a simple example but the same methods can be used to log any data you want from the CPX.%
\end{introduction}%
%
%
\typeout{************************************************}
\typeout{Section 9.1 Initial Setup}
\typeout{************************************************}
%
\begin{sectionptx}{Section}{Initial Setup}{}{Initial Setup}{}{}{start}
For starters, let's get the CPX to print button presses. The code I'm using is shown below and can also be found on \href{https://github.com/cmontalvo251/Microcontrollers/blob/master/Circuit_Playground/CircuitPython/Data_Logging/button_presses.py}{\linkhilite{Github}}\hyperlink{Github}{\linkhilite{[{\xreffont 33}]}}.%
\begin{figureptx}{Figure}{CircuitPython code that prints button presses}{start-3}{}%
\begin{image}{0.135}{0.73}{0.135}{}%
\includegraphics[width=\linewidth]{external/button_presses.png}
\end{image}%
\tcblower
\end{figureptx}%
The code is pretty similar to what I had in the past. I import board, digitalio, and time. I create a buttonA object using the digitalio library to record button presses. I then enter into a while loop print the buttonA.value. The difference here is that I use the int() function to convert the buttonA.value to an integer. The reason why I do this is because buttonA.value is a boolean. It is either True or False. An integer though is a number and thus a value of False is 0 and True is 1. If you open the serial monitor and push the A button down a few times you'll see some zeros and 1's.%
\begin{figureptx}{Figure}{Serial monitor showing button presses in Mu}{start-5}{}%
\begin{image}{0.135}{0.73}{0.135}{}%
\includegraphics[width=\linewidth]{external/MuButtonPress.png}
\end{image}%
\tcblower
\end{figureptx}%
Mu also has a really neat builtin plotter. You'll see next to the Serial button there is a button called Plotter. If you click that button now nothing will pop up on the screen. Unfortunately in order to plot using the Plotter you need to modify the print() statement to this:%
\begin{preformatted}
print((int(buttonA.value),))
\end{preformatted}
Notice the extra parentheses and the comma. Now if you click Plotter you'll see something like this. You'll notice that the print statement now has commas in it and the Plotter is recording button presses.%
\begin{figureptx}{Figure}{Plotter open showing button presses in Mu}{start-9}{}%
\begin{image}{0.135}{0.73}{0.135}{}%
\includegraphics[width=\linewidth]{external/MuPlotter.png}
\end{image}%
\tcblower
\end{figureptx}%
The problem with this is we still can't save the recorded data anywhere. Before we get into saving data let's first edit the print statement again to get rid of the Plotter by removing the extra parentheses and add time.monotonic() that way we can keep track of when a button was pressed. My print statement looks like this now:%
\begin{preformatted}
print(time.monotonic(),int(buttonA.value))
\end{preformatted}
Looking at the serial monitor now you'll see that time is being printed alongside the button presses.%
\begin{figureptx}{Figure}{Serial monitor showing button presses and time in Mu}{start-13}{}%
\begin{image}{0.135}{0.73}{0.135}{}%
\includegraphics[width=\linewidth]{external/MuPlotterTime.png}
\end{image}%
\tcblower
\end{figureptx}%
Now we are in a position where we can record some data and save it to our computer. There are 4 ways to record data. I call the first, Method1 and you basically just copy and paste from the serial monitor, Method2 where you have the CPX\slash{}CPB type data into a spreadsheet and Method3 where you log data internally onto the CPX\slash{}CPB itself. The 4th method called Method4 utilizes the Bluetooth Module. Since that has it's own issues there is a completely separate section on how to explain Bluetooth (See \hyperref[bluetooth]{\linkhilite{Chapter~{\xreffont\ref{bluetooth}}}}). Note you can only do Bluetooth if you have the Circuit Playground Bluefruit (CPB).%
\end{sectionptx}
%
%
\typeout{************************************************}
\typeout{Section 9.2 Method 1 - Copying Serial Monitor Data}
\typeout{************************************************}
%
\begin{sectionptx}{Section}{Method 1 - Copying Serial Monitor Data}{}{Method 1 - Copying Serial Monitor Data}{}{}{daq-4}
If you open up the serial monitor you can see the data output. If you unplug the CPX while it's taking data the code will stop. \alert{Note: In newer versions, unplugging your CPX will result in a loss of data. If this happens try pressing CTRL+C after you click the REPL window.} With the code stopped you can select all the data in the Serial monitor and then copy and paste the data into a text file on your computer. You can actually copy this into a new file on Thonny and just save it as a \textasteriskcentered{}.txt file. Once you have the data in a text file you can proceed to plotting on your desktop which I discuss in the last section. Here's some example data in Gedit which is a simple text editing program.%
\begin{figureptx}{Figure}{Example button data in a text editor}{daq-4-3}{}%
\begin{image}{0.135}{0.73}{0.135}{}%
\includegraphics[width=\linewidth]{external/Gedit_Data.png}
\end{image}%
\tcblower
\end{figureptx}%
\end{sectionptx}
%
%
\typeout{************************************************}
\typeout{Section 9.3 Method 2 - Automatically Populate a Spreadsheet}
\typeout{************************************************}
%
\begin{sectionptx}{Section}{Method 2 - Automatically Populate a Spreadsheet}{}{Method 2 - Automatically Populate a Spreadsheet}{}{}{daq-5}
The downside with the above method of course is if you have a ton of data to record you could lose the data or run into a massive copy and paste issue. The second option is to use this module called keyboard which takes control of your keyboard on your desktop computer and actively types your data into a spreadsheet. The code is very extensive but I'll include the simple one here so we can discuss it. Below are the first 30 lines of code. The first 6 lines of code are just comments since I heavily adopted this code from the \href{https://learn.adafruit.com/make-it-a-keyboard/circuitpython}{\linkhilite{Adafruit Learn System}}\hyperlink{Adafruit}{\linkhilite{[{\xreffont 32}]}}. \href{https://github.com/cmontalvo251/Microcontrollers/blob/master/Circuit_Playground/CircuitPython/Data_Logging/typing/record_button_presses_typing.py}{\linkhilite{My version of this code}} can be found on my Github \hyperlink{Github}{\linkhilite{[{\xreffont 33}]}}. Lines 8 - 14 are import commands as we've seen previously. The regular import modules board, time and digitalio are imported but we are also importing the Keyboard module so that the CPX can takeover our keyboard. In order to import that module properly you need to \href{https://circuitpython.org}{\linkhilite{install modules}} onto you \emph{lib} folder (See \hyperref[troubleshooting]{\linkhilite{Chapter~{\xreffont\ref{troubleshooting}}}}) \hyperlink{CircuitPython}{\linkhilite{[{\xreffont 30}]}}. Lines 16-22 create two buttons. First we create buttonA attached to pin D4 and then a switch attached to pin D7. If you look on the CPX there is a switch labeled D7. Before you copy this code onto the CPX make sure you move the switch towards the ear looking symbol. Lines 26-28 created the keyboard object. We are going to call it layout for this example code.%
\begin{figureptx}{Figure}{Snippet of Typing code in CircuitPython}{daq-5-3}{}%
\begin{image}{0.135}{0.73}{0.135}{}%
\includegraphics[width=\linewidth]{external/Typing1.png}
\end{image}%
\tcblower
\end{figureptx}%
The next 30 lines are shown below. Lines 32-35 define a function. Functions in CircuitPython have a pretty standard structure. The keyword def is used to denote that the next line is a definition for a function. The name of the function is \emph{slow\textunderscore{}write()}. The input to the function is string which ironically enough is a string object. Line 33-35 define what the function does. Line 33 sets up a for loop where the code loops through each character in the string. Everytime it gets to a new character it will use your keyboard to type that character using the layout.write(c) command. The time.sleep(0.02) is just to slow down the keyboard so your computer can keep up. That function is defined above the standard while True: statement on line 37 but is called on line 42. You'll see there is a \emph{slow\textunderscore{}write(output)} on line 42. In this case output is a string and it's sent to the function \emph{slow\textunderscore{}write()}. So in this case we have a function that can write a string so we just need to take data and then write it using our keyboard. Line 38 is an if statement that will only be true if the switch on pin D7 is pushed towards the music note on the CPX. If the switch is not thrown the code will move to the else statement on line 52 and tell the user that you need to flip the switch. If the switch is thrown line 40 will take data for us. First it will record the time.monotonic() and store it as a floating point number using the \%0.1f designation which means that it will store 1 decimal as a \%floating point number for f.%
\begin{figureptx}{Figure}{Remainder of Typing code in CircuitPython}{daq-5-5}{}%
\begin{image}{0.135}{0.73}{0.135}{}%
\includegraphics[width=\linewidth]{external/Typing2.png}
\end{image}%
\tcblower
\end{figureptx}%
The second number in the string is an integer or a base 10 (decimal) integer designated by the \%d part of the format. The integer is int(buttonA.value). You'll see a \textbackslash{}t in between the formatted numbers which is a tab. The tab is there to tab between cells in a spreadsheet. Line 41 will print the output string to the Serial monitor and it will also type the contents of the string. Very important here. When you flip the switch on the CPX your keyboard will start typing in whatever active window is selected. If you don't have a spreadsheet opened and active (selected), the keyboard will just begin typing in whatever window is open. Make sure you have a spreadsheet program open and ready to go. Lines 44-51 tell the keyboard to hit the \emph{DOWN\textunderscore{}ARROW} on your keyboard to move to the next row and the \emph{LEFT\textunderscore{}ARROW} twice to move back to the first column. Line 55 is a sleep to only log data once a second. I ran this code for a bit and had it type into LibreOffice Calc which is a free spreadsheet program. Google Sheets or Microsoft Excel will also work just fine.%
\begin{figureptx}{Figure}{Example data in a spreadsheet from Typing code}{daq-5-7}{}%
\begin{image}{0.135}{0.73}{0.135}{}%
\includegraphics[width=\linewidth]{external/Typing3.png}
\end{image}%
\tcblower
\end{figureptx}%
You'll see that the first column is time with 1 decimal point and the second column is the button press values. At this point you must click \emph{Save As...} and save the document as a CSV which stands for Comma Separated Value. Once you have the file saved you can proceed to plotting on your Desktop.%
\end{sectionptx}
%
%
\typeout{************************************************}
\typeout{Section 9.4 Method 3 - Logging Data Directly to on board memory}
\typeout{************************************************}
%
\begin{sectionptx}{Section}{Method 3 - Logging Data Directly to on board memory}{}{Method 3 - Logging Data Directly to on board memory}{}{}{daq-6}
The problem with the above 2 methods is that you need a laptop to log data in the field. It would be nice if you could use the optional battery pack and just have the CPX log data on the CPX itself. This is the most complex way but in my opinion the best way. In order to get this to work you need to allow the drive on the CPX to have read\slash{}write permissions. This requires you to load a piece of software called boot.py and put it on the CPX. I have this \href{https://github.com/cmontalvo251/Microcontrollers/blob/master/Circuit_Playground/CircuitPython/Data_Logging/boot.py}{\linkhilite{software on my Github}}\hyperlink{Github}{\linkhilite{[{\xreffont 33}]}}. The software is shown below. The first 10 lines are probably very familiar. Import some modules and then create a switch object. Line 13 is where all of the storage permissions are changed. If the flip is switched towards the A button, the storage module is used to allow you to write to the CPX. The problem here is that if you do this, you won't be able to edit code. I'll explain the procedure here in a minute. As always, the relevant Adafruit tutorial is on the \href{https://learn.adafruit.com/adafruit-circuit-playground-express/circuitpython-storage}{\linkhilite{Adafruit Learn System}}\hyperlink{Adafruit}{\linkhilite{[{\xreffont 32}]}} if you want to read more about it. Again make sure you store this file onto the CIRCUITPY drive and save it as boot.py%
\begin{figureptx}{Figure}{boot.py script in CircuitPython}{daq-6-3}{}%
\begin{image}{0.135}{0.73}{0.135}{}%
\includegraphics[width=\linewidth]{external/boot.png}
\end{image}%
\tcblower
\end{figureptx}%
In addition to storing the file boot.py you'll need to edit your main.py script to only log data when the switch is moved towards the B button. The software to record button presses on disk is shown below and as always \href{https://github.com/cmontalvo251/Microcontrollers/blob/master/Circuit_Playground/CircuitPython/Data_Logging/write_button_presses_disk.py}{\linkhilite{on my Github}}\hyperlink{Github}{\linkhilite{[{\xreffont 33}]}}. In this software we again see the standard commands. Lines 1-3 import all the modules we need and then 5-15 create a switch, a button and an LED. In this case we're using the LED soldered to the board. Line 17-20 check to see if the user has flipped the switch. If the switch is False the storage module on boot.py will allow the drive to act like a data logger and it will open a file called \emph{Test\textunderscore{}Data.txt} for writing \emph{(\textquotedblleft{}w\textquotedblright{})}. If the switch is True then the user will be notified that the file has not been opened for writing. Lines 22 through 33 include the infinite while loop. Line 23 turns the LED on and line 24 prints out the current time and the button value in integer form. If the switch value is False the program will create an output string by converting all numbers to strings using the str function. Notice that there is a \emph{str(\textquotedblleft{}\textbackslash{}textbackslash n\textquotedblright{})} at the end of the output variable which tells the computer to write a new line of data to the file. Lines 28 and 29 write the output to the file from line 18 and then flush the output which means the CPX waits for the data to be fully written before moving on. It also turns the LED off so we know the CPX took data even when we aren't looking at the Serial monitor. If the switch value is true it means that the we never opened the data file and thus we tell the user we aren't logging data and it's time to flip the switch and hit reset. \alert{NOTE: The figure below is from an old version of the code. The newest version produces the same outcome. The only difference is that the D13 LED blinks when the code is running and the first neopixel toggles 3 different colors when the system is logging data. This allows for extra user information when operating Method 3 without a computer and the serial monitor. Note that these additions require the use of neopixel.mpy module.}%
\begin{figureptx}{Figure}{Snippet of writing data to a text file in CircuitPython}{daq-6-5}{}%
\begin{image}{0.135}{0.73}{0.135}{}%
\includegraphics[width=\linewidth]{external/method3_1.png}
\end{image}%
\tcblower
\end{figureptx}%
So here is the flow of what you want to do for method 3.%
%
\begin{enumerate}
\item{}Unplug the CPX%
\item{}Flip the switch towards the A button.%
\item{}Plug in the CPX and save the boot.py and main.py files. Remember you can only save CircuitPython scripts when the switch is flipped towards the A button.%
\item{}When you are ready to start recording data, flip the switch towards the B button. If you're looking at the Serial monitor, the software will throw an error. Just ignore it and hit the reset button. When your computer recognizes the CPX you can turn the Serial monitor on and off.%
\item{}When you are done taking data simply slide the switch over towards the A button and hit reset again. This is what my Serial monitor looks like when I do this. You'll see that I was writing to disk for like 25 seconds and then I flipped the switch back towards the A button.%
\end{enumerate}
\begin{figureptx}{Figure}{Serial monitor open in Mu showing Method 3 for data logging}{daq-6-8}{}%
\begin{image}{0.135}{0.73}{0.135}{}%
\includegraphics[width=\linewidth]{external/method3_2.png}
\end{image}%
\tcblower
\end{figureptx}%
With the switch flipped and data taken, open your folder manager and take a look at the CIRCUITPY drive. This is what mine looks like. You'll see I have two CircuitPython files and a file \emph{Test\textunderscore{}Data.txt} with all my data in it.%
\begin{figureptx}{Figure}{Text file shown in file manager on computer}{daq-6-10}{}%
\begin{image}{0.135}{0.73}{0.135}{}%
\includegraphics[width=\linewidth]{external/method3_3.png}
\end{image}%
\tcblower
\end{figureptx}%
If you open the \emph{Test\textunderscore{}Data.txt} file you will hopefully see data in it.%
\begin{figureptx}{Figure}{Example text file from Method 3}{daq-6-12}{}%
\begin{image}{0.135}{0.73}{0.135}{}%
\includegraphics[width=\linewidth]{external/Gedit_Data.png}
\end{image}%
\tcblower
\end{figureptx}%
At this point you can copy this text file over to your desktop computer and proceed to the plotting portion. Ok so let's recap method 3 one more time.%
\begin{enumerate}
\item{}Unplug CPX (or remove power)%
\item{}Slide switch to A%
\item{}Plug in CPX (or provide battery power) and wait for system to fully boot up%
\item{}Slide switch to B%
\item{}Hit Reset and wait for system to fully boot up%
\item{}Take data for however long you want%
\item{}Slide switch to A when you're ready to stop taking data%
\item{}Hit Reset and wait for system to fully boot up%
\item{}Remove power if you're on battery power%
\item{}Plug CPX into computer if not already connected%
\item{}Transfer data file to computer%
\end{enumerate}
%
\end{sectionptx}
%
%
\typeout{************************************************}
\typeout{Section 9.5 Method 4 - Logging Data on a Cell Phone using Bluetooth (CPB Only)}
\typeout{************************************************}
%
\begin{sectionptx}{Section}{Method 4 - Logging Data on a Cell Phone using Bluetooth (CPB Only)}{}{Method 4 - Logging Data on a Cell Phone using Bluetooth (CPB Only)}{}{}{daq-7}
As mentioned in the introduction it's possible to have the Circuit Playground send data wirelessly to a cell phone using Bluetooth provided you have Bluetooth setup and a smart phone with the Adafruit Connect App. Bluetooth is explained in detail in its own chapter (See \hyperref[bluetooth]{\linkhilite{Chapter~{\xreffont\ref{bluetooth}}}}). Method 4 is a valid form of logging data it just requires a cell phone to be powered the entire time and it must be within 30 feet of the Circuit Playground at all times. This also only works on the CPB since the CPX does not have a Bluetooth transmitter.%
\end{sectionptx}
%
%
\typeout{************************************************}
\typeout{Section 9.6 Method 5 - Logging Data to an SD Card}
\typeout{************************************************}
%
\begin{sectionptx}{Section}{Method 5 - Logging Data to an SD Card}{}{Method 5 - Logging Data to an SD Card}{}{}{daq-8}
Another to log data is using a standalone \href{https://www.adafruit.com/product/254}{\linkhilite{micro sd card breakout board}} (See \hyperref[sdcard]{\linkhilite{Chapter~{\xreffont\ref{sdcard}}}})\hyperlink{AdafruitMAIN}{\linkhilite{[{\xreffont 34}]}}. Method 1 and 2 are simple to use but they require a laptop and Method 4 can only be done with a CPB. Method 3 is light-weight in the sense that you don't need any extra components but this method can be somewhat unreliable. It's possible to brick the CPX which requires formatting (See \hyperref[troubleshooting]{\linkhilite{Chapter~{\xreffont\ref{troubleshooting}}}}) or data corruption in general to the file. I've seen so many students not follow Method 3 correctly and either not get data or corrupt their entire CPX which requires them to reformat. Although some students have been successful, the complexity of getting \emph{boot.py} to work properly requires great care and this is where Method 5 excels. In this method rather than logging data directly to the CPX\slash{}CPB and potentially corrupting the internal memory, the data is logged to an external SD card. After logging data the card can be removed and uploaded to a computer using any SD card to USB converter.%
\end{sectionptx}
%
%
\typeout{************************************************}
\typeout{Section 9.7 Method 6 - Logging Data via Serial}
\typeout{************************************************}
%
\begin{sectionptx}{Section}{Method 6 - Logging Data via Serial}{}{Method 6 - Logging Data via Serial}{}{}{daq-9}
Yet another method to log data is using the same serial communication that is used in Mu (See \hyperref[serial]{\linkhilite{Chapter~{\xreffont\ref{serial}}}}). In this case you need your laptop just like method 1 and 2 but you forgo having to copy and paste the Serial monitor yourself. In this method you would have a program on your computer that automatically reads the serial output from the CPX and saves that to a text file. It reduces the complexity of the CircuitPython code and allows you to focus on the experiment while your computer handles logging data.%
\end{sectionptx}
%
%
\typeout{************************************************}
\typeout{Section 9.8 Plotting Logged Data}
\typeout{************************************************}
%
\begin{sectionptx}{Section}{Plotting Logged Data}{}{Plotting Logged Data}{}{}{daq-10}
Alright so there you have it. I have explained 4 methods to datalogging. Here are the methods again in summary.%
\begin{enumerate}
\item{}Print data to Serial and copy and paste%
\item{}Use the Keyboard module to save data to a spreadsheet%
\item{}Access the storage of your CPX and write data to a text file on the CPX%
\item{}Send data wirelessly to a Cell Phone using Bluetooth - CPB Only (See \hyperref[bluetooth]{\linkhilite{Chapter~{\xreffont\ref{bluetooth}}}})%
\end{enumerate}
%
\par
All methods will work but some will obviously have their pros and cons. I suggest you get comfortable with 1 method and use that for the remainder of the semester. Whatever option you choose though will provide you with a data file that you can open on your desktop computer to plot. The simplest way to import data is by using the loadtxt function from the module numpy. Here is some very simple code to plot data from a text file.%
\par
When you plot make sure your \emph{Test\textunderscore{}Data.txt} file is in the same folder as your plotting script in Thonny or Spyder. Here's my example code (this code is not on Github but you only need 3 or 4 lines of code to plot).%
\begin{figureptx}{Figure}{Plotted data shown in Thonny}{daq-10-5}{}%
\begin{image}{0.135}{0.73}{0.135}{}%
\includegraphics[width=\linewidth]{external/plotdata.png}
\end{image}%
\tcblower
\end{figureptx}%
In this example lines 1 and 2 import numpy and matplotlib. Line 4 imports data from the \emph{Test\textunderscore{}Data.txt} file and then 6 and 7 save the first and second columns into time and button. The remaining lines plot the data and create x and y labels as well as a grid. Hopefully now you are well versed in taking data and plotting on your desktop computer.%
\end{sectionptx}
%
%
\typeout{************************************************}
\typeout{Section 9.9 Extra Help}
\typeout{************************************************}
%
\begin{sectionptx}{Section}{Extra Help}{}{Extra Help}{}{}{daq-11}
This is a pretty hard project to do with multiple different methods. After you've read through this document I suggest you watch a youtube I created on \href{https://youtu.be/yX0zeIfn2j4}{\linkhilite{Logging Data with the Circuit Playground Express}}. I have also made video where I \href{https://www.youtube.com/watch?v=Vc7_KxRnORI}{\linkhilite{log temperature and accelerometer data using on board memory with the CircuitPlayground Bluefruit}}. Finally, I also have a \href{https://www.youtube.com/watch?v=tJOz-ty-2ec\&list=PL_D7_GvGz-v1RsDs_OdNW65qRjEjmpfQx\&index=12}{\linkhilite{Youtube video explaining how to plot a text file}}\hyperlink{youtube}{\linkhilite{[{\xreffont 50}]}}.%
\end{sectionptx}
%
%
\typeout{************************************************}
\typeout{Section 9.10 Assignment}
\typeout{************************************************}
%
\begin{sectionptx}{Section}{Assignment}{}{Assignment}{}{}{daq-12}
\begin{introduction}{}%
For this project you must use method 1, 2, 3 or 4 to save time and button presses to a text file. You must then plot the button presses as a function of time. Remember to add x and y labels to all figures.%
%
\begin{enumerate}
\item{}Title page with name,title and date - 5\%%
\item{}Introduction - 5\%%
\item{}All figures must have appropriate annotations, figure captions, figure labels and a paragraph of supporting text - 10\%%
\item{}Include a snippet (5 lines) of your data file - 40\%%
\item{}Include a plot of your button presses with time on the x-axis and button presses on the y-axis (no screenshots) - 40\%%
\item{}Appendix A - Video Details - File yourself and your entire screen pressing the button and showing the number change on the screen, then show yourself plotting your data on your computer - Pass\slash{}Fail%
\item{}Appendix B - Code used to take data and code used to plot in plain text (no screenshots) - Pass\slash{}Fail%
\item{}Appendix C - Screenshot of Excel or Sheets if applicable - Pass\slash{}Fail%
\item{}\alert{If you are asked to fix something in the Appendices such as missing content in your video or screenshots of your code, you will lose 10\% of your grade everytime you submit an update after the deadline.}%
\end{enumerate}
\end{introduction}%
\end{sectionptx}
\end{chapterptx}
%
%
\typeout{************************************************}
\typeout{Chapter 10 Bluetooth Logging - Method 4}
\typeout{************************************************}
%
\begin{chapterptx}{Chapter}{Bluetooth Logging - Method 4}{}{Bluetooth Logging - Method 4}{}{}{bluetooth}
\renewcommand*{\chaptername}{Chapter}
\begin{introduction}{}%
\begin{listptx}{List}{\textbf{Parts List}}{bluetooth-2-1}{}%
%
\begin{enumerate}
\item{}Smart Phone%
\item{}Adafruit BLE Connect App (\href{https://play.google.com/store/apps/details?id=com.adafruit.bluefruit.le.connect}{\linkhilite{Play Store}}\slash{}\href{https://apps.apple.com/us/app/bluefruit-connect/id830125974}{\linkhilite{App Store}})\hyperlink{AdafruitBLE}{\linkhilite{[{\xreffont 36}]}}%
\item{}CircuitPlayground Bluefruit%
\item{}USB Cable%
\item{}Laptop%
\end{enumerate}
\end{listptx}%
\begin{listptx}{List}{\textbf{Learning Objectives}}{bluetooth-2-2}{}%
%
\begin{enumerate}
\item{}Understand the bluetooth module on the CircuitPlayground Bluefruit%
\item{}Learn how to send data via the Bluefruit to your smart phone%
\item{}Understand how to plot data sent via UART%
\end{enumerate}
\end{listptx}%
I mentioned in the DAQ project that Method 4 can be used to log data by using bluetooth to send data from your phone to your Circuitplayground Bluefruit (CPB) and vice versa. Once the data is on your phone you can export the data to a text file. That basically means you can use the bluetooth module as another method to save data from the CPB. There is a lot you can do with bluetooth but for starters all code required for this module is on \href{https://github.com/cmontalvo251/Microcontrollers/tree/master/Circuit_Playground/CircuitPython/CPB_DataLogger}{\linkhilite{Github}}. This chapter walks through how to write some example bluetooth code and get the app to receive the data on your smart phone.%
\end{introduction}%
%
%
\typeout{************************************************}
\typeout{Section 10.1 Setup}
\typeout{************************************************}
%
\begin{sectionptx}{Section}{Setup}{}{Setup}{}{}{bluetooth-3}
First we're going to run the \href{https://github.com/cmontalvo251/Microcontrollers/blob/master/Circuit_Playground/CircuitPython/CPB_DataLogger/bluefruit_uart_button.py}{\linkhilite{bluetooth\textunderscore{}uart\textunderscore{}button.py}} script which sends button data to your smart phone via something called UART which is a type of serial communication. It's beyond the scope of this lesson but serial is digital as opposed to analog which is done using the AnalogIn functions (See \hyperref[voltage]{\linkhilite{Chapter~{\xreffont\ref{voltage}}}}) \hyperlink{Github}{\linkhilite{[{\xreffont 33}]}}.%
\begin{figureptx}{Figure}{Example Bluetooth code in CircuitPython}{bluetooth-3-3}{}%
\begin{image}{0.135}{0.73}{0.135}{}%
\includegraphics[width=\linewidth]{external/bluetooth_code.png}
\end{image}%
\tcblower
\end{figureptx}%
Lines 3-11 import a ton of modules. You'll recognize many of them like analogio, and time but the new ones are the ones that say ble. These are the bluetooth modules required for the CPB. Lines 21-24 setup the button so we can log the button via bluetooth and Lines 14-17 kick of the BLERadio object and the UARTService() to send data. Line 25 also grabs the current time before the infinite while loop that way the timer starts closer to zero.%
\begin{figureptx}{Figure}{Remainder of Bluetooth code in CircuitPython}{bluetooth-3-5}{}%
\begin{image}{0.135}{0.73}{0.135}{}%
\includegraphics[width=\linewidth]{external/bluetooth_code1.png}
\end{image}%
\tcblower
\end{figureptx}%
The code above is the infinite while loop which actually contains 2 while loops. Lines 28-33 prints the name of your bluetooth sensor and then starts advertising bluetooth to whoever is listening. It will then enter a while loop from 32-33 until bluetooth is connected. Once bluetooth is connected it will enter into the second while loop from line 39-52. In those lines 40-46 is responsible for taking all the necessary measurements and printing them to the serial monitor in Mu. In this case it's only printing the current time and the value of the button as an integer. \emph{buttonA.value} is either True or False and the \emph{int} function converts that to a 0 or a 1. Line 49 then sends the data over bluetooth using the UART server. You'll notice in this case the code is sending t,b by using the format variable and the 2 empty  brackets. If you want to send more data you need to add more empty brackets and more variables to the format function. When you first save this script your CPB will not be connected and enter into an infinite while loop where it waits for your smart phone to connect. If you open your smart phone and open the Bluefruit Connect App the following screen will pop up.%
\begin{figureptx}{Figure}{Screenshot of Bluetooth devices on a smart phone}{bluetooth-3-7}{}%
\begin{image}{0.35}{0.3}{0.35}{}%
\includegraphics[width=\linewidth]{external/phoneapp1.png}
\end{image}%
\tcblower
\end{figureptx}%
In this case there are numerous different bluetooth modules can be seen but the one you need to click is the one that says CIRCUITPYf8e8. You will have a different code after CIRCUITPY and you can figure out what your 4 digit code is by making sure you have the \emph{print('Look for',ble.name)} in your code. Once you do that the CPB will begin sending time and the button press to your smart phone.%
\begin{figureptx}{Figure}{Interfaces in Bluetooth App on smart phone}{bluetooth-3-9}{}%
\begin{image}{0.35}{0.3}{0.35}{}%
\includegraphics[width=\linewidth]{external/phoneapp2.png}
\end{image}%
\tcblower
\end{figureptx}%
There are numerous items you can click. The Controller is very fun for creating a remote control robot but we're only going to go over the UART and Plotter tabs. If you click the plotter tab you will be greeted with a live screen of the data being sent.%
\begin{figureptx}{Figure}{Example plotter in Bluetooth App on smart phone}{bluetooth-3-11}{}%
\begin{image}{0.35}{0.3}{0.35}{}%
\includegraphics[width=\linewidth]{external/phoneapp3.png}
\end{image}%
\tcblower
\end{figureptx}%
In the photo above you can see three data streams that is coming directly from the CPB. The red line is time and the blue line is the button value. Notice the blue line goes from 0 to 1 which means I pressed the button a few times. The red line is always increasing which kind of messes up the plotter so you can always go back to your code in Mu and just send your data. This is great for live demonstrations and for debugging if you need to see data from an experiment and you don't have access to a laptop with Mu. If you hit the back arrow and then click UART you will see the raw data come in as text.%
\begin{figureptx}{Figure}{Example UART data in Bluetooth App on smart phone}{bluetooth-3-13}{}%
\begin{image}{0.35}{0.3}{0.35}{}%
\includegraphics[width=\linewidth]{external/phoneapp4.png}
\end{image}%
\tcblower
\end{figureptx}%
Again here you can see the 3 data streams separated by commas. The very neat thing with the UART tab is that you can click the three vertical dots in the upper right hand corner and click export to TXT. The easiest thing for me was to export the data to google drive and then download the data to my computer. Once I downloaded the TXT file to my computer and opened it the data file looked like this.%
\begin{figureptx}{Figure}{Example data exported to CSV file}{bluetooth-3-15}{}%
\begin{image}{0.1}{0.8}{0.1}{}%
\includegraphics[width=\linewidth]{external/csv_fileapp.png}
\end{image}%
\tcblower
\end{figureptx}%
If you export the file as a CSV the data file will look completely different and it's much more complicated to plot. If you export the data as a TXT file you just need to use the np.loadtxt command to read in the data. Note you might have commas in your data file. If there are commas just use the CTRL+H command and replace all commas with spaces or use the np.loadtxt('buttonble.txt',delimiter=',') command. Plotting your button presses should be as simple as the previous lab thus plotting the button is left as an exercise to the reader.%
\end{sectionptx}
%
%
\typeout{************************************************}
\typeout{Section 10.2 Extra Help}
\typeout{************************************************}
%
\begin{sectionptx}{Section}{Extra Help}{}{Extra Help}{}{}{bluetooth-4}
You might find plotting data via bluetooth to be rather difficult and it was pretty difficult for me until I learned that you can export data as a txt file rather than a csv file. Before I learned how to do that I put together a \href{https://www.youtube.com/playlist?list=PL_D7_GvGz-v0-U3JACRMgldvqQqn2mje9}{\linkhilite{4 part series}} describing everything in this module. Worst case you can just watch the \href{https://www.youtube.com/watch?v=9EHNFdVX9O8\&list=PL_D7_GvGz-v0-U3JACRMgldvqQqn2mje9\&index=3}{\linkhilite{third video in the series}}. The video is 30 minutes but the first 8 minutes goes through setting up the bluetooth module and the rest of the video is just on plotting the exported csv data which took me some time. Note that exporting data as a txt file is the preferred method as parsing the file is way easier \hyperlink{youtube}{\linkhilite{[{\xreffont 50}]}}.%
\end{sectionptx}
%
%
\typeout{************************************************}
\typeout{Section 10.3 Assignment}
\typeout{************************************************}
%
\begin{sectionptx}{Section}{Assignment}{}{Assignment}{}{}{bluetooth-5}
\begin{introduction}{}%
This project is similar to the Data Acquisition project only you must use method 4 to save time and button presses to a text file. You must then plot the button presses as a function of time. Remember to add x and y labels to all figures.%
%
\begin{enumerate}
\item{}Title page with name,title and date - 5\%%
\item{}Introduction - 5\%%
\item{}All figures must have appropriate annotations, figure captions, figure labels and a paragraph of supporting text - 10\%%
\item{}Include a screenshot (entire screen) from your phone showing the raw UART output - 20\%%
\item{}Include a snippet (5 lines) of your data file - 20\%%
\item{}Include a plot of your button presses with time on the x-axis and button presses on the y-axis (no screenshots) - 40\%%
\item{}Appendix A - Video Details - Film yourself and your entire screen showing data broadcast to your phone while you press and release the button, then show yourself plotting the data on your computer - Pass\slash{}Fail%
\item{}Appendix B - Code used to take data and code used to plot in plain text (no screenshots) - Pass\slash{}Fail%
\item{}Appendix C - Screenshot of Excel or Sheets if applicable - Pass\slash{}Fail%
\item{}\alert{If you are asked to fix something in the Appendices such as missing content in your video or screenshots of your code, you will lose 10\% of your grade everytime you submit an update after the deadline.}%
\end{enumerate}
\end{introduction}%
\end{sectionptx}
\end{chapterptx}
%
%
\typeout{************************************************}
\typeout{Chapter 11 SD Card Logging - Method 5}
\typeout{************************************************}
%
\begin{chapterptx}{Chapter}{SD Card Logging - Method 5}{}{SD Card Logging - Method 5}{}{}{sdcard}
\renewcommand*{\chaptername}{Chapter}
\begin{introduction}{}%
\begin{listptx}{List}{\textbf{Parts List}}{sdcard-2-1}{}%
%
\begin{enumerate}
\item{}Laptop%
\item{}USB Cable%
\item{}Micro SD Card%
\item{}Micro SD Card USB Reader%
\item{}\href{https://www.adafruit.com/product/254}{\linkhilite{Micro SD Card Breakout Board}}%
\item{}Arduino or CPX\slash{}CPB%
\end{enumerate}
\end{listptx}%
\begin{listptx}{List}{\textbf{Learning Objectives}}{sdcard-2-2}{}%
%
\begin{enumerate}
\item{}Understand how to use a micro SD card to log data from a microcontroller%
\item{}Learn to install libraries for the Arduino%
\item{}Learn yet another way to use a microcontroller as a data acquisition system%
\end{enumerate}
\end{listptx}%
This chapter explores writing data directly to a micro SD card. This is helpful because when the datalogger is done taking data, it is easy to power down the system, take the SD card out and pop it into a laptop for post-processing. For the CPX\slash{}CPB setp it helps because you don't have to have data writing to the internal storage which can potentially brick the system. You have code and data storage in two different places. Setting up this system requires a few steps for both CPX\slash{}CPB and Arduino microcontrollers. The Arduino UNO is capable of writing to the micro SD card but the only difference is that the pins setup is different. There is a very indepth \href{https://learn.adafruit.com/adafruit-micro-sd-breakout-board-card-tutorial}{\linkhilite{SD card tutorial on Adafruit}} that can be used in conjunction with this tutorial if extra help is needed \hyperlink{AdafruitMAIN}{\linkhilite{[{\xreffont 34}]}}.%
\begin{figureptx}{Figure}{Micro SD card Breakout Board without header pins soldered}{sdcard-2-4}{}%
\begin{image}{0.135}{0.73}{0.135}{}%
\includegraphics[width=\linewidth]{external/microSDcard.jpg}
\end{image}%
\tcblower
\end{figureptx}%
No matter what for both microcontrollers you need to have a micro SD card formatted to FAT32. In Linux that means you need to navigate to the Disk utility app and format the disk. On Windows or Mac you need to use a similar program to format the disk.%
\begin{figureptx}{Figure}{Using the Disk Utility to Format a Micro SD Card on Linux}{sdcard-2-6}{}%
\begin{image}{0.135}{0.73}{0.135}{}%
\includegraphics[width=\linewidth]{external/SDCARD_FORMAT.png}
\end{image}%
\tcblower
\end{figureptx}%
Also, when you open the packaging of this breakout board you'll notice that the header pins are missing. First you'll need to cut a row of 8 to \alert{solder the header pins to the board}. If you're taking my class I can solder this for you or teach everyone about soldering during a lecture session of class. If you are taking this class elsewhere you have two options: try and find someone who can solder this real quick (only takes about 5 minutes) or buy your own soldering iron and try to solder yourself. Once the device is soldered you can "plug" it into a breadboard and then move on to either the Arduino or CPX\slash{}CPB sections.%
\end{introduction}%
%
%
\typeout{************************************************}
\typeout{Section 11.1 Setting up the CPX\slash{}CPB}
\typeout{************************************************}
%
\begin{sectionptx}{Section}{Setting up the CPX\slash{}CPB}{}{Setting up the CPX\slash{}CPB}{}{}{sdcard-3}
Before wiring up the system I recommend setting up the CPX\slash{}CPB first before wiring up the system. The example code we're going to use for this chapter is on \href{https://github.com/cmontalvo251/Microcontrollers/blob/master/Circuit_Playground/CircuitPython/SD/sdcard_cpb_button_presses.py}{\linkhilite{Github}} as usual.%
\begin{figureptx}{Figure}{Snippet of SD card Software for the CPB\slash{}CPX}{sdcard-3-3}{}%
\begin{image}{0.135}{0.73}{0.135}{}%
\includegraphics[width=\linewidth]{external/SDcard_Snippet_CPB.png}
\end{image}%
\tcblower
\end{figureptx}%
The figure above shows the preamble of the code which has a few things that you need to do before getting this to work. \alert{First, you need to have an \emph{sd} folder in the your \emph{CIRCUITPY} drive.} That means, when you plug in the CPB\slash{}CPX to your computer you must add a solder called \emph{sd} into the main folder alongside the \emph{lib} folder which needs the \emph{adafruit\textunderscore{}sdcard} module. Before installing extra modules into the \emph{lib} folder you want to make sure you have your \href{https://circuitpython.org/downloads}{\linkhilite{Circuit Python UF2}} up to date\hyperlink{CircuitPython}{\linkhilite{[{\xreffont 30}]}}. In this example I'm using the 10.X version. Once I updated my UF2 I also updated my \href{https://circuitpython.org/libraries}{\linkhilite{Circuit Python Libraries}} (See \hyperref[installing-modules]{\linkhilite{Section~{\xreffont\ref{installing-modules}}}} for help on installing extra modules on your CPX\slash{}CPB)\hyperlink{CircuitPython}{\linkhilite{[{\xreffont 30}]}}. Again, for this lab you just need the \emph{adafruit\textunderscore{}sdcard} module in the lib folder. After that you can copy the software from Github into your \emph{code.py} file and save it. The code will not run right away and will probably through an \emph{SD Card Not Found} error. At this point you can unplug the CPX\slash{}CPB and begin wiring the breakout board making to plug in the micro SD card of course.%
\begin{figureptx}{Figure}{CPB wired to a micro SD card breakout board}{sdcard-3-5}{}%
\begin{image}{0.135}{0.73}{0.135}{}%
\includegraphics[width=\linewidth]{external/CPB_Method5_Wiring.png}
\end{image}%
\tcblower
\end{figureptx}%
In the figure above moving counter-clockwise, the red alligator clip from \emph{5V} to a green male-male wire to \emph{5V}, the first yellow alligator clip goes directly from \emph{A1(CLK)} to \emph{CLK}, the first green alligator clip goes from \emph{A2(MISO)} to a white male-male wire to \emph{DO}, the second green alligator clip goes from \emph{A3(MOSI)} to a blue male-male wire to \emph{DI}, the second yellow alligator clip is connected to \emph{A4} then to a black male-male wire connected to \emph{CS} and finally the black alligator clip goes directly from \emph{GND} to \emph{GND}.%
\begin{figureptx}{Figure}{Example output when writing to the SD card}{sdcard-3-7}{}%
\begin{image}{0.135}{0.73}{0.135}{}%
\includegraphics[width=\linewidth]{external/CPB_Method5_DatFileOutput.png}
\end{image}%
\tcblower
\end{figureptx}%
Once you plug in the CPB\slash{}CPX and press the A button to start recording you can press the A button over and over again to stop logging and then restart logging. Everytime you close the file and open a new file, the code will find a new file to create so that you don't continue to overwrite the file that is already on the sd card. In this code the A button is used to start and stop recording and the B button along with a dummy counter is logged to the SD card.%
\begin{figureptx}{Figure}{Example Data file after writing to SD Card}{sdcard-3-9}{}%
\begin{image}{0.135}{0.73}{0.135}{}%
\includegraphics[width=\linewidth]{external/DataFile_CPB_Data.png}
\end{image}%
\tcblower
\end{figureptx}%
When you are done recording make sure to press the A button to the close the current file and then remove power from the CPX\slash{}CPB. You can then place the SD card into your computer and you should see multiple \emph{DataFile\#.txt} text files in your computer. You can then copy those to your computer and plot the data similar to \hyperref[daq]{\linkhilite{Chapter~{\xreffont\ref{daq}}}}. This code can easily be augmented to accommodate any other sensor such as GPS (\hyperref[gps]{\linkhilite{Chapter~{\xreffont\ref{gps}}}}), accelerometer (\hyperref[imu]{\linkhilite{Chapter~{\xreffont\ref{imu}}}}), pressure (\hyperref[pressure]{\linkhilite{Chapter~{\xreffont\ref{pressure}}}}) or any other sensor you can think of. Although you add the complexity of having an extra piece of hardware, you remove the software headache and potential bricking of your CPX\slash{}CPB by attempting to use Method 3 (\hyperref[daq]{\linkhilite{Chapter~{\xreffont\ref{daq}}}}) while still removing the need for a laptop.%
\end{sectionptx}
%
%
\typeout{************************************************}
\typeout{Section 11.2 Setting up the Arduino}
\typeout{************************************************}
%
\begin{sectionptx}{Section}{Setting up the Arduino}{}{Setting up the Arduino}{}{}{sdcard-4}
There is no setup for the Arduino other than making sure the SD card is formatted to \emph{FAT32} and that the breakout board is soldered properly to header pins. The wiring for an Arduino also requires 5 wires that I recommend doing while the system is powered off.%
\begin{figureptx}{Figure}{Micro SD card Breakout Board wired to an Arduino MEGA}{sdcard-4-3}{}%
\begin{image}{0.135}{0.73}{0.135}{}%
\includegraphics[width=\linewidth]{external/MEGA_SD_2.png}
\end{image}%
\tcblower
\end{figureptx}%
In the figure above the dark green male-male wire goes from \emph{5V} to \emph{5V}, the orange male-male wire goes from \emph{GND} to \emph{GND}, the light green wire goes from \emph{52} to \emph{CLK}, the white wire goes from \emph{50} to \emph{DO}, the purple male-male wire goes from \emph{51} to \emph{DI} and the blue male-male wire goes from \emph{53} to \emph{CS}. The software to get this to work is on \href{https://github.com/cmontalvo251/Microcontrollers/blob/master/Arduino/MEGA/SD_Example/SD_Example.ino}{\linkhilite{Github}} as usual and a snippet is shown in the figure below.%
\begin{figureptx}{Figure}{Arduino MEGA Software Snippet to send data to a micro SD card breakout board}{sdcard-4-5}{}%
\begin{image}{0.135}{0.73}{0.135}{}%
\includegraphics[width=\linewidth]{external/SDCardSoftware_Arduino.png}
\end{image}%
\tcblower
\end{figureptx}%
You'll notice in this code that there is a \emph{\#include \textless{}SD.h\textgreater{}} line of code which imports the SD header file. In C++ and other compilation languages you have a \emph{.cpp} or \emph{.c} file which has all the code necessary to run somthing and a \emph{.h} file which declares all the variables. In Python you don't have to declare a variable the language figures it out on its own but in Arduino you must tell the computer what type of variable you are using. In this case if you need an extra module you don't import it you include it. To include extra module you must click \emph{Sketch\textgreater{}Include Library\textgreater{}Manage Libraries...}.%
\begin{figureptx}{Figure}{Installing New Arduino Libraries}{sdcard-4-7}{}%
\begin{image}{0.135}{0.73}{0.135}{}%
\includegraphics[width=\linewidth]{external/Arduino_Libraries_1.png}
\end{image}%
\tcblower
\end{figureptx}%
After clicking that a new window will pop up where you can then search for "SD.h" and after scrolling down a bit you'll find the SD library to install. After installing it you will be able to Compile the software without any errors and Upload the script to your Arduino. The Serial monitor will work with this script as \emph{Serial.begin(115200)} is in the setup routine and there are many \emph{Serial.print} calls in the loop routine. If you open the serial monitor you should see the \emph{millis()} time printed on the screen provided you set the baud rate to 115200.%
\begin{figureptx}{Figure}{Arduino Serial Monitor Output when writing to an SD Card}{sdcard-4-9}{}%
\begin{image}{0.135}{0.73}{0.135}{}%
\includegraphics[width=\linewidth]{external/SD_Arduino_Serial.png}
\end{image}%
\tcblower
\end{figureptx}%
In this case the serial monitor also states that is is writing the \emph{DATA004.TXT}. Everytime the code runs it will look for a new file that doesn't exist so that everytime you boot up the code it will create a new file. It ends up creating a lot of files but at least your data isn't overwritten. It also only outputs the internal timer because the Arduino does not have a built-in button but once we do other labs with pressure and acceleration it will be easy to send that data to the SD card as is shown in the code snippet below.%
\begin{figureptx}{Figure}{Code to send acceleration data to an SD card}{sdcard-4-11}{}%
\begin{image}{0.135}{0.73}{0.135}{}%
\includegraphics[width=\linewidth]{external/Acceleration_SD_Code.png}
\end{image}%
\tcblower
\end{figureptx}%
The code to get an external accelerometer to work on an Arduino is dicussed in more detail in \hyperref[imu]{\linkhilite{Chapter~{\xreffont\ref{imu}}}}. Still, for the simple example of just writing time to an SD card, by simply unplugging the Arduino and placing the SD card into an SD card reader will show all the files on your laptop and opening the txt file itself will show all the data you just took.%
\begin{figureptx}{Figure}{Data files on a micro SD Card}{sdcard-4-13}{}%
\begin{image}{0.135}{0.73}{0.135}{}%
\includegraphics[width=\linewidth]{external/SD_Arduino_Data.png}
\end{image}%
\tcblower
\end{figureptx}%
At this point it is identical to \hyperref[daq]{\linkhilite{Chapter~{\xreffont\ref{daq}}}} to plot the data in whatever plotting environment you like.%
\end{sectionptx}
%
%
\typeout{************************************************}
\typeout{Section 11.3 Assignment}
\typeout{************************************************}
%
\begin{sectionptx}{Section}{Assignment}{}{Assignment}{}{}{sdcard-5}
This project is similar to the Data Acquisition project only you must use method 5 to save time and button presses to a text file. You must then plot the button presses as a function of time. Remember to add x and y labels to all figures.%
%
\begin{enumerate}
\item{}Title page with name,title and date - 5\%%
\item{}Introduction with wiring diagram - 5\%%
\item{}All figures must have appropriate annotations, figure captions, figure labels and a paragraph of supporting text - 10\%%
\item{}Include a screenshot (entire computer screen) of your SD card open and mounted on your computer with your data file in it - 20\%%
\item{}Include a snippet (5 lines) of your data file - 20\%%
\item{}Include a plot of your button presses with time on the x-axis and button presses on the y-axis (no screenshots) - 40\%%
\item{}Appendix A - Video Details - Film yourself and your entire screen with the Serial monitor open as you press a button on the microcontroller if a button exists. Then show yourself removing the SD card and plotting the data on your computer - Pass\slash{}Fail%
\item{}Appendix B - Code used to take data and code used to plot in plain text (no screenshots) - Pass\slash{}Fail%
\item{}Appendix C - Screenshot of Excel or Sheets if applicable - Pass\slash{}Fail%
\item{}\alert{If you are asked to fix something in the Appendices such as missing content in your video or screenshots of your code, you will lose 10\% of your grade everytime you submit an update after the deadline.}%
\end{enumerate}
\end{sectionptx}
\end{chapterptx}
%
%
\typeout{************************************************}
\typeout{Chapter 12 Logging with Serial - Method 6}
\typeout{************************************************}
%
\begin{chapterptx}{Chapter}{Logging with Serial - Method 6}{}{Logging with Serial - Method 6}{}{}{serial}
\renewcommand*{\chaptername}{Chapter}
\begin{introduction}{}%
\begin{listptx}{List}{\textbf{Parts List}}{serial-2-1}{}%
%
\begin{enumerate}
\item{}USB Cable%
\item{}Computer with Python installed%
\item{}CircuitPlayground Express or Bluefruit%
\end{enumerate}
\end{listptx}%
\begin{listptx}{List}{\textbf{Learning Objectives}}{serial-2-2}{}%
%
\begin{enumerate}
\item{}Understand how to use serial communication to log data from a microcontroller%
\item{}Learn to read and interpret serial data in Python%
\item{}Create a simple data logging script using the serial library%
\end{enumerate}
\end{listptx}%
Method 6 is another "automatic" way to log data from your microcontroller. This method uses the serial library in Python to read data sent from your CircuitPlayground Express or Bluefruit. The idea is to have the Circuit Playground send data via USB to your laptop just like the REPL in Mu but instead of just viewing the output in Mu a Python script on a laptop can read and process the data. Serial communication is a fundamental method for microcontrollers to send data to a computer. In this chapter, we will explore how to use the serial library in Python to read data from your CircuitPlayground Express or Bluefruit. This method allows you to log data in real-time and can be a powerful tool for data acquisition and analysis. The code on the CPX\slash{}CPB is the same button press code used for the last few chapters and again can be found on \href{https://github.com/cmontalvo251/Microcontrollers/blob/master/Circuit_Playground/CircuitPython/Data_Logging/button_presses.py}{\linkhilite{Github}}. The code itself is discussed in detail in Chapter \hyperref[daq]{\linkhilite{Chapter~{\xreffont\ref{daq}}}}.%
\end{introduction}%
%
%
\typeout{************************************************}
\typeout{Section 12.1 Serial Setup}
\typeout{************************************************}
%
\begin{sectionptx}{Section}{Serial Setup}{}{Serial Setup}{}{}{serial-3}
In order to read serial data from a microcontroller, the serial library must be installed in Python. The serial library allows the computer to communicate with the serial port on a computer and read data sent from the microcontroller. The steps to install the serial module is the same as installing \emph{numpy} or \emph{matplotlib} explained in Chapter \hyperref[plotting]{\linkhilite{Chapter~{\xreffont\ref{plotting}}}}. If using Spyder or command line the command is shown below.%
\begin{preformatted}
!pip install serial
\end{preformatted}
or%
\begin{preformatted}
!pip3 install serial
\end{preformatted}
if using Python3. If using Thonny, the process is the same as installing any other module by going to Tools \textgreater{} Manage Packages. The code to read serial data from the CPX\slash{}CPB is shown below and can also be found on \href{https://github.com/cmontalvo251/Microcontrollers/blob/master/Circuit_Playground/CircuitPython/Serial_Handshake/read_serial_cpb.py}{\linkhilite{Github}}.%
\begin{figureptx}{Figure}{Example Serial code in Python}{serial-3-7}{}%
\begin{image}{0.135}{0.73}{0.135}{}%
\includegraphics[width=\linewidth]{external/Serial_Code.png}
\end{image}%
\tcblower
\end{figureptx}%
The code above is pretty advanced and uses a class called \emph{CPX} to handle all of the serial communication as well as saving the text file and plotting. The \emph{\textunderscore{}\textunderscore{}init\textunderscore{}\textunderscore{}} function creates a filename and opens a serial port on "\slash{}dev\slash{}ttyACM0" which is a serial port for Linux. If using Windows the port will be something like "COM3" and if using Mac it will be something like "\slash{}dev\slash{}tty.usbmodem14101". The \emph{read} function reads the serial data and saves it to a text file by first stripping the binary information, line feeds and converts the number of a 64 bit number. The \emph{extractdata} function plots the data after the serial reading is done. The main part of the code is in the while loop at the bottom which continuously reads serial data for 10 seconds in this case. So basically, the CPX\slash{}CPB is coded using Mu to send serial data to the computer just like the REPL in Mu and then the laptop program is run to save 10 second of data to a text file and then plot the data. The figure below shows the result of running the example script. \begin{figureptx}{Figure}{Example Serial Plot of Button Presses vs Time}{serial-3-8-5}{}%
\begin{image}{0.135}{0.73}{0.135}{}%
\includegraphics[width=\linewidth]{external/Plot_Serial.png}
\end{image}%
\tcblower
\end{figureptx}%
%
\end{sectionptx}
%
%
\typeout{************************************************}
\typeout{Section 12.2 Extra Help}
\typeout{************************************************}
%
\begin{sectionptx}{Section}{Extra Help}{}{Extra Help}{}{}{serial-4}
There are two great youtube videos on this topic. The first video discusses \href{https://www.youtube.com/watch?v=ULAeE9UxtHc}{\linkhilite{how to use the "screen" program on Ubuntu (Linux) to read serial data from a microcontroller without the use of Mu}}. This is a great starting point to see how serial communication is sent to a laptop via USB. The second video is a bit more advanced and shows \href{https://www.youtube.com/watch?v=UG6aolK8p3I}{\linkhilite{how to use the serial library in Python to read serial data from a microcontroller}}. The video walks through the code explained in the previous sections and is a great starting point if video lectures are preferred to text%
\end{sectionptx}
%
%
\typeout{************************************************}
\typeout{Section 12.3 Assignment}
\typeout{************************************************}
%
\begin{sectionptx}{Section}{Assignment}{}{Assignment}{}{}{serial-5}
This project is similar to the Data Acquisition project only you must use method 6 to save time and button presses to a text file. You must then plot the button presses as a function of time. Remember to add x and y labels to all figures.%
%
\begin{enumerate}
\item{}Title page with name,title and date - 5\%%
\item{}Introduction with wiring diagram - 5\%%
\item{}All figures must have appropriate annotations, figure captions, figure labels and a paragraph of supporting text - 10\%%
\item{}Include a screenshot (entire computer screen) of your computer actively reading serial data from your CPX\slash{}CPB - 20\%%
\item{}Include a snippet (5 lines) of your data file - 20\%%
\item{}Include a plot of your button presses with time on the x-axis and button presses on the y-axis (no screenshots) - 40\%%
\item{}Appendix A - Video Details - Film yourself and your entire screen with the Serial monitor open in Mu as you press a button on the CPX\slash{}CPB. Then show your computer actively reading and saving the data via serial. Finall,y show yourself plotting the data on your computer - Pass\slash{}Fail%
\item{}Appendix B - Code used to take data and code used to plot in plain text (no screenshots) - Pass\slash{}Fail%
\item{}Appendix C - Screenshot of Excel or Sheets if applicable - Pass\slash{}Fail%
\item{}\alert{If you are asked to fix something in the Appendices such as missing content in your video or screenshots of your code, you will lose 10\% of your grade everytime you submit an update after the deadline.}%
\end{enumerate}
\end{sectionptx}
\end{chapterptx}
%
%
\typeout{************************************************}
\typeout{Chapter 13 Potentiometer}
\typeout{************************************************}
%
\begin{chapterptx}{Chapter}{Potentiometer}{}{Potentiometer}{}{}{voltage}
\renewcommand*{\chaptername}{Chapter}
\begin{introduction}{}%
\begin{listptx}{List}{\textbf{Parts List}}{voltage-2-1}{}%
%
\begin{enumerate}
\item{}Laptop%
\item{}CPX\slash{}CPB%
\item{}USB Cable%
\item{}Potentiometer%
\item{}Resistor (the Ohms depends on how large your potentiometer is)%
\item{}Breadboard%
\end{enumerate}
\end{listptx}%
\begin{listptx}{List}{\textbf{Learning Objectives}}{voltage-2-2}{}%
%
\begin{enumerate}
\item{}Understand voltage division of resistors in series%
\item{}Measure an analog signal on the CircuitPlayground%
\item{}Understand the binary measurement done by the analog to digital conversion (ADC)%
\end{enumerate}
\end{listptx}%
At this point in a standard Instrumentation course, you've learned about analog to digital converters (ADC). It turns out that the CPX has 8 analog ports hooked up to a 3.3V logic 16 bit ADC. The input range on the ADC is 0 to 3.3V and the output range is 0 to 65536 which is \(2^{16}\) hence 16 bits. In order to get accustomed to the ADC on the CPX, we're going to do a simple example where we measure the voltage drop across a potentiometer. You can read about \href{https://www.build-electronic-circuits.com/potentiometer/}{\linkhilite{potentiometers online if you wish}}\hyperlink{potentiometerlearn}{\linkhilite{[{\xreffont 60}]}}. Basically though, a potentiometer is a variable resistance resistor that changes resistance by turning a knob. The knob changes the connection point of a wire and thus the length of the wire. This in turn changes the resistance. Potentiometers come in all shapes and sizes. Here are some \href{https://uk.rs-online.com/web/content/discovery/ideas-and-advice/potentiometers-guide}{\linkhilite{potentiometer examples}}. I've done this lab with a few potentiometer. Ideally you'd like to have the potentiometer hooked up in series with another resistor so that you end up building a voltage divider but it' possible you can do it without it as shown in the \href{https://learn.adafruit.com/sensor-plotting-with-mu-and-circuitpython/potentiometer}{\linkhilite{figure below}} (Courtesy of \href{https://learn.adafruit.com/u/kattni}{\linkhilite{Kattni Rembor}}) \hyperlink{potentiometer}{\linkhilite{[{\xreffont 38}]}}.%
\begin{figureptx}{Figure}{Wiring Diagram of a Circuit Playground Express hooked up to a Potentiometer\hyperlink{potentiometer}{\linkhilite{[{\xreffont 38}]}}}{voltage-2-4}{}%
\begin{image}{0.1}{0.8}{0.1}{}%
\includegraphics[width=\linewidth]{external/sensors_PlottingMuPotentiometer_bb.jpg}
\end{image}%
\tcblower
\end{figureptx}%
In this chapter you're going to wire up a potentiometer to the CPX and read the voltage across it as you turn the knob.%
\end{introduction}%
%
%
\typeout{************************************************}
\typeout{Section 13.1 Setup}
\typeout{************************************************}
%
\begin{sectionptx}{Section}{Setup}{}{Setup}{}{}{voltage-3}
Here's my circuit all hooked up without a resistor in series. Two legs are connected to 3.3V and GND while the middle leg of the potentiometer is connected to pin A2.%
\begin{figureptx}{Figure}{Circuit Playground Express hooked up to a Potentiometer}{voltage-3-3}{}%
\begin{image}{0.1}{0.8}{0.1}{}%
\includegraphics[width=\linewidth]{external/potentiometer1.jpeg}
\end{image}%
\tcblower
\end{figureptx}%
\alert{As I said before, some potentiometers do not have enough resistance when turned all the way down. I suggest that you put a resistor in between the third leg and ground. Some experimenters have melted plastic or gotten really hot. One student even blew up a potentiometer.} Here is my circuit with a resistor in series.%
\begin{figureptx}{Figure}{Potentiometer with a blue resistor in series}{voltage-3-5}{}%
\begin{image}{0.1}{0.8}{0.1}{}%
\includegraphics[width=\linewidth]{external/potentiometer2.png}
\end{image}%
\tcblower
\end{figureptx}%
There is a relevant \href{https://learn.adafruit.com/adafruit-circuit-playground-express/circuitpython-analog-in}{\linkhilite{Adafruit Learn Tutorial}} to help with the \emph{analogio} module but I'll explain the minimum required here to get some analog values plotted in \emph{Plotter} and on your computer\hyperlink{Adafruit}{\linkhilite{[{\xreffont 32}]}}. First let's take a look at some \href{https://github.com/cmontalvo251/Microcontrollers/blob/master/Circuit_Playground/CircuitPython/Analog/analog_simple.py}{\linkhilite{simple example code}} to read an analog signal and plot it using the \emph{Plotter}\hyperlink{Github}{\linkhilite{[{\xreffont 33}]}}.%
\begin{figureptx}{Figure}{Code to read analog signal in CircuitPython}{voltage-3-7}{}%
\begin{image}{0.25}{0.5}{0.25}{}%
\includegraphics[width=\linewidth]{external/analogio.png}
\end{image}%
\tcblower
\end{figureptx}%
In the example code above, lines 1-3 again import the necessary modules with \emph{analogio} being the new module here. Line 5 creates the analog object by attaching pin A2 to the analog function. Lines 7-9 then simple read the analog value and print it to \emph{Serial} and the \emph{Plotter}. Running this code on my laptop and turning the knob on the potentiometer produces this output. My potentiometer has a very large knob on the front and is easy to turn. Some potentiometers have a small screw on top that you need to turn with a screwdriver. Turning the screw or the knob results in chaning the resistance and therefore changing the voltage read by the CPX.%
\begin{figureptx}{Figure}{Reading an analog signal in Mu with Serial monitor and Plotter open}{voltage-3-9}{}%
\begin{image}{0.135}{0.73}{0.135}{}%
\includegraphics[width=\linewidth]{external/analogio_mu.png}
\end{image}%
\tcblower
\end{figureptx}%
For this lab I want you to spin the potentiometer all the way to one side and then the other while recording time and the analog value. I then want you to plot the data with time on the x-axis and voltage on the y-axis. Remember to convert a digital output to voltage you just need to use the equation below where D is the raw value from the analog port. 3.3V is the range of the ADC and \(2^{16}\) is the maximum value the ADC can represent.%
\begin{equation}
V = \frac{3.3D}{2^{16}}\label{voltage-3-10-2}
\end{equation}
%
\par
After doing this experiment myself, this is the plot I obtain. The code is not provided as reading data and plotting has been discussed in a previous lab (See \hyperref[daq]{\linkhilite{Chapter~{\xreffont\ref{daq}}}}). From the screenshot though you can see how I convert the digital output to an analog signal.%
\par
\alert{NOTE THAT ON LINE 6 IT READS}%
\begin{preformatted}
time -= time[0]
\end{preformatted}
\alert{Notice the minus sign in front of the equal sign. That effects a lot.}%
\begin{figureptx}{Figure}{Plotting analog signal in volts in Thonny}{voltage-3-15}{}%
\begin{image}{0.135}{0.73}{0.135}{}%
\includegraphics[width=\linewidth]{external/plotting_analogio.png}
\end{image}%
\tcblower
\end{figureptx}%
\end{sectionptx}
%
%
\typeout{************************************************}
\typeout{Section 13.2 Extra Help}
\typeout{************************************************}
%
\begin{sectionptx}{Section}{Extra Help}{}{Extra Help}{}{}{voltage-4}
I've made some youtube videos on first just \href{https://youtu.be/_gnDvPOvPqk}{\linkhilite{creating the circuit and plotting the data}} and then another video where I \href{https://youtu.be/9UF3OVUIYjU}{\linkhilite{write data to the CPX using method 3}}\hyperlink{youtube}{\linkhilite{[{\xreffont 50}]}}.%
\end{sectionptx}
%
%
\typeout{************************************************}
\typeout{Section 13.3 Assignment}
\typeout{************************************************}
%
\begin{sectionptx}{Section}{Assignment}{}{Assignment}{}{}{voltage-5}
\begin{introduction}{}%
Your assignment for this lab is wire up the potentiometer, read the analog signal using the analog to digital converter on the CPX and plot it on your desktop computer. Remember to record both time and digital output as you rotate the potentiometer \alert{through the full range} and then plot that as a function of time. Specific requirements are shown below.%
%
\begin{enumerate}
\item{}Title page with name,title and date - 5\%%
\item{}Introduction with wiring diagram - 5\%%
\item{}All figures must have appropriate annotations, figure captions, figure labels and a paragraph of supporting text - 10\%%
\item{}Include a photo of your circuit showing the potentiometer wired up to an analog pin on your CPX\slash{}CPB - 10\%%
\item{}Include a screenshot (entire computer screen) of Mu with the Plotter open showing the digital output of the potentiometer. The code in Mu also needs to also show the same analog pin as your potentiometer. - 10\%%
\item{}Based on the digital output from the potentiometer, compute the minimum and maximum voltage across the potentiometer. - 10\%%
\item{}Relate the voltage across your potentiometer to angle of the potentiometer in degrees and plot degrees on the x-axis and voltage across the potentiometer on the y-axis. Be sure to include the equation relating angle to voltage - 10\%%
\item{}Plot your digital output (raw potentiometer analog value) vs time \alert{make sure to rotate the potentiometer through the full range as when you take data} - 10\%%
\item{}Then convert your digital output (Do) to voltage and plot that vs time - 10\%%
\item{}Finally convert your voltage to angle in degrees and plot that vs time - 20\%%
\item{}Appendix A - Video Details - Film yourself and your entire screen with the Plotter open in Mu as you rotate the potentiometer from it's maximum to minimum position. Then show yourself plotting the raw data on your computer - Pass\slash{}Faill%
\item{}Appendix B - Code used to take data and code used to plot in plain text (no screenshots) - Pass\slash{}Fail%
\item{}\alert{If you are asked to fix something in the Appendices such as missing content in your video or screenshots of your code, you will lose 10\% of your grade everytime you submit an update after the deadline.}%
\end{enumerate}
\end{introduction}%
\end{sectionptx}
\end{chapterptx}
%
%
\typeout{************************************************}
\typeout{Chapter 14 Pitot Probe}
\typeout{************************************************}
%
\begin{chapterptx}{Chapter}{Pitot Probe}{}{Pitot Probe}{}{}{pitot}
\renewcommand*{\chaptername}{Chapter}
\begin{introduction}{}%
\begin{listptx}{List}{\textbf{Parts List}}{pitot-2-1}{}%
%
\begin{enumerate}
\item{}Laptop%
\item{}CPX\slash{}CPB%
\item{}USB Cable%
\item{}Alligator Clips (x3)%
\item{}\href{https://www.amazon.com/Pressure-MPXV7002DP-Breakout-Transducer-Components/dp/B08C5CF7CB}{\linkhilite{Pressure Transducer}} + \href{https://www.amazon.com/Jadeshay-Differential-Speedometer-Precision-Instrument/dp/B07Q6126B6/ref=pd_bxgy_thbs_d_sccl_1/131-3143700-2176464?pd_rd_w=7zZaR\&content-id=amzn1.sym.9bef5913-5870-4504-8883-3ba89d7f8e39\&pf_rd_p=9bef5913-5870-4504-8883-3ba89d7f8e39\&pf_rd_r=KSZ04S642T1YZ2QXTHWK\&pd_rd_wg=MrPsd\&pd_rd_r=63018516-b34d-43f9-9369-36d494eda7a3\&pd_rd_i=B07Q6126B6\&psc=1}{\linkhilite{Pitot Probe}}%
\item{}Breadboard%
\end{enumerate}
\end{listptx}%
\begin{listptx}{List}{\textbf{Learning Objectives}}{pitot-2-2}{}%
%
\begin{enumerate}
\item{}Understand how pitot probes works%
\item{}Understand the relationship between a voltage signal from a pitot probe to a pressure value%
\item{}Understand the relationship between pressure and windspeed%
\end{enumerate}
\end{listptx}%
Although a CPX has numerous sensors built in, you can easily augment the capabilities of the CPX using either I2C or just the ADC on board the CPX. In this lab, if you purchased a \href{https://www.amazon.com/Hobbypower-Airspeed-MPXV7002DP-Differential-controller/dp/B00WSFWO36/ref=sr_1_3?dchild=1\&keywords=Airspeed+sensor+kit\&qid=1590532161\&sr=8-3}{\linkhilite{pitot probe}} you will be able to do this assignment\hyperlink{pitot-amazon}{\linkhilite{[{\xreffont 40}]}}. Since you don't need the pitot probe for very long you can always borrow one from some other team or from your instructor if they have one. Let's talk about the hardware and the wiring to get this to work. First, pitot probes work by mechanically changing the dynamic pressure of the incoming airflow as shown in the \href{https://makersportal.com/blog/2019/02/06/arduino-pitot-tube-wind-speed-theory-and-experiment}{\linkhilite{figure below}} (Courtesy of \href{https://makersportal.com/blog?author=59b036fc6073554c1cfffef7}{\linkhilite{Joshua Hrisko}}) \hyperlink{pitot-author}{\linkhilite{[{\xreffont 39}]}}.%
\begin{figureptx}{Figure}{Schematic of Pitot probe showing static and dynamic pressures\hyperlink{pitot-author}{\linkhilite{[{\xreffont 39}]}}}{pitot-2-4}{}%
\begin{image}{0.135}{0.73}{0.135}{}%
\includegraphics[width=\linewidth]{external/pitot_tube_drawing_transparent_cleaner.png}
\end{image}%
\tcblower
\end{figureptx}%
The pitot probe has two pressure taps that measure static (ambient) pressure and dynamic (stagnation) pressure. These taps move through two silicon tubes to a pressure transducer that has a strain gauge that separates both pressures. When the pressure on one side of the transducer is larger than the other, it will flex the membrane and create strain. This strain runs through a wheatstone bridge with a voltage offset to the pin labeled analog. The analog signal from the analog pin will be denoted as \(V_{pitot}\). The pressure transducer used in this lab is the MPXV7002 which is a differential pressure sensor. The \href{https://www.nxp.com/docs/en/data-sheet/MPXV7002.pdf}{\linkhilite{datasheet}} indicates that the voltage offset of the pitot probe is 2.5V and the change in voltage is proportional to the change in pressure in units of kPa\hyperlink{pitot-datasheet}{\linkhilite{[{\xreffont 41}]}}. This means that \(\Delta P\) is given by the equation below.%
\begin{equation}
\Delta P_{kPa} = V_{pitot} - 2.5\label{pitot-2-5-5}
\end{equation}
remembering that \(V_{pitot}=3.3D_o/65535\). Also note that 2.5V is the nominal voltage even though your sensor might have something slightly different like 2.4V or 2.6V. To measure wind speed you can use a variation of Bernoulli's principle. Note it is necessary to convert the pressure \(P\) from kilopascals (kPa) to Pascals (Pa) by multiplying by \(1000\) before using it in this formula.%
\begin{equation}
U = \sqrt{\frac{2 \Delta P_{Pa}}{\rho}}\label{pitot-2-5-9}
\end{equation}
where \(U\) is the wind speed in \(m/s\), \(\Delta P\) is the differential pressure (in Pascals), and \(\rho=1.225~kg/m^3\) is the density of air. As for wiring up the circuit itself, the transducer has 3 pins, +5V, GND and Analog. The figure below shows the transducer connected to an Arduino (Courtesy of \href{https://makersportal.com/blog/2019/02/06/arduino-pitot-tube-wind-speed-theory-and-experiment?srsltid=AfmBOoqrja9WFwnJNLPOLxPEG5mRa1XqM8iP86s0ZYy5Qb1aSA-GpU0D}{\linkhilite{Joshua Hrisko}})\hyperlink{pitot-author}{\linkhilite{[{\xreffont 39}]}}. \begin{figureptx}{Figure}{Arduino connected to a pressure transducer\hyperlink{pitot-author}{\linkhilite{[{\xreffont 39}]}}}{pitot-2-5-16}{}%
\begin{image}{0.135}{0.73}{0.135}{}%
\includegraphics[width=\linewidth]{external/Arduino+MPXV7002DP+Wiring+for+Differential+Pressure+Measurement.png}
\end{image}%
\tcblower
\end{figureptx}%
%
\par
It is pretty straightforward to wire this up to a CPX\slash{}CPB as is shown in the wiring diagram below. \begin{figureptx}{Figure}{Pitot probe wired to a CPX}{pitot-2-6-1}{}%
\begin{image}{0.05}{0.9}{0.05}{}%
\includegraphics[width=\linewidth]{external/pitot_wiring.png}
\end{image}%
\tcblower
\end{figureptx}%
 Remember that +5V needs to go to VOUT, GND to GND and Analog to any analog pin. Continue to read on to setup your own CPX\slash{}CPB and complete this assignment%
\end{introduction}%
%
%
\typeout{************************************************}
\typeout{Section 14.1 Setup}
\typeout{************************************************}
%
\begin{sectionptx}{Section}{Setup}{}{Setup}{}{}{setup}
To complete this lab, I chose pin A2 as shown in the Figure below.%
\begin{figureptx}{Figure}{CircuitPlayground Express connected to Pitot Probe and Transducer}{setup-3}{}%
\begin{image}{0.05}{0.9}{0.05}{}%
\includegraphics[width=\linewidth]{external/pitot_probe_circuit.jpeg}
\end{image}%
\tcblower
\end{figureptx}%
At that point it's very simple to just print the analog signal in bits to Serial. I've done this below. The code is the same \href{https://github.com/cmontalvo251/Microcontrollers/blob/master/Circuit_Playground/CircuitPython/Analog/analog_simple.py}{\linkhilite{analog code}} that we've used in the past\hyperlink{Github}{\linkhilite{[{\xreffont 33}]}}.%
\begin{figureptx}{Figure}{Analogio code in CircuitPython}{setup-5}{}%
\begin{image}{0.25}{0.5}{0.25}{}%
\includegraphics[width=\linewidth]{external/analogio.png}
\end{image}%
\tcblower
\end{figureptx}%
The goal of the experiment is to take pitot probe data for 15 seconds with no wind, then 15 seconds of data with a fan on and then 15 seconds of no wind data. You'll need to use one of the datalogging methods (See \hyperref[daq]{\linkhilite{Chapter~{\xreffont\ref{daq}}}}) to log both time and pitot probe analog value. Once you have that data, import the data on your desktop computer and convert the signal to windspeed as explained above. Using your data, create a plot of windspeed with time on the x-axis and windspeed on the y-axis. Some steps that might help you as you complete this project. First, have Mu plot the voltage coming from the pitot probe. If you've done everything right it will not be zero. The data sheet says there's an offset voltage of 2.5V so you will hopefully get something around 50,000 when you don't blow into the pitot probe. 50,000 multiplied by 3.3\slash{}\(2^{16}\) is around 2.5V. Make a note of that average value you get so you can subtract it off later. Once you've verified you're reading the pitot probe correctly, blow into the pitot probe and using the Plotter or Serial, verify that the analog signal increases. If the signal decreases, it means the pressure taps on the pressure transducer are backwards and you need to flip them. Either that or just flip the sign in your plotting routine on your computer but flipping the tubes might be easier for you. Hopefully when you do this lab you will get some data that looks like this. In this Figure you'll see that when the fan wasn't running the signal was something around 49,800 which is fine. It means your bias is around 2.5 volts. Every pitot probe and circuit will be different. You can then convert this signal to voltage then and then pressure and then finally wind speed.%
\begin{figureptx}{Figure}{Example pitot probe data showing fan off\slash{}on\slash{}off scenario}{setup-7}{}%
\begin{image}{0.135}{0.73}{0.135}{}%
\includegraphics[width=\linewidth]{external/pitot_probe_data.png}
\end{image}%
\tcblower
\end{figureptx}%
The code to accomplish this is relatively simple and a portion of the code is shown below. You'll see that when I subtracted the bias from the voltage I also zeroed out any negative values. That is, any delta voltage less than zero was set to zero. A couple of things about this chart. The data from the pitot probe is super noisy which means attaching a complementary filter is probably a good idea provided you don't over filter the signal and run into aliasing issues (See \hyperref[aliasing]{\linkhilite{Section~{\xreffont\ref{aliasing}}}}). You can see that I implemented an offline complementary filter and plotted it in the orange line which helps the noise issue quite a bit. You'll also notice that the noise is about 2 m\slash{}s. It turns out that pitot probes are actually not very accurate lower than about 2 m\slash{}s. They would be great for an airplane or you driving down the highway but they wouldn't be very good to take wind data outside on a calm day.%
\begin{figureptx}{Figure}{Example windspeed data showing complementary filter in Thonny}{setup-9}{}%
\begin{image}{0.135}{0.73}{0.135}{}%
\includegraphics[width=\linewidth]{external/pitot_probe_final.png}
\end{image}%
\tcblower
\end{figureptx}%
\end{sectionptx}
%
%
\typeout{************************************************}
\typeout{Section 14.2 Extra Help}
\typeout{************************************************}
%
\begin{sectionptx}{Section}{Extra Help}{}{Extra Help}{}{}{pitot-4}
I explain the process of substracting off bias in this \href{https://www.youtube.com/watch?v=e4xs9Ky7_YI\&feature=youtu.be}{\linkhilite{accelerometer video}}. That may help you figure out how to subtract off the bias of your initial pitot probe data. I've also done this pitot project before and have posted a video on Youtube about \href{https://youtu.be/jSLIRC1cfvE}{\linkhilite{Converting Pitot Probe Data to Windspeed}}. \alert{There is a typo in the video though. V1 is supposed to have a sqrt())}. I also have a video about \href{https://youtu.be/zTGa4sk6UZE}{\linkhilite{complementary filters}} and another about \href{https://youtu.be/8F_8st_8MmA}{\linkhilite{aliasing issues}} \hyperlink{youtube}{\linkhilite{[{\xreffont 50}]}}.%
\end{sectionptx}
%
%
\typeout{************************************************}
\typeout{Section 14.3 Assignment}
\typeout{************************************************}
%
\begin{sectionptx}{Section}{Assignment}{}{Assignment}{}{}{pitot-5}
\begin{introduction}{}%
For this assignment you are to wire up a pitot probe and record time and raw analog signal from the pitot probe as you turn a fan on and off. My suggestion is you record at least 30 seconds of data with the fan off and then 30 seconds with the fan on and then again 30 seconds with the fan off. Specific requirements are shown below.%
%
\begin{enumerate}
\item{}Title page with name,title and date - 5\%%
\item{}Introduction with wiring diagram - 5\%%
\item{}All figures must have appropriate annotations, figure captions, figure labels and a paragraph of supporting text - 10\%%
\item{}\alert{In addition to the standard format above, you must also return the pitot probe you borrowed in class - Pass\slash{}Fail}%
\item{}Include a photo of your circuit with your pitot probe wired up to an analog pin - 10\%%
\item{}Include a screenshot (entire computer screen) of Mu with the Plotter open showing the raw analog signal. The Mu code also needs to show the same analog pin as the circuit above - 10\%%
\item{}Include a plot of the raw analog signal vs time that clearly shows when the fan is on and off - 20\%%
\item{}Convert your raw signal to windspeed and filter your signal using an offline complementary filter. Plot both unfiltered and filtered windspeed on the same plot and include a legend. Make sure you add the saturation filter to prevent a negative in the square root - 20\%%
\item{}What is the maximum windspeed that the CPX\slash{}CPB can read? - 10\%%
\item{}The data sheet also suggests that you use a \(470~pF\) capacitor to filter the output. Select two resistors for a low pass filter such that the DC Gain would be equal to 1. Assume the total impedance of both resistors is \(2~k\Omega\). Also compute the cutoff frequency. - 10\%%
\item{}Appendix A - Video Details - Include 2 vidoes: First, film your entire team and your screen as you show the digital output change in the Plotter in Mu as the freestream airflow changes. Then, include a video of you plotting your raw digital output and converting that to windspeed on your computer - Pass\slash{}Fail%
\item{}Appendix B - Code used to take data and code used to plot in plain text (no screenshots) - Pass\slash{}Fail%
\item{}Appendix C - Screenshot of Excel or Sheets if applicable - Pass\slash{}Fail%
\item{}\alert{If you are asked to fix something in the Appendices such as missing content in your video or screenshots of your code, you will lose 10\% of your grade everytime you submit an update after the deadline.}%
\end{enumerate}
\end{introduction}%
\end{sectionptx}
\end{chapterptx}
%
%
\typeout{************************************************}
\typeout{Chapter 15 Built-In Sensors}
\typeout{************************************************}
%
\begin{chapterptx}{Chapter}{Built-In Sensors}{}{Built-In Sensors}{}{}{modules}
\renewcommand*{\chaptername}{Chapter}
\begin{introduction}{}%
\begin{listptx}{List}{\textbf{Parts List}}{modules-2-1}{}%
%
\begin{enumerate}
\item{}Laptop%
\item{}CPX\slash{}CPB%
\item{}USB Cable%
\end{enumerate}
\end{listptx}%
\begin{listptx}{List}{\textbf{Learning Objectives}}{modules-2-2}{}%
%
\begin{enumerate}
\item{}Understand the different sensors on the Circuitplayground%
\item{}Learn the difference between high level and low level control%
\item{}Get more practice plotting data from onboard sensors%
\end{enumerate}
\end{listptx}%
The CPX has numerous built-in sensors. These include a light sensor, an IR sensor, an accelerometer, a microphone, a speaker, some neopixels, a temperature sensor and 8 analog inputs with ADCs and even I2C (pronounced I squared C - it's a kind of serial communication) that you can use to easily hook up more sensors to it. We're not going to utilize all of these sensors since that would be a rather large project. Instead we're going to learn how to use the temperature, light and sound sensors as well as the accelerometer. For each of these examples there is a relatively easy way to access the sensors using a built-in module called \emph{adafruit\textunderscore{}circuitplayground.express} if you are using the CPX. If you are using the CPB you need to type \emph{adafruit\textunderscore{}circuitplayground.bluefruit}. It's a very nice module because it imports everything on the board. Section \hyperref[installing-modules]{\linkhilite{Section~{\xreffont\ref{installing-modules}}}} goes over the steps required to install different modules on the CPX\slash{}CPB. The problem is you can run into module conflicts. This happens when two different modules try to access the same pins on the CP. Sometimes if you import \emph{adafruit\textunderscore{}circuitplayground} you won't be able to import some other modules. \alert{Note you might need to add the \emph{adafruit\textunderscore{}circuitplayground} library to your lib\slash{} folder on your CIRCUITPY drive}. Due to this module conflict issue, there are some low level control commands you can use to access each of the sensors on board. We're obviously going to learn the low level control method first and then I'll show you how to access the sensors using the \emph{adafruit\textunderscore{}circuitplayground} module. \alert{If you get a \emph{currently in use} error it means you have a module conflict. Hence why I'm showing you the low level control method}.%
\end{introduction}%
%
%
\typeout{************************************************}
\typeout{Section 15.1 Low Level Control}
\typeout{************************************************}
%
\begin{sectionptx}{Section}{Low Level Control}{}{Low Level Control}{}{}{modules-3}
%
%
\typeout{************************************************}
\typeout{Subsection 15.1.1 Light}
\typeout{************************************************}
%
\begin{subsectionptx}{Subsection}{Light}{}{Light}{}{}{light-sensor}
The light sensor on the CPX is just a simple photocell wired in series with a resistor. A photocell (or Light Dependent Resistor, LDR) is a sensor whose resistance changes significantly based on the intensity of light falling on its sensitive surface. In dimmer light, its resistance is high, and in brighter light, its resistance is low. By integrating it into a simple voltage divider circuit, Adafruit makes it possible to read the voltage across the photocell and correlate that to light level. There is a relevant \href{https://learn.adafruit.com/photocells/circuitpython}{\linkhilite{Adafruit Tutorial on Photocells}} and the code required to measure the voltage if you'd like to read more about it\hyperlink{Adafruit}{\linkhilite{[{\xreffont 32}]}}.%
\par
The GND leg of the photocell is connected to pin A8. You can check the pin by looking at the graphic of an eye on the CPX and taking a look at the digital pin next to it. Since this circuit is connected to a voltage divider, the equations below can be used to relate the resistance in the photocell to the voltage being read by the ADC.%
\begin{equation}
V_{photocell} = V_{out}\frac{R_{photocell}}{R_{series}+R_{photocell}}\label{light-sensor-3-1}
\end{equation}
%
\par
In this case \(V_{out}\) is 3.3V and \(R_{series}=10000 \Omega\) is a resistor value chosen by Adafruit.  So, if you measure \(V_{photocell}\) using the ADC on the CPX\slash{}CPB you'll be able to solve for the resistance in the photocell \(R_{photocell}\). The resistance can then be used to determine the Lux value. Some \href{https://www.researchgate.net/publication/375038565_Measuring_the_Voltage_Current_and_Resistance_of_the_LDR_Sensor_through_the_Arduino_UNO}{\linkhilite{sources}} suggest that Lux is a power law given by the equation below \hyperlink{b_photocell}{\linkhilite{[{\xreffont 48}]}}\hyperlink{lux}{\linkhilite{[{\xreffont 61}]}}.%
\begin{equation}
Lux = 500000R_{photocell}^{-0.7}\label{light-sensor-4-8}
\end{equation}
where \(R_{photocell}\) is the resistance of the photocell in Ohms. We've already learned how to access analog pins (See \hyperref[voltage]{\linkhilite{Chapter~{\xreffont\ref{voltage}}}}) in a previous lab so just use the \href{https://github.com/cmontalvo251/Microcontrollers/blob/master/Circuit_Playground/CircuitPython/Analog/analog_simple.py}{\linkhilite{code from that lab}} and change the pin to A8\hyperlink{Github}{\linkhilite{[{\xreffont 33}]}}. Here's what my code looks like when I change the pin to A8. I also brought the Plotter up and moved my finger in front of the light to make sure the light was working. Verify that your CPX responds the same way before moving on.%
\begin{figureptx}{Figure}{Serial monitor and Plotter open in Mu showing digital output of pin A8 (light sensor)}{light-sensor-5}{}%
\begin{image}{0.135}{0.73}{0.135}{}%
\includegraphics[width=\linewidth]{external/modules_light_mu.png}
\end{image}%
\tcblower
\end{figureptx}%
\end{subsectionptx}
%
%
\typeout{************************************************}
\typeout{Subsection 15.1.2 Sound}
\typeout{************************************************}
%
\begin{subsectionptx}{Subsection}{Sound}{}{Sound}{}{}{modules-3-3}
The sound sensor uses the audiobusio library and creates a mic object using the (Pulse Density Modulation) PDM library. You have to set the sample rate and the number of bits to use to capture the data. We're going to set the bits to 16 to utilize the whole spectrum and then set the sample rate to 16 kHz. It's not quite 44.1 kHz like most modern microphones but it will do. After creating the mic object we have to compute some root mean squared values and thus two functions are defined before the while true loop in the code. The code itself is shown below. The code starts on line 22 because the first 22 lines are copyright from \href{https://learn.adafruit.com/adafruit-pdm-microphone-breakout/circuitpython}{\linkhilite{Dan Halbert, Kattni Rembor, and Tony DiCola}} from Adafruit Industries\hyperlink{Adafruit}{\linkhilite{[{\xreffont 32}]}}. I have edited the code a bit to fit my needs and uploaded \href{https://github.com/cmontalvo251/Microcontrollers/blob/master/Circuit_Playground/CircuitPython/Audio/record_sound_simple.py}{\linkhilite{my version to Github}}\hyperlink{Github}{\linkhilite{[{\xreffont 33}]}}. In the code line 23-27 import standard modules as well as some new ones. The array module is used to create array like matrices. The math module is used to compute functions like cos, sin, and sqrt. Then of course the audiobusio module is used to create the mic object on line 42. Notice the two functions defined on 33 and 39 which create a function for computing the mean and for computing the normalized root mean square value of the data stream. Basically what's going to happen is we're going to record 160 samples as defined on line 160. So on line 43 we create a hexadecimal array (hexademical: base 16 hence the num\textunderscore{}bits set to 16 on line 31) with 160 zeros. In the while true loop we're going to sleep for 0.01 seconds and then record some samples. Since we're sampling at 16 kHz the time it takes to record 160 samples is 160\slash{}16000 = 16\slash{}1600 = 1\slash{}100 = 0.01 seconds. Since we're taking 160 samples we need to compute some sort of average which is why the normalized root mean square value is computed on line 48.%
\begin{figureptx}{Figure}{Reading sound code in CircuitPython}{modules-3-3-3}{}%
\begin{image}{0.135}{0.73}{0.135}{}%
\includegraphics[width=\linewidth]{external/sound_code.png}
\end{image}%
\tcblower
\end{figureptx}%
When I run this code and talk normally into the microphone, I get this output in the Plotter. You'll notice that the data is pretty noisy in the beginning but then there are noticeable humps in the data. This is me saying something random into the microphone at normal volume. It's possible we could increase the number of samples we take each loop by editing line 30 but that would slow down our code. So there's a tradeoff between filtering here and speed. That's something will investigate in some later labs.%
\begin{figureptx}{Figure}{Serial monitor and Plotter open in Mu showing digital output sound sensor}{modules-3-3-5}{}%
\begin{image}{0.135}{0.73}{0.135}{}%
\includegraphics[width=\linewidth]{external/sound_mu.png}
\end{image}%
\tcblower
\end{figureptx}%
\end{subsectionptx}
%
%
\typeout{************************************************}
\typeout{Subsection 15.1.3 Temperature}
\typeout{************************************************}
%
\begin{subsectionptx}{Subsection}{Temperature}{}{Temperature}{}{}{temperature-sensor}
The temperature sensor is actually a \href{https://en.wikipedia.org/wiki/Thermistor}{\linkhilite{thermistor}}\hyperlink{thermistor-wikipedia}{\linkhilite{[{\xreffont 63}]}}. A thermistor is basically a \textquotedblleft{}thermometer resistor\textquotedblright{} which means the resistance depends on temperature. Since this thermistor on the CPX\slash{}CPB is connected to an analog pin, you can read the analog signal coming from the thermistor just by reading the analog signal from pin A9. If you look for the thermometer symbol on the CPX you'll see pin A9. Therefore, it is possible to just use the analogio library and just read in the analog voltage but in order to convert to celsius and then fahrenheit you need to use some heat transfer equations to convert the analog signal to celsius.%
\par
In this case, the thermistor is wired in a voltage divider circuit in series with a \(10000~\Omega\) resistor. For this circuit though, it turns out that the thermistor is actually wired to the high side of the voltage divider. That means when you measure pin A9 you're actually measuring the voltage across the series resistor rather than the thermistor. This means that the voltage across the series resistor must be inverted to compute the resistance across the thermistor. This is done using the equation below similar to the equation used for the photocell.%
\begin{equation}
V_{measured} = V_{out}\frac{R_{series}}{R_{series}+R_{thermistor}}\label{temperature-sensor-3-2}
\end{equation}
%
\par
Where \(V_{out}=3.3V\) and \(R_{series}=10000~\Omega\). The equation above can be inverted to obtain the resistance in the thermistor . This means that the digital output from the analog to digital converter can be converted to voltage and then to resistance as has been done for various ADC labs in this textbook. Remember that the ADC on boad the CPX\slash{}CPB is actually measuring the digital output \(D_o\). This can be converted to measured voltage as shown in the equation below%
\begin{equation}
V_{measured} = V_{out}\frac{D_o}{65535}\label{temperature-sensor-4-4}
\end{equation}
If you absolutely must have the voltage across the thermistor you simply use Kirchhoff's Voltage Law given in the equation below%
\begin{equation}
V_{thermistor} = V_{out}-V_{measured}\label{temperature-sensor-4-5}
\end{equation}
So the process is to first read the digital output from the ADC, convert that to voltage, then use the voltage to compute the resistance in the thermistor and finally use the resistance to compute the temperature in Kelvin. In order to convert measured voltage to resistance, we can rearrange the equations above to get the equation below%
\begin{equation}
R_{thermistor} = R_{series}\left(\frac{V_{out}}{V_{measured}}-1\right)\label{temperature-sensor-4-6}
\end{equation}
Note that \(V_{out}/V_{measured}=65535/D_o\). Once you have the resistance from the thermistor, you can use a modified version of the \href{https://en.wikipedia.org/wiki/Steinhart}{\linkhilite{Stein-Hart}}\hyperlink{stein-hart}{\linkhilite{[{\xreffont 64}]}}\hyperlink{steinhart-wikipedia}{\linkhilite{[{\xreffont 65}]}} to convert the resistance to temperature in Celsius.%
%
\begin{equation}
T_{thermistor}=\left(\frac{1}{T_0}+\frac{1}{\beta}log(R_{thermistor}/R_{series})\right)^{-1}\label{temperature-sensor-5}
\end{equation}
where \(\beta=3950\) is a heat transfer coefficient specific to the bulk semiconductor material over a given temperature range of interest and \(T_0=298.15K\) is the nominal temperature of the semiconductor in Kelvin. Note that in the equation above, the output is in Kelvin.%
\par
If this all seems complex, the folks at Adafruit have done it again with an \emph{adafruit\textunderscore{}thermistor} module. If you head over to their \href{https://github.com/adafruit/Adafruit_CircuitPython_Thermistor/blob/master/adafruit_thermistor.py}{\linkhilite{github on this module}} you'll see the relevant conversion under the definition temperature which at the time of this writing is on line 127\hyperlink{Github}{\linkhilite{[{\xreffont 33}]}}. \href{https://learn.adafruit.com/thermistor/circuitpython}{\linkhilite{The Adafruit Learn system}} also does a bit of work to explain the conversion from voltage to temperature\hyperlink{Adafruit}{\linkhilite{[{\xreffont 32}]}} but understand that the tutorial on that page assumes that the thermistor is on the low (grounded) side and thus the voltage divider equation is different. If you look at the \emph{adafruit\textunderscore{}thermistor} module you'll see that there is a \emph{self.high\textunderscore{}side} boolean that changes between either a thermistor on the high side or the low side. In this case, for the built-in A9 thermistor the thermistor is on the high side. Given all this complexity, we will just appreciate the simplicity of the code below which uses the \emph{adafruit\textunderscore{}thermistor} module. I've also created a simple version of the code to read the analog value of the thermistor and posted it \href{https://github.com/cmontalvo251/Microcontrollers/blob/master/Circuit_Playground/CircuitPython/Temp/record_temperature_thermistor.py}{\linkhilite{on my Github}}\hyperlink{Github}{\linkhilite{[{\xreffont 33}]}}. \alert{Note the code below is from a much earlier and simpler version of my code which only reads the temperature in Celsius. The code on my Github has more functionality such as printing the raw ADC value and the voltage.}%
\begin{figureptx}{Figure}{CircuitPython code to read the onboard thermistor}{temperature-sensor-8}{}%
\begin{image}{0.135}{0.73}{0.135}{}%
\includegraphics[width=\linewidth]{external/thermistor_code.png}
\end{image}%
\tcblower
\end{figureptx}%
For the code above, lines 1-3 import the relevant modules and then line 8 create the thermistor object. You'll notice the input arguments are the pin which is A9 as well as the resistor values which are in series with the thermistor. These resistors are soldered to the PCB so they are fixed at \(10 k\Omega\). The 25 is for the nominal resistance temperature in celsius of the thermistor and 3950 is the b coefficient which is a heat transfer property. Running this code and then placing my finger on the A9 symbol causes the temperature to rise just a bit. You'll notice the temperature rise quite quickly when I place my finger on the sensor but when I remove the sensor it takes some time before the sensor cools off. This has to do with the dynamic response of the sensor. We'll discuss this in some future labs on dynamic measurements. For now you can move on to the accelerometer.%
\begin{figureptx}{Figure}{Serial monitor and Plotter open in Mu showing temperature in Celsius}{temperature-sensor-10}{}%
\begin{image}{0.135}{0.73}{0.135}{}%
\includegraphics[width=\linewidth]{external/thermistor_mu.png}
\end{image}%
\tcblower
\end{figureptx}%
\end{subsectionptx}
%
%
\typeout{************************************************}
\typeout{Subsection 15.1.4 Accelerometer}
\typeout{************************************************}
%
\begin{subsectionptx}{Subsection}{Accelerometer}{}{Accelerometer}{}{}{accelerometer}
The accelerometer is a 3-axis sensor. As such, it is going to spit out not just 1 value but 3 values. Accelerations in x,y and z or North, East, Down or Forward, Side to Side, Up and Down. Since it's reading 3 values we can't just read 3 analog signals (we can but the accelerometer chip design didn't want to do that) so instead we're going to use the I2C (I like "Indigo" and 2C like "Squared C". So "I Squared C". Not "12C" or "one two C". It's "I squared C") functionality. I2C is a type of serial communication that allows computers to send strings rather than numbers. I2C is beyond the scope of this course but just know it's a type of serial communication that uses hexadecimal addresses. It's a much more complex form of communication but since it's a standard form of communication, we can just use the busio module which contains the I2C function.%
\begin{figureptx}{Figure}{CircuitPython code to read the onboard accelerometer}{accelerometer-3}{}%
\begin{image}{0.135}{0.73}{0.135}{}%
\includegraphics[width=\linewidth]{external/accelerometer_code.png}
\end{image}%
\tcblower
\end{figureptx}%
In this code we see alot more imports than normal. In addition to the standard time, board and digitalio modules we need the busio module and the \emph{adafruit\textunderscore{}lis3dh}. You might think that LIS3DH is a very weird name for an accelerometer but it's actually the name of the chip on your CPX. The chip itself is very standard and is well documented on multiple websites. \href{https://www.st.com/en/mems-and-sensors/lis3dh.html}{\linkhilite{Here's one from ST}}\hyperlink{ST}{\linkhilite{[{\xreffont 66}]}}. You can also buy the \href{https://www.adafruit.com/product/2809?gclid=EAIaIQobChMI98qfjrC86gIVEI_ICh2JtQ-JEAAYAiAAEgJ2_fD_BwE}{\linkhilite{chip on a breakout board from Adafruit}} and then of course the Adafruit Learn site has plenty of tutorials on \href{https://learn.adafruit.com/circuitpython-hardware-lis3dh-accelerometer/software}{\linkhilite{reading Accelerometer data in CircuitPython}}\hyperlink{Adafruit}{\linkhilite{[{\xreffont 32}]}}\hyperlink{AdafruitMAIN}{\linkhilite{[{\xreffont 34}]}}. As always I've learned what I can from the relevant tutorials and created \href{https://github.com/cmontalvo251/Microcontrollers/blob/master/Circuit_Playground/CircuitPython/Accelerometer/low_level_accel.py}{\linkhilite{my own simple version to read the accelerometer data}} and posted it to Github\hyperlink{Github}{\linkhilite{[{\xreffont 33}]}}. I digress, lines 8-11 of the code do alot. It first uses the SCL and SDA pins to set up an I2C object which establishes serial communication to the accelerometer. Line 9 creates an interrupt which is beyond the scope of this course. Finally, line 10 creates the actual accelerometer object by sending it the I2C pins, the hexadecimal address in the I2C protocol and finally the interrupt pin. Line 11 then sets the range. Line 14 in the while loop is where the x,y and z values of the accelerometer are read and then promptly printed to Serial on line 15. If I run this code and shake the sensor a bit I can get all the values to vary. If you put the CPX on a flat surface, the Z axis will measure something close to 9.81. The units of the accelerometer are clearly in \(m/s^2\). \begin{figureptx}{Figure}{Serial monitor and Plotter open in Mu showing accelerometer output in \(m/s^2\)}{accelerometer-4-11}{}%
\begin{image}{0.135}{0.73}{0.135}{}%
\includegraphics[width=\linewidth]{external/accelerometer_mu.png}
\end{image}%
\tcblower
\end{figureptx}%
%
\par
Note that the accelerometer can be used to obtain pitch and roll angles in benign environments. This is done via trigonometry. You'll notice that when you point the CPX directly at you with the cable pointing up the x axis reads around 0 and the y axis reads around gravity. If you then rotate the sensor so that the cable is coming out of the left side of the CPX, the x axis is reading about gravity while the y axis is zero. Similarly, rotating the CPX in the X\slash{}Z plane can be done by rotating the USB cable where it plugs into the CPX. In this case we can ignore the Y axis data. If you print the raw accelerometer data, you'll notice that when you place the CPX directly onto a flat surface, the x and y axes read a value around 0, while the z axis reads around gravity (9.81\(m/s^2\)). If you then rotate the sensor clockwise 90 degrees, the x axis is reading about gravity while the z axis is now zero. This means we can form two triangles (one for the X\slash{}Z plane and another for the Y\slash{}Z plane) and get the angles of the CPX\slash{}CPB using the equations below.%
\begin{equation}
\theta = -sin^{-1}(\bar{a}_x)\label{accelerometer-5-2}
\end{equation}
%
\begin{equation}
\phi = tan^{-1}\left(\frac{\bar{a}_y}{\bar{a}_z}\right)\label{accelerometer-5-3}
\end{equation}
In the equations above, \(a\) is the accelerations along all 3 axes. Shorthand is used for \(cos(\eta)=c_{\eta}\) and \(sin(\eta)=s_{\eta}\). The \(\bar{\eta}\) notation is used to indicate a normalization of the vector. That is \(\bar{\eta} = \eta/||\eta||\) where \(||\eta||\) is the norm of a 3 dimensional vector \(||\eta||=\sqrt{\eta_x^2+\eta_y^2+\eta_z^2}\). The derivation for these angles from accelerometers is quite involved and requires the knowledge of rotation matrices. That derivation is in another textbook called \href{https://github.com/cmontalvo251/LaTeX/tree/master/Aerospace_Mechanics}{\linkhilite{Aerospace Mechanics}}\hyperlink{Github}{\linkhilite{[{\xreffont 33}]}}.%
\end{subsectionptx}
\end{sectionptx}
%
%
\typeout{************************************************}
\typeout{Section 15.2 High Level Control}
\typeout{************************************************}
%
\begin{sectionptx}{Section}{High Level Control}{}{High Level Control}{}{}{modules-4}
Alright so we've learned the hard way for all the sensors using low level control of the various sensors. Let's now import the simple \emph{adafruit\textunderscore{}circuitplayground.express} module. \href{https://learn.adafruit.com/circuitpython-made-easy-on-circuit-playground-express/circuit-playground-express-library}{\linkhilite{The Adafruit Learn site}} offers pretty much every example code snippet you'd ever need for all the different push buttons and sensors on the CPX\hyperlink{Adafruit}{\linkhilite{[{\xreffont 32}]}}. Head over there if you ever need something outside of the scope of this text. As I said before, the main module you need to import is done by adding the following to the top of your code%
\begin{preformatted}
from adafruit_circuitplayground.express import cpx
\end{preformatted}
\alert{Note that if you have a CPB you need to change that line to adafruit\textunderscore{}circuitplayground.bluefruit import cpb. Then everywhere you see cpx you replace with cpb.} Also, keep in mind that you will need to install modules onto your CPX\slash{}CPB (See \hyperref[installing-modules]{\linkhilite{Section~{\xreffont\ref{installing-modules}}}}). Assuming you have the right module, the line above will import the cpx module into your working code. From here the commands to read different things are relatively simple. Here are the commands for all the various sensors%
\begin{preformatted}
light = cpx.light

x,y,z = cpx.acceleration

temperature = cpx.temperature
\end{preformatted}
There unfortunately is no simple module for the sound sensor on the CPX. You'll still need to use the low level control no matter what. However, if you get the Circuit Playground Bluefruit there is a \href{https://learn.adafruit.com/circuitpython-made-easy-on-circuit-playground-express/sound}{\linkhilite{simple way to read the sound level}} using \emph{cpb.sound\textunderscore{}level}\hyperlink{Adafruit}{\linkhilite{[{\xreffont 32}]}}. Implementing the various sensors into a while loop on my CPX looks like this.%
\begin{figureptx}{Figure}{Serial monitor and Plotter open in Mu showing light, temperature and acceleration with high level contorl}{modules-4-7}{}%
\begin{image}{0.135}{0.73}{0.135}{}%
\includegraphics[width=\linewidth]{external/high_level_mu.png}
\end{image}%
\tcblower
\end{figureptx}%
I left out the low level sound sensor stuff just because it kind of messes with the simplicity of the code above. The \emph{adafruit\textunderscore{}circuitplayground.express} module outputs just as before except for the light sensor. In the low level control we simply computed the voltage across the photocell but the \emph{adafruit\textunderscore{}circuitplayground.express} module outputs data in \href{https://en.wikipedia.org/wiki/Lux}{\linkhilite{Lux}}\hyperlink{lux-wikipedia}{\linkhilite{[{\xreffont 62}]}}. Recall, if you have a CPB you can use the high level \emph{cpb.sound\textunderscore{}level} command as shown in the Figure below. Notice that the import command has changed to \emph{adafruit\textunderscore{}circuitplayground.bluefruit import cpb}. \begin{figureptx}{Figure}{Example High Level Sound Level on the CPB}{cpb-sound-level}{}%
\begin{image}{0.135}{0.73}{0.135}{}%
\includegraphics[width=\linewidth]{external/sound_level.png}
\end{image}%
\tcblower
\end{figureptx}%
%
\end{sectionptx}
%
%
\typeout{************************************************}
\typeout{Section 15.3 Extra Help}
\typeout{************************************************}
%
\begin{sectionptx}{Section}{Extra Help}{}{Extra Help}{}{}{modules-5}
If you need extra help on this assignment I have uploaded a youtube video where I \href{https://www.youtube.com/watch?v=ceI0uw3ICE0}{\linkhilite{read the temperature and accelerometer}} from the CircuitPlayground Bluefruit\hyperlink{youtube}{\linkhilite{[{\xreffont 50}]}}.%
\end{sectionptx}
%
%
\typeout{************************************************}
\typeout{Section 15.4 Assignment}
\typeout{************************************************}
%
\begin{sectionptx}{Section}{Assignment}{}{Assignment}{}{}{modules-6}
\begin{introduction}{}%
Using either \emph{low or high level control}, take at least 60 seconds of data using the microphone, photocell, accelerometer and thermistor on your CPX\slash{}CPB. Make sure to log time and the raw sensor value at 1Hz or faster. To make the project more challenging, try and log all sensor data all at once (although this isn't a strict requirement).%
%
\begin{enumerate}
\item{}Title page with name,title and date - 5\%%
\item{}Introduction - 5\%%
\item{}All figures must have appropriate annotations, figure captions, figure labels and a paragraph of supporting text - 10\%%
\item{}Include 4 screenshots (entire computer screen) of Mu with the plotter open showing example data from each sensor. - 10\% per photo%
\item{}Include 4 plots with time on the x-axis and sensor data on the y-axis (No screenshots) - 10\% per plot%
\item{}Appendix A - Video Details - Filmd yourself and your entire screen with the Plotter open in Mu as you demonstrate the output of all 4 sensors. The values on the Plotter must change due to your inputs - Pass\slash{}Fail%
\item{}Appendix B - Code used to take data and code used to plot in plain text (no screenshots) - Pass\slash{}Fail%
\item{}Appendix C - Screenshot of Excel or Sheets if applicable - Pass\slash{}Fail%
\item{}\alert{If you are asked to fix something in the Appendices such as missing content in your video or screenshots of your code, you will lose 10\% of your grade everytime you submit an update after the deadline.}%
\end{enumerate}
\end{introduction}%
\end{sectionptx}
\end{chapterptx}
%
%
\typeout{************************************************}
\typeout{Chapter 16 Integrating Acceleration}
\typeout{************************************************}
%
\begin{chapterptx}{Chapter}{Integrating Acceleration}{}{Integrating Acceleration}{}{}{acceleration}
\renewcommand*{\chaptername}{Chapter}
\begin{introduction}{}%
\begin{listptx}{List}{\textbf{Parts List}}{acceleration-2-1}{}%
%
\begin{enumerate}
\item{}CPX\slash{}CPB%
\item{}USB Cable%
\item{}Laptop%
\item{}Some sort of temporary adhesive%
\item{}Automobile%
\end{enumerate}
\end{listptx}%
\begin{listptx}{List}{\textbf{Learning Objectives}}{acceleration-2-2}{}%
%
\begin{enumerate}
\item{}Taking acceleration data of real systems%
\item{}Numerical Integration%
\item{}The dangers of integration bias%
\end{enumerate}
\end{listptx}%
The main goal for this lab is for you to secure your CPX\slash{}CPB on the dashboard of your car and take acceleration data while you accelerate and decelerate. You will then plot that data and integrate it to get velocity and position. You will also compare your results to GPS data taken on your phone using PhyPhox. This is a pretty fun lab and it's really cool to see the acceleration data coming from your car. It's also really interesting to see how the velocity and position plots look after integrating the acceleration data. You will probably be surprised at how bad the velocity and position plots look after integrating the acceleration data. This is because of bias in the accelerometer data which gets integrated as well and causes the velocity and position plots to drift over time. \alert{Note that you will be using a vehicle to do this lab and you should be very careful when gathering data. Make sure to do this in a safe location and only accelerate to speeds you are comfortable with. I also require you to have a partner in the car with you while you take data so that you are not driving and taking data at the same time. I am not responsible for any damage done to your vehicle or anyone else because you are doing this project.}%
\end{introduction}%
%
%
\typeout{************************************************}
\typeout{Section 16.1 Setup}
\typeout{************************************************}
%
\begin{sectionptx}{Section}{Setup}{}{Setup}{}{}{acceleration-3}
The code for this lab is to have the CPX log acceleration. So when you're done with this lab you will hopefully have a data file with 4 columns of data: time, acceleration x, acceleration y, acceleration z. The code I'm using is the same as the lab on accelerometers (See \hyperref[modules]{\linkhilite{Chapter~{\xreffont\ref{modules}}}}). I'm using method 1 for datalogging so I'm just having it print to Serial (See \hyperref[daq]{\linkhilite{Chapter~{\xreffont\ref{daq}}}}).%
\begin{figureptx}{Figure}{Accelerometer code used print acceleration to Serial}{acceleration-3-3}{}%
\begin{image}{0.1}{0.8}{0.1}{}%
\includegraphics[width=\linewidth]{external/accelerometer_code1.png}
\end{image}%
\tcblower
\end{figureptx}%
The code on Github has a sleep of 0.1 seconds but make sure to have the CPX take data as fast as possible. A sleep of 0.01 is probably good. You will probably get a lot of data points for this experiment. Once your code is working, place the CPX on your dashboard with one of the axes of the accelerometer pointing towards the nose of your car. Try and place the CPX on as flat a surface as possible. You can use 3M tape or duct tape or hot glue. Just make sure you don't damage your car and make sure the CPX is well anchored to the dashboard. This way when the car accelerates, the CPX will measure that acceleration. Note, if you'd like to do this with a bike or some other motor vehicle that is just fine. Just make sure to take pictures and videos when you do the experiment. I suggest you do this in a parking lot for safety reasons. I am not responsible for any damage done to your vehicle or anyone else because you are doing this project. Once you have the CPX anchored, accelerate your vehicle to 20 mph (or however fast you are comfortable driving) and then slam on the brakes. Once your data is logged, plot your acceleration on your desktop computer. After doing the experiment myself, this is what my acceleration plot looks like. I had to clip the time series to only include the part from where I accelerated and decelerated quickly. I also subtracted the first data point from each accelerometer axis to zero it out and subtract off the bias. Since I took some data for a bit before I started moving I could have averaged the first few data points to obtain the bias. Instead just to get something working properly I went ahead and just used the first data point.%
\begin{figureptx}{Figure}{Example acceleration curves coming from accelerating and decelerating in a car}{acceleration-3-5}{}%
\begin{image}{0.135}{0.73}{0.135}{}%
\includegraphics[width=\linewidth]{external/accelerometer_plots.png}
\end{image}%
\tcblower
\end{figureptx}%
It's clear from the plots that the z axis was oriented towards the nose of the car. In this case I am going to have to flip the z-axis since the beginning is acceleration and the end is deceleration. I also through the acceleration in the z-axis through a complementary filter with a filter value of 0.25. I think it makes the acceleration profile a bit less jumpy. I then used a Reimann sum and integrated the acceleration data points to get velocity. The equation itself looks like this:%
\begin{equation}
V_i = \sum\limits_{n=0}^{i}{(a_i-a_0)\Delta t}\label{acceleration-3-6-1}
\end{equation}
%
\par
This of course assumes the initial velocity is zero. Notice that I take the individual acceleration points and subtract off the bias. Computing that summation by hand is pretty trivial but getting the code to work is another story. For a Reimann sum we're going to use a for loop where we loop through all the data points. The good news is that the time between data points is the same so we can just treat that as a constant. Once you have acceleration integrated you can plot velocity. This is what mine looks like after I did the experiment. According to my plot I accelerated to about 45 mph. I guess I can't lie in this instance. I said to accelerate to 20 mph but I really wanted to see a large change in acceleration so I punched it. Notice though that at the end of the time series the velocity is negative. This is because as time goes on you are integrating error and the error just gets worse.%
\begin{figureptx}{Figure}{Integrating Acceleration and plotting Velocity in Thonny}{acceleration-3-8}{}%
\begin{image}{0.135}{0.73}{0.135}{}%
\includegraphics[width=\linewidth]{external/accelerometer_integration.png}
\end{image}%
\tcblower
\end{figureptx}%
This is why speedometers are used. They are just much more accurate than integrating acceleration which is prone to bias and drift. \href{https://github.com/cmontalvo251/Microcontrollers/blob/master/Circuit_Playground/CircuitPython/cpx_assignments/Velocity_from_Acceleration/velocity_from_acceleration.py}{\linkhilite{This folder on Github has some codes that will help with your project}}\hyperlink{Github}{\linkhilite{[{\xreffont 33}]}}. \alert{Note, that some of those codes have a bias filter, truncation filter and complementary filter. That code may not work for you and you may need to tune the filters for your specific data set. Make sure to understand what each filter does and think about how it applies to your data set otherwise your code may throw an error due to the differences in your data set.}%
\end{sectionptx}
%
%
\typeout{************************************************}
\typeout{Section 16.2 Extra Help}
\typeout{************************************************}
%
\begin{sectionptx}{Section}{Extra Help}{}{Extra Help}{}{}{acceleration-4}
I explain how to get \href{https://www.youtube.com/watch?v=e4xs9Ky7_YI\&feature=youtu.be}{\linkhilite{mean and standard deviation on Youtube}} if you need help subtracting the bias from your data\hyperlink{youtube}{\linkhilite{[{\xreffont 50}]}}.%
\end{sectionptx}
%
%
\typeout{************************************************}
\typeout{Section 16.3 Assignment}
\typeout{************************************************}
%
\begin{sectionptx}{Section}{Assignment}{}{Assignment}{}{}{acceleration-5}
\begin{introduction}{}%
For this assignment you are to find a safe place to accelerate and decelerate your vehicle while recording acceleration data on the Circuit Playground. I also want you to take GPS data using PhyPhox so you can compare data. I suggest using a temporary adhesive to secure your CPX\slash{}CPB to your dashboard (making sure it's level) and have a passenger in the car to help you record data as well as operate Phyphox and verify that you have a GPS lock. \alert{For safety's sake I suggest finding an empty parking lot and only accelerating to 20 mph or less.}%
\begin{enumerate}
\item{}Title page with name,title and date - 5\%%
\item{}Introduction - 5\%%
\item{}All figures must have appropriate annotations, figure captions, figure labels and a paragraph of supporting text - 10\%%
\item{}Include a picture of your car and your passenger that is helping you record data. In your description write what speed you achieved in your experiment \alert{(If your passenger is in this class you must do the experiment twice and get two different individual sets of data)}. - 10\%%
\item{}Include a photo of your CPX\slash{}CPB mounted to your vehicle indicating which axis of acceleration is pointing forward - 10\%%
\item{}Plot one axis of your accelerometer data vs time which clearly indicates when you accelerated and decelerated. In your description be sure to explain which axis you are plotting and any signal conditioners you applied to get your clean signal \alert{(Again, if your passenger is also taking this class, you must do the experiment twice and get two different individual\slash{}unique sets of data).}. - 20\%%
\item{}Integrate acceleration and plot velocity as a function of time. Comment on whether or not the maximum velocity is the same as what you did in your actual car. Make sure to superimpose your Phyphox GPS speed on top of your plot to compare. You may need to have your raw GPS coordinates converted to velocity if it doesn't log speed natively. - 20\%%
\item{}Integrate the velocity and compute position. Plot your position as a function of time and include that in your report. Also include the position of your car using your GPS coordinates keeping in mind that you may need to convert your Lat\slash{}Lon coordinates to meters. Although you didn't measure how far you went, comment on the accuracy of the plot and whether or not you think you traveled that far especially given the superimposed data from your GPS. - 20\%%
\item{}Appendix A - Video Details - Include 2 videos: The first is of yourself and your passenger in the car taking data showing the Plotter open in Mu, the second is plotting the acceleration data and converting that to velocity \alert{(Note, if your partner is in this class they must submit their own videos with their own data)}. - Pass\slash{}Fail%
\item{}Appendix B - Code used to take data and code used to plot in plain text (no screenshots) - Pass\slash{}Fail%
\item{}Appendix C - Screenshot of Excel or Sheets if applicable - Pass\slash{}Fail%
\item{}\alert{If you are asked to fix something in the Appendices such as missing content in your video or screenshots of your code, you will lose 10\% of your grade everytime you submit an update after the deadline.}%
\end{enumerate}
%
\end{introduction}%
\end{sectionptx}
\end{chapterptx}
%
%
\typeout{************************************************}
\typeout{Chapter 17 Pedometer}
\typeout{************************************************}
%
\begin{chapterptx}{Chapter}{Pedometer}{}{Pedometer}{}{}{pedometer}
\renewcommand*{\chaptername}{Chapter}
\begin{introduction}{}%
\begin{listptx}{List}{\textbf{Parts List}}{pedometer-2-1}{}%
%
\begin{enumerate}
\item{}Laptop%
\item{}CPX\slash{}CPB%
\item{}USB Cable%
\item{}External Battery Pack%
\end{enumerate}
\end{listptx}%
\begin{listptx}{List}{\textbf{Learning Objectives}}{pedometer-2-2}{}%
%
\begin{enumerate}
\item{}Understand how to run CPB\slash{}CPX while not tethered to a computer%
\item{}Reinforce bluetooth tech for data transfer%
\item{}Understand post-processing for debugging to be used for online calculations%
\item{}Understand the fundamentals of how a pedometer works%
\end{enumerate}
\end{listptx}%
A pedometer is a device that counts the number of steps. Typically these are worn as watches with popular brands like Garmin or Fitit owning the market share at the time of this writing. It turns out though that your phone and apps like Google Fit can also track steps just by being inside your pocket all day. The way they do this is by measuring the acceleration, angular velocity and potentially even the angles (using a magnetometer) to count steps. In this lab this we will just focus on attempting to get steps using accelerometer data.%
\end{introduction}%
%
%
\typeout{************************************************}
\typeout{Section 17.1 Gathering Accelerometer Data}
\typeout{************************************************}
%
\begin{sectionptx}{Section}{Gathering Accelerometer Data}{}{Gathering Accelerometer Data}{}{}{pedometer-3}
First we need to make sure we can gather accelerometer data. The low level accelerometer code is relatively simple and is explained in the Modules lab (\hyperref[modules]{\linkhilite{Chapter~{\xreffont\ref{modules}}}}). In order to gather the accelerometer data while running you'll need to be able to operate the CPB\slash{}CPX untethered from a computer. This means you have to use Method 3,4 or 5 (See \hyperref[daq]{\linkhilite{Chapter~{\xreffont\ref{daq}}}}, \hyperref[bluetooth]{\linkhilite{Chapter~{\xreffont\ref{bluetooth}}}} or \hyperref[sdcard]{\linkhilite{Chapter~{\xreffont\ref{sdcard}}}}. Remember that Method 1 and 2 require a computer and running with a computer would be difficult if not impossible. Method 3 requires a lot of setup to log data directly to the disk and Method 5 needs extra hardware. So, for this lab we will just use the Bluetooth module to send data directly to your phone. Note that if you have a CPX you will have to use Method 3 or 5. The best way to do this experiment is with a partner. Have the CPB\slash{}CPX measure acceleration and place the entire device with a battery pack inside the runners pocket. Then have your partner connect to the CPB with the Adafruit Connect app and log data using the UART and Export to txt function. Remember not to run too far because the Bluetooth signal distance is only about 30 feet. See if you can combine the Bluetooth code and the acceleration code into one code to send time and the 3-axis accelerometer data. If you're still having trouble, code for this lab can be found on \href{https://github.com/cmontalvo251/Microcontrollers/tree/master/Circuit_Playground/CircuitPython/pedometer}{\linkhilite{Github}}\hyperlink{Github}{\linkhilite{[{\xreffont 33}]}}. Note if you have a CPX you will need to combine the accelerometer code with the Method 3 version of data logging.%
\par
Running the CPX\slash{}CPB untethered does require a few extra steps besides writing the code. The first step is obviously to write the software that you want to run on the CPX\slash{}CPB. I recommend testing the code extensively while tethered to the computer so you can debug using the REPL. Once you're certain the code works you can disconnect the CPX\slash{}CPB and connect it to a battery pack. Once again I recommend testing the code with the battery back before you perform the experiment. In this experiment I used an external USB battery bank as shown in the photo below. My code also utilizes the neopixel library to turn on some LEDs.%
\begin{figureptx}{Figure}{CircuitPlayground Bluefruit connected to an external battery pack}{pedometer-3-4}{}%
\begin{image}{0.1}{0.8}{0.1}{}%
\includegraphics[width=\linewidth]{external/External_Battery_Pack.jpg}
\end{image}%
\tcblower
\end{figureptx}%
\end{sectionptx}
%
%
\typeout{************************************************}
\typeout{Section 17.2 Computing Number of Steps: Post-Processing}
\typeout{************************************************}
%
\begin{sectionptx}{Section}{Computing Number of Steps: Post-Processing}{}{Computing Number of Steps: Post-Processing}{}{}{pedometer-4}
Using the hardware and software defined above I had my partner run down the hallway after I ensured there was a solid connection between the CPB and my smart phone. I then exported the data to a text file and plotted the raw data using Thonny.%
\begin{figureptx}{Figure}{Counting steps in Thonny}{pedometer-4-3}{}%
\begin{image}{0.135}{0.73}{0.135}{}%
\includegraphics[width=\linewidth]{external/Pedometer3.png}
\end{image}%
\tcblower
\end{figureptx}%
Upon inspecting the raw data it seems as though my partner began running around 68 seconds. At about 80 seconds my partner reached the end of the hallway and the CPB got a bit out of range. As such there is a gap in the data. The accelerometer streams return once my partner begins running back down the hallway. In order to simply look at the data of one run the data was truncated from 69 seconds to 79 seconds as shown in the Figure below.%
\begin{figureptx}{Figure}{Counting step in truncated data in Thonny}{pedometer-4-5}{}%
\begin{image}{0.135}{0.73}{0.135}{}%
\includegraphics[width=\linewidth]{external/Pedometer2.png}
\end{image}%
\tcblower
\end{figureptx}%
In order to count steps the algorithm is fairly simple and not very robust but it does at least give you a sense of how data can be analyzed to obtain steps. First, the norm of the accelerometer data is computed and then the norm is substracted by 9.81 \(m/s^2\). When glancing at the data the steps seem to be taken when the result of the norm-9.81 goes from positive to negative. It is possible that there is some aliasing in the data but for a simple experiment like this a rudimentary algorithm can be created. First the STEP counter is set to zero and then a for loop is created to loop through the data. When the data goes from positive to negative a STEP is created. This is done by using a RESET flag and checking whether or not the data becomes positive. This algorithm computes 26 steps which seems reasonable for the length of the hallway.%
\end{sectionptx}
%
%
\typeout{************************************************}
\typeout{Section 17.3 Computing Number of Steps: Online}
\typeout{************************************************}
%
\begin{sectionptx}{Section}{Computing Number of Steps: Online}{}{Computing Number of Steps: Online}{}{}{pedometer-5}
The benefit of post-processing in Python is that the CPX\slash{}CPB run a almost identical derivative of Python called CircuitPython. This means that almost any line of code used in Python can be copied directly onto the CPX\slash{}CPB as shown in the Figure below.%
\begin{figureptx}{Figure}{CircuitPython code to count steps on the CircuitPlayground Express}{pedometer-5-3}{}%
\begin{image}{0.135}{0.73}{0.135}{}%
\includegraphics[width=\linewidth]{external/Pedometer4.png}
\end{image}%
\tcblower
\end{figureptx}%
You can see that the STEPS counter is set to zero before the infinite while loop and then after connection the RESET flag and bias value are checked for a switch from positive to negative. Some code is added to change the pixels on the CPB so that the user can see a step change. The raw data is still transmitted via bluetooth but it would be a neat exercise to have the CPB transmit the number of steps to a cell phone for instant feedback of the number of steps.%
\end{sectionptx}
%
%
\typeout{************************************************}
\typeout{Section 17.4 Assignment}
\typeout{************************************************}
%
\begin{sectionptx}{Section}{Assignment}{}{Assignment}{}{}{pedometer-6}
\begin{introduction}{}%
For this assignment you are to run down the hallway while wearing your CPX\slash{}CPB and record acceleration. The acceleration signal will then be processed to compute the number of steps you took. While running, be sure to make a note on the number of steps you are taking.%
%
\begin{enumerate}
\item{}Title page with name,title and date - 5\%%
\item{}Introduction - 5\%%
\item{}All figures must have appropriate annotations, figure captions, figure labels and a paragraph of supporting text - 10\%%
\item{}Include a photo of your CPX\slash{}CPB with its own battery pack - 40\%%
\item{}Include a plot of accelerometer data vs time showing stars where the peaks are. In your description be sure to comment on how many steps the code outputted vs how many steps you took during your experiment - 40\%%
\item{}Appendix A - Video Details - Include 2 videos: First, include a cell phone video of your partner placing the system in a pocket and running or walking, then show a video with yourself and your screen as you plot the acceleration vs time with peaks shown on your computer - Pass\slash{}Fail%
\item{}Appendix B - Code used to take data and code used to plot in plain text (no screenshots) - Pass\slash{}Fail%
\item{}Appendix C - Screenshot of Excel or Sheets if applicable - Pass\slash{}Fail%
\item{}\alert{If you are asked to fix something in the Appendices such as missing content in your video or screenshots of your code, you will lose 10\% of your grade everytime you submit an update after the deadline.}%
\end{enumerate}
\end{introduction}%
\end{sectionptx}
\end{chapterptx}
%
%
\typeout{************************************************}
\typeout{Chapter 18 Statistics using Photocells}
\typeout{************************************************}
%
\begin{chapterptx}{Chapter}{Statistics using Photocells}{}{Statistics using Photocells}{}{}{photocell}
\renewcommand*{\chaptername}{Chapter}
\begin{introduction}{}%
\begin{listptx}{List}{\textbf{Parts List}}{photocell-2-1}{}%
%
\begin{enumerate}
\item{}Laptop%
\item{}CPX\slash{}CPB%
\item{}USB Cable%
\item{}Photocell%
\item{}Resistor (Any \(\Omega\) value will do)%
\item{}Alligator Clips (x3)%
\item{}Breadboard%
\end{enumerate}
\end{listptx}%
\begin{listptx}{List}{\textbf{Learning Objectives}}{photocell-2-2}{}%
%
\begin{enumerate}
\item{}Understand how a photocell responds to light level by measuring the voltage across the photocell in different light conditions%
\item{}Learn how to create histograms of a noisy data signal%
\item{}Understand mean and standard deviation and how that applies to Normal distributions.%
\end{enumerate}
\end{listptx}%
In order to prepare for this lab, be sure you are familiar with photocells and how they work (See \hyperref[light-sensor]{\linkhilite{Subsection~{\xreffont\ref{light-sensor}}}}). The only difference between \hyperref[light-sensor]{\linkhilite{Subsection~{\xreffont\ref{light-sensor}}}} and this lab is that you're going to use an external photocell which means you need to build your own voltage divider circuit. Once you have the photocell wired up to the CPX\slash{}CPB you can use the same code as the potentiometer lab to take data and then plot that data on your desktop computer. The second part of this lab is to take 1000 data points at three different light levels and create a histogram of the data as well as compute the mean, median, and standard deviation of the data. Finally, you are going to superimpose a Gaussian distribution on top of your histogram to see how close your data is to a normal distribution. This will be a good exercise in understanding the noise in your system and how that noise is distributed.%
\par
Wiring a photocell is similar to a potentiometer (See \hyperref[voltage]{\linkhilite{Chapter~{\xreffont\ref{voltage}}}}) and remember that you need a resistor in series with the photocell to create the voltage divider circuit. The photo below also shows a voltage divider circuit of a photocell using an \href{https://learn.adafruit.com/photocells/arduino-code}{\linkhilite{Arduino}} (Courtesy of \href{https://learn.adafruit.com/u/adafruit2}{\linkhilite{Lady Ada}} herself!)\hyperlink{Adafruit}{\linkhilite{[{\xreffont 32}]}}. \begin{figureptx}{Figure}{Photocell wired up to an Arduino\hyperlink{photocell-wiring}{\linkhilite{[{\xreffont 44}]}} \alert{Note that any \(\Omega\) value will do for the resistor.}}{photocell-2-4-5}{}%
\begin{image}{0.135}{0.73}{0.135}{}%
\includegraphics[width=\linewidth]{external/arduino_photocell.png}
\end{image}%
\tcblower
\end{figureptx}%
 Wiring this circuit to a CPX is almost identical to the Arduino. The only difference is that you need to use alligator clips instead of using jumper wires and a breadboard. You also need to use 3.3V instead of 5V since the CPX is running on 3.3V logic.The wiring diagram is shown below. Again, the diagram below shows a 150 \(\Omega\) resistor but you can use any value you like. \begin{figureptx}{Figure}{Wiring diagram for photocell circuit on CPX}{photocell-2-4-7}{}%
\begin{image}{0.135}{0.73}{0.135}{}%
\includegraphics[width=\linewidth]{external/photocell_wiring.png}
\end{image}%
\tcblower
\end{figureptx}%
 \alert{Note that in this case the resistor is before the photocell but it doesn't matter which way you wire it as long as the resistor and photocell are in series and you measure the voltage across the photocell. If you end up measuring the voltage across the resistor as is shown in the figure above, the voltage divider equations will change and thus effect your light level measurement. Notice the difference in equations below depending on where you measure the voltage}.%
\begin{equation}
V_{resistor} = V_{out}\frac{R_{resistor}}{R_{series}+R_{photocell}}\label{photocell-2-4-9}
\end{equation}
%
\begin{equation}
V_{photocell} = V_{out}\frac{R_{photocell}}{R_{series}+R_{photocell}}\label{photocell-2-4-10}
\end{equation}
%
\end{introduction}%
%
%
\typeout{************************************************}
\typeout{Section 18.1 Taking Data}
\typeout{************************************************}
%
\begin{sectionptx}{Section}{Taking Data}{}{Taking Data}{}{}{photocell-3}
For this lab you need to wire up the circuit, take data at the low and high value of the photocell by covering the sensor with your finger and then shining a light on it and plotting the entire data set on your desktop computer. The wiring diagram was already presented and when I did the lab my example circuit is shown below. An alligator clip is connected to 3.3V on the CPX and the other end is connected to either end of the photocell. The photocell is then in series with a resistor and the other end of the resistor is connected to GND via another alligator clip. Finally, another alligator clip from pin A2 is plugged into the same row on the breadboard as the resistor and photocell. \alert{Again any \(\Omega\) value will do for the resistor. I used a 10k resistor but you can use a 1k or even a 330 resistor and it will work just fine. The only thing that will change is the range of voltages you get across the photocell. Recall that the equations for the photocell in a voltage divider are in \hyperref[light-sensor]{\linkhilite{Subsection~{\xreffont\ref{light-sensor}}}}.}%
\begin{figureptx}{Figure}{CircuitPlayground Express connected to a Photocell}{photocell-3-3}{}%
\begin{image}{0.135}{0.73}{0.135}{}%
\includegraphics[width=\linewidth]{external/photocell_circuit.jpeg}
\end{image}%
\tcblower
\end{figureptx}%
Once you have the circuit wired properly you can use the \href{https://github.com/cmontalvo251/Microcontrollers/blob/master/Circuit_Playground/CircuitPython/Analog/analog_simple.py}{\linkhilite{same code as the potentiometer lab}}\hyperlink{Github}{\linkhilite{[{\xreffont 33}]}}. The example screenshot below shows the analog signal below showing a high spike where I placed a flashlight over the photocell and then a low spot where I covered the photocell with my finger. Remember that you can use any Analog pin on the CPX provided you change line 5 to the same pin.%
\begin{figureptx}{Figure}{Serial monitor open showing Photocell output in Mu}{photocell-3-5}{}%
\begin{image}{0.135}{0.73}{0.135}{}%
\includegraphics[width=\linewidth]{external/photocell_mu.png}
\end{image}%
\tcblower
\end{figureptx}%
Once you've gotten some example data you can plot the result as you did for the potentiometer lab. Here's what your plot may look like.%
\begin{figureptx}{Figure}{Photocell data plotted in Thonny}{photocell-3-7}{}%
\begin{image}{0.135}{0.73}{0.135}{}%
\includegraphics[width=\linewidth]{external/photocell_plots.png}
\end{image}%
\tcblower
\end{figureptx}%
\end{sectionptx}
%
%
\typeout{************************************************}
\typeout{Section 18.2 Statistics of your Data}
\typeout{************************************************}
%
\begin{sectionptx}{Section}{Statistics of your Data}{}{Statistics of your Data}{}{}{photocell-4}
If you allow the light source to be constant you'll notice that the data is quite noisy. What I'd like you to do for part 2 of this lab is take 1000 data points with the photocell with as constant of a light source as possible. Do this for three different light ranges. Low Light, ambient light and then sunlight if you're outside. If you're completing this assignment late you can use a flashlight as your third light source. With the three different data streams, create a histogram of the data with appropriate labels and compute the mean, median, and standard deviation of the data stream. Here is my example code showing code to get mean, median, and standard deviation as well as create the histogram. Notice in my code I imported the statistics module to compute the mode. Although it worked in my code, it's not typical to compute the mode of a continuous variable because often times you will not ever get the same value twice. Still, feel free to compute the mode if you so desire.%
\begin{figureptx}{Figure}{Histogram of photocell data}{photocell-4-3}{}%
\begin{image}{0.135}{0.73}{0.135}{}%
\includegraphics[width=\linewidth]{external/histogram.png}
\end{image}%
\tcblower
\end{figureptx}%
Again make sure to convert to voltage and then Lux before you plot that way you can see what the noise level is in volts and Lux. When I ran this experiment for a second time my CPX started and stopped 3 separate times. You'll see in the time series plot below that the voltage dipped in the first set and the second data set had some weird bumps probably from me changing tabs on my chrome tab. The photocell was close to my computer so that effected it. Thankfully the 3rd data set looked pretty good.%
\begin{figureptx}{Figure}{Example of bad photocell data}{photocell-4-5}{}%
\begin{image}{0.135}{0.73}{0.135}{}%
\includegraphics[width=\linewidth]{external/photocell_bad.png}
\end{image}%
\tcblower
\end{figureptx}%
The only problem with the 3rd data set is that I put my hand over it for testing purposes. Because of that I had to remove those outliers. To do that I computed the current mean and standard deviation and then threw out all data points that were 3 standard deviations away from the mean. The code looks like this.%
\begin{preformatted}
##COMPUTE CURRENT MEAN AND DEV 

mean = np.mean(voltage) 

dev = np.std(voltage) 

print(mean,dev) 

time = time[voltage > mean - 3*dev]

voltage = voltage[voltage > mean - 3*dev] 

time = time[voltage < mean + 3*dev] 

voltage = voltage[voltage < mean + 3*dev]

###COMPUTE NEW MEAN,STD 

mean = np.mean(voltage) 

dev = np.std(voltage) print(mean,dev)
\end{preformatted}
Once I did all that clean up I was able to get a nice time series plot of my data.%
\begin{figureptx}{Figure}{Photocell data plotted as a histogram with Gaussian Distribution plotted on top}{photocell-4-9}{}%
\begin{image}{0.135}{0.73}{0.135}{}%
\includegraphics[width=\linewidth]{external/histogram_good.png}
\end{image}%
\tcblower
\end{figureptx}%
I also was able to plot the Normal Gaussian Distribution on top of the histogram. You can see that in the left plot in orange. The code to do that is shown below where the 72 in the plot is the "height" of the histogram. Note that your histogram will have a different height and you will need to get that specifically from your plot.%
\begin{preformatted}
###COMPUTE THE NORMAL DISTRIBUTION 

x = np.linspace(-3*s+mu,3*s+mu,100) 

pdf = 1.0/(s*np.sqrt(2*np.pi))*np.exp((-(x-mu)**2)/(2.0*s**2)) * (s*np.sqrt(2*np.pi)) * 72
\end{preformatted}
The equations above make a time series from +-3 standard deviations from the mean and then plot the PDF of a normal Gaussian distribution. The only extra thing you have to do is multiply by (s\textasteriskcentered{}np.sqrt(2\textasteriskcentered{}np.pi)) \textasteriskcentered{} 72 which first causes the height of the PDF to be 1 and then multiply by 72 which again is the height of the histogram which will be different for your system.%
\end{sectionptx}
%
%
\typeout{************************************************}
\typeout{Section 18.3 Extra Help}
\typeout{************************************************}
%
\begin{sectionptx}{Section}{Extra Help}{}{Extra Help}{}{}{photocell-5}
\href{https://www.youtube.com/watch?v=bfeJfAWTqzY\&list=PL_D7_GvGz-v1RsDs_OdNW65qRjEjmpfQx\&index=22\&t=0s}{\linkhilite{Creating a histogram in Python}} is fairly simple as explained above but I have a Youtube Vide to supplement this tutorial. I also have another video where I get \href{https://www.youtube.com/watch?v=e4xs9Ky7_YI\&list=PL_D7_GvGz-v1RsDs_OdNW65qRjEjmpfQx\&index=20\&t=0s}{\linkhilite{mean and median values for accelerometer data}}\hyperlink{youtube}{\linkhilite{[{\xreffont 50}]}}.%
\end{sectionptx}
%
%
\typeout{************************************************}
\typeout{Section 18.4 Assignment}
\typeout{************************************************}
%
\begin{sectionptx}{Section}{Assignment}{}{Assignment}{}{}{photocell-6}
\begin{introduction}{}%
For this assignment you are to build another voltage divider circuit using a photocell and measure the voltage across the photocell using the analog to digital converter on the CPB\slash{}CPX. First take 30 seconds of data where you cover the photocell with your hand and also shine a flashlight on it. This will represent the min and max values of Lux as well as the min and max values of the resistance in your photocell. After that you are to run some statistics on the photocell by taking 1000 data points three separate times at three different light levels. The three light levels are low light (hand covered), ambient light (office light or shade light), and high light (flashlight or sunlight). For each set of 1000 data points, create a histogram, compute the mean, median, standard deviation and plot a Gaussian distribution curve on top of the histogram.%
%
\begin{enumerate}
\item{}Title page with name,title and date - 5\%%
\item{}Introduction with wiring diagram - 5\%%
\item{}All figures must have appropriate annotations, figure captions, figure labels and a paragraph of supporting text - 10\%%
\item{}Include a photo of your circuit and write the value of the resistor that is in series with your photocell. Again, the value of the resistor can be whatever you want. Also, explain how you varied the light levels - 10\%%
\item{}Plot voltage, resistance (converted from voltage) and Lux (converted from resistance) of the photocell vs time for your first 30 seconds of data where you randomly vary light conditions. - 20\%%
\item{}Include the mean, median, and standard deviation of all light levels in Lux - 10\%%
\item{}Using the mean, standard deviation and number of samples, compute your 90\% confidence interval for your 3 light levels - 10\%%
\item{}Similarly, using the mean and standard deviation of your low light level, compute the probability that you would get a reading greater than the mean plus 1 standard deviation of your low light level data set \(P(L>\mu+\sigma)\). Remember to show your work - 10\%%
\item{}Include 3 histogram plots of low light, ambient light and high light levels in units of Lux. On top of the histogram plot the normal Gaussian distribution to see how close your histogram is to a Gaussian distribution - 20\%%
\item{}Appendix A - Video Details - Film yourself and your entire screen with the Plotter open in Mu as you change the light level by putting your hand over the photocell and then shining a light on the photocell. Then show yourself plotting a histogram - Pass\slash{}Fail%
\item{}Appendix B - Code used to take data and code used to plot in plain text (no screenshots) - Pass\slash{}Fail%
\item{}Appendix C - Screenshot of Excel or Sheets if applicable - Pass\slash{}Fail%
\item{}\alert{If you are asked to fix something in the Appendices such as missing content in your video or screenshots of your code, you will lose 10\% of your grade everytime you submit an update after the deadline.}%
\end{enumerate}
\end{introduction}%
\end{sectionptx}
\end{chapterptx}
%
%
\typeout{************************************************}
\typeout{Chapter 19 Servo Control}
\typeout{************************************************}
%
\begin{chapterptx}{Chapter}{Servo Control}{}{Servo Control}{}{}{servo}
\renewcommand*{\chaptername}{Chapter}
\begin{introduction}{}%
\begin{listptx}{List}{\textbf{Parts List}}{servo-2-1}{}%
%
\begin{enumerate}
\item{}CPX\slash{}CPB%
\item{}USB Cable%
\item{}Laptop%
\item{}protractor or piece of paper%
\item{}servo%
\item{}Alligator clips (x3)%
\end{enumerate}
\end{listptx}%
\begin{listptx}{List}{\textbf{Learning Objectives}}{servo-2-2}{}%
%
\begin{enumerate}
\item{}Understand what a PWM signal is and how it affects a servo%
\item{}Understand the inner workings of a servo and how to make it move%
\item{}Practice first order regression and calibration techniques%
\item{}Understand how to compute roll and pitch from an accelerometer%
\item{}Learn the fundamental concepts of feedback control%
\item{}Learn how to Combine two codes (servo code and accelerometer code into one)%
\end{enumerate}
\end{listptx}%
For this lab we are going to learn how to drive a \href{https://www.adafruit.com/product/1143}{\linkhilite{servo}}\hyperlink{AdafruitMAIN}{\linkhilite{[{\xreffont 34}]}}. Servos accept what are called Pulse-Width-Modulation(PWM) signals which are basically square waves of varying frequency. The servo itself has a microprocessor on board that turns a DC motor based on the incoming PWM frequency which is typically called a duty cycle. The DC motor runs through some gears to rotate a shaft. Because of this rotation, you can make a number of things move! Servos are used for all sorts of things, opening doors, deflecting control surfaces on RC aircraft and many more. The neat thing about PWM signals is that they can not only move servos but they can also drive electronic speed controllers (ESCs) to turn 3 phase motors and even change the light intensity of LEDs. Servos typically come in two different color schemes as shown below.%
\begin{figureptx}{Figure}{Example servos}{servo-2-4}{}%
\begin{image}{0.1}{0.8}{0.1}{}%
\includegraphics[width=\linewidth]{external/servo.jpeg}
\end{image}%
\tcblower
\end{figureptx}%
The servo on the right is a micro servo and the one on the left is a standard servo. The micro servo is smaller and typically has less torque than the standard servo. For this lab you can use either one but I recommend using the micro servo since it draws less current and thus is less likely to cause your CPX to brown out. If you have a standard servo you may want to consider powering it with an external power supply instead of the CPX. The wiring for both servos is the same so it doesn't matter which one you use.%
\par
For this lab we are going to first learn how to change the PWM frequency on the CPX\slash{}CPB and turn the servo. Then we are going to experimentally find the limits of the servo and create a calibration curve to relate servo angle (deg) to PWM signal (ms). Finally, we are going to use the accelerometer on the CPX\slash{}CPB to measure pitch and roll (See \hyperref[accelerometer]{\linkhilite{Subsection~{\xreffont\ref{accelerometer}}}}) and use that information to feedback to the servo. This is a very common control problem in aerospace engineering and is a great example of how to combine different pieces of code together to create a more complex piece of software.%
\par
As can be easily seen when you look at the servo, servos have 3 pins (GND, PWR, SIGNAL). A wiring diagram of a servo connected to the \href{https://learn.adafruit.com/adafruit-circuit-playground-express/circuitpython-servo}{\linkhilite{CPX\slash{}CPB}} is shown below (Couresy of \href{https://learn.adafruit.com/u/kattni}{\linkhilite{Kattni Rembor}})\hyperlink{Adafruit}{\linkhilite{[{\xreffont 32}]}}.%
\begin{figureptx}{Figure}{Wiring diagram of a servo connected to a CircuitPlayground Express\hyperlink{servo-wiring}{\linkhilite{[{\xreffont 45}]}}}{servo-2-8}{}%
\begin{image}{0.1}{0.8}{0.1}{}%
\includegraphics[width=\linewidth]{external/circuitpython_CPXServo_bb.jpg}
\end{image}%
\tcblower
\end{figureptx}%
The brown or black wire is GND, the red wire needs to go to a 5V signal so for the CPX it needs to go to the VOUT pin and the yellow or white wire is the signal wire which need to go to an analog port on the CPX that supports PWM signals. Which ones support PWM signals? \href{https://learn.adafruit.com/adafruit-circuit-playground-express/pinouts}{\linkhilite{Adafruit Learn}} has a great description of which pins support PWM signals but as a quick check you can use the following analog pins: A1, A2, A3, A6 and A7.\hyperlink{Adafruit}{\linkhilite{[{\xreffont 32}]}}.%
\end{introduction}%
%
%
\typeout{************************************************}
\typeout{Section 19.1 Moving the Servo}
\typeout{************************************************}
%
\begin{sectionptx}{Section}{Moving the Servo}{}{Moving the Servo}{}{}{servo-3}
To start, make sure you follow the wiring diagram in the introduction and wire up your servo to the CPX\slash{}CPB. In my circuit I just picked pin A2 since I've been using it so much in the past. \alert{Very important: It is recommended to power a servo through an external power supply instead of the CPX. Servos can draw a lot of current and the CPX although it supports 5V can only provide so much power (P). Remember that P = VI so if P is low it means current is low. If the servo pulls more current than the CPX can provide the CPX will \textquotedblleft{}brown out\textquotedblright{} which means it will go into a safe-mode setting. If you have the AA power supply you may consider doing that. For small servos you hopefully won't have any issues. However, on the off chance your servo does put you into safe mode I recommend making small movements in your code to drive your servo. For example, if moving the servo directly from 0 to 180 puts your system into safe mode, try moving from 0 to 45 degrees instead.} \begin{figureptx}{Figure}{Photo of a servo connected to a CircuitPlayground Express}{servo-3-2-2}{}%
\begin{image}{0.1}{0.8}{0.1}{}%
\includegraphics[width=\linewidth]{external/servo_circuit.jpeg}
\end{image}%
\tcblower
\end{figureptx}%
 As I said before, a servo takes in a square wave. The square wave has a duty cycle in units of microseconds. If you send a roughly 500 us square wave to a servo it will rotate all the way to the left. If you send a roughly 2500 us signal to the servo, it will turn all the way to the right. The code to send PWM signals has been thoroughly explained in the \href{https://learn.adafruit.com/using-servos-with-circuitpython/overview}{\linkhilite{Adafruit Learn system}}\hyperlink{Adafruit}{\linkhilite{[{\xreffont 32}]}}. I also have a \href{https://github.com/cmontalvo251/Microcontrollers/blob/master/Circuit_Playground/CircuitPython/Servo/servo.py}{\linkhilite{simple servo.py script}} on Github\hyperlink{Github}{\linkhilite{[{\xreffont 33}]}}. \begin{figureptx}{Figure}{CircuitPython code to move a servo}{servo-3-2-7}{}%
\begin{image}{0.1}{0.8}{0.1}{}%
\includegraphics[width=\linewidth]{external/servo_code.png}
\end{image}%
\tcblower
\end{figureptx}%
 In this code as usual the top 3 lines are used to import the necessary modules. The pulseio module is used here to create a servo object on line 6 by connecting to pin A2. Make sure to change the pin to whatever pin you have the signal wire hooked up to. Lines 9-12 create a function that pulse in milliseconds and compute the duty cycle of PWM signal. Lines 16-19 then kick off an infinite while loop where a 800 us signal is sent to the servo and then a 2000 us signal is sent using a for loop which starts on line 16. You'll see servo command on line 18 which is responsible for sending the microsecond signal to the servo. The function servo\textunderscore{}duty\textunderscore{}cycle converts the pulse in milliseconds to a duty cycle. The value is then passed to the attribute of the servo object servo.duty\textunderscore{}cycle. If you put this code on the CPX and run the code you will hopefully see your servo turning left and right in 1 second intervals.%
\par
Besides making the servo move back and forth I'd like you to vary the pulses on line 16 \emph{SLOWLY} until the servo can't move any farther. This line of code is a for loop which loops through the array currently showing \emph{[0.8,2.0].} If you change that array to \emph{[0.9,1.2,1.5,1.8]} the servo will move to a pulse in milliseconds of \emph{0.9, 1.2, 1.5 and then 1.8}. The for loop is a great way to loop through multiple commands. Using this array, determine the minimum pulse you can send to the servo and the maximum pulse you can send to the servo. If the servo makes a funny noise it means you sent a signal outside the bounds so try a different signal. Hence the need for moving the pulse signal slowly. If you change the array to just 1 number \emph{[0.8]} the servo will just move to 1 angle and stay there forever. \alert{NOTE THAT IN CIRCUITPYTHON VERSION 7.0.0 YOU HAVE TO USE THE PWMIO LIBRARY INSTEAD OF THE PULSEIO LIBRARY}.%
\end{sectionptx}
%
%
\typeout{************************************************}
\typeout{Section 19.2 Servo Calibration}
\typeout{************************************************}
%
\begin{sectionptx}{Section}{Servo Calibration}{}{Servo Calibration}{}{}{servo-calibration}
If you notice, sending a pulse signal in microseconds moved the servo to a specific angle. Thing is I would like to be able to move the servo to a specific angle rather than having to just guess and check like we did in the last lab. So what we're going to do is start at the minimum pulse signal you computed in the last project and then change the servo pulse in equal increments until we reach the maximum servo pulse signal. When I did this lab I found that 0.6 ms was just about the smallest I could get the servo to move. We're going to call this 0 degrees. My maximum pulse signal ended up being about 2.4 ms. So I want you to test 10 different points between your specific maximum and minimum value which will hopefully be different for all of you. Everytime you test a pulse I want to measure the angle the servo makes with the minimum value being 0 degrees. Create a table of data with two columns. In the first column put pulse in milliseconds and in the second column put angle of servo in degrees. Use a protractor to measure the angle. If you don't have a protractor you can actually build one or just pull up a protractor on your computer screen and use that. I did this project with just 3 data points and here are my data points. Again you need to have around 10 data points.%
\begin{tableptx}{Table}{\textbf{Example Data taken for Calibration of Servo}}{table-servo-calibration}{}%
\centering%
{\tabularfont%
\begin{tabular}{AcAcA}\hrulethin
Pulse (ms)&Angle (Degrees)\tabularnewline\hrulethin
0.6&0\tabularnewline\hrulethin
1.5&90\tabularnewline\hrulethin
2.4&180\tabularnewline\hrulethin
\end{tabular}
}%
\end{tableptx}%
Take your table of data and put it into a spreadsheet and save the data as a CSV or simply put your data into a text file. Since I only had 3 data points I just put them into a text file. Plot the data on your desktop computer with servo angle on the x-axis and duty cycle on the y axis. Using the data, determine if the data set is linear, quadratic or cubic. Fit a trend line to the data and plot your trend line on top of the data. I made a helpful \href{https://github.com/cmontalvo251/Python/blob/master/instrumentation/least_squares_regression.py}{\linkhilite{python script with some fictitious data on Github}} that fits the data with linear and quadratic fits\hyperlink{Github}{\linkhilite{[{\xreffont 33}]}}. Here is my data plotted alongside the trendline in Python. I sort of made up the data and made it perfect on purpose so my trendline is perfect. Yours will not be so perfect.%
\begin{figureptx}{Figure}{Example trendline fit to example servo calibration data}{servo-calibration-5}{}%
\begin{image}{0.1}{0.8}{0.1}{}%
\includegraphics[width=\linewidth]{external/trendline.png}
\end{image}%
\tcblower
\end{figureptx}%
You'll notice that I first import the data from the text file using the np.loadtxt function and then I use the polyfit and polyval functions to create the trendline. The polyfit function requires you to give it the X and Y axes and the order of the trendline which since the trend line is linear I sent it a 1 but you could easily do 2 for quadratic or 3 for cubic. I then print the coefficients which are [0.01 0.6]. This means my trend line looks like this.%
\begin{equation}
Pulse = 0.6 + 0.01*Angle\label{servo-calibration-6-1}
\end{equation}
%
\par
Where Pulse is in ms and Angle is in degrees. Now you have an equation where you can use angle in degrees to compute the pulse in milliseconds. In the code I've posted I then use np.linspace to create 1000 data points from 0 to 180 degrees and then use the polyval function to compute the pulse for all 1000 angles I created using the linspace command. I finally plot the trend line in red and use the remaining part of the script to create labels and legends. Once you have this plot write down your coefficients and create an equation like I did above. Again your numbers won't be so neat. Once you have this equation, return to Mu and create a function using the def keyword that takes in an angle as an input and then returns a pulse signal in milli seconds. It will look sometime like this%
\begin{preformatted}
def angle2pulse(angle):
        return 0.6 + 0.01*angle
\end{preformatted}
\alert{Note: Functions in python must be after all your imports but before your while loop. If you put this function inside your while loop the code will not work}. Using that equation, modify your servo.py script to have the servo move through the following angles, 0, 45,90,135 (your servo may not travel to 180 degrees). Verify that your equation is working correctly by placing your protractor below the servo. Much of the code required for this project is not included because it is left as an exercise for the student.%
\end{sectionptx}
%
%
\typeout{************************************************}
\typeout{Section 19.3 Feedback Control}
\typeout{************************************************}
%
\begin{sectionptx}{Section}{Feedback Control}{}{Feedback Control}{}{}{servo-5}
\begin{introduction}{}%
Feedback control can be its own course or multiple courses but can be broken down into a few simple steps. The goal of feedback control is to drive the \textquotedblleft{}state\textquotedblright{} of a \textquotedblleft{}system\textquotedblright{} to a desired \textquotedblleft{}command\textquotedblright{} by sending a \textquotedblleft{}control\textquotedblright{} signal to the \textquotedblleft{}system\textquotedblright{}. The figure below shows a standard block diagram for a dynamic system.%
\begin{figureptx}{Figure}{Simple feedback control block diagram}{servo-5-2-2}{}%
\begin{image}{0.2}{0.6}{0.2}{}%
\includegraphics[width=\linewidth]{external/feedback.png}
\end{image}%
\tcblower
\end{figureptx}%
In this figure, R is the reference signal or the commanded signal. The reference signal feeds into a summation block where the reference signal is substracted from the measured signal. The measured signal is the output of the H block which represents the sensor. The output of the summation block is the error signal E. The error signal feeds into the controller C and outputs the control signal U to the plant system G. The output of the plant G is Y or the state. The state Y feeds into the sensor block H to measure this signal and out the measured signal to the summation block.%
\end{introduction}%
%
%
\typeout{************************************************}
\typeout{Subsection 19.3.1 Example Feedback Control System}
\typeout{************************************************}
%
\begin{subsectionptx}{Subsection}{Example Feedback Control System}{}{Example Feedback Control System}{}{}{servo-5-3}
An example feedback control system could be employed for the longitudinal pitch dynamics of an airplane. In this case the aircraft pitches through the angle \(\theta\) and the elevator \(\delta_e\) is used to control the pitch angle. Typically the pitch angle is measured with an angle sensor or even an accelerometer. The elevator is then typically controlled with a servo. Thus, for this lab we are going to make the bare bones circuitry required for pitch angle control of an airplane. The \textquotedblleft{}system\textquotedblright{} then is the aircraft. The \textquotedblleft{}state\textquotedblright{} is the \textquotedblleft{}pitch\textquotedblright{} angle (measured by the CPX), the \textquotedblleft{}command\textquotedblright{} is the desired pitch angle (programmed by you the pilot in command) and the \textquotedblleft{}control\textquotedblright{} signal is the elevator command (actuated by a servo). I'm using the same circuit I created in parts 1 and 2.%
\par
To measure the pitch and roll angles of the CPX\slash{}CPB we're going to use the accelerometer values and convert them to angles using the same equations discussed in the accelerometer lab (See \hyperref[accelerometer]{\linkhilite{Subsection~{\xreffont\ref{accelerometer}}}}).%
\par
The measurement of the pitch angle would represent the H block in the figure above. The elevator on an airplane is a control surface responsible for pitching the aircraft up and down. For this example we are going to assume that the desired pitch angle is 0\(^o\). This means our \textquotedblleft{}error\textquotedblright{} signal is going to be 0 minus the pitch angle \(e = \theta_C - \theta = 0 - \theta = -\theta\) . Our \textquotedblleft{}control\textquotedblright{} signal will be the angle of the elevator. As I said, we are going to use a servo to control the elevator so this just means we need a way to relate our \textquotedblleft{}error\textquotedblright{} signal to servo pulse width. There are a few steps here before we can move on. First, we need to relate our \textquotedblleft{}error\textquotedblright{} signal to the \textquotedblleft{}control\textquotedblright{} signal which will be the elevator pitch angle. I've made a table below to explain what I mean.%
\begin{tableptx}{Table}{\textbf{Table of relationship between pitch angle, error signal and elevator deflection angle}}{servo-5-3-5}{}%
\centering%
{\tabularfont%
\begin{tabular}{AcAcAcA}\hrulethin
Pitch (deg)&Error (deg)&Elevator (deg)\tabularnewline\hrulethin
-90&90&+90\tabularnewline\hrulethin
0&0&0\tabularnewline\hrulethin
+90&-90&-90\tabularnewline\hrulethin
\end{tabular}
}%
\end{tableptx}%
This table basically says that if the aircraft is level with a pitch angle of zero I want the elevator to be zero as well. If the aircraft pitches down, I want the elevator to pitch up and counteract that rotation. Using these three data points I can create a simple equation to relate elevator angle to pitch angle.%
\begin{equation}
\delta_e = -\theta\label{e_kpeasy}
\end{equation}
%
\par
This is a simplified version of proportional control however. In reality the control signal should be given by the equation below.%
\begin{equation}
\delta_e = Ce\label{servo-5-3-7-1}
\end{equation}
where C is the controller and e is the error signal. If we replace the error signal with \(e=\theta_c-\theta\) and the controller with just proportional gain \(k_p\) we arrive at the equation below.%
\begin{equation}
\delta_e = kp(\theta_c-\theta)\label{servo-5-3-7-4}
\end{equation}
%
\par
It's easy to see that if \(k_p=1\) and \(\theta_c=0\) the equation for the elevator would simplify to \hyperref[e_kpeasy]{\linkhilite{({\xreffont\ref{e_kpeasy}})}}. Now that we have the elevator pitch angle we need to relate this to the servo angle. Servos can only move from 0 to 180 degrees which means we can't have the servo go negative. Thus we need to offset the elevator angle to the servo angle. Again we can make a table here.%
\begin{tableptx}{Table}{\textbf{Table showing relationship between Elevator deflection angle and Servo angle}}{servo-5-3-9}{}%
\centering%
{\tabularfont%
\begin{tabular}{AcAcA}\hrulethin
Elevator (deg)&Servo (deg)\tabularnewline\hrulethin
-90&0\tabularnewline\hrulethin
0&90\tabularnewline\hrulethin
+90&180\tabularnewline\hrulethin
\end{tabular}
}%
\end{tableptx}%
This also results in a simple equation to relate servo angle \((s)\) to elevator angle \((\delta_e)\).%
\begin{equation}
s = \delta_e + 90\label{servo-5-3-10-3}
\end{equation}
%
\par
Finally, we can then use our calibration coefficients (See \hyperref[servo-calibration]{\linkhilite{Section~{\xreffont\ref{servo-calibration}}}}) to relate servo angle to pulse width. When I calibrated my servo I obtained the following equation where \(\mu\) is the pulse in PWM.%
\begin{equation}
\mu = 0.6 + 0.01s\label{servo-5-3-11-3}
\end{equation}
%
\par
So now I have an equation that relates the pitch angle from the airplane \(\theta\) to the elevator deflection angle \(\delta_e\). I can then relate the servo deflection angle to servo deflection angle (\(s\)) and then finally the servo deflection angle to PWM signal \(\mu\). With these 3 equations I can now program my servo to respond to changes in the pitch angle of the CPX.%
\par
Using the accelerometer to measure pitch and the servo to deflect the elevator I can put the entire system together. Pictured below is my code which again is also \href{https://github.com/cmontalvo251/Microcontrollers/blob/master/Circuit_Playground/CircuitPython/Servo/feedback_control_servo_pwmio.py}{\linkhilite{online on Github}}\hyperlink{Github}{\linkhilite{[{\xreffont 33}]}}. \alert{Note that the version on Github is constantly edited and as such will be slightly different than the version below.}%
\begin{figureptx}{Figure}{Start of feedback control code in CircuitPython}{servo-5-3-14}{}%
\begin{image}{0.1}{0.8}{0.1}{}%
\includegraphics[width=\linewidth]{external/feedback1.png}
\end{image}%
\tcblower
\end{figureptx}%
The first 22 lines here will hopefully seem familiar. Line 1-7 are import commands of all the various modules needed. Lines 10-13 create the definition that converts pulse width to duty cycle. Line 16 creates the servo and lines 19-22 create the accelerometer. Hopefully this is a good example of combining different codes together to get a more complex piece of software. Lines 24-45 include a very long while loop. I will try and go through each line. Line 26 grabs the accelerometer data on the CPX. Line 28 uses the x and z axis accelerometer data and converts the values to pitch angle using the atan2 function in the math module which was imported on line 4. Line 30 computes the elevator pitch angle and line 32 computes the servo deflection angle. Line 34-37 is a type of signal conditioner called a saturation filter. Basically, I don't want the servo to break because I tried to make the servo rotate more than 180 degrees or less than 0 degrees. So I created two if statements that restrict the servo to be within these two values. If the servo angle is less than 0 as stated on line 34, the servo angle is set to 0 on line 35. If the servo angle is greater than 180 as stated in line 36 the servo angle is set to 180.%
\begin{figureptx}{Figure}{Second part of feedback cntrol code in CircuitPython}{servo-5-3-16}{}%
\begin{image}{0.2}{0.6}{0.2}{}%
\includegraphics[width=\linewidth]{external/feedback2.png}
\end{image}%
\tcblower
\end{figureptx}%
Line 39 uses the calibration equation from the previous experiment to convert servo angle to pulse width. You'll need to replace these numbers with your servo since all servos are different. Line 41 uses the definition created on lines 10-13 to convert pulse width to duty cycle. Line 43 makes the servo move. Line 44 prints everything to Serial for debugging purposes and line 45 pauses the script for 0.1 seconds which helps with some twitchiness in the servo. When I did this I didn't have to program a complementary filter so I guess the servo may have it's own low pass filter. Either way this circuit is ready to be placed on an aircraft. Whether or not it is effective is a completely different story.%
\end{subsectionptx}
%
%
\typeout{************************************************}
\typeout{Subsection 19.3.2 Second Order Pitch Dynamics}
\typeout{************************************************}
%
\begin{subsectionptx}{Subsection}{Second Order Pitch Dynamics}{}{Second Order Pitch Dynamics}{}{}{second-order-dynamics}
We could take an entire aircraft design course or an undergraduate controls and systems dynamics course but to start let's just simulate this control system and see how it does. First we need to write the dynamic equation for pitch of an aircraft. That is shown below.%
\begin{equation}
\frac{J}{q_{\infty}S\bar{c}}\ddot{\theta} - \frac{\bar{c}}{2u_0}C_{mq}\dot{\theta} - C_{m\alpha}\theta = C_{m\delta_e}\delta_e + C_{m0}\label{e_second-order}
\end{equation}
%
\par
In the equation above \(J\) is the inertia of the aircraft, \(C_{m0},C_{mq},C_{m\alpha}\) and \(C_{m\delta_e}\) are non-dimensional coefficients, \(\bar{c}\) is the mean aerodynamic chord, \(u_0\) is the equilibrium velocity, \(\bar{c}\) is the mean aerodynamic chord, S is the planform area of the wing and \(q_{\infty}\) is the dynamic pressure. The dynamic pressure is given by the equation below%
\begin{equation}
q_{\infty} = \frac{1}{2}\rho {u_0}^2\label{second-order-dynamics-3-8}
\end{equation}
where \(\rho\) is the air density. To simulate this system use the table below.%
\begin{tableptx}{Table}{\textbf{Table of parameters used to simulate second order pitch response of an example radio controlled aircraft}}{second-order-dynamics-4}{}%
\centering%
{\tabularfont%
\begin{tabular}{AcAcAcAcAcAcA}\hrulethin
\(C_{m\alpha}\)&-2.19&\(C_{mq}\)&-24.45&\(C_{m\delta_e}\)&-1.15\tabularnewline\hrulethin
\(\bar{c}\)&0.2286~m&S&0.34~\(m^2\)&J&0.1232~\(kg-m^2\)\tabularnewline\hrulethin
\(u_0\)&15~m\slash{}s&\(\rho\)&1.225~\(kg/m^3\)&\(C_{m0}\)&-0.076\tabularnewline\hrulethin
\end{tabular}
}%
\end{tableptx}%
Now that the dynamics are written, the second order equation can be integrated using a standard numerical tool. In this case it's safe to assume the angular velocity is zero and the initial pitch angle is also zero. You can then use proportional control to command your pitch angle to whatever you want. Note that proportional control is not sufficient to stabilize your aircraft but I'll leave that discussion to your controls professor.%
\end{subsectionptx}
%
%
\typeout{************************************************}
\typeout{Subsection 19.3.3 First Order Velocity Dynamics}
\typeout{************************************************}
%
\begin{subsectionptx}{Subsection}{First Order Velocity Dynamics}{}{First Order Velocity Dynamics}{}{}{servo-5-5}
In order to gain some practice with simulating dynamic systems it's possible to approximate the velocity of an aircraft as a first order system. In order to accelerate a radio controlled aircraft, the motor must spin faster. It turns out that the PWM commands sent to servos to cause them to rotate back and forth is actually the same PWM signal used to control an Electronic Speed Controller (ESC) which controls a motor. So if you send around 1.1 ms to an ESC, the motor will be off and 1.9 ms will be full throttle. In this case we can make a simple equation that relates the throttle command to the velocity of the aircraft. The equation is shown below.%
\begin{equation}
\dot{u} = \frac{1}{m}[C_{x\delta t} (\delta_t-\delta_0) - 0.5\rho u_0 S C_x u]\label{e_first-order}
\end{equation}
where \(m=1.23 kg\), \(u\) is the velocity of the aircraft, \(\delta_t\) is the throttle command in \(ms\), and \(\delta_0=1.1 ms\) is the nominal PWM pulse to an ESC. The coefficients \(C_{x\delta t}=13.78 \frac{N}{ms}\) and \(C_x=0.017\) are coefficients that relate throttle to acceleration and velocity to drag respectively. The rest of the variables are explained above in the second order system. For this simulation you can assume that proportional control works just fine and thus%
\begin{equation}
\delta_t = k_p(u_c-u)\label{servo-5-5-2-9}
\end{equation}
where \(u_c\) is the commanded velocity and \(k_p\) is the proportional gain.%
\end{subsectionptx}
\end{sectionptx}
%
%
\typeout{************************************************}
\typeout{Section 19.4 Extra Help}
\typeout{************************************************}
%
\begin{sectionptx}{Section}{Extra Help}{}{Extra Help}{}{}{servo-6}
If you're having trouble with this lab, I did the calibration part of the project myself and \href{https://youtu.be/ynlGiPZk5VM}{\linkhilite{posted a YouTube video about it}}.If you need help with trend lines in Python you can watch this video where I  \href{https://www.youtube.com/watch?v=4vYYPHRMdqM\&feature=youtu.be}{\linkhilite{create a regression line}}. There is a video about \href{https://www.youtube.com/watch?v=XSLzcwTOsWk}{\linkhilite{making a protractor with paper}} if you don't have one.  I explain feedback control in much more detail in a \href{https://www.youtube.com/watch?v=PAK5V8wzVXY\&list=PL_D7_GvGz-v30U58EUUOdJGgO4u75DXoB\&index=1}{\linkhilite{controls overview Youtube video}} if you'd like to learn more about Control Theory \hyperlink{youtube}{\linkhilite{[{\xreffont 50}]}}.%
\end{sectionptx}
%
%
\typeout{************************************************}
\typeout{Section 19.5 Assignment}
\typeout{************************************************}
%
\begin{sectionptx}{Section}{Assignment}{}{Assignment}{}{}{servo-7}
\begin{introduction}{}%
This servo project can be split into 2 projects at your instructors discretion. The first part involves learning how to rotate the servo and then calibrating the servo so that degrees of rotation can be related to PWM in microseconds while the second part deals with feedback control. In part 1 ensure that you can rotate the servo and then experimentally determine the minimum and maximum PWM signal you can send to the servo before it stops rotating anymore. Then create a calibration equation and plot with degrees on the x-axis and PWM on the y-axis. Take at least 10 data points for your calibration equation and use a protractor on your computer screen to obtain angle in degrees of your servo. Your calibration equation needs to be linear. I'm also asking you to simulate a first order system to look at a simple feedback control problem. In part 2 you are to use the angle obtained by the accelerometer to feedback to the servo using proportional control. You are to do this for a \(k_p\) value of 1 and 2. In addition, you are to simulate the second order system above for 10 seconds with varying values of \(k_p\) and plotting the results. \alert{Note, for this lab, once you have the minimum and maximum values of your servo, you need to add a saturation filter into your code so that the value sent to your servo is always within these bounds. Otherwise you could damage or even break your servo.}%
\end{introduction}%
%
%
\typeout{************************************************}
\typeout{Subsection 19.5.1 Part 1}
\typeout{************************************************}
%
\begin{subsectionptx}{Subsection}{Part 1}{}{Part 1}{}{}{servo-7-3}
%
\begin{enumerate}
\item{}Title page with name,title and date - 5\%%
\item{}Introduction with wiring diagram - 5\%%
\item{}All figures must have appropriate annotations, figure captions, figure labels and a paragraph of supporting text - 10\%%
\item{}Include 2 photos of your servo with one photo rotated to the maximum and the other photo rotated to the minimum degrees. Also, include a screenshot (entire computer screen) of Mu showing the PWM pulse in the serial monitor for both rotations. - 10\%%
\item{}Report the minimum and maximum values of your servo in a table. Be sure to report the PWM signals in milliSeconds as well as the angle of the servo in degrees. - 10\%%
\item{}Include your raw calibration data with PWM signal in one column and degrees in the other column - 10\%%
\item{}Include a Figure of your calibration data plotted with your trend line on top. Put your \(r^2\) value in the title of the plot - 20\%%
\item{}Write your regression equation as (\emph{servo\textunderscore{}pulse = m\textasteriskcentered{}servo\textunderscore{}angle + b}) being sure to replace m and b with your actual values. Remember to round to appropriate significant figures - 10\%%
\item{}Using your calibration equation imbedded in your Mu code, compute the pulse required to rotate your servo to 90 degrees. Report on how close it is to 1.5 ms. Include a photo of your servo rotated to 90 degrees and a screenshot (entire computer screen) of your calibration equation placed into Mu with your pulse value in the Serial monitor. The code must also show your saturation filter as well - 10\%%
\item{}Simulate the first order system in Equation \hyperref[e_first-order]{\linkhilite{({\xreffont\ref{e_first-order}})}} for 10 seconds with \(k_p=1\) and \(k_p=10\) for a \(u_c=20 m/s\). Plot \(u\) vs time for both values of proportional gain. - 10\%%
\item{}Appendix A - Video Details - Film yourself and your entire screen with the Serial monitor open in Mu as you rotate the servo from 0 to roughly 180 degrees. Then show yourself plotting the regression line on top of your data on your laptop - Pass\slash{}Fail%
\item{}Appendix B - Code used to rotate servo and code used to plot data in plain text (no screenshots) - Pass\slash{}Fail%
\item{}\alert{If you are asked to fix something in the Appendices such as missing content in your video or screenshots of your code, you will lose 10\% of your grade everytime you submit an update after the deadline.}%
\end{enumerate}
\end{subsectionptx}
%
%
\typeout{************************************************}
\typeout{Subsection 19.5.2 Part 2}
\typeout{************************************************}
%
\begin{subsectionptx}{Subsection}{Part 2}{}{Part 2}{}{}{servo-7-4}
%
\begin{enumerate}
\item{}Title page with name,title and date - 5\%%
\item{}Introduction - 5\%%
\item{}All figures must have appropriate annotations, figure captions, figure labels and a paragraph of supporting text - 10\%%
\item{}Include an annotated photo of your circuit showing the accelerometer, wiring, and servo. - 20\%%
\item{}Include a photo of your CPX rotated to 45 degrees and your servo rotated with \(k_p=1\) and \(k_p=2\). Have your servo held up to your protractor when rotated so readers can see the angle at which your servo is rotated. - 40\%%
\item{}Simulate the second order system in Equation \hyperref[e_second-order]{\linkhilite{({\xreffont\ref{e_second-order}})}} for 10 seconds with \(k_p=-1\) and \(k_p=-10\) for a \(\theta_c=45^o\). Plot \(\theta\) vs time for both values of proportional gain. - 20\%%
\item{}Appendix A - Video Details - Film yourself and your entire screen with the Serial monitor open in Mu as you rotate the CPX\slash{}CPB +\slash{}-45 degrees and also showing the servo rotate IN SYNC (+\slash{}-45 degrees) with your CPX\slash{}CPB. - Pass\slash{}Fail%
\item{}Appendix B - Code used to rotate servo as well as simulation code in plain text (no screenshots) - Pass\slash{}Fail%
\item{}Appendix C - Screenshot of Excel or Sheets if applicable - Pass\slash{}Fail%
\item{}\alert{If you are asked to fix something in the Appendices such as missing content in your video or screenshots of your code, you will lose 10\% of your grade everytime you submit an update after the deadline.}%
\end{enumerate}
\end{subsectionptx}
\end{sectionptx}
\end{chapterptx}
%
%
\typeout{************************************************}
\typeout{Chapter 20 Time Constant of a Thermistor}
\typeout{************************************************}
%
\begin{chapterptx}{Chapter}{Time Constant of a Thermistor}{}{Time Constant of a Thermistor}{}{}{thermistor}
\renewcommand*{\chaptername}{Chapter}
\begin{introduction}{}%
\begin{listptx}{List}{\textbf{Parts List}}{thermistor-2-1}{}%
%
\begin{enumerate}
\item{}CPX\slash{}CPB%
\item{}USB Cable%
\item{}Laptop%
\end{enumerate}
\end{listptx}%
\begin{listptx}{List}{\textbf{Learning Objectives}}{thermistor-2-2}{}%
%
\begin{enumerate}
\item{}Learn the basic form of a first order system%
\item{}Learn the basic solution of a first order system%
\item{}Understand settling time and how that effects engineering%
\item{}Applied Estimation of a First Order system%
\end{enumerate}
\end{listptx}%
For this lab we're going to use the thermistor located on Pin A9 of the CPX\slash{}CPB. In a previous lab we measured the analog signal from the pin and also used the \emph{adafruit\textunderscore{}thermistor.py} module that converts the voltage from the thermistor to celsius (See \hyperref[temperature-sensor]{\linkhilite{Subsection~{\xreffont\ref{temperature-sensor}}}}).%
\par
For this lab though, we're not only going to measure the temperature from the thermistor, we're going to analyze the dynamics of the thermistor and how it responds to an abrupt change in temperature. It is said that a themistor can be modeled as a first order system. The basic form of a first order system can be shown below.%
\begin{equation}
\dot{T} = \sigma (T_f - T)\label{thermistor-2-4-1}
\end{equation}
where \(T\) is our state variable and \(\dot{T}\) is the derivative of our state. For this example let's assume that \(T\) is temperature. In this case \(T_f\) is an external temperature that is either higher or lower than current temperature causing the derivative of temperature to be non-zero. You can see in this case that once the temperature of the system is equal to the external forcing temperature, the derivative of the temperature of the system goes to zero. This implies that the system has reached equilibrium. The variable \(\sigma\) has units of \(Hz\) or \(1/sec\). The inverse of \(\sigma\) is \(\tau\) the time constant.%
\begin{equation}
\tau = \frac{1}{\sigma}\label{thermistor-2-4-11}
\end{equation}
%
\par
The time constant of the system is related to how quickly the system responds to change or even how long it takes for the system to reach equilibrium. Reaching equilibrium quantitatively happens when the state has changed 98\%. In other words the state of the system is within 2\% of the equilibrium state. This is called the settling time \(t_s\). The time constant is related to the settling time using the equation below.%
\begin{equation}
\tau = \frac{t_s}{4}\label{thermistor-2-5-2}
\end{equation}
%
\par
The only question now of course is what is the solution \(T(t)\) or rather the temperature as a function of time. In this case, the solution to the above dynamics equation can be solved using \href{https://www.youtube.com/watch?v=VOv2HI3i7oo}{\linkhilite{standard first order differential equation techniques}} to obtain the solution below\hyperlink{youtube}{\linkhilite{[{\xreffont 50}]}}.%
\begin{equation}
T(t) = (T_0-T_f)e^{-\sigma t} + T_f\label{e_thermistor}
\end{equation}
%
\par
Looking at the equation now you can see some properties right away. When \(t=0\), the first term reduces to \((T_0-T_f)\) which means the initial temperature is \(T_0\) or the initial temperature. When \(t\rightarrow\infty\), the first term drops to zero and thus the temperature is \(T_f\) or final temperature. In this lab we're going to get the time constant of the \href{https://en.wikipedia.org/wiki/Thermistor}{\linkhilite{thermistor}} on board the CPX\hyperlink{thermistor-wikipedia}{\linkhilite{[{\xreffont 63}]}}. If you take data on the CPX and walk outside or place the thermistor directly into a fridge, the temperature change will not be immediate. In order to estimate the time constant of a thermistor you need to change the temperature somehow. Blelow are some ideas for how to change the temperature of the CPX\slash{}CPB. You can do one or more of these methods to get data for this lab but you're only required to pick one.%
\begin{enumerate}
\item{}Start logging data and then place the CPX into the fridge%
\item{}Put the CPX into the fridge and then pull the CPX out of the fridge and watch the temperature return to ambient%
\item{}Walk outside on a hot (or cold) day and watch the CPX change temperature due to your A\slash{}C.%
\item{}Walk inside and watch the CPX get warmer (or colder) as your HVAC changes the CPX temperature.%
\item{}Use a hair dryer to heat up the CPX and then watch it cool down.%
\end{enumerate}
%
\end{introduction}%
%
%
\typeout{************************************************}
\typeout{Section 20.1 Initial Setup}
\typeout{************************************************}
%
\begin{sectionptx}{Section}{Initial Setup}{}{Initial Setup}{}{}{thermistor-3}
There is no wiring diagram for this project because we're using the thermistor that's already on the CPX\slash{}CPB. The code to read the thermistor is pretty simple and is shown below. Remember that the thermistor is a resistor that changes with temperature. The ADC on the CPX converts the voltage across the thermistor to temperature. \href{https://learn.adafruit.com/thermistor/circuitpython}{\linkhilite{The Adafruit Learn system}} does a bit of work to explain the conversion from voltage to temperature and is also explained in \hyperref[temperature-sensor]{\linkhilite{Subsection~{\xreffont\ref{temperature-sensor}}}}\hyperlink{Adafruit}{\linkhilite{[{\xreffont 32}]}}. Just remember that the thermistor on pin A9 is connected to the high side of the voltage divider which effects the equations to compute the resistance and voltage in the thermistor. My version of the code is also on my \href{https://github.com/cmontalvo251/Microcontrollers/blob/master/Circuit_Playground/CircuitPython/Temp/record_temperature_thermistor.py}{\linkhilite{Github}}\hyperlink{Github}{\linkhilite{[{\xreffont 33}]}}. \alert{A simple version of the code is shown below but note that the code is from a much earlier and simpler version of my code which only reads the temperature in Celsius. The code on my Github has more functionality such as printing the raw ADC value and the voltage.}%
\begin{figureptx}{Figure}{Read thermistor data code in CircuitPython}{thermistor-3-3}{}%
\begin{image}{0.135}{0.73}{0.135}{}%
\includegraphics[width=\linewidth]{external/thermistor.png}
\end{image}%
\tcblower
\end{figureptx}%
\end{sectionptx}
%
%
\typeout{************************************************}
\typeout{Section 20.2 Estimating the Time Constant}
\typeout{************************************************}
%
\begin{sectionptx}{Section}{Estimating the Time Constant}{}{Estimating the Time Constant}{}{}{thermistor-4}
I did the first two examples (using a fridge) and plotted both data sets in the same script as shown below. I opted to use method 1 (see \hyperref[daq]{\linkhilite{chapter~{\xreffont\ref{daq}}}}) from the datalogging project and just have the data print to Serial and then unplug the CPX when I'm done taking data and copy and paste the data into a text file.%
\begin{figureptx}{Figure}{Heating and cooling examples plotted in Thonny}{thermistor-4-3}{}%
\begin{image}{0.135}{0.73}{0.135}{}%
\includegraphics[width=\linewidth]{external/thermistor_plots.png}
\end{image}%
\tcblower
\end{figureptx}%
At this point it's possible to get the time constant by remembering that the settling time (time it takes the temperature to settle out) is equal to 4 times the time constant (\(\tau\)) and thus the time constant is the settling time divided by 4. After computing the settling time for both data sets and overlaying the equations on the measured data I get these two plots.%
\begin{figureptx}{Figure}{Exponential regression curves fit to heating and cooling data in Thonny}{thermistor-4-5}{}%
\begin{image}{0.135}{0.73}{0.135}{}%
\includegraphics[width=\linewidth]{external/thermistor_plots2.png}
\end{image}%
\tcblower
\end{figureptx}%
What was interesting was that the time constant for heating up was 62.5 seconds and for cooling down it was 155 seconds. The time to get cold was way slower than heating up. I'm not a heat transfer expert so I won't comment as to why this happened. One other import note I'd like to mention is the cool down phase was much more accurate than the heat up phase. This is most likely because when I pulled the thermistor out of the fridge I touched it with my hands and then moved it to a table. There was also alot more airflow outside the fridge which would change the overall dynamics. Still, the fitted data matches up pretty well and I hope yours does too.%
\end{sectionptx}
%
%
\typeout{************************************************}
\typeout{Section 20.3 Assignment}
\typeout{************************************************}
%
\begin{sectionptx}{Section}{Assignment}{}{Assignment}{}{}{thermistor-5}
\begin{introduction}{}%
For this assignment, you are to change the ambient temperature of the CPX\slash{}CPB through heating or cooling and then record time and voltage while waiting for the CPX\slash{}CPB to return to ambient temperature. You are to then plot voltage, resistance and temperature as a function of time and then using the graph of temperature (not using a curve fit), estimate \(T_0\), \(T_f\),\(\sigma\),\(t_s\) and \(\tau\). Finally, using \hyperref[e_thermistor]{\linkhilite{equation~({\xreffont\ref{e_thermistor}})}}, plot the estimate of \(T(t)\).%
%
\begin{enumerate}
\item{}Title page with name,title and date - 5\%%
\item{}Introduction - 5\%%
\item{}All figures must have appropriate annotations, figure captions, figure labels and a paragraph of supporting text - 10\%%
\item{}Take a picture of you actively heating or actively cooling the CPX\slash{}CPB - 10\%%
\item{}Plot the raw voltage and resistance in the thermistor vs time - 20\%%
\item{}Plot your temperature data with time on the x-axis and temperature on the y-axis and overlay your fitted data on top of your measured data. Points will be given based on how well your fit is - 20\%%
\item{}Include a table of your estimated parameters - 20\%%
\item{}Using the time constant you computed, assume you accidentally put your CPX\slash{}CPB in the oven at 400F. Simulate your system heating up from room temperatue (25C) to 400F and tell me how long I would have before I destroy my CPX\slash{}CPB. Make sure to show your work or you receive no credit - 10\%%
\item{}Appendix A - Video Details - Film yourself and your screen as you explain how you heated or cooled the CPX\slash{}CPB. Then show the thermistor running on your computer with the Serial Monitor open in Mu. Finally plot the regression line over your plotted data on your computer - Pass\slash{}Fail%
\item{}Appendix B - Code used to take data and code used to plot in plain text (no screenshots) - Pass\slash{}Fail%
\item{}Appendix C - Screenshot of Excel or Sheets if applicable - Pass\slash{}Fail%
\item{}\alert{If you are asked to fix something in the Appendices such as missing content in your video or screenshots of your code, you will lose 10\% of your grade everytime you submit an update after the deadline.}%
\end{enumerate}
\end{introduction}%
\end{sectionptx}
\end{chapterptx}
%
%
\typeout{************************************************}
\typeout{Chapter 21 Second Order Systems and Aliasing}
\typeout{************************************************}
%
\begin{chapterptx}{Chapter}{Second Order Systems and Aliasing}{}{Second Order Systems and Aliasing}{}{}{second-order}
\renewcommand*{\chaptername}{Chapter}
\begin{introduction}{}%
\begin{listptx}{List}{\textbf{Parts List}}{second-order-2-1}{}%
%
\begin{enumerate}
\item{}Laptop%
\item{}CPX\slash{}CPB%
\item{}USB Cable%
\item{}Some sort of oscillating system like a swinging pendulum or vibrating ruler%
\end{enumerate}
\end{listptx}%
\begin{listptx}{List}{\textbf{Learning Objectives}}{second-order-2-2}{}%
%
\begin{enumerate}
\item{}Learn the basic form of a second order system%
\item{}Understand the difference between underdamped, critically damped, and overdamped%
\item{}Understand natural frequency and damping ratio%
\item{}Applied Estimation of a Second Order system%
\item{}Understand the pitfalls of aliasing%
\end{enumerate}
\end{listptx}%
A second order system undergoing free motion will have dynamics that look like this%
\begin{equation}
\ddot{\theta} + 2\zeta \omega_n \dot{\theta} + {\omega_n}^2\label{second-order-2-3-1}
\end{equation}
%
\par
In this case, the solution to the above equation can be solved using \href{https://www.youtube.com/watch?v=VOv2HI3i7oo}{\linkhilite{standard second order differential equation techniques}} to obtain the solution below \hyperlink{youtube}{\linkhilite{[{\xreffont 50}]}}.%
\begin{equation}
\theta(t) = \theta_0e^{-\sigma t}cos(\omega_d t)\label{e_secondorder}
\end{equation}
%
\begin{equation}
\omega_n = \sqrt{\sigma^2 + {\omega_d}^2}\label{second-order-2-4-4}
\end{equation}
%
\begin{equation}
\zeta = \frac{\sigma}{\omega_n}\label{second-order-2-4-5}
\end{equation}
%
\par
Looking at the equations you can see that if the time series of the oscillations are known, the damping constant and damped natural frequency can be obtained. These two values can be combined to obtain the natural frequency and the damping ratio. So let's get some data.%
\par
What I'd like you to do for this lab is to find some sort of parameter that varies in some sort of sinusoidal way. I've come up with a few ideas below. You may pick anyone you want although some are easier than others.%
\begin{enumerate}
\item{}Drive over a speed bump - If you drive over a speed bump slowly and place your CPX on the dashboard your car will hopefully vibrate for a few seconds after you drive over it. Your acceleration will look somewhat like a sine wave.%
\item{}Vibrate a ruler - If you place a ruler on the edge of a desk and deflect it, it will vibrate. If you place the CPX on the ruler you'll be able to measure the vibration of the ruler.%
\item{}Attach CPX to a spring or a weight to the end of rubber bands - Take the CPX and attach it to a weight of some kind and attach the weight to a spring or a set of rubber bands. This is a mass spring damper simple and will vibrate at the natural frequency of square root of stiffness divided by mass. I chose to do this example.%
\item{}Build a Pendulum: If you decide to build a pendulum, you need to hang a weight on the CPX or attach the CPX to something heavy. This will make the ratio between cable to end mass much smaller and thus better for data and fitting. An example includes  duct taping your CPX to a water bottle or something. It would also be better to use a string and mount the CPX to the string with an external battery pack and have the CPX log data internally. This way most of the mass would be concentrated at the tip of the pendulum. Another idea is to take a paper towel cardboard tube and tape the cpx inside it with the cable running through it. Then duct a large weight on the end of the paper towel tube and then hinge the top of the tube by skewering a screwdriver through it. This will allow the tube to swing like a pendulum rather than the string.%
\end{enumerate}
%
\par
There are most likely many other options so try and get creative and find something oscillatory or dynamic in some way. Try to find something that changes relatively quickly. The temperature outside changes in a sinusoidal fashion but it's so slow it would take you days to do this experiment.%
\end{introduction}%
%
%
\typeout{************************************************}
\typeout{Section 21.1 Setup Pendulum Example}
\typeout{************************************************}
%
\begin{sectionptx}{Section}{Setup Pendulum Example}{}{Setup Pendulum Example}{}{}{second-order-3}
In this example I'm going to swing a pendulum in the X\slash{}Y plane of the accelerometer sensor so that I can ignore the Z axis data. I'm going to get the actual angle of the pendulum but if you're building something else you can ignore this part. In order to convert the accelerometer values to pitch angle, I'm going to use the same equations discussed in the accelerometer lab (See \hyperref[accelerometer]{\linkhilite{Subsection~{\xreffont\ref{accelerometer}}}})%
\par
Remember, on the CPX\slash{}CPB specifically you'll want to import the math module and use the atan2 function. When I swing the pendulum, this is the result I get from Plotter in Mu.%
\begin{figureptx}{Figure}{Oscillations shown from atan2 measurement of accelerometer in Mu}{second-order-3-4}{}%
\begin{image}{0.135}{0.73}{0.135}{}%
\includegraphics[width=\linewidth]{external/oscillation_mu.png}
\end{image}%
\tcblower
\end{figureptx}%
If I then bring this into Python I get the following plot below. In my data set I only logged the angle. Since the time step between each point was 0.01 seconds I was able to create a time series. It's pretty clear that there is some nonlinearity in the data so I chose to start the data at 0.5 seconds. Another thing I noticed was that the angle settled out to around 8 degrees so I chose to subtract off that bias from the angle data.%
\begin{figureptx}{Figure}{Oscillation data plotted in Thonny}{second-order-3-6}{}%
\begin{image}{0.135}{0.73}{0.135}{}%
\includegraphics[width=\linewidth]{external/oscillation_plots.png}
\end{image}%
\tcblower
\end{figureptx}%
\end{sectionptx}
%
%
\typeout{************************************************}
\typeout{Section 21.2 Estimating Parameters}
\typeout{************************************************}
%
\begin{sectionptx}{Section}{Estimating Parameters}{}{Estimating Parameters}{}{}{second-order-4}
After trimming the data and removing some bias it was time to get my damping constant and damped natural frequency. There are a few equations that can help you obtain these parameters. First, the settling time is the length of time it takes for the oscillations to settle. The settling time can be used to find the damping constant. This is equal to:%
\begin{equation}
\sigma = \frac{4}{T_s}\label{second-order-4-2-1}
\end{equation}
%
\par
For my data set the settling time was about 1.25 seconds which gave a damping constant of 3.2. Once I had the damping constant I could obtain the damped natural frequency. This was done by measuring the distance between two peaks in the data set. There is a peak at around 0.5 seconds and another at around 0.95 seconds. I can use this to compute a period T. Period can be computed to angular frequency using the equation below.%
\begin{equation}
\omega_d = \frac{2\pi}{T}\label{second-order-4-3-1}
\end{equation}
%
\par
Using the period in my wave form I obtained a damped natural frequency of about 14.8 rad\slash{}s. Using these values I can plot the simulated data on top of the measured data noting that my initial angle was about 35 degrees minus the bias of 8 degrees. When I first plotted the data I noticed that my fit wasn't entirely perfect. My period was correct but my damping rate was too high. I realized it was because my settling time was too big. I increased the settling time to 1.8 seconds and got this plot here. You can see that my fitted data lined up almost perfectly with my measured data.%
\begin{figureptx}{Figure}{Sinusoidal regression curve plotted on top of oscillation data in Thonny}{second-order-4-5}{}%
\begin{image}{0.135}{0.73}{0.135}{}%
\includegraphics[width=\linewidth]{external/oscillation_plots2.png}
\end{image}%
\tcblower
\end{figureptx}%
\end{sectionptx}
%
%
\typeout{************************************************}
\typeout{Section 21.3 Aliasing}
\typeout{************************************************}
%
\begin{sectionptx}{Section}{Aliasing}{}{Aliasing}{}{}{aliasing}
Aliasing occurs when you don't take data fast enough. The speed of taking data to avoid aliasing depends on the natural frequency of the system itself. To show this I changed the length of the pendulum which changed the natural frequency. I then proceeded to measure the angle as it oscillated at its natural frequency. I repeated this at sampling frequencies of 1, 10 and 100 Hz. The way I changed the sampling frequency was by changing the time.sleep value in the while True loop. The \href{https://github.com/cmontalvo251/Microcontrollers/blob/master/Circuit_Playground/CircuitPython/Accelerometer/low_level_accel.py}{\linkhilite{accelerometer code}} I used can be found on Github\hyperlink{Github}{\linkhilite{[{\xreffont 33}]}}. After I finished the experiment I had 3 data files that I plotted on top of each other. This what I got with the longest string. It's easy to see in the photo that sampling at 1 Hz was way too slow to capture the natural oscillations of the water bottle. However, 100 Hz and even 10 Hz was plenty fast to sample the oscillations. According to the recorded data there was above 19 cycles in 9 seconds which is about 2 Hz. In this case, as long as we sample at 4 Hz the signal will be captured properly which is why 10 Hz and 100 Hz is able to capture the signal correctly.%
\begin{figureptx}{Figure}{Aliasing shown with oscillatory data in Thonny}{aliasing-3}{}%
\begin{image}{0.135}{0.73}{0.135}{}%
\includegraphics[width=\linewidth]{external/aliasing.png}
\end{image}%
\tcblower
\end{figureptx}%
Once I sampled my waveform at 3 different sampling rates I elected to shrink the length of the pendulum to change the natural frequency and see if that affected any aliasing I saw. The results are shown in the figure below.%
\begin{figureptx}{Figure}{Aliasing for shorter pendulum shown in Thonny}{aliasing-5}{}%
\begin{image}{0.135}{0.73}{0.135}{}%
\includegraphics[width=\linewidth]{external/aliasing2.png}
\end{image}%
\tcblower
\end{figureptx}%
In this case you can see that there were 33 cycles in 10 seconds which is 3.3 Hz. The Nyquist criteria states that I need to sample at 6.6 Hz. The Nyquist criteria is very specific though in that if you sample at twice the frequency, you will just obtain the correct frequency. This does not mean that you will capture every data point properly. Hence in the chart above, the blue line at 100 Hz is perfect, the green line at 1 Hz is too slow (less than 6.6 Hz) and the orange line at 10 Hz captures the frequency correctly but between 5 and 8 seconds does not adequately capture the waveform. In my opinion, in order to sample the data effectively and not just obtain the frequency of the waveform, you need to sample at 4 times the natural frequency. So for the waveform of 3.3 Hz, the Montalvo frequency would be 13.2 Hz which is higher than the 10 Hz. This explains why the 10 hz sample is not perfect while the 100 hz sample rate is much better.%
\end{sectionptx}
%
%
\typeout{************************************************}
\typeout{Section 21.4 Extra Help}
\typeout{************************************************}
%
\begin{sectionptx}{Section}{Extra Help}{}{Extra Help}{}{}{second-order-6}
I have an accompanying youtube video where I discuss \href{https://youtu.be/8F_8st_8MmA}{\linkhilite{aliasing issues}} which might help you better understand aliasing in general\hyperlink{youtube}{\linkhilite{[{\xreffont 50}]}}.%
\end{sectionptx}
%
%
\typeout{************************************************}
\typeout{Section 21.5 Assignment}
\typeout{************************************************}
%
\begin{sectionptx}{Section}{Assignment}{}{Assignment}{}{}{second-order-7}
\begin{introduction}{}%
For this assignment you are to create a system that oscillates and capture data while oscillating. You are to sample the system at 1, 10 and 100 Hz and plot the shifted and cleaned data for all 3 sample rates on the same graph. For the data sampled at 100 Hz you are then to estimate the initial amplitude, the settling time and the period using the graph. Then compute the natural frequency and damping ratio. Finally, using \hyperref[e_secondorder]{\linkhilite{equation~({\xreffont\ref{e_secondorder}})}}, plot the amplitude as a function of time on top of your 100hz sampled data. \alert{Note, if you decide to make a pendulum you must ensure that the CPX\slash{}CPB does not spin when oscillating otherwise you won't be able to use your acceleration data easily.}%
%
\begin{enumerate}
\item{}Title page with name,title and date - 5\%%
\item{}Introduction - 5\%%
\item{}All figures must have appropriate annotations, figure captions, figure labels and a paragraph of supporting text - 10\%%
\item{}Include a photo of your experiment and explain the state that you are measuring - 10\%%
\item{}Plot the results of your data (make sure to connect your data points) vs time for all 3 sampling frequencies (1,10, 100Hz) and include a legend - 10\%%
\item{}Include a table of your estimated parameters as well as the natural frequency and damping ratio - 20\%%
\item{}Plot your measured data sampled at 100hz with your fitted equation on top - 20\%%
\item{}Based on your results and your computed natural frequency, comment on the Nyquist frequency and the Montalvo frequency of your system and whether or not you encountered aliasing in your experiment based on the figures and the Nyquist\slash{}Montalvo frequency - 10\%%
\item{}Using the natural frequency and damping ratio you computed, simulate your system for an initial amplitude of 1 for two times the settling time. - 10\%%
\item{}Appendix A - Video Details - Film yourself and your entire screen showing the camera the experiment you used to create oscillations in the accelerometer. Then plot your accelerometer data and show the regression line fit over your oscillations - Pass\slash{}Fail%
\item{}Appendix B - Code used to take data and code used to plot in plain text (no screenshots) - Pass\slash{}Fail%
\item{}Appendix C - Screenshot of Excel or Sheets if applicable - Pass\slash{}Fail%
\item{}\alert{If you are asked to fix something in the Appendices such as missing content in your video or screenshots of your code, you will lose 10\% of your grade everytime you submit an update after the deadline.}%
\end{enumerate}
\end{introduction}%
\end{sectionptx}
\end{chapterptx}
%
%
\typeout{************************************************}
\typeout{Chapter 22 Inertial Measurement Unit}
\typeout{************************************************}
%
\begin{chapterptx}{Chapter}{Inertial Measurement Unit}{}{Inertial Measurement Unit}{}{}{imu}
\renewcommand*{\chaptername}{Chapter}
\begin{introduction}{}%
\begin{listptx}{List}{\textbf{Parts List}}{imu-2-1}{}%
%
\begin{enumerate}
\item{}Laptop%
\item{}CPX\slash{}CPB or Arduino%
\item{}USB Cable%
\item{}\href{https://www.adafruit.com/product/4485}{\linkhilite{LSM6DS33+LIS3MDL}} (Note that other Inertial Measurement or "DOF" Sensors will work for this lab you will just have to change the code to accommodate the change in hardware)\hyperlink{AdafruitMAIN}{\linkhilite{[{\xreffont 34}]}}%
\item{}Alligator Clips (x4)%
\item{}Bread Board%
\item{}Soldering iron%
\end{enumerate}
\end{listptx}%
\begin{listptx}{List}{\textbf{Learning Objectives}}{imu-2-2}{}%
%
\begin{enumerate}
\item{}See and understand the concept of soldering%
\item{}Understand the I2C protocol at a high level%
\item{}Learn the components of an Inertial Measurement Unit (IMU) (Accelerometer, Rate Gyro, Magnetometer)%
\item{}Read IMU data and plot%
\end{enumerate}
\end{listptx}%
In this lab we're going to use an Inertial Measurement Unit (IMU) sensor to measure angular velocity and the magnetic field of the surrounding environment.%
\begin{figureptx}{Figure}{Inertial Measurement Unit (IMU) without header pins soldered}{imu-2-4}{}%
\begin{image}{0.135}{0.73}{0.135}{}%
\includegraphics[width=\linewidth]{external/imu.jpeg}
\end{image}%
\tcblower
\end{figureptx}%
The sensor above is the \href{https://www.adafruit.com/product/4485}{\linkhilite{LSM6DS33+LIS3MDL}}\hyperlink{AdafruitMAIN}{\linkhilite{[{\xreffont 34}]}}. You basically have 2 separate microchips in one. The first (LSM6DS33) is a 3-axis accelerometer and 3-axis rate gyro. The first measures acceleration and the second measures angular velocity. The LIS3MDL is a 3-axis magnetometer which measures magnetic fields. The three of these sensors put together (accelerometer, rate gyro, magnetometer) is called an IMU (Inertial Measurement Unit). It's actually possible to measure roll and pitch of an airplane using the accelerometer and heading using the magnetometer. Combining the angular velocity of the rate gyro can create a complete attitude estimation algorithm for six-degree-of-freedom vehicles (6DOF). This sensor does not come standard in the current iteration of the kit but you can purchase one on Adafruit for only \textdollar{}10 at the time of this writing. The interesting thing about this device is that you can actually purchase the LSM6DS33 and LIS3MDL separately but this breakout board has both chips on board. The goal of this lab is not only to use this specific sensor but to understand IMUs and I2C protocol.Most if not all breakout boards on the Adafruit website use I2C communication. You'll know if the breakout board uses I2C if you find SDA\slash{}SCL pins on the board. It will also say it in the quick description of the sensor.%
\begin{figureptx}{Figure}{IMU on Adafruit}{imu-2-6}{}%
\begin{image}{0.135}{0.73}{0.135}{}%
\includegraphics[width=\linewidth]{external/imu2.png}
\end{image}%
\tcblower
\end{figureptx}%
The nice thing about this sensor is that there is a magnetometer on board. Recall that the accelerometer can be used to measure roll and pitch (See \hyperref[accelerometer]{\linkhilite{Subsection~{\xreffont\ref{accelerometer}}}}). By combining the magnetometer in this IMU, it's possible to get the compass heading of the sensor.%
%
\begin{equation}
\psi = tan^{-1}\left(\frac{\bar{\beta}_z s_{\phi} - \bar{\beta}_y c_{\phi}}{\bar{\beta}_x c_{\phi} + \bar{\beta}_y s_{\theta}s_{\phi} + \bar{\beta}_z c_{\phi}s_{\theta}}\right)\label{imu-2-8}
\end{equation}
In the equations above, \(\beta\) is the magnetic field along all 3 axes. Again shorthand is used for cosine and sine functions and the derivation is left for another course entirely (See \href{https://github.com/cmontalvo251/LaTeX/tree/master/Aerospace_Mechanics}{\linkhilite{Aerospace Mechanics}})\hyperlink{Github}{\linkhilite{[{\xreffont 33}]}}. When you open the packaging of this breakout board you'll notice that the header pins are missing. First you'll need to cut a row of 4 and 6 for the top and bottom side of the board and \alert{solder the header pins to the sensor}. If you're taking my class I can solder this for you or teach everyone about soldering during a lecture session of class. If you are taking this class elsewhere you have two options: try and find someone who can solder this real quick (only takes about 5 minutes) or buy your own soldering iron and try to solder yourself. Once the device is soldered you can "plug" it into a breadboard.%
\end{introduction}%
%
%
\typeout{************************************************}
\typeout{Section 22.1 CPX\slash{}CPB Wiring and Software}
\typeout{************************************************}
%
\begin{sectionptx}{Section}{CPX\slash{}CPB Wiring and Software}{}{CPX\slash{}CPB Wiring and Software}{}{}{imu-3}
The wiring for this system requires 4 wires. The figure below shows the IMU connected to \href{https://learn.adafruit.com/lis3mdl-triple-axis-magnetometer/python-circuitpython}{\linkhilite{Feather M4}} (Courtesy of \href{https://learn.adafruit.com/u/siddacious}{\linkhilite{Bryan Siepert}})\hyperlink{Adafruit}{\linkhilite{[{\xreffont 32}]}}.%
\begin{figureptx}{Figure}{Wiring diagram of IMU connected to a FeatherWing\hyperlink{imu-wiring}{\linkhilite{[{\xreffont 42}]}}}{imu-3-3}{}%
\begin{image}{0.135}{0.73}{0.135}{}%
\includegraphics[width=\linewidth]{external/sensors_python_feather_wiring_breadboard.png}
\end{image}%
\tcblower
\end{figureptx}%
The only difference between the wiring diagram above and the CPX\slash{}CPB is that you will be using 4 alligator clips. The rest is straightforward. You need \emph{3.3V} to run to \emph{VIN}, \emph{GND} to \emph{GND} and then \emph{SDA} to \emph{SDA} and \emph{SCL} to \emph{SCL}.%
\begin{figureptx}{Figure}{Photo of IMU Connected to a CircuitPlayground Express}{imu-3-5}{}%
\begin{image}{0.135}{0.73}{0.135}{}%
\includegraphics[width=\linewidth]{external/imu_correct.jpeg}
\end{image}%
\tcblower
\end{figureptx}%
The photo above uses a combination of alligator clips and male to male wires into a breadboard. The red alligator clip is connected to the \emph{3.3V} pin on the CPB and then also connected to the red male-male wire which is connected to the \emph{VIN} pin on the IMU. The black alligator clip is connected to the \emph{GND} pin on the CPB and also connected to the brown male-male wire which is connected to the \emph{GND} pin on the IMU. The blue alligator clip is connected directly to the \emph{SDA} pins on both the CPB and the IMU. The white alligator clip connected to the yellow male-male wire which is connected to the \emph{SCL} pins on both the CPB and the IMU. Recall that SDA and SCL are external pins for I2C. In the accelerometer lab we used the internal I2C pins to access the accelerometer (See \hyperref[accelerometer]{\linkhilite{Subsection~{\xreffont\ref{accelerometer}}}}). In this lab we're going to use the external I2C pads to connect different I2C sensors.%
\par
Once you have the circuit wired and soldered it's time to work on software. First, you want to make sure you have your \href{https://circuitpython.org/downloads}{\linkhilite{Circuit Python UF2}} up to date\hyperlink{CircuitPython}{\linkhilite{[{\xreffont 30}]}}. In this example I'm using the 6.X version. Once I updated my UF2 I also updated my \href{https://circuitpython.org/libraries}{\linkhilite{Circuit Python Libraries}} (See \hyperref[installing-modules]{\linkhilite{Section~{\xreffont\ref{installing-modules}}}} for help on installing extra modules on your CPX\slash{}CPB)\hyperlink{CircuitPython}{\linkhilite{[{\xreffont 30}]}}. The specific modules I needed for this lab were:%
\begin{enumerate}
\item{}adafruit\textunderscore{}bus\textunderscore{}device%
\item{}adafruit\textunderscore{}register%
\item{}lis3mdl%
\item{}lsm6ds33%
\end{enumerate}
%
\par
Once you have the necessary modules you can run some example code. The Adafruit Learn page has a \href{https://learn.adafruit.com/lsm6ds33-6-dof-imu=accelerometer-gyro/python-circuitpython}{\linkhilite{tutorial for the LSM6DS33}}\hyperlink{Adafruit}{\linkhilite{[{\xreffont 32}]}}. The problem with the tutorial is that it seems like it was written for the Raspberry or some other microcontroller. As such I had to find some example code on \href{https://github.com/adafruit/Adafruit_CircuitPython_LSM6DS/blob/master/examples/lsm6ds_lsm6ds33_simpletest.py}{\linkhilite{Adafruit's Github}}\hyperlink{Github}{\linkhilite{[{\xreffont 33}]}}. After following both tutorials I was able to make my own script and upload it to my \href{https://github.com/cmontalvo251/Microcontrollers/blob/master/Circuit_Playground/CircuitPython/Accelerometer/external_lis3mdl_lsm6dss.py}{\linkhilite{Github}}\hyperlink{Github}{\linkhilite{[{\xreffont 33}]}}. \alert{Note that it's a good exercise when attaching I2C devices to run a scan on all connected I2C devices to make sure you can see the device and that you have the correct address (See \hyperref[I2C_scanner]{\linkhilite{Section~{\xreffont\ref{I2C_scanner}}}}).}%
\begin{figureptx}{Figure}{CircuitPython code to read an IMU}{imu-3-9}{}%
\begin{image}{0.1}{0.8}{0.1}{}%
\includegraphics[width=\linewidth]{external/imu_code.png}
\end{image}%
\tcblower
\end{figureptx}%
In the code above lines 1-6 import all the modules with line 5 importing the accelerometer on board the CPX and line 6 importing the external sensor wired up to SDA and SCL. Line 8 creates an I2C object using the SDA and SCL pins from the alligator clips and line 11 creates the sensor object. I also include lines 14-17 to include the onboard accelerometer. Notice I can access both sensors no problem. In the while loop line 20 checks the accelerometer on the CPX, line 21 checks the accelerometer on the breakout board and line 22 checks the angular velocity on the breakout board. Lines 23 and 24 print to serial and output to the plotter. Note that some lines are commented out because I wanted to try one thing at a time. With both accelerometers printing to the Plotter I could move the CPX and the breakout board in unison and get the following output.%
\begin{figureptx}{Figure}{Serial monitor open showing IMU data in Mu}{imu-3-11}{}%
\begin{image}{0.135}{0.73}{0.135}{}%
\includegraphics[width=\linewidth]{external/imu_mu.png}
\end{image}%
\tcblower
\end{figureptx}%
Notice that there are 6 numbers printed and 6 lines on the plotter. Both the CPX and the breakout board have little XYZ cartesian coordinate systems. I had to line them up properly before I started moving them. My suggestion would be for you to get some hot glue or 3M tape and place both breakout board and CPX on some sort of hard material like plywood, masonite, or even a cutting board. Anything to keep everything together.%
\par
Once you've done this, try uncommenting the line of code that prints the angular velocity. When I do that and move the breakout board around I can measure the angular velocity of each axis. The units are in radians per second but it's pretty obvious just from the magnitude of the graph.%
\begin{figureptx}{Figure}{Angular Velocity data of IMU shown in Serial Monitor or Mu}{imu-3-14}{}%
\begin{image}{0.135}{0.73}{0.135}{}%
\includegraphics[width=\linewidth]{external/imu_mag.png}
\end{image}%
\tcblower
\end{figureptx}%
The final part is to get the LIS3MDL (magnetometer) to work. The starting point for me was the \href{https://learn.adafruit.com/lis3mdl-triple-axis-magnetometer/python-circuitpython}{\linkhilite{Adafruit Learn page}}, along with the simple example from \href{https://github.com/adafruit/Adafruit_CircuitPython_LIS3MDL/blob/master/examples/lis3mdl_simpletest.py}{\linkhilite{Adafruit's Github}}. After that I was able to create \href{https://github.com/cmontalvo251/Microcontrollers/blob/master/Circuit_Playground/CircuitPython/Accelerometer/external_lis3mdl.py}{\linkhilite{my own code}}\hyperlink{Github}{\linkhilite{[{\xreffont 33}]}}\hyperlink{Adafruit}{\linkhilite{[{\xreffont 32}]}}. The only difference in your code is that the address will be 0x1c instead of 0x1e.%
\begin{figureptx}{Figure}{Magnetometer code in CircuitPython}{imu-3-16}{}%
\begin{image}{0.1}{0.8}{0.1}{}%
\includegraphics[width=\linewidth]{external/imu_code_mag.png}
\end{image}%
\tcblower
\end{figureptx}%
The code is almost identical to the code before except all the LIS3DH and LSM6DS33 code is commented out. Instead I have code to grab the magnetometer (LIS3MDL) at address 0x1E. Line 25 then calls the magnetometer and prints it.%
\end{sectionptx}
%
%
\typeout{************************************************}
\typeout{Section 22.2 Arduino Wiring and Software}
\typeout{************************************************}
%
\begin{sectionptx}{Section}{Arduino Wiring and Software}{}{Arduino Wiring and Software}{}{}{imu-4}
Wiring this to the Arduino requires the same setup as the CPX\slash{}CPB and also uses I2C. From the Arduino to the accelerometer, the \emph{3.3V} pin runs to \emph{VIN}, \emph{GND} to \emph{GND}, \emph{SDA} to \emph{SDA} and \emph{SCL} to \emph{SCL}. This is shown in the figure below. Note that the SD card breakout board from \hyperref[sdcard]{\linkhilite{Chapter~{\xreffont\ref{sdcard}}}} is also used in this example to showcase how easy it is to send accelerometer data to an SD card%
\begin{figureptx}{Figure}{Accelerometer connected to an Arduino}{imu-4-3}{}%
\begin{image}{0.135}{0.73}{0.135}{}%
\includegraphics[width=\linewidth]{external/Accelerometer_Arduino.png}
\end{image}%
\tcblower
\end{figureptx}%
Moving on to software, the \emph{.ino} file for both the \href{https://github.com/cmontalvo251/Microcontrollers/blob/master/Arduino/MEGA/Accelerometer_LSM6DS33/Accelerometer_LSM6DS33.ino}{\linkhilite{accelerometer}} by itself and the \href{https://github.com/cmontalvo251/Microcontrollers/blob/master/Arduino/MEGA/Accelerometer_LSM6DS33_SD/Accelerometer_LSM6DS33_SD.ino}{\linkhilite{accelerometer writing to an SD card}} are both on Github. For both of the versions of the code you need to install some libraries. You'll notice that at the top of the code there is a \#include statement which adds a specific library.%
\begin{preformatted}
#include <Adafruit_LSM6DS33.h>
Adafruit_LSM6DS33 lsm6ds33;
\end{preformatted}
In this case you need the \emph{Adafruit\textunderscore{}LSM6DS33} library which also requires other libraries in addition to that single \#include call. Also, since the libraries are from Adafruit rather than Arduino you need to navigate to Github and download a few zip libraries. The list of libraries you need and links to those websites are shown below%
%
\begin{enumerate}
\item{}\href{https://github.com/adafruit/adafruit_lsm6ds}{\linkhilite{Adafruit\textunderscore{}LSM6DS}}%
\item{}\href{https://github.com/adafruit/Adafruit_BusIO}{\linkhilite{Adafruit\textunderscore{}BusIO}}%
\item{}\href{https://github.com/adafruit/adafruit_sensor}{\linkhilite{Adafruit\textunderscore{}Sensor}}%
\end{enumerate}
Once you have the zip folders downloaded you need to navigate to Sketch\textgreater{}Include Library\textgreater{}Add .ZIP Library... and then add the ZIP file you just downloaded. This is different than adding a library directly from the Arduino website as is done for the SD library in \hyperref[sdcard]{\linkhilite{Chapter~{\xreffont\ref{sdcard}}}}. Once you've got the code downloaded and the libraries installed you should be able to compile without seeing any errors. After uploading the code the serial monitor should produce something similar to the figure below.%
\begin{figureptx}{Figure}{Serial Monitor Output of Accelerometer Code on Arduino}{imu-4-9}{}%
\begin{image}{0.135}{0.73}{0.135}{}%
\includegraphics[width=\linewidth]{external/Accel_Serial_Output.png}
\end{image}%
\tcblower
\end{figureptx}%
In the figure above, the sensor is outputtting temperature, acceleration x,y and z as well as the angular velocity along the x,y, and z axes. If you then compile the code which contains the SD card setup and writing you'll be able to write the data to an external SD card which would be super helpful for a standalone project. Again setting up the SD card is done in \hyperref[sdcard]{\linkhilite{Chapter~{\xreffont\ref{sdcard}}}}%
\begin{figureptx}{Figure}{Accelerometer Data File on SD Card}{imu-4-11}{}%
\begin{image}{0.135}{0.73}{0.135}{}%
\includegraphics[width=\linewidth]{external/Accelerometer_Data_File.png}
\end{image}%
\tcblower
\end{figureptx}%
In the figure above, the code only writes the millis timer and acceleration data to the SD card but it would be very simple to add temperature and angular velocity to the code.%
\end{sectionptx}
%
%
\typeout{************************************************}
\typeout{Section 22.3 Assignment}
\typeout{************************************************}
%
\begin{sectionptx}{Section}{Assignment}{}{Assignment}{}{}{imu-5}
\begin{introduction}{}%
For this assignment you are to wire up the external IMU and get data from it. You need to mount the CPX\slash{}CPB and the IMU to some sort of hard surface so that when you move the CPX\slash{}CPB the IMU moves as well. Make sure the axis of the IMU and the CPX\slash{}CPB are oriented in the same direction. With system mounted to a hard surface, perform doublet manuevers on each axis for a total of 3 doublets. A doublet is where you rotate the system to +90 degrees and then -90 degrees and then back to zero typically taking around 3 seconds for the entire maneuever. Using the accelerometer, compute the roll and pitch angle in degrees. Then use the magnetometer to compute the yaw angle. Using the roll, pitch and yaw angles, take a derivative to compute the angular velocity. Finally, take the angular velocity data and integrate it to obtain the pitch, roll and yaw angles.%
%
\begin{enumerate}
\item{}Title page with name,title and date - 5\%%
\item{}Introduction with wiring diagram - 5\%%
\item{}All figures must have appropriate annotations, figure captions, figure labels and a paragraph of supporting text - 10\%%
\item{}Include a photo of your CPX\slash{}CPB mounted to a hard surface with the IMU on a breadboard. Be sure to explain in your description about your axis system for both sensors. - 10\%%
\item{}Include a screenshot (entire computer screen) of Mu showing the Plotter open and all 3 angular velocity axes - 10\%%
\item{}Include a plot of both accelerometers for the 3 doublet manuevers. - 10\%%
\item{}Include a plot of angular velocity data for the 3 doublet manuevers. Also plot the derivative of the pitch, roll and yaw angles on top of this plot and add a legend to clearly indicate which is which. - 20\%%
\item{}Include a plot of magnetometer data for the 3 doublet maneuevers - 10\%%
\item{}Plot the roll, pitch and yaw angles in degrees from the accelerometer\slash{}magnetometer as well as the integrated rate gyro angles. Add a legend to clearly indicate which line is which. - 20\%%
\item{}Appendix A - Video Details - Film yourself and your entire screen showing the acceleration and angular velocity values change in the Plotter in Mu - Pass\slash{}Fail%
\item{}Appendix B - Code used to take data and code used to plot in plain text (no screenshots) - Pass\slash{}Fail%
\item{}Appendix C - Screenshot of Excel or Sheets if applicable - Pass\slash{}Fail%
\item{}\alert{If you are asked to fix something in the Appendices such as missing content in your video or screenshots of your code, you will lose 10\% of your grade everytime you submit an update after the deadline.}%
\end{enumerate}
\end{introduction}%
\end{sectionptx}
\end{chapterptx}
%
%
\typeout{************************************************}
\typeout{Chapter 23 Barometers}
\typeout{************************************************}
%
\begin{chapterptx}{Chapter}{Barometers}{}{Barometers}{}{}{pressure}
\renewcommand*{\chaptername}{Chapter}
\begin{introduction}{}%
\begin{listptx}{List}{\textbf{Parts List}}{pressure-2-1}{}%
%
\begin{enumerate}
\item{}Laptop%
\item{}CPX\slash{}CPB or Arduino%
\item{}USB Cable%
\item{}\href{https://www.adafruit.com/product/4633}{\linkhilite{LPS22}} (Note that other barometers will work for this lab you will just have to change the code to accommodate the change in hardware)\hyperlink{AdafruitMAIN}{\linkhilite{[{\xreffont 34}]}}%
\item{}Alligator Clips (x4)%
\item{}Bread Board%
\item{}Soldering iron%
\end{enumerate}
\end{listptx}%
\begin{listptx}{List}{\textbf{Learning Objectives}}{pressure-2-2}{}%
%
\begin{enumerate}
\item{}Continue to understand the concept of soldering%
\item{}Reinforce the I2C protocol at a high level%
\item{}Learn the components of a Pressure and temperature Sensor (Barometer)%
\item{}Read Pressure data and plot%
\end{enumerate}
\end{listptx}%
In this lab we're going to use a pressure and temperature sensor (Barometer) to measure atmospheric pressure and temperature. Although \hyperref[temperature-sensor]{\linkhilite{Subsection~{\xreffont\ref{temperature-sensor}}}} already discussed temperature sensors, this external device (LPS22) has a pressure sensor as well adding the ability to measure more atmospheric conditions. Assuming a standard atmosphere model it's also possible to convert pressure to altitude.%
\begin{figureptx}{Figure}{Pressure Sensor (Barometer) without header pins soldered}{pressure-2-4}{}%
\begin{image}{0.135}{0.73}{0.135}{}%
\includegraphics[width=\linewidth]{external/lps22.jpg}
\end{image}%
\tcblower
\end{figureptx}%
The sensor above is the \href{https://www.adafruit.com/product/4633}{\linkhilite{LPS22}}\hyperlink{AdafruitMAIN}{\linkhilite{[{\xreffont 34}]}}. The first sensor on board is a pressure sensor which uses a pressure transducer (something that creates a voltage when pressure is applied). The second is a temperature sensor which is more than likely a thermistor. The two of these sensors put together (pressure, temperature) is called a Barometer. This can bed purchased on Adafruit for only \textdollar{}7 at the time of this writing. The interesting thing about this device is that you can actually purchase the LPS22 and the temperature sensor separately but this breakout board has both chips on board. The goal of this lab is not only to use this specific sensor but to understand Barometers and reinforce I2C protocol like has been discussed in previous chapters (\hyperref[imu]{\linkhilite{Chapter~{\xreffont\ref{imu}}}} and \hyperref[accelerometer]{\linkhilite{Subsection~{\xreffont\ref{accelerometer}}}}).%
\end{introduction}%
%
%
\typeout{************************************************}
\typeout{Section 23.1 LPS Setup}
\typeout{************************************************}
%
\begin{sectionptx}{Section}{LPS Setup}{}{LPS Setup}{}{}{pressure-3}
When you open the packaging of this breakout board you'll notice that the header pins are missing. First you'll need to cut a row of 8 header pins and \alert{solder the header pins to the sensor}. If you're taking my class I can solder this for you or teach everyone about soldering during a lecture session of class. If you are taking this class elsewhere you have two options: try and find someone who can solder this real quick (only takes about 5 minutes) or buy your own soldering iron and try to solder yourself. Once the device is soldered you can "plug" it into a breadboard. The wiring for this system requires 4 wires. The figure below shows the pressure sensor connected to a \href{https://learn.adafruit.com/lis3mdl-triple-axis-magnetometer/python-circuitpython}{\linkhilite{Feather M4}} (Courtesy of \href{https://learn.adafruit.com/u/siddacious}{\linkhilite{Bryan Siepert}})\hyperlink{Adafruit}{\linkhilite{[{\xreffont 32}]}}. Notice that the \emph{SDA\slash{}SCL} pins on the LPS22 are listed \emph{SDI\slash{}SCK} respectively.%
\begin{figureptx}{Figure}{Pressure Sensor Connected to a Feather M4}{pressure-bryan}{}%
\begin{image}{0.135}{0.73}{0.135}{}%
\includegraphics[width=\linewidth]{external/adafruit_products_d_cp_feather_wiring_bb.png}
\end{image}%
\tcblower
\end{figureptx}%
\end{sectionptx}
%
%
\typeout{************************************************}
\typeout{Section 23.2 CPB Wiring and Software}
\typeout{************************************************}
%
\begin{sectionptx}{Section}{CPB Wiring and Software}{}{CPB Wiring and Software}{}{}{pressure-4}
The only difference between the wiring diagram shown previously and the CPX\slash{}CPB is that you will be using 4 alligator clips. The rest is straightforward. You need \emph{3.3V} to run to \emph{VIN}, \emph{GND} to \emph{GND} and then \emph{SDA} to \emph{SDI} and \emph{SCL} to \emph{SCK}. The two figures below show a CPB wired to the LPS22.%
\begin{figureptx}{Figure}{Pressure Sensor Connected to a CPB}{pressure-4-3}{}%
\begin{image}{0.135}{0.73}{0.135}{}%
\includegraphics[width=\linewidth]{external/lps22_CPB.jpeg}
\end{image}%
\tcblower
\end{figureptx}%
In the photos above and below, the red alligator clip is connected to \emph{3.3V} and is then connected to a red male-male wire which is then connected to the \emph{VIN} pin on the LPS22. The black alligator clip is connected to the \emph{GND} pin which is connected to a brown male-male wire connected to the \emph{GND} pin on the LPS22. The white alligator clip is connected to the \emph{SCL} pin connected to an orange male-male wire connected to the \emph{SCK} pin on the LPS22. Finally, the yellow alligator clip is connected to the \emph{SDA} pin connected to a yellow male-male wire connected to the \emph{SDI} pin on the LPS22.%
\begin{figureptx}{Figure}{Pressure Sensor Connected to a CPB with a Breadboard}{pressure-4-5}{}%
\begin{image}{0.135}{0.73}{0.135}{}%
\includegraphics[width=\linewidth]{external/lps22.jpeg}
\end{image}%
\tcblower
\end{figureptx}%
Once you have the circuit wired and soldered it's time to work on software. First, you want to make sure you have your \href{https://circuitpython.org/downloads}{\linkhilite{Circuit Python UF2}} up to date\hyperlink{CircuitPython}{\linkhilite{[{\xreffont 30}]}}. Once I updated my UF2 I also updated my \href{https://circuitpython.org/libraries}{\linkhilite{Circuit Python Libraries}} (See \hyperref[installing-modules]{\linkhilite{Section~{\xreffont\ref{installing-modules}}}} for help on installing extra modules on your CPX\slash{}CPB)\hyperlink{CircuitPython}{\linkhilite{[{\xreffont 30}]}}. In this lab I purposefully made the code a bit more elaborate by adding, bluetooth, and method 3 data logging. The reason is that my students in Aerospace Propulsion and Spacecraft Design have to build a hobbyist level rocket and get pressure data during the flight. This way they know how high their rocket went. As such they need to build a standalone datalogger that can be placed into a rocket. The specific modules I needed for this lab for all of that functionality is as follows:%
\begin{enumerate}
\item{}adafruit\textunderscore{}lis3dh%
\item{}adafruit\textunderscore{}lps2x%
\item{}adafruit\textunderscore{}thermistor%
\item{}neopixel%
\item{}adafruit\textunderscore{}ble%
\end{enumerate}
%
\par
Notice that I also added the thermistor and lis3dh libraries to get acceleration and temperature from the CPB. Once you have the necessary modules you can run some example code. \alert{Note that it's a good exercise when attaching I2C devices to run a scan on all connected I2C devices to make sure you can see the device and that you have the correct address (See \hyperref[I2C_scanner]{\linkhilite{Section~{\xreffont\ref{I2C_scanner}}}}).} Once you're positive the I2C address of the LPS22 is setup properly you can run the example code to \href{https://github.com/cmontalvo251/Microcontrollers/tree/master/Circuit_Playground/CircuitPython/Pressure_Accel_Datalogging}{\linkhilite{log pressure and acceleration data}}. Note that the link previously shown is to a folder of many many different pieces of software. The reason is because some students have a CPX that doesn't have bluetooth, some students want their CPX\slash{}CPB to beep and some students use a BME280 instead of an LPS22 to log data. Make sure to select the correct version for your hardware. Once you get it working the output will look something like this.%
\begin{figureptx}{Figure}{Pressure Data in Mu}{pressure-4-8}{}%
\begin{image}{0.135}{0.73}{0.135}{}%
\includegraphics[width=\linewidth]{external/Pressure_Logging.png}
\end{image}%
\tcblower
\end{figureptx}%
The data shown have multiple columns of data which are time,x accel,y accel,z accel,pressure,relative humidity,lps22 temperature, CPX\slash{}CPB temperature. Notice that the relative humidity is 0.0 because the LPS22 doesn't have a humidity sensor. If you were to use a BME280 instead of an LPS22 you would get relative humidity data as well. It's possible to convert the pressure data to altitude and is shown in \hyperref[pressure-altitude]{\linkhilite{Section~{\xreffont\ref{pressure-altitude}}}}.%
\end{sectionptx}
%
%
\typeout{************************************************}
\typeout{Section 23.3 Arduino Wiring and Software}
\typeout{************************************************}
%
\begin{sectionptx}{Section}{Arduino Wiring and Software}{}{Arduino Wiring and Software}{}{}{pressure-5}
The wiring for the Arduino is similar to any I2C setup like the external IMU (\hyperref[imu]{\linkhilite{Chapter~{\xreffont\ref{imu}}}}). In this case \emph{3.3V} and \emph{GND} are straightforward while \emph{SDI} and \emph{SCK} pins are connected to the dedicated \emph{SDA} and \emph{SCL} pins on the Arduino. The good news is that the entire Arduino family has dedicated \emph{SDA} and \emph{SCL} pins for I2C. The Figure below shows the pressure sensor connected to an Arduino MEGA.%
\begin{figureptx}{Figure}{Arduino MEGA hooked up to an LPS22, IMU and SD Card breakout board}{pressure-5-3}{}%
\begin{image}{0.135}{0.73}{0.135}{}%
\includegraphics[width=\linewidth]{external/Arduino_LPS_Accel_Hardware.jpg}
\end{image}%
\tcblower
\end{figureptx}%
The Figure above also has an external IMU (\hyperref[imu]{\linkhilite{Chapter~{\xreffont\ref{imu}}}}) and micro sd card breakout board (\hyperref[sdcard]{\linkhilite{Chapter~{\xreffont\ref{sdcard}}}}) connected to the Arduino to build a standalone datalogger capable of measuring acceleration and pressure and log that to a micro SD card. This is so students can place this whole setup into a rocket and measure pressure and acceleration while logging that to an SD card. The LPS22 sensor is the sensor on the left but unfortunately has header pins soldered upside down. As such the LPS22 sensor is upside down. Since you can't see the pins I placed another identical LPS22 right below it so the pins are shown. In this case, the blue wire is connected to the same \emph{3.3V} pin as the IMU, the left most green wire is connected to \emph{GND} on the Arduino, the next green wire is the \emph{SCK} pin connected to the same \emph{SCL} pin as the IMU, and the last dark green wire is connected to the \emph{SDI} pin on the LPS22 while the other end is connected to the same \emph{SDA} pin as the IMU. Note that I2C has multiple \textquotedblleft{}addresses\textquotedblright{}. This means you can have multiple different I2C sensors connected to the same \emph{SDA\slash{}SCL} pins and the code\slash{}hardware will know which sensor is where just like addresses on a street. Moving on to software, the code is similar to any I2C device in that it is necessary to get the LP22 library but also \emph{Adafruit\textunderscore{}BusIO} and \emph{Adafruit\textunderscore{}Sensor}. Since these libraries are from Adafruit rather than Arduino I find it easier to navigate to Github and download the zip libraries directly. The list of libraries you need and links to those websites are shown below%
%
\begin{enumerate}
\item{}\href{https://github.com/adafruit/Adafruit_LPS2X}{\linkhilite{Adafruit\textunderscore{}LPS2X}}%
\item{}\href{https://github.com/adafruit/Adafruit_BusIO}{\linkhilite{Adafruit\textunderscore{}BusIO}}%
\item{}\href{https://github.com/adafruit/adafruit_sensor}{\linkhilite{Adafruit\textunderscore{}Sensor}}%
\end{enumerate}
Once you have the zip folders downloaded you need to navigate to Sketch\textgreater{}Include Library\textgreater{}Add .ZIP Library... and then add the ZIP file you just downloaded. This is different than adding a library directly from the Arduino website as is done for the SD library in \hyperref[sdcard]{\linkhilite{Chapter~{\xreffont\ref{sdcard}}}}. The code for just reading the pressure sensor can be found in the example set of \emph{Adafruit\textunderscore{}LPS2X} and is done by navigating the toolbar as shown below.%
\begin{preformatted}
File>Examples>Adafruit_LPS2X>adafruit_lps2x_test.
\end{preformatted}
Just in case you're having trouble finding the example script, I've placed a \href{https://github.com/cmontalvo251/Microcontrollers/blob/master/Arduino/MEGA/LPS22_Test/LPS22_Test.ino}{\linkhilite{streamlined LPS22 copy}} of the example on Github. A snippet of that code is shown below.%
\begin{figureptx}{Figure}{Arduino MEGA Code for LPS22 Only}{pressure-5-9}{}%
\begin{image}{0.135}{0.73}{0.135}{}%
\includegraphics[width=\linewidth]{external/LPS22_Code.png}
\end{image}%
\tcblower
\end{figureptx}%
Notice that lines 2-4 of the software have the \emph{\#include} calls to the libraries just installed. \emph{Wire.h} is a pre-installed header file and \emph{Adafruit\textunderscore{}BusIO} is called from \emph{Adafruit\textunderscore{}Sensor} so you the code doesn't require an excplicit include. Once you've got the code opened and the libraries installed you should be able to compile without seeing any errors. After uploading the code the serial monitor should produce something similar to the figure below.%
\begin{figureptx}{Figure}{Serial Monitor Output for LPS22 Only}{pressure-5-11}{}%
\begin{image}{0.135}{0.73}{0.135}{}%
\includegraphics[width=\linewidth]{external/LPS22_Arduino_Output.png}
\end{image}%
\tcblower
\end{figureptx}%
Again, it's possible to convert pressure to altitude and is done in \hyperref[pressure-altitude]{\linkhilite{Section~{\xreffont\ref{pressure-altitude}}}}. Along with getting pressure data from the LPS22, it's possible to augment the software to also read the external IMU from I2C and log all data to the micro SD card. A snippet of the software is shown below while the software itself for a \href{https://github.com/cmontalvo251/Microcontrollers/blob/master/Arduino/MEGA/Accelerometer_LSM6DS33_Pressure_LPS22_SD/Accelerometer_LSM6DS33_Pressure_LPS22_SD.ino}{\linkhilite{standalone Pressure+IMU datalogger}} is on Github us usual.%
\begin{figureptx}{Figure}{Arduino MEGA Code of IMU, Pressue and SD Card Breakout boards}{pressure-5-13}{}%
\begin{image}{0.135}{0.73}{0.135}{}%
\includegraphics[width=\linewidth]{external/Accel_LPS22_Arduino_MEGA_Code.png}
\end{image}%
\tcblower
\end{figureptx}%
In that snippet of code, you can see the loop routine where the IMU and the Pressure sensors are polled and printed to the serial monitor together. Once that code is compiled and flashed to the Arduino the serial monitor output will create something similar to the figure below.%
\begin{figureptx}{Figure}{Serial Monitor Output for LPS22 and IMU}{pressure-5-15}{}%
\begin{image}{0.135}{0.73}{0.135}{}%
\includegraphics[width=\linewidth]{external/Pressure_Accel_Output_Arduino.png}
\end{image}%
\tcblower
\end{figureptx}%
\end{sectionptx}
%
%
\typeout{************************************************}
\typeout{Section 23.4 Converting Pressure to Altitude}
\typeout{************************************************}
%
\begin{sectionptx}{Section}{Converting Pressure to Altitude}{}{Converting Pressure to Altitude}{}{}{pressure-altitude}
The pressure data obtained from the CPX\slash{}CPB is in hectopascals (hPa) which is the same as millibars (mbar). The standard atmospheric pressure at sea level is 1013.25 hPa. As you go up in altitude the pressure decreases. The equation to convert pressure to altitude assuming a standard atmosphere is as follows:%
\begin{equation}
h = \frac{1.0-(\frac{p}{p_0}^{1.0/5.25588})}{2.2557 \times 10^{-5.0}}\label{pressure-altitude-2-1}
\end{equation}
where \emph{p} is pressure and \(p_0\) is the standard atmospheric pressure at sea level (1013.25 hPa). This equation assumes a standard atmosphere and is only accurate up to about 11 km in altitude. Beyond that the equation becomes less accurate and you would need to use a more complex model of the atmosphere to get accurate altitude data. However, for our purposes of launching a hobbyist rocket, this equation will work just fine. It's also possible to convert altitude to velocity using a first order derivative.%
%
\begin{equation}
v = \frac{h_{i+1}-h_i}{t_{i+1}-t_i}\label{pressure-altitude-3}
\end{equation}
In order to get that equation to have reasonable results, students typically have to heavily filter their data given the amount of noise in the data. The equation to convert to velocity is shown below but the results and filtering are left to the reader as an assignment. I think students have had a lot of success using moving average and complimentary filters. I'm sure Google and AI tools will be able to help get that working.%
\end{sectionptx}
%
%
\typeout{************************************************}
\typeout{Section 23.5 Data from a Rocket Launch}
\typeout{************************************************}
%
\begin{sectionptx}{Section}{Data from a Rocket Launch}{}{Data from a Rocket Launch}{}{}{pressure-7}
My students have placed this entire data logger into a rocket and launched it. The figure below shows the rocket on the launch pad in Samson,AL.%
\begin{figureptx}{Figure}{Rocket on the Launch Pad}{pressure-7-3}{}%
\begin{image}{0.135}{0.73}{0.135}{}%
\includegraphics[width=\linewidth]{external/rocket.png}
\end{image}%
\tcblower
\end{figureptx}%
The datalogger itself is attached to a piece of wood that is called a sled and then placed into the rocket prior to launch.%
\begin{figureptx}{Figure}{Datalogger Attached to Sled}{pressure-7-5}{}%
\begin{image}{0.135}{0.73}{0.135}{}%
\includegraphics[width=\linewidth]{external/pressure_sled.png}
\end{image}%
\tcblower
\end{figureptx}%
In the figure above the small gray piece below the CPB is a small 1S Lipo battery and the LPS22 is on the other side of the sled. The figure below shows pressure as a function of time.%
\begin{figureptx}{Figure}{Pressure Data from a Rocket Launch}{pressure-7-7}{}%
\begin{image}{0.135}{0.73}{0.135}{}%
\includegraphics[width=\linewidth]{external/pressure_rocket.png}
\end{image}%
\tcblower
\end{figureptx}%
It's clear from the data that pressure drops as the flight continues showing the boost phase of the rocket. There is a small oscillation near the apogee of the flight. This is due to the delay charge to blow the nose cone off the rocket body to deploy the parachute. The pressure sensor is measuring local atmospheric pressure. So if there is a small explosion inside the rocket to deploy the parachute, the data will show a slight oscillation.%
\begin{figureptx}{Figure}{Altitude Data from a Rocket Launch}{pressure-7-9}{}%
\begin{image}{0.135}{0.73}{0.135}{}%
\includegraphics[width=\linewidth]{external/altitude_rocket.png}
\end{image}%
\tcblower
\end{figureptx}%
The figure above shows the pressure data converted to altitude using the equation shown above. Note that in order to get an altitude of zero at time (t=0) it is important to set \(p_0\) equal to the pressure at time (t=0). This way the altitude will be zero at the start of the flight and will increase as the rocket goes up. Again the altitude plot shows a small oscillation in altitude which isn't true but indicative of the explosion to deploy the parachute.%
\end{sectionptx}
%
%
\typeout{************************************************}
\typeout{Section 23.6 Extra Help}
\typeout{************************************************}
%
\begin{sectionptx}{Section}{Extra Help}{}{Extra Help}{}{}{pressure-8}
There is one youtube video on \href{https://www.youtube.com/watch?v=MKq2glxVkMI}{\linkhilite{logging pressure data with the LPS22}}. If you need a refresher on Method 3 be sure to go back to \hyperref[daq]{\linkhilite{Chapter~{\xreffont\ref{daq}}}} for more help on that.%
\end{sectionptx}
%
%
\typeout{************************************************}
\typeout{Section 23.7 Assignment}
\typeout{************************************************}
%
\begin{sectionptx}{Section}{Assignment}{}{Assignment}{}{}{pressure-9}
\begin{introduction}{}%
For this assignment you are to wire up the external pressure sensor and get data from it. You need to mount the CPX\slash{}CPB and the pressure sensor to some sort of hard surface so that when you move the CPX\slash{}CPB the pressure moves as well. In order to get the pressure to change significantly you'll need to change your altitude by a lot. My recommendation is to ride the elevator of a tall building or drive up a really big hill. If you have the ability to do so, place the sensor in a rocket or an airplane!%
%
\begin{enumerate}
\item{}Title page with name,title and date - 5\%%
\item{}Introduction with wiring diagram - 5\%%
\item{}All figures must have appropriate annotations, figure captions, figure labels and a paragraph of supporting text - 10\%%
\item{}Include a photo of your CPX\slash{}CPB mounted to a hard surface with the pressure sensor on a breadboard. - 10\%%
\item{}Include a selfie of you right before you change altitude either by going up a big hill or by riding the elevator. Explain how you plan on changing altitude significantly in the supporting text. - 10\%%
\item{}Include a screenshot (entire computer screen) of Mu showing the Plotter open showing pressure - 20\%%
\item{}Plot pressure vs time for your experiment where pressure significantly changes. - 20\%%
\item{}Convert pressure to altitude and plot altitude vs time for your experiment where altitude significantly changes. - 20\%%
\item{}Appendix A - Video Details - Film yourself and your entire screen showing the pressure value in the serial Monitor in Mu. Also explain what you plan on doing to significantly change the altitude.- Pass\slash{}Fail%
\item{}Appendix B - Code used to take data and code used to plot in plain text (no screenshots) - Pass\slash{}Fail%
\item{}Appendix C - Screenshot of Excel or Sheets if applicable - Pass\slash{}Fail%
\item{}\alert{If you are asked to fix something in the Appendices such as missing content in your video or screenshots of your code, you will lose 10\% of your grade everytime you submit an update after the deadline.}%
\end{enumerate}
\end{introduction}%
\end{sectionptx}
\end{chapterptx}
%
%
\typeout{************************************************}
\typeout{Chapter 24 Global Positioning System (GPS)}
\typeout{************************************************}
%
\begin{chapterptx}{Chapter}{Global Positioning System (GPS)}{}{Global Positioning System (GPS)}{}{}{gps}
\renewcommand*{\chaptername}{Chapter}
\begin{introduction}{}%
\begin{listptx}{List}{\textbf{Parts List}}{gps-2-1}{}%
%
\begin{enumerate}
\item{}Laptop%
\item{}CPX\slash{}CPB or Arduino%
\item{}USB Cable%
\item{}\href{https://www.adafruit.com/product/746}{\linkhilite{GPS}} (Note that other GPS modules will work for this lab provided they output standard NMEA sentences)\hyperlink{AdafruitMAIN}{\linkhilite{[{\xreffont 34}]}}%
\item{}Alligator Clips (x4)%
\item{}Bread Board%
\item{}Soldering iron%
\end{enumerate}
\end{listptx}%
\begin{listptx}{List}{\textbf{Learning Objectives}}{gps-2-2}{}%
%
\begin{enumerate}
\item{}Continue to understand the concept of soldering%
\item{}Understand serial communication using TX\slash{}RX%
\item{}Learn the components of a GPS module%
\item{}Learn about standard NMEA sentences output by GPS modules%
\item{}Read GPS data and plot%
\end{enumerate}
\end{listptx}%
In this lab we're going to use a GPS module to obtain location information, convert it to distance in meters and velocity.%
\begin{figureptx}{Figure}{GPS Module without header pins soldered}{gps-2-4}{}%
\begin{image}{0.135}{0.73}{0.135}{}%
\includegraphics[width=\linewidth]{external/gps.jpg}
\end{image}%
\tcblower
\end{figureptx}%
The sensor above is the \href{https://www.adafruit.com/product/746}{\linkhilite{GPS}}\hyperlink{AdafruitMAIN}{\linkhilite{[{\xreffont 34}]}}. The large square is the GPS receiver and the breakout board sends serial data via TX\slash{}RX pins to any microcontroller listening. So this is different than I2C as has been done in the previous labs with accelerometers (\hyperref[accelerometer]{\linkhilite{Subsection~{\xreffont\ref{accelerometer}}}}), IMUs (\hyperref[imu]{\linkhilite{Chapter~{\xreffont\ref{imu}}}}), or pressure sensors (\hyperref[pressure]{\linkhilite{Chapter~{\xreffont\ref{pressure}}}}) but due to the high level coding it's still just as simple to parse the data . This can be purchased on Adafruit for \textdollar{}30 at the time of this writing. The interesting thing about this device is that you can actually purchase the GPS separately but this breakout board has simple \emph{TX\slash{}RX} outputs for ease of communication. The goal of this lab is not only to use this specific sensor but to understand GNSS and reinforce UART protocol.%
\end{introduction}%
%
%
\typeout{************************************************}
\typeout{Section 24.1 GPS Setup}
\typeout{************************************************}
%
\begin{sectionptx}{Section}{GPS Setup}{}{GPS Setup}{}{}{gps-3}
When you open the packaging of this breakout board you'll notice that the header pins are missing. First you'll need to cut a row of 9 pins and \alert{solder the header pins to the sensor}. If you're taking my class I can solder this for you or teach everyone about soldering during a lecture session of class. If you are taking this class elsewhere you have two options: try and find someone who can solder this real quick (only takes about 5 minutes) or buy your own soldering iron and try to solder yourself. Once the device is soldered you can "plug" it into a breadboard. The wiring for this system requires 4 wires. The figure below shows the GPS sensor connected to an \href{https://learn.adafruit.com/adafruit-ultimate-gps/circuitpython-parsing}{\linkhilite{Adafruit Metro M0}}. (Courtesy of \href{https://learn.adafruit.com/u/adafruit2}{\linkhilite{Lady Ada}} herself!)\hyperlink{Adafruit}{\linkhilite{[{\xreffont 32}]}}.%
\begin{figureptx}{Figure}{GPS Wiring Diagram for Adafruit Metro M0\hyperlink{gps-wiring}{\linkhilite{[{\xreffont 46}]}}}{gps-3-3}{}%
\begin{image}{0.135}{0.73}{0.135}{}%
\includegraphics[width=\linewidth]{external/adafruit_products_m0_express_gps_bb.png}
\end{image}%
\tcblower
\end{figureptx}%
\end{sectionptx}
%
%
\typeout{************************************************}
\typeout{Section 24.2 CPB Wiring and Software}
\typeout{************************************************}
%
\begin{sectionptx}{Section}{CPB Wiring and Software}{}{CPB Wiring and Software}{}{}{gps-4}
The only difference between the wiring diagram shown previously and the CPX\slash{}CPB is that you will be using 4 alligator clips. The rest is straightforward. You need \emph{3.3V} to run to \emph{VIN}, \emph{GND} to \emph{GND} and then \emph{TX} to \emph{RX} and \emph{RX} to \emph{TX}. That is the \emph{TX} and \emph{RX} pins are swapped so that the transmit wire goes to the receive wire on the other device.%
\begin{figureptx}{Figure}{GPS Module Connected to a CPB}{gps-4-3}{}%
\begin{image}{0.135}{0.73}{0.135}{}%
\includegraphics[width=\linewidth]{external/gps_CPB.jpeg}
\end{image}%
\tcblower
\end{figureptx}%
In the photos above and below, the red alligator clip is connected to \emph{3.3V} and is then connected to a red male-male wire which is then connected to the \emph{VIN} pin on the GPS. The black alligator clip is connected to the \emph{GND} pin which is connected to a brown male-male wire connected to the \emph{GND} pin on the GPS. The white alligator clip is connected to the \emph{TX} pin connected to an orange male-male wire connected to the \emph{RX} pin on the GPS. Finally, the yellow alligator clip is connected to the \emph{RX} pin connected to a yellow male-male wire connected to the \emph{TX} pin on the GPS.%
\begin{figureptx}{Figure}{GPS Module Connected to a CPB with a Breadboard}{gps-4-5}{}%
\begin{image}{0.135}{0.73}{0.135}{}%
\includegraphics[width=\linewidth]{external/gps.jpeg}
\end{image}%
\tcblower
\end{figureptx}%
Once you have the circuit wired and soldered it's time to work on software. First, you want to make sure you have your \href{https://circuitpython.org/downloads}{\linkhilite{Circuit Python UF2}} up to date\hyperlink{CircuitPython}{\linkhilite{[{\xreffont 30}]}}. Once I updated my UF2 I also updated my \href{https://circuitpython.org/libraries}{\linkhilite{Circuit Python Libraries}} (See \hyperref[installing-modules]{\linkhilite{Section~{\xreffont\ref{installing-modules}}}} for help on installing extra modules on your CPX\slash{}CPB)\hyperlink{CircuitPython}{\linkhilite{[{\xreffont 30}]}}. In this lab I purposefully made the code a bit more elaborate by adding, bluetooth, and method 3 data logging. The reason is that my students in Flight Dynamics and Aircraft Design have to build a hobbyist radio controlled aircraft and get GPS data during the flight. This way they know where they flew and how fast. As such they need to build a standalone datalogger that can be placed into an airplane. The specific modules I needed for this lab for all of that functionality is as follows:%
\begin{enumerate}
\item{}adafruit\textunderscore{}lis3dh%
\item{}adafruit\textunderscore{}gps%
\item{}adafruit\textunderscore{}thermistor%
\item{}neopixel%
\item{}adafruit\textunderscore{}ble%
\end{enumerate}
%
\par
Notice that I also added the thermistor and lis3dh libraries to get acceleration and temperature from the CPB. Once you have the necessary modules you can run some example code. Once you have all the modules, you can run the example code to \href{https://github.com/cmontalvo251/Microcontrollers/tree/master/Circuit_Playground/CircuitPython/GPS}{\linkhilite{log gps and acceleration data}}. Note that the link previously shown is to a folder of a few different pieces of software. The reason is because some students have a CPX that doesn't have bluetooth. Make sure to select the correct version for your hardware. Once you get it working the output will look something like this.%
\begin{figureptx}{Figure}{GPS Data in Mu}{gps-4-8}{}%
\begin{image}{0.135}{0.73}{0.135}{}%
\includegraphics[width=\linewidth]{external/GPS_No_Lock_Mu.PNG}
\end{image}%
\tcblower
\end{figureptx}%
Notice that in the photo above the multiple columns of data are time,-99,-99,-99,0.0,x accel,y accel, z accel. The -99 values are dummy values for latitude, longitude and altitude because the GPS doesn't have a lock on any satellites. The 0.0 is the speed in knots which is also zero because there is no lock. The x,y,z accel values are from the accelerometer on the CPB. Move to \hyperref[gps-explanation]{\linkhilite{Section~{\xreffont\ref{gps-explanation}}}} to learn about getting a lock and how to process the GPS information once you get a lock.%
\end{sectionptx}
%
%
\typeout{************************************************}
\typeout{Section 24.3 Arduino Wiring and Software}
\typeout{************************************************}
%
\begin{sectionptx}{Section}{Arduino Wiring and Software}{}{Arduino Wiring and Software}{}{}{gps-5}
The wiring for the Arduino depends on which type of Arduino you're using. For every Arduino the \emph{Serial0} pins are hardwired to the USB port so you can't use them which means you need to use \emph{Serial1} or similar. The problem is that the Arduino UNO does not have any more serial ports so you have to use \emph{SoftwareSerial}. When using \emph{SoftwareSerial} you can plug the \emph{TX\slash{}RX} pins into any two digital ports and then create a \emph{SoftwareSerial} instance. If you're using an Arduino MEGA you can use any of the hardware Serial ports. In the figure below the GPS is plugged into \emph{Serial1} and thus \emph{TX} on the GPS is connected to \emph{RX1} on the Arduino MEGA and \emph{RX} on the GPS is connected to \emph{TX1} on the Arduino MEGA. The Figure below also has an IMU (\hyperref[imu]{\linkhilite{Chapter~{\xreffont\ref{imu}}}}) and micro sd card breakout board (\hyperref[sdcard]{\linkhilite{Chapter~{\xreffont\ref{sdcard}}}}) connected to the Arduino to build a standalone datalogger capable of measuring acceleration and GPS and log that to a micro SD card.%
\begin{figureptx}{Figure}{Arduino MEGA Hooked up to a GPS and external IMU}{gps-5-3}{}%
\begin{image}{0.135}{0.73}{0.135}{}%
\includegraphics[width=\linewidth]{external/Arduino_GPS_Accel_Hardware.png}
\end{image}%
\tcblower
\end{figureptx}%
In the Figure above, the blue wire goes from \emph{VIN} to \emph{5V} on the Arduino, the left most green wire goes from \emph{GND} to \emph{GND} while the next two green wires connect \emph{TX\slash{}RX} to \emph{RX1\slash{}TX1} respectively on the Arduino MEGA. Moving on to software also depends on the Arduino MEGA or UNO. On the UNO since you're using \emph{SoftwareSerial} you need to ensure that you have \emph{\#include \textless{}SoftwareSerial.h\textgreater{}} at the top of your code and that the \emph{SoftwareSerial} library is installed in Arduino IDE. Thankfully \emph{SoftwareSerial} is already built-in to the Arduino IDE so there's no need to download it. The code for an \href{https://github.com/cmontalvo251/Microcontrollers/blob/master/Arduino/UNO/GPS/GPS_Shield_Print/GPS_Shield_Print.ino}{\linkhilite{Arduino UNO example}} can be found on Github but will not be shown in this chapter. Notice though that on line 33 upon following that link, pins 8 and 7 are used for the \emph{SoftwareSerial} port given by this line of code: \emph{SoftwareSerial mySerial(8, 7);}. For the \href{https://github.com/cmontalvo251/Microcontrollers/blob/master/Arduino/MEGA/GPS_Only/GPS_Only.ino}{\linkhilite{Arduino MEGA}}, the code is also on Github but the \emph{Serial} port is different as shown in the code below.%
\begin{figureptx}{Figure}{Arduino MEGA Code for GPS Only}{gps-5-5}{}%
\begin{image}{0.135}{0.73}{0.135}{}%
\includegraphics[width=\linewidth]{external/Arduino_MEGA_GPS_Only.png}
\end{image}%
\tcblower
\end{figureptx}%
Notice, in this code the call to create the \emph{GPS} object is done using \emph{\&Serial1} rather than \emph{\&mySerial} which is an instance of the \emph{SoftwareSerial} library. While \emph{\&Serial1} is simply built-in to the Arduino MEGA software suite \emph{\&mySerial} is a specific instance of \emph{SoftwareSerial} that is only required on the Arduino UNO or similar boards that don't have Hardware Serial capabilities. However, for either the Arduino UNO or MEGA it is still required to get the \href{https://github.com/adafruit/Adafruit_GPS}{\linkhilite{Adafruit\textunderscore{}GPS library}} which can be downloaded from Github and installed by navigating to Sketch\textgreater{}Include Library\textgreater{}Add .ZIP Library... and then add the ZIP file you just downloaded. This library is needed to parse the specific raw GPS sentences (\hyperref[gps-explanation]{\linkhilite{Section~{\xreffont\ref{gps-explanation}}}}) coming in. Once you've installed all the libraries, compiled the code, wired up the Arduino with the GPS and flashed the code you should be able to open the Serial Monitor with a baud rate of 115200 and see the GPS output.%
\begin{figureptx}{Figure}{Serial monitor output for GPS Only}{gps-5-7}{}%
\begin{image}{0.135}{0.73}{0.135}{}%
\includegraphics[width=\linewidth]{external/GPS_Arduino_Output.png}
\end{image}%
\tcblower
\end{figureptx}%
The Figure above shows the GPS output when you don't have a lock on any satellites which is why Latitude and Longitude have zeros. Still, this is a good check to make sure you're getting information correctly from the GPS since the date and time is typically accurate even without a lock. Move to \hyperref[gps-explanation]{\linkhilite{Section~{\xreffont\ref{gps-explanation}}}} to learn about getting a lock and how to process the GPS information once you get a lock. Along with getting a lock on the GPS, it's possible to augment the software to also read the external IMU from I2C and log all data to the micro SD card. A snippet of the software is shown below while the software itself for a \href{https://github.com/cmontalvo251/Microcontrollers/blob/master/Arduino/MEGA/Accelerometer_GPS_LSM6DS33_SD/Accelerometer_GPS_LSM6DS33_SD.ino}{\linkhilite{standalone GPS+IMU datalogger}} is on Github us usual.%
\begin{figureptx}{Figure}{Arduino MEGA Code for GPS,IMU and SD card breakout boards}{gps-5-9}{}%
\begin{image}{0.135}{0.73}{0.135}{}%
\includegraphics[width=\linewidth]{external/Arduino_GPS_IMU_Code.png}
\end{image}%
\tcblower
\end{figureptx}%
The code shows the loop routine polling the GPS using the \emph{useInterrupt} function as well as the \emph{sensors.getEvent()} routine to poll the IMU. After compiling the code, verifying the hardware connections and flashing the software the Serial monitor will produce something similar to the code below. Again in this case the latitude and longitude are shown as zeros because there is no lock.%
\begin{figureptx}{Figure}{Serial Monitor Output for LPS22 and IMU}{gps-5-11}{}%
\begin{image}{0.135}{0.73}{0.135}{}%
\includegraphics[width=\linewidth]{external/Arduino_GPS_Accel_Output.png}
\end{image}%
\tcblower
\end{figureptx}%
\end{sectionptx}
%
%
\typeout{************************************************}
\typeout{Section 24.4 GPS Receiver and Sentence Explanation}
\typeout{************************************************}
%
\begin{sectionptx}{Section}{GPS Receiver and Sentence Explanation}{}{GPS Receiver and Sentence Explanation}{}{}{gps-explanation}
The GPS constellation of satellites orbiting planet Earth are around 26 satellites. Typically there are about 13 satellites above you from horizon to horizon. However, with cloud cover and buildings and trees you can typically only receive data from about 7 of them. You need at least 3 but 4 is preferred to get a GPS lock. If you go outside on a clear data and wait your data will return coordinates for your current position. Latitude is your degree value above or below the equator, longitude is your degree value east or west of the prime meridian and then altitude is your height above sea level. Note that if you're code is printing raw NMEA sentences, you may get something like this\begin{preformatted}
$GPGGA,123519,4807.038,N,01131.000,E,1,08,0.9,545.4,M,46.9,M,,*47
\end{preformatted}
This is the raw NMEA sentence output by the GPS. The code in the example folder parses this data to get the latitude, longitude, altitude and speed. The NMEA sentence is a command that contains all the GPS data. The first part is the type of sentence (GPGGA in this case) and then the rest of the data is separated by commas. The format is type of sentence, time, latitude, N\slash{}S, longitude, E\slash{}W, fix quality, number of satellites, horizontal dilution of position, altitude, M (meters), height of geoid above WGS84 ellipsoid, M (meters), time since last DGPS update, DGPS reference station id, checksum. The code in either Arduino (\emph{Adafruit\textunderscore{}GPS}) or CPX\slash{}CPB (\emph{adafruit\textunderscore{}gps}) parses this data to get the latitude, longitude, altitude and speed. The code then prints this data in a more readable format. The code also converts the latitude and longitude from the NMEA format to decimal degrees which is easier to work with. Again, once you go outside and get a lock from more than 3 satellites, the latitude and longitude values will be replaced with your actual coordinates and speed. It's possible to convert GPS position to X and Y coordinate provided you use the first GPS coordinate as the origin. The equations to convert GPS coordinates to position are shown below.%
\begin{equation}
x = R cos(\lambda)(\phi - \phi_0) \label{gps-explanation-2-4}
\end{equation}
%
\begin{equation*}
y = R (\lambda - \lambda_0)
\end{equation*}
where \(R\) is the radius of the Earth, \(\phi\) is the longitude, \(\lambda\) is the latitude and \(\lambda_0\) and \(\phi_0\) are the latitude and longitude of the origin point. In this case the x-axis is East-West and the y-axis is North-South. Once you have position you can take the derivative to get velocity typically doing a heavily filtered moving average derivative to get a smooth velocity plot. To start though a simple first order derivative will at least get you started.%
\begin{equation}
v_x = \frac{x_{i+1}-x_i}{t_{i+1}-t_i} \label{gps-explanation-2-11}
\end{equation}
%
\begin{equation*}
v_y = \frac{y_{i+1}-y_i}{t_{i+1}-t_i}
\end{equation*}
%
\end{sectionptx}
%
%
\typeout{************************************************}
\typeout{Section 24.5 Results from an aircraft flight}
\typeout{************************************************}
%
\begin{sectionptx}{Section}{Results from an aircraft flight}{}{Results from an aircraft flight}{}{}{gps-7}
My students have placed this entire data logger into an airplane and flown it. The aircraft itself was a P-51 clone made out of balsa and clear monokote.%
\begin{figureptx}{Figure}{Aircraft with my students after their maiden flight}{gps-7-3}{}%
\begin{image}{0.135}{0.73}{0.135}{}%
\includegraphics[width=\linewidth]{external/Aircraft.png}
\end{image}%
\tcblower
\end{figureptx}%
The datalogger itself was soldered to a custom printed circuit board courtesy of Aker Solutions and permanently soldered to a \href{https://www.adafruit.com/product/4600}{\linkhilite{QtPy}}\hyperlink{Adafruit}{\linkhilite{[{\xreffont 32}]}}. The QtPy is a super small microcontroller without any extra sensors so the students only measured GPS and nothing else.%
\begin{figureptx}{Figure}{Datalogger inside balsa airplane}{gps-7-5}{}%
\begin{image}{0.135}{0.73}{0.135}{}%
\includegraphics[width=\linewidth]{external/GPS_Datalogger_Aircraft.png}
\end{image}%
\tcblower
\end{figureptx}%
After the flight at the ACMA (Azalea City of Model Aeronautics) airfield in Irvington, AL, the students were able to plot the GPS position of the flight. The color code indicates speed throughout the flight.%
\begin{figureptx}{Figure}{Latitude vs Longitude for Aircraft Flight}{gps-7-7}{}%
\begin{image}{0.135}{0.73}{0.135}{}%
\includegraphics[width=\linewidth]{external/GPS_Flight.png}
\end{image}%
\tcblower
\end{figureptx}%
It's clear from the data that the pilot took off and landed on the runway after flying a few loops around the radio controlled flying field.%
\end{sectionptx}
%
%
\typeout{************************************************}
\typeout{Section 24.6 Extra Help}
\typeout{************************************************}
%
\begin{sectionptx}{Section}{Extra Help}{}{Extra Help}{}{}{gps-8}
There is one youtube video on \href{https://www.youtube.com/watch?v=h1u0EiZqVLI}{\linkhilite{logging GPS data with the Adafruit Ultimate GPS}}. If you need a refresher on Method 3 be sure to go back to \hyperref[daq]{\linkhilite{Chapter~{\xreffont\ref{daq}}}} for more help on that.%
\end{sectionptx}
%
%
\typeout{************************************************}
\typeout{Section 24.7 Assignment}
\typeout{************************************************}
%
\begin{sectionptx}{Section}{Assignment}{}{Assignment}{}{}{gps-9}
\begin{introduction}{}%
For this assignment you are to wire up the external GPS and get data from it. You need to mount the CPX\slash{}CPB and the GPS sensor to some sort of hard surface so that when you move the CPX\slash{}CPB the GPS moves as well. In order to get a significant change in GPS coordinates I recommend going for a bike ride or a walk. If you decide to use a car I strongly recommend having a partner drive the car so you can focus on data acquisition and plotting. Once you have your GPS coordinates, convert them to position and velocity.%
%
\begin{enumerate}
\item{}Title page with name,title and date - 5\%%
\item{}Introduction with wiring diagram - 5\%%
\item{}All figures must have appropriate annotations, figure captions, figure labels and a paragraph of supporting text - 10\%%
\item{}Include a photo of your CPX\slash{}CPB mounted to a hard surface with the GPS sensor on a breadboard. - 10\%%
\item{}Include a selfie of you right before you change location either by going for a bike ride or a walk. Explain how you plan on changing location significantly in the supporting text. - 10\%%
\item{}Include a screenshot (entire computer screen) of Mu showing the Plotter open showing GPS data - 10\%%
\item{}Plot GPS coordinates (Latitude vs Longitude) for your experiment where coordinates significantly change. - 10\%%
\item{}Convert GPS coordinates to position and plot (X vs Y) for your experiment where position significantly changes. - 20\%%
\item{}Convert your position to velocity and plot velocity vs time for your experiment where velocity significantly changes. Comment on whether your velocity is accurate to the speed you were traveling - 20\%%
\item{}Appendix A - Video Details - Film yourself and your entire screen showing your GPS getting a lock with the Serial monitor open in Mu. Also explain what you plan on doing to significantly change the GPS coordinates. - Pass\slash{}Fail%
\item{}Appendix B - Code used to take data and code used to plot in plain text (no screenshots) - Pass\slash{}Fail%
\item{}Appendix C - Screenshot of Excel or Sheets if applicable - Pass\slash{}Fail%
\item{}\alert{If you are asked to fix something in the Appendices such as missing content in your video or screenshots of your code, you will lose 10\% of your grade everytime you submit an update after the deadline.}%
\end{enumerate}
\end{introduction}%
\end{sectionptx}
\end{chapterptx}
%
%
\typeout{************************************************}
\typeout{Chapter 25 Load Cells}
\typeout{************************************************}
%
\begin{chapterptx}{Chapter}{Load Cells}{}{Load Cells}{}{}{load-cell}
\renewcommand*{\chaptername}{Chapter}
\begin{introduction}{}%
\begin{listptx}{List}{\textbf{Parts List}}{load-cell-2-1}{}%
%
\begin{enumerate}
\item{}Laptop%
\item{}CPX\slash{}CPB or Arduino%
\item{}USB Cable%
\item{}\href{https://www.adafruit.com/product/5974}{\linkhilite{HX711 Load Cell Amplifier}} (Note that other amplifiers will work)\hyperlink{AdafruitMAIN}{\linkhilite{[{\xreffont 34}]}}%
\item{}\href{https://www.adafruit.com/product/5974}{\linkhilite{Strain Gauge Load Cell}} \hyperlink{AdafruitMAIN}{\linkhilite{[{\xreffont 34}]}}%
\item{}Alligator Clips (x4)%
\item{}Bread Board%
\item{}Soldering iron%
\item{}Digital Scale%
\end{enumerate}
\end{listptx}%
\begin{listptx}{List}{\textbf{Learning Objectives}}{load-cell-2-2}{}%
%
\begin{enumerate}
\item{}Continue to understand the concept of soldering%
\item{}Load Cell calibration%
\item{}Read load cell data and plot it%
\end{enumerate}
\end{listptx}%
In this lab we're going to use a strain gauge load cell and load cell amplifier to make a digital scale.%
\begin{figureptx}{Figure}{Strain Gauge Load Cell}{load-cell-2-4}{}%
\begin{image}{0.135}{0.73}{0.135}{}%
\includegraphics[width=\linewidth]{external/load_cell.jpg}
\end{image}%
\tcblower
\end{figureptx}%
The sensor above is the \href{https://www.adafruit.com/product/4543}{\linkhilite{a strain gauge load cell}}\hyperlink{AdafruitMAIN}{\linkhilite{[{\xreffont 34}]}}. The load cell itself is mounted to a metal beam with large holes in it so it can deflect when a load is placed on to it. The strain gauge itself changes its resistance when a load is applied but the output voltage is so low that an amplifier is needed.%
\begin{figureptx}{Figure}{Load Cell Amplifier}{load-cell-2-6}{}%
\begin{image}{0.135}{0.73}{0.135}{}%
\includegraphics[width=\linewidth]{external/load_cell_amplifier.jpg}
\end{image}%
\tcblower
\end{figureptx}%
The figure above shows the \href{https://www.adafruit.com/product/5974}{\linkhilite{HX711 load cell amplifier}} which takes the signal from the load cell and amplifies it and converts it to a nice clean digital signal similar to I2C but not quite. In this chapter we're going to wire up the load cell and then read the data from amplifier after calibrating it properly.%
\end{introduction}%
%
%
\typeout{************************************************}
\typeout{Section 25.1 Load Cell Setup}
\typeout{************************************************}
%
\begin{sectionptx}{Section}{Load Cell Setup}{}{Load Cell Setup}{}{}{load-cell-3}
When you open the packaging of the amplifier and strain gauge, you'll notice that the header pins on the amplifier are missing. First you'll need to cut a row of 6 pins and \alert{solder the header pins to the sensor}. If you're taking my class I can solder this for you or teach everyone about soldering during a lecture session of class. If you are taking this class elsewhere you have two options: try and find someone who can solder this real quick (only takes about 5 minutes) or buy your own soldering iron and try to solder yourself. Once the device is soldered you can "plug" it into a breadboard. The wiring for this system requires 4 wires from the microcontroller to the amplifier and then 4 wires from the strain gauge to the amplifier. The figure below shows the load cell and load cell amplifier connected to a \href{https://learn.adafruit.com/adafruit-nau7802-24-bit-adc-stemma-qt-qwiic/python-circuitpython}{\linkhilite{Feather M4}}. (Courtesy of \href{https://learn.adafruit.com/u/BlitzCityDIY}{\linkhilite{Liz Clark}})\hyperlink{Adafruit}{\linkhilite{[{\xreffont 32}]}}.%
\begin{figureptx}{Figure}{Load Cell Wiring Diagram for Adafruit Feather M4\hyperlink{load-cell-wiring}{\linkhilite{[{\xreffont 47}]}}}{load-cell-3-3}{}%
\begin{image}{0.135}{0.73}{0.135}{}%
\includegraphics[width=\linewidth]{external/adafruit_products_featherBB_revC_bb.png}
\end{image}%
\tcblower
\end{figureptx}%
In the wiring diagram above, the strain gauge is connected to the amplifier with 4 wires. The red wire is connected to \emph{E+}, the black wire is connected to \emph{E-}, the white wire is connected to \emph{A-} and the green wire is connected to \emph{A+}.%
\end{sectionptx}
%
%
\typeout{************************************************}
\typeout{Section 25.2 CPX\slash{}CPB Wiring and Software}
\typeout{************************************************}
%
\begin{sectionptx}{Section}{CPX\slash{}CPB Wiring and Software}{}{CPX\slash{}CPB Wiring and Software}{}{}{load-cell-4}
For the CPX\slash{}CPB the strain gauge and amplifier are the same but you will be using 4 alligator clips as shown in the photo below. Although very non-standard, there are two yellow alligator clips. The first yellow alligator clip is connected to \emph{3.3V} and is then connected to a purple male-male wire which is then connected to the \emph{VIN} pin on the amplifier. The white alligator clip is connected to the \emph{GND} pin which is connected to a brown male-male wire connected to the \emph{GND} pin on the amplifier. The green alligator clip is connected to the \emph{A2} pin connected to an white male-male wire connected to the \emph{SCK} pin on the amplifier. Finally, the second yellow alligator clip is connected to the \emph{A1} pin connected to a green male-male wire connected to the \emph{DATA} pin on the amplifier.%
\begin{figureptx}{Figure}{Load Cell Amplifier connected to a CPX}{load-cell-4-3}{}%
\begin{image}{0.135}{0.73}{0.135}{}%
\includegraphics[width=\linewidth]{external/CPX_HX711_Wiring.png}
\end{image}%
\tcblower
\end{figureptx}%
This isn't as straightforward because the amplifier in the parts list has an \emph{SCK} and \emph{DATA} line instead of \emph{SDA\slash{}SCL}. The good news is that it still uses digital output so the CPX\slash{}CPB can read it. So not quite I2C but still digital communication.%
\begin{figureptx}{Figure}{Strain Gauge Connected to a Load Cell Amplifier}{load-cell-4-5}{}%
\begin{image}{0.135}{0.73}{0.135}{}%
\includegraphics[width=\linewidth]{external/HX711_Amplifier.png}
\end{image}%
\tcblower
\end{figureptx}%
Once you have the circuit wired there are no extra modules to install since we're just going to be using the digitalio library. The example code on Github \href{https://github.com/cmontalvo251/Microcontrollers/blob/master/Circuit_Playground/CircuitPython/HX711/hx711_mytest.py}{\linkhilite{reads and calibrates the load cell data}} and converts it to kg for ease of use. Note that kg is not a unit of weight but since we're on Earth we're assuming that we're under 1g acceleration and thus kg can be used as weight even though it's actually a mass. \begin{figureptx}{Figure}{Load Cell Output in Kg in Mu}{load-cell-4-6-2}{}%
\begin{image}{0.135}{0.73}{0.135}{}%
\includegraphics[width=\linewidth]{external/HX711_Screenshot.png}
\end{image}%
\tcblower
\end{figureptx}%
 The photo above shows the example output of the load cell in kg by applying a load to the strain gauge. Using this simple script it is easy to weigh certain objects and get the value from the load cell. Notice however, that on line 73 at the time of this writing, the scale value is set to 12,000. That value will need to be adjusted based on your own calibration measurements which is left as an assignment to the reader.%
\end{sectionptx}
%
%
\typeout{************************************************}
\typeout{Section 25.3 Arduino Wiring and Software}
\typeout{************************************************}
%
\begin{sectionptx}{Section}{Arduino Wiring and Software}{}{Arduino Wiring and Software}{}{}{load-cell-5}
The wiring for the Arduino family is extremely similar to the CPX\slash{}CPB. The figure below shows an Arduino connect to a load cell as well as many other external components that my students used for their final project.%
\begin{figureptx}{Figure}{Arduino Connected to an HX711 plus many other components}{load-cell-5-3}{}%
\begin{image}{0.135}{0.73}{0.135}{}%
\includegraphics[width=\linewidth]{external/Arduino_HX711.png}
\end{image}%
\tcblower
\end{figureptx}%
The \emph{DATA} line connects to pin 3 on the Arduino while \emph{SCK} connects to pin 2 and then \emph{5V} and \emph{GND} pins are straightforward. The \href{https://github.com/cmontalvo251/Microcontrollers/blob/master/Arduino/UNO/Load_Cell/HX711_Arduino_Example.ino}{\linkhilite{Arduino code}} is also on Github and starts with the example script from the \href{https://github.com/bogde/HX711}{\linkhilite{HX711 library}} which is also on Github. To install downloaded ZIP libraries, navigate to Sketch\textgreater{}Include Library\textgreater{}Add .ZIP Library... and then add the ZIP file downloaded from Github.%
\begin{figureptx}{Figure}{Arduino Code for the HX711}{load-cell-5-5}{}%
\begin{image}{0.125}{0.75}{0.125}{}%
\includegraphics[width=\linewidth]{external/HX711_Arduino_Code.png}
\end{image}%
\tcblower
\end{figureptx}%
The screenshot above shows the \emph{\#include \textless{}HX711.h\textgreater{}} on line 34 which imports the HX711 library and then the rest of the setup sets up pins 2 and 3 for the strain gauge. There are also lines of code used to setup Serial communication so that you can see the output of the scale and also change the calibration factor by pressing a and z or + and - to either inrease or decrease the scale factor.%
\end{sectionptx}
%
%
\typeout{************************************************}
\typeout{Section 25.4 Assignment}
\typeout{************************************************}
%
\begin{sectionptx}{Section}{Assignment}{}{Assignment}{}{}{load-cell-6}
\begin{introduction}{}%
For this assignment you are to wire up a strain gauge and load cell amplifier and calibrate the sensor using known weights and another digital scale. First, set the scale value to zero and then place known weights on to the load cell and record the output. Place the known weight on a different calibrated scale and record those measurements as well. Repeat this experiment for 10 different known weights and plot the actual weight on the y-axis from the known digital scale and the output from the load cell on the x-axis. Use this data to compute a regression curve similar to \hyperref[servo]{\linkhilite{Chapter~{\xreffont\ref{servo}}}} and write the equation relating load cell output to actual "weight" (mass) in kg and report the scale value necessary for your code.%
%
\begin{enumerate}
\item{}Title page with name,title and date - 5\%%
\item{}Introduction with wiring diagram - 5\%%
\item{}All figures must have appropriate annotations, figure captions, figure labels and a paragraph of supporting text - 10\%%
\item{}Include a photo of your CPX\slash{}CPB connected to your amplifier and load cell. - 10\%%
\item{}Include a screenshot (entire computer screen) of Mu showing the Plotter open showing load cell data as you press on the strain gauge - 10\%%
\item{}Include a table of your known weights and load cell output in kg. - 20\%%
\item{}Include a Figure of your calibration data plotted with your trend line on top. Put your \(r^2\) value in the title of the plot - 20\%%
\item{}Write your regression equation as (\emph{actual weight = load cell output \slash{} scale}) being sure to replace scale with your actual values. Remember to round to appropriate significant figures - 10\%%
\item{}Using your calibration equation imbedded in your Mu code, place an unknown weight on your load cell and see what the output is and compare to the actual weight once you place it on a scale. Include a photo and screenshot of this experiment - 10\%%
\item{}Appendix A - Video Details - Film yourself and your entire screen showing the load cell value changing as you press on the strain gauge. - Pass\slash{}Fail%
\item{}Appendix B - Code used to take data and code used to plot in plain text (no screenshots) - Pass\slash{}Fail%
\item{}Appendix C - Screenshot of Excel or Sheets if applicable - Pass\slash{}Fail%
\item{}\alert{If you are asked to fix something in the Appendices such as missing content in your video or screenshots of your code, you will lose 10\% of your grade everytime you submit an update after the deadline.}%
\end{enumerate}
\end{introduction}%
\end{sectionptx}
\end{chapterptx}
%
%
\typeout{************************************************}
\typeout{Chapter 26 Results and Discussion}
\typeout{************************************************}
%
\begin{chapterptx}{Chapter}{Results and Discussion}{}{Results and Discussion}{}{}{results-discussion}
\renewcommand*{\chaptername}{Chapter}
\begin{introduction}{}%
The chapters above detailed about 20 projects with multiple subsections that can be used for an instrumentation course. These projects have iterated and been deployed in a real classroom environment since the Fall of 2019. The results and discussion below represent multiple years of analysis and \textquotedblleft{}in the field\textquotedblright{} iteration of courses. That is, the authors and faculty have used feedback from students to make the course better and better every semester. This chapter here represents the results and discussion of feedback from the students as well as conclusions drawn from faculty and the authors of this text.%
\end{introduction}%
%
%
\typeout{************************************************}
\typeout{Section 26.1 Data Collection and Analysis}
\typeout{************************************************}
%
\begin{sectionptx}{Section}{Data Collection and Analysis}{}{Data Collection and Analysis}{}{}{results-discussion-3}
Qualitative data collection for this study was done using end of the semester surveys. Course surveys are delivered by the University Office of Institutional Effectiveness. The instructor of the course only receives results of the survey. Eight questions in the survey are restricted to a 1-5 Likert scale while four questions are open ended text responses\hyperlink{likert}{\linkhilite{[{\xreffont 49}]}}. The surveys are sent to every student taking the course via email. The surveys are optional and as such not all students completed the survey. The data for the course was compiled for all questions and compared against the Departmental average. This provides a base of comparison for this course across the Department of Mechanical Engineering. It is worth noting that the authors do reflect on how their own positionality may affect interpretation of data. This reflectivity is a core characteristic in qualitative research and is important in establishing credibility and showing transparency. A more in depth analysis of the results are examined in the discussion section.%
\end{sectionptx}
%
%
\typeout{************************************************}
\typeout{Section 26.2 Results}
\typeout{************************************************}
%
\begin{sectionptx}{Section}{Results}{}{Results}{}{}{results-discussion-4}
The main goal of implementing this survey was to determine the answer to the following question: \alert{What response does implementation of the kit yield from students?} To answer this question, the 4th question in the course survey will be examined. This question states: \alert{"What did you like best about how the instructor taught the course?"} \hyperref[tab_quotes1]{\linkhilite{Table~{\xreffont\ref{tab_quotes1}}}} and \hyperref[tab_quotes2]{\linkhilite{Table~{\xreffont\ref{tab_quotes2}}}} includes examples of questionnaire responses which is in the appendix.%
\begin{tableptx}{Table}{\textbf{Questionnaire responses to \textquotedblleft{}What did you like best about how the instructor taught the course?\textquotedblright{} (2019-2022)}}{tab_quotes1}{}%
\centering%
{\tabularfont%
\begin{tabular}{AcAcA}\hrulethin
Semester&Quotes\tabularnewline\hrulemedium
Fall 2019&\multicolumn{1}{m{0.7\linewidth}A}{\raggedright\setlength{\parskip}{0.5\baselineskip}%
\textquotedblleft{}The course was taught expertly with proper use of in class equipment available and hands on design problems with experiments relative to the material being taught.\textquotedblright{}%
\par
\textquotedblleft{}I like how he introduced us to the Circuit Playground which allowed to be hands-on since none of my other classes are.\textquotedblright{}%
}\tabularnewline\hrulethin
Spring 2020 (Remote Mid-Way)&\multicolumn{1}{m{0.7\linewidth}A}{\raggedright\setlength{\parskip}{0.5\baselineskip}%
\textquotedblleft{}Hands on experience\textquotedblright{}%
\par
\textquotedblleft{}I especially enjoyed the hands-on Fridays portion of the course. The practical knowledge of creating something with the book-learning we gained during the lecture was priceless.\textquotedblright{}%
\par
\textquotedblleft{}Fridays were my favorite days. I got to experience hands on projects where I learned the most from.\textquotedblright{}%
}\tabularnewline\hrulethin
Fall 2020 (Hybrid)&\multicolumn{1}{m{0.7\linewidth}A}{\raggedright\setlength{\parskip}{0.5\baselineskip}%
\textquotedblleft{}I loved the hands on approach to the course. It is the first course, outside senior design, where we had the opportunity to exercise our engineering knowledge on an actual project.\textquotedblright{}%
\par
\textquotedblleft{}The projects! this [taught] a lot of people hands-on skills\textquotedblright{}%
}\tabularnewline\hrulethin
Spring 2021 (Hybrid)&\multicolumn{1}{m{0.7\linewidth}A}{\raggedright\setlength{\parskip}{0.5\baselineskip}%
\textquotedblleft{}I liked the hands on aspect of the class the most\textquotedblright{}%
\par
\textquotedblleft{}I wish all other class were taught as well\textquotedblright{}%
\par
\textquotedblleft{}I like that we actually had a hands on project, and several projects that were fun to do\textquotedblright{}%
}\tabularnewline\hrulethin
Fall 2021&\multicolumn{1}{m{0.7\linewidth}A}{\raggedright%
\textquotedblleft{}I love how he tries to engage his students and help them learn real world engineering, not just textbook stuff. He really wants to watch his students succeed inside and outside of the classroom. We need more classes and professors like this in engineering! It was a fantastic class and instructor\textquotedblright{}%
}\tabularnewline\hrulethin
Spring 2022&\multicolumn{1}{m{0.7\linewidth}A}{\raggedright%
\textquotedblleft{}I had no interest in coding or robotics like things but now I do. I grew an interest in coding after absolutely hating it. I love seeing the application of what I've learned so far in my degree path. I am a hands-on learner, so the labs and the project were my favorite.\textquotedblright{}%
}\tabularnewline\hrulethin
Fall 2022&\multicolumn{1}{m{0.7\linewidth}A}{\raggedright\setlength{\parskip}{0.5\baselineskip}%
\textquotedblleft{}The hands on projects and labs\textquotedblright{}%
\par
\textquotedblleft{}Very engaging with the students and cared for each student and wanted to help when needed. Gave good real world applications as well as fun, interesting assignments that helped us students learn.\textquotedblright{}%
}\tabularnewline\hrulethin
\end{tabular}
}%
\end{tableptx}%
\begin{tableptx}{Table}{\textbf{Questionnaire responses to \textquotedblleft{}What did you like best about how the instructor taught the course?\textquotedblright{} (2023-2025)}}{tab_quotes2}{}%
\centering%
{\tabularfont%
\begin{tabular}{AcAcA}\hrulethin
Spring 2023&\multicolumn{1}{m{0.7\linewidth}A}{\raggedright%
\textquotedblleft{}I love how hands on everything was we would learn about a concept on either Monday or Wednesday and then Friday use that concept on the mini projects\textquotedblright{}%
}\tabularnewline\hrulethin
Fall 2023&\multicolumn{1}{m{0.7\linewidth}A}{\raggedright\setlength{\parskip}{0.5\baselineskip}%
\textquotedblleft{}Dr. Montalvo gives a very calm atmosphere to learn in. He made this subject so much better and enjoyable than I ever would of thought of before starting this class. It has gotten me interested in perusing my own projects using what I learned outside of class as well. Dr. Montalvo also showed us a handful of resources for learning and getting components for our projects, which was fascinating as well.\textquotedblright{}%
\par
\textquotedblleft{}His method of teaching was genuinely stimulating, and encouraged students to learn. This was the only class I actually looked forward to attending\textquotedblright{}%
\par
\textquotedblleft{}Engaging and hands-on\textquotedblright{}%
}\tabularnewline\hrulethin
Spring 2024&\multicolumn{1}{m{0.7\linewidth}A}{\raggedright\setlength{\parskip}{0.5\baselineskip}%
\textquotedblleft{}He makes things that seemed super hard and complicated in circuits, very easy. The labs are very engaging and he is very available and willing to help any time you have questions or issues.\textquotedblright{}%
\par
\textquotedblleft{}Engaging and Funday fridays were actually fun\textquotedblright{}%
}\tabularnewline\hrulethin
Fall 2024&\multicolumn{1}{m{0.7\linewidth}A}{\raggedright\setlength{\parskip}{0.5\baselineskip}%
\textquotedblleft{}How he engaged with the class.\textquotedblright{}%
\par
\textquotedblleft{}His in class projects\textquotedblright{}%
}\tabularnewline\hrulethin
Spring 2025&\multicolumn{1}{m{0.7\linewidth}A}{\raggedright%
\textquotedblleft{}I like how he does a big project with numerous steps throughout the semester. This is very valuable in the design process learning how to keep moving forward but at the same time learning how to troubleshoot when issues do arise.\textquotedblright{}%
}\tabularnewline\hrulethin
Fall 2025&\multicolumn{1}{m{0.7\linewidth}A}{\raggedright%
\emph{No quotes this semester due to low response}%
}\tabularnewline\hrulethin
\end{tabular}
}%
\end{tableptx}%
Note that certain semesters were during the Pandemic of COVID-19. In the Spring of 2020 the University transitioned to fully online lectures on March 23rd. That was mid-way through the semester. In Fall 2020 and Spring 2021 the students participated in a hybrid lecture format where students were both allowed to come to class as well as remain remote and log on via Zoom. In the Fall 2021 semester, the University returned to fully in person lectures without masks.%
\par
To further answer this question, results from the question \textquotedblleft{}How effective was the instructor in helping students achieve course objectives\textquotedblright{} is shown in \hyperref[tab_scores]{\linkhilite{Table~{\xreffont\ref{tab_scores}}}} also in the appendix. This table shows the response rate per semester as well as average score for the question The last column shows the average score across the department. Spring 2020 has an asterisk because the University of South Alabama created a different course survey due to the onset of COVID-19. Instead the question \textquotedblleft{}Please tell us your perception of the quality of the class prior to the transition to [online] on March 23rd\textquotedblright{} was used during Spring 2020. As such, the departmental scores are not listed.%
\begin{tableptx}{Table}{\textbf{Average Score for \textquotedblleft{}Instructor Effectiveness in Helping Students Achieve Course Objectives\textquotedblright{} (\textasteriskcentered{}Alternate question during COVID-19 Pandemic)}}{tab_scores}{}%
\centering%
{\tabularfont%
\begin{tabular}{AcAcAcAcA}\hrulethin
Semester&Response Rate&Average Score (1-5)&Department Avg Score (1-5)\tabularnewline\hrulemedium
Fall 2019&22\%&4.20&3.63\tabularnewline\hrulethin
Spring 2020 (Remote Mid-Way) \textasteriskcentered{}&30\%&4.66\textasteriskcentered{}&N\slash{}A\tabularnewline\hrulethin
Fall 2020 (Hybrid)&46\%&4.64&3.57\tabularnewline\hrulethin
Spring 2021 (Hybrid)&45\%&4.75&3.73\tabularnewline\hrulethin
Fall 2021&42\%&4.70&4.14\tabularnewline\hrulethin
Spring 2022&47\%&4.94&3.91\tabularnewline\hrulethin
Fall 2022&57\%&5.0&4.0\tabularnewline\hrulethin
Spring 2023&35\%&4.94&3.93\tabularnewline\hrulethin
Fall 2023&92\%&5.0&4.09\tabularnewline\hrulethin
Spring 2024&55\%&4.81&4.12\tabularnewline\hrulethin
Fall 2024&40\%&5.0&4.07\tabularnewline\hrulethin
Spring 2025&37\%&4.94&4.08\tabularnewline\hrulethin
Fall 2025&11\%&5.0&4.01\tabularnewline\hrulethin
\end{tabular}
}%
\end{tableptx}%
From analysis of the surveys, 4 themes emerged which capture student perceptions of the newly developed course. These include: Teaching Style, Course Design and; Comparison, Balance between Theory and; Application and Connected Learning. \hyperref[t_themes]{\linkhilite{Table~{\xreffont\ref{t_themes}}}} shows these four themes with definitions of each. These themes overall represent the perspectives students had on the course based on the survey results.%
\begin{tableptx}{Table}{\textbf{Resultant themes and definitions of themes}}{t_themes}{}%
\centering%
{\tabularfont%
\begin{tabular}{AcAcA}\hrulethin
\terminology{Theme}&\terminology{Definition}\tabularnewline\hrulemedium
Teaching Style&\multicolumn{1}{m{0.6\linewidth}A}{\raggedright%
Instructor teaching style and personality has an impact on course which is sometimes hard to distinguish from the course design itself%
}\tabularnewline\hrulethin
Course Design \& Comparison&\multicolumn{1}{m{0.6\linewidth}A}{\raggedright%
Student evaluation of different aspects of the course, including the flipped-classroom design, structure of lecture and class time, as well as specifics of the book. Course is compared to others.%
}\tabularnewline\hrulethin
Balance between Theory \& Application&\multicolumn{1}{m{0.6\linewidth}A}{\raggedright%
Students enjoyed the application aspect of the PBL Instrumentation course, but suggested adding more theory to the GitBook provided for the course. Students wanted more harmony between the textbook and the kit.%
}\tabularnewline\hrulethin
Connected Learning&\multicolumn{1}{m{0.6\linewidth}A}{\raggedright%
Students extend learning outside of the classroom and make connections to their personal life%
}\tabularnewline\hrulethin
\end{tabular}
}%
\end{tableptx}%
\end{sectionptx}
%
%
\typeout{************************************************}
\typeout{Section 26.3 Discussion}
\typeout{************************************************}
%
\begin{sectionptx}{Section}{Discussion}{}{Discussion}{}{}{results-discussion-5}
This lab at home kit has been implemented since the Fall of 2019. By the time of this writing, the kit has been implemented for 7 semesters, three of which were during COVID-19. The kit was very useful during the COVID-19 pandemic (Spring 2020, Fall 2020 and Spring 2021) because students could order the kit and complete the projects remotely. If students had trouble they could email issues or take a video and share it or even set up a quick video conferencing call. The students were able to see their course objectives hands on which provided a unique way of learning during the pandemic. When the students used the kit in the classroom once the University returned to in person lectures, the students still enjoyed the projects as shown in the survey results.%
\par
In regards to the second research question about student satisfaction of the course through the use of the kit, \hyperref[tab_quotes1]{\linkhilite{Table~{\xreffont\ref{tab_quotes1}}}} and \hyperref[tab_quotes2]{\linkhilite{Table~{\xreffont\ref{tab_quotes2}}}} clearly show a positive response in regards to the projects and the lab at home kit. For brevity only two to three quotes were shown but for some semesters at least 25 quotes were generated with many of them reflecting on the "hands on learning" component of the course.%
\par
Furthermore, \hyperref[tab_scores]{\linkhilite{Table~{\xreffont\ref{tab_scores}}}} shows an on average course rating that is 22\% higher than the departmental average. It can be said with confidence that the quotes in \hyperref[tab_quotes1]{\linkhilite{Table~{\xreffont\ref{tab_quotes1}}}} and \hyperref[tab_quotes2]{\linkhilite{Table~{\xreffont\ref{tab_quotes2}}}} are in direct response to the implementation of the kit and that the students would not have written these words had the kit not been an addition to the course. This proves the first research question that a kit can be implemented as well as answers the second research question.%
\par
When these quotes are taken into consideration with the higher course ratings and higher student responses, it is evident that the take-home lab kit promoted more active student engagement in the course and, by extension, more active student learning. This data shows, in conjunction with student commentary, that the projects associated with kits fostered a positive learning experience for students while also covering required content for the course. Projects like this one are a strong strategy professors can add to their repertoire as they teach content. Per Gardner's and Bloom's theories of learning, this would indicate the students likely mastered content and skills in this setting just as well as, if not better than they would in a more traditional lecture format\hyperlink{Armstrong}{\linkhilite{[{\xreffont 19}]}}\hyperlink{Gardner}{\linkhilite{[{\xreffont 11}]}}.%
\end{sectionptx}
%
%
\typeout{************************************************}
\typeout{Section 26.4 Limitations of Findings}
\typeout{************************************************}
%
\begin{sectionptx}{Section}{Limitations of Findings}{}{Limitations of Findings}{}{}{results-discussion-6}
Although survey analysis resulted in themes that gave perspectives on the course, there are clear limitations in the resultant findings. The authors are well aware that this data does not conclusively show that implementing a kit in a course increases achievement of course objectives. Many other variable are present in this study, such as the instructor, the course and even the classroom the course is taught in. This study is limited in generalizability based on the small response rate in some of the semesters. Findings are also not generalizable to all engineering Instrumentation courses across the world or even the United States given the influence of the instructor's teaching style and personality in the course. Future research should expand on the study by recruiting individual interview participants and survey respondents to gain a better understanding of the impact of the course on students.%
\par
However, findings are useful for instructors seeking to improve their own course by adopting the at-home kit and documentation developed. This theme appeared multiple times among students and it is very possible that students were confused about whether they liked the course or liking the professor. For students, the instructor had a big influence on making the course enjoyable; Preliminary results from this qualitative study can be used to inform course design and teaching practice in future courses.%
\end{sectionptx}
\end{chapterptx}
%
%
\typeout{************************************************}
\typeout{Chapter 27 Conclusions}
\typeout{************************************************}
%
\begin{chapterptx}{Chapter}{Conclusions}{}{Conclusions}{}{}{conclusions}
\renewcommand*{\chaptername}{Chapter}
\begin{introduction}{}%
A lab at home kit has been described here. This kit is part of the curriculum for Instrumentation and Experimental Methods and uses the CircuitPlayground Express (CPX) from Adafruit as well as other electronic components. The CPX uses CircuitPython, a derivative of Python. Python like all scripting languages allows students to focus on writing concise code to use the CPX as a tool rather than getting bogged down in syntax, compilation and runtime errors. All software is free, open source and has numerous examples on the Adafruit Learn page as well as the internet at large and Github \hyperlink{Github}{\linkhilite{[{\xreffont 33}]}}.%
\par
As with all other courses in engineering it is possible to teach this course without any hands on material. Instead, a project based learning (PBL) style is utilized where students use the CPX and other components of the lab at home kit to complete guided projects requiring data acquisition, analysis and post-processing all geared towards enhancing learning of course objectives. The final project in the class allows the students to be creative, take ownership, collaborate with others, and use critical thinking to develop something new.%
\par
The responses from the students is positive with a 22\% increase over the department average over 7 semesters. The qualitative survey responses from students are also positive and show that the kit is effective in enhancing student satisfaction through interaction of tangible objects. Student satisfaction correlates to student performance as shown in the literature. Hopefully, similar kits can be created for other classes at USA as well as other universities in STEM across the world.%
\end{introduction}%
\end{chapterptx}
%
%
\typeout{************************************************}
\typeout{Chapter 28 Acknowledgements}
\typeout{************************************************}
%
\begin{chapterptx}{Chapter}{Acknowledgements}{}{Acknowledgements}{}{}{acknowledgements}
\renewcommand*{\chaptername}{Chapter}
\begin{introduction}{}%
Dr. Carlos Montalvo would like to acknowledge a few key members who made this textbook possible. First and foremost I would like to thank Adafruit for their entire ecosystem of electronics, tutorials, blogs and forums. Much of what I have learned here to teach Instrumentation was from \href{https://www.adafruit.com/}{\linkhilite{Adafruit}} and the \href{https://learn.adafruit.com/}{\linkhilite{Adafruit Learn}} system and specifically people like Lady Ada and John Park who have helped shape \href{https://circuitpython.org/}{\linkhilite{CircuitPython}} and the \href{https://www.adafruit.com/product/3333}{\linkhilite{CircuitPlayground Express}} to what it is today\hyperlink{Adafruit}{\linkhilite{[{\xreffont 32}]}}\hyperlink{AdafruitMAIN}{\linkhilite{[{\xreffont 34}]}}\hyperlink{CircuitPython}{\linkhilite{[{\xreffont 30}]}}. I would also like to thank Dr. Saami Yazdani for creating the blueprint for Instrumentation at my university by creating a laboratory environment for an otherwise totally theoretical course. His course was the foundation for this textbook and for that I thank him for showing the way. I’d like to also thank and acknowledge Tangibles that Teach for giving me the opportunity to morph this loose set of projects into a textbook that can be used for multiple universities and classrooms and of course help students learn and acquire knowledge through creating. I'd also like to thank Steven Clontz for first showing me PreTeXt and for Oscar Levin and all of his helpful replies via email and Zoom. I also like to thank Justin Dyer, Landon Freeman, Noah Broadus, and Logan Allen for their rocket pressure results and Judson Duke, Justin Dyer, and Jason Brignac for their GPS aircraft flight results. I'd like to thank my graduate students Nghia Hyunh and Prisha Chanana for their help on a few chapters. Finally, I'd like to thank the co-authors Lisa and Marine for helping me add some statistical survey data to this textbook and being all around great academic colleagues and friends%
\end{introduction}%
\end{chapterptx}
%
\backmatter%
%
\clearpage\phantomsection%
\addcontentsline{toc}{part}{Back Matter}%
%
%
\typeout{************************************************}
\typeout{References  References}
\typeout{************************************************}
%
\begin{references-chapter-numberless}{References}{References}{}{References}{}{}{backmatter-1}
%% If this is a top-level references
%%   you can replace the "referencelist" environment
%%   with a standard "thebibliography" environment
\begin{referencelist}
\bibitem[1]{copilot}\hypertarget{copilot}{}GitHub. GitHub Copilot. https:\slash{}\slash{}copilot.github.com\slash{}, 2021. Accessed on December 5, 2025.%

\bibitem[2]{gemini}\hypertarget{gemini}{}Google DeepMind. Gemini: Googles multimodal ai model. https:\slash{}\slash{}gemini.google.com\slash{}, 2023. Accessed on December 5, 2025.%

\bibitem[3]{vandv}\hypertarget{vandv}{}Wikipedia. V-model. https:\slash{}\slash{}en.wikipedia.org\slash{}wiki\slash{}V-Model. Accessed December 21st, 2021.%

\bibitem[4]{drawio}\hypertarget{drawio}{}Draw.io. JGraph Ltd, 2026, www.drawio.com\slash{}. Accessed 29 Mar. 2026%

\bibitem[5]{blender}\hypertarget{blender}{}Blender. Blender Foundation, 2026, www.blender.org\slash{}. Accessed 29 Mar. 2026%

\bibitem[6]{fritzing}\hypertarget{fritzing}{}Fritzing. Fritzing, 2026, fritzing.org\slash{}. Accessed 29 Mar. 2026%

\bibitem[7]{solidworks}\hypertarget{solidworks}{}SolidWorks. Dassault Systèmes, 2026, www.solidworks.com\slash{}. Accessed 29 Mar. 2026%

\bibitem[8]{autodesk}\hypertarget{autodesk}{}Autodesk. Autodesk, 2026, www.autodesk.com\slash{}. Accessed 29 Mar. 2026%

\bibitem[9]{cirkit}\hypertarget{cirkit}{}Cirkit Designer. Cirkit, 2026, https:\slash{}\slash{}app.cirkitdesigner.com\slash{}. Accessed 29 Mar. 2026%

\bibitem[10]{inkscape}\hypertarget{inkscape}{}Inkscape. Inkscape, 2026, www.inkscape.org\slash{}. Accessed 29 Mar. 2026%

\bibitem[11]{Gardner}\hypertarget{Gardner}{}Howard gardner’s theory of multiple intelligences. in instructional guide for university faculty and teaching assistants. https:\slash{}\slash{}www.niu.edu\slash{}citl\slash{}resources\slash{}guides\slash{}instructional-guide, 2020. Northern Illinois University Center for Innovative Teaching and Learning.%

\bibitem[12]{Ralph}\hypertarget{Ralph}{}R.A. Ralph. Post secondary project-based learning in science, technology, engineering and mathematics. Journal of Technology and Science Education (JOTSE), 6(1):26-35., 2015.%

\bibitem[13]{Learning-Styles}\hypertarget{Learning-Styles}{}Linda K. Silverman Richard M. Felder. Learning and teaching styles in engineering education. Engineering Education, 78(7):674–681, 1988.%

\bibitem[14]{Yong}\hypertarget{Yong}{}Tang Yong. The mode of theoretical knowledge and practical knowledge combination: The significance of internship. World Journal of Education, 2(4):55–63., 2012.%

\bibitem[15]{PBL1}\hypertarget{PBL1}{}J. C. Perrenet, P. A. J. Bouhuijs, and J. G. M. M. Smits. The suitability of problem-based learning for engineering education: Theory and practice. Teaching in Higher Education, 5(3):345–358, 2000.%

\bibitem[16]{PBL-Book}\hypertarget{PBL-Book}{}Krajcik Joseph S. and Phyllis C. Blumenfeld. Project-based learning. Academia, 2006.%

\bibitem[17]{TANGIBLEINTERFACES}\hypertarget{TANGIBLEINTERFACES}{}Paul Marshall. Do tangible interfaces enhance learning? In Proceedings of the 1st International Conference on Tangible and Embedded Interaction, TEI ’07, page 163–170, New York, NY, USA, 2007. Association for Computing Machinery.%

\bibitem[18]{StudentSatisfaction}\hypertarget{StudentSatisfaction}{}BARBARA M. OLDS and RONALD L. MILLER. The effect of a first-year integrated engineering curriculum on graduation rates and student satisfaction: A longitudinal study. Journal of Engineering Education, 93(1):23–35, 2004.%

\bibitem[19]{Armstrong}\hypertarget{Armstrong}{}P. Armstrong. Bloom’s taxonomy. Accessed 2\slash{}2\slash{}22 https:\slash{}\slash{}cft.vanderbilt.edu\slash{}guides-sub-pages\slash{}blooms-taxonomy\slash{}, 2010. Vanderbilt University Center for Teaching.%

\bibitem[20]{Labs}\hypertarget{Labs}{}Lyle D. Feisel and Albert J. Rosa. The role of the laboratory in undergraduate engineering education. Journal of Engineering Education, 94(1):121–130, 2005.%

\bibitem[21]{LabAtHome}\hypertarget{LabAtHome}{}Juan P. Oliver and Fiorella Haim. Lab at home: Hardware kits for a digital design lab. IEEE Transactions on Education, 52(1):46–51, 2009.%

\bibitem[22]{LabAtHome1}\hypertarget{LabAtHome1}{}F. G. Martin. Integrating hardware experiences into a computer architecture core course. Journal of Computer Science College, 21:39–52, 6 2006.%

\bibitem[23]{LabAtHome2-EE}\hypertarget{LabAtHome2-EE}{}N. Manjikian and S. Simmons. Evolution and enhancements of a microprocessor systems course. IEEE Transactions Education, 42(4):360, 11 1999.%

\bibitem[24]{LabAtHome3}\hypertarget{LabAtHome3}{}P. Li W. Durfee and D. Waletzko. Take-home lab kits for system dynamics and controls courses. Proceedings of the American Control Conference, pages 1319–1322, 2004.%

\bibitem[25]{LabAtHome4-EE}\hypertarget{LabAtHome4-EE}{}D. Nicoletti D. Cyganski and J. A. Orr. A new introductory electrical engineering curriculum for the first-year student. IEEE Transactions on Education, 37(2):171–177, 5 1994.%

\bibitem[26]{codetype}\hypertarget{codetype}{}Sebastian Nanz and Carlo A. Furia. A comparative study of programming languages in rosetta code. In 2015 IEEE\slash{}ACM 37th IEEE International Conference on Software Engineering, volume 1, pages 778–788, 2015.%

\bibitem[27]{IDE}\hypertarget{IDE}{}Integrated development environments. IntegratedDevelopmentEnvironments, 2021. Accessed: 2021-12-09. https:\slash{}\slash{}wiki.python.org\slash{}moin\slash{}%

\bibitem[28]{Tiobe}\hypertarget{Tiobe}{}Tiobe index. https:\slash{}\slash{}www.tiobe.com\slash{}tiobe-index\slash{}, 2022. Accessed: 2021-3-25.%

\bibitem[29]{CodeComparison}\hypertarget{CodeComparison}{}Hans Fangohr. A comparison of c, matlab, and python as teaching languages in engineering. In Marian Bubak, Geert Dick van Albada, Peter M. A. Sloot, and Jack Dongarra, editors, Computational Science - ICCS 2004, pages 1210–1217, Berlin, Heidelberg, 2004. Springer Berlin Heidelberg.%

\bibitem[30]{CircuitPython}\hypertarget{CircuitPython}{}CircuitPython. https:\slash{}\slash{}circuitpython.org\slash{}, 2021. Accessed: 12-09-2021.%

\bibitem[31]{Gitbook}\hypertarget{Gitbook}{}Carlos Montalvo, Lisa Schibelius, and Marine Leabeater. Project Based Engineering Instrumentation with CircuitPython. Github, 2023. https:\slash{}\slash{}cmontalvo251.github.io\slash{}project-based-instrumentation\slash{}.%

\bibitem[32]{Adafruit}\hypertarget{Adafruit}{}Adafruit. Adafruit learn, 2021. Adafruit https:\slash{}\slash{}learn.adafruit.com\slash{}.%

\bibitem[33]{Github}\hypertarget{Github}{}Carlos Montalvo. Facility for aerial systems and technology github repository. cmontalvo251, 2022. https:\slash{}\slash{}github.com\slash{}cmontalvo251%

\bibitem[34]{AdafruitMAIN}\hypertarget{AdafruitMAIN}{}Adafruit. Adafruit, 2023. Adafruit https:\slash{}\slash{}adafruit.com\slash{}.%

\bibitem[35]{Arduino}\hypertarget{Arduino}{}Arduino. Arduino, 2023. https:\slash{}\slash{}www.arduino.cc\slash{}.%

\bibitem[36]{AdafruitBLE}\hypertarget{AdafruitBLE}{}Adafruit. Adafruit ble connect, 2023. https:\slash{}\slash{}learn.adafruit.com\slash{}bluefruit-le-connect?gclid=CjwKCAiA3KefBhByEiwAi2LDHOVV9IbeAFTEsKBx285M-B7o3VSCeW7uoU0gbDqJ3D0KDX\textunderscore{}xnKD-uxoCrrAQAvD\textunderscore{}BwE.%

\bibitem[37]{Metro}\hypertarget{Metro}{}DiCola, Tony. "Digital Outputs." CircuitPython Basics: Digital Inputs and Outputs, Adafruit Learning System, 28 Aug. 2017, learn.adafruit.com\slash{}circuitpython-digital-inputs-and-outputs\slash{}digital-outputs. Accessed 16 Mar. 2026.%

\bibitem[38]{potentiometer}\hypertarget{potentiometer}{}Ada (Fried), Lady (Limor), and Kattni Rembor. "CircuitPython Analog In." Adafruit Circuit Playground Express, Adafruit Learning System, 12 Oct. 2017, learn.adafruit.com\slash{}adafruit-circuit-playground-express\slash{}circuitpython-analog-in. Accessed 18 Mar. 2026.%

\bibitem[39]{pitot-author}\hypertarget{pitot-author}{}Hrisko, Joshua. "Arduino Pitot Tube Wind Speed and Airspeed Indicator - Theory and Experiments." Maker Portal, 8 Feb. 2019, makersportal.com\slash{}blog\slash{}2019\slash{}02\slash{}06\slash{}arduino-pitot-tube-wind-speed-theory-and-experiment. Accessed 18 Mar. 2026.%

\bibitem[40]{pitot-amazon}\hypertarget{pitot-amazon}{}“Hobbypower Airspeed Sensor Kit MPXV7002DP Differential Pressure Sensor Board for APM Pixhawk Flight Controller.” Amazon, 2026, www.amazon.com\slash{}Hobbypower-Airspeed-MPXV7002DP-Differential-controller\slash{}dp\slash{}B00WSFWO36. Accessed 27 Mar. 2026.%

\bibitem[41]{pitot-datasheet}\hypertarget{pitot-datasheet}{}NXP Semiconductors. MPXV7002: Integrated Silicon Pressure Sensor On-Chip Signal Conditioned, Temperature Compensated and Calibrated. Rev. 4, 2017, www.nxp.com\slash{}docs\slash{}en\slash{}data-sheet\slash{}MPXV7002.pdf. Accessed 27 Mar. 2026.%

\bibitem[42]{imu-wiring}\hypertarget{imu-wiring}{}Siepert, Bryan. "LIS3MDL Triple-axis Magnetometer", Adafruit Learning System January 17th, 2020, https:\slash{}\slash{}learn.adafruit.com\slash{}lis3mdl-triple-axis-magnetometer\slash{}python-circuitpython. Accessed 20 Mar. 2026%

\bibitem[43]{pressure-wiring}\hypertarget{pressure-wiring}{}Siepert, Bryan. "LPS22 Pressure Sensor", Adafruit Learning System January 17th, 2020, https:\slash{}\slash{}learn.adafruit.com\slash{}lps22-pressure-sensor\slash{}python-circuitpython. Accessed 20 Mar. 2026%

\bibitem[44]{photocell-wiring}\hypertarget{photocell-wiring}{}Ada (Fried), Lady (Limor), "Photocells", Adafruit Learning System, July 29th, 2012, https:\slash{}\slash{}learn.adafruit.com\slash{}photocells\slash{}using-a-photocell, Accessed Mar 20. 2026%

\bibitem[45]{servo-wiring}\hypertarget{servo-wiring}{}Ada (Fried), Lady (Limor), Rembor, Kattni, "Adafruit Circuit Playground Express - CircuitPython Servo", 12 Oct. 2017, https:\slash{}\slash{}learn.adafruit.com\slash{}adafruit-circuit-playground-express\slash{}circuitpython-servo. Accessed 20 Mar. 2026%

\bibitem[46]{gps-wiring}\hypertarget{gps-wiring}{}Ada (Fried), Lady (Limor), Rembor, Kattni, "Adafruit Ultimate GPS - CircuitPython Parsing", 12 Oct. 2017, https:\slash{}\slash{}learn.adafruit.com\slash{}adafruit-ultimate-gps\slash{}circuitpython-parsing. Accessed 20 Mar. 2026%

\bibitem[47]{load-cell-wiring}\hypertarget{load-cell-wiring}{}Clark, Liz "Adafruit NAU7802 24-Bit ADC - STEMMA QT \slash{} Qwiic" - June 24, 2025, https:\slash{}\slash{}learn.adafruit.com\slash{}adafruit-nau7802-24-bit-adc-stemma-qt-qwiic\slash{}python-circuitpython Accessed 11 Jun, 2026%

\bibitem[48]{b_photocell}\hypertarget{b_photocell}{}Kujtim Mustafa, Ragmi Mustafa, and Refik Ramadani. Measuring the voltage, current and resistance of the ldr sensor through the arduino uno. Asian Journal of Research in Computer Science, 16:211–222, 10 2023.%

\bibitem[49]{likert}\hypertarget{likert}{}Andrew T. Jebb, Vincent Ng, and Louis Tay. A review of key likert scale development advances: 1995–2019. Frontiers in Psychology, 12, 2021%

\bibitem[50]{youtube}\hypertarget{youtube}{}Montalvo, Carlos. YouTube, 2026, www.youtube.com\slash{}@CarlosMontalvo251.%

\bibitem[51]{python}\hypertarget{python}{}Python Software Foundation. Python, 2026. https:\slash{}\slash{}www.python.org\slash{}.%

\bibitem[52]{thonny}\hypertarget{thonny}{}Aasumets, Aivar. Thonny: Python IDE for Beginners. 2026, thonny.org\slash{}.%

\bibitem[53]{spyder}\hypertarget{spyder}{}Spyder: The Scientific Python Development Environment. Spyder Project Contributors, 2026, www.spyder-ide.org\slash{}.%

\bibitem[54]{colab}\hypertarget{colab}{}Google Colaboratory. Google, 2026. https:\slash{}\slash{}colab.research.google.com\slash{}.%

\bibitem[55]{pydroid}\hypertarget{pydroid}{}IIEC \slash{} Teamux LLC Pydroid 3 - IDE for Python 3. IIEC, 2026. https:\slash{}\slash{}play.google.com\slash{}store\slash{}apps\slash{}details?id=ru.iiec.pydroid3\&hl=en\textunderscore{}US\&gl=US.%

\bibitem[56]{pythonista}\hypertarget{pythonista}{}Zorn, Ole. Pythonista 3. Pythonista, 2026. https:\slash{}\slash{}apps.apple.com\slash{}us\slash{}app\slash{}pythonista-3\slash{}id1085978097.%

\bibitem[57]{mu}\hypertarget{mu}{}Mu Editor. Mu, 2026. https:\slash{}\slash{}codewith.mu\slash{}en\slash{}.%

\bibitem[58]{CPXchip}\hypertarget{CPXchip}{}Microchip Technology. “ATSAMD21G18 - Microcontrollers and Processors.” Microchip Technology, 2026, www.microchip.com\slash{}wwwproducts\slash{}en\slash{}ATsamd21g18. Accessed 27 Mar. 2026.%

\bibitem[59]{sparkfun}\hypertarget{sparkfun}{}SparkFun Tutorials. SparkFun Electronics, 2026, learn.sparkfun.com\slash{}tutorials\slash{}. Accessed 27 Mar. 2026.%

\bibitem[60]{potentiometerlearn}\hypertarget{potentiometerlearn}{}Platzer, Øyvind Nydal. “How to Use a Potentiometer.” Build Electronic Circuits, 2026, www.build-electronic-circuits.com\slash{}potentiometer\slash{}. Accessed 27 Mar. 2026.%

\bibitem[61]{lux}\hypertarget{lux}{}Nofrizal, et al. “Measuring the Voltage, Current and Resistance of the LDR Sensor through the Arduino UNO.” Journal of Physics: Conference Series, vol. 2596, no. 1, 2023, doi:10.1088\slash{}1742-6596\slash{}2596\slash{}1\slash{}012051. Accessed 27 Mar. 2026.%

\bibitem[62]{lux-wikipedia}\hypertarget{lux-wikipedia}{}Lux. Wikipedia, 2026, en.wikipedia.org\slash{}wiki\slash{}Lux. Accessed 27 Mar. 2026.%

\bibitem[63]{thermistor-wikipedia}\hypertarget{thermistor-wikipedia}{}Thermistor. Wikipedia, 2026, en.wikipedia.org\slash{}wiki\slash{}Thermistor. Accessed 27 Mar. 2026.%

\bibitem[64]{stein-hart}\hypertarget{stein-hart}{}Steinhart, John S., and Stanley R. Hart. “Calibration Curves for Thermistors.” Deep Sea Research and Oceanographic Abstracts, vol. 15, no. 4, 1968, pp. 497–503., doi:10.1016\slash{}0011-7471(68)90057-0. Accessed 27 Mar. 2026.%

\bibitem[65]{steinhart-wikipedia}\hypertarget{steinhart-wikipedia}{}Steinhart-Hart equation. Wikipedia, 2026, en.wikipedia.org\slash{}wiki\slash{}Steinhart\%E2\%80\%93Hart\textunderscore{}equation. Accessed 27 Mar. 2026.%

\bibitem[66]{ST}\hypertarget{ST}{}STMicroelectronics. “LIS3DH: MEMS Digital Output Motion Sensor: Ultra-Low-Power High-Performance 3-Axis ‘Nano’ Accelerometer.” STMicroelectronics, 2026, www.st.com\slash{}en\slash{}mems-and-sensors\slash{}lis3dh.html. Accessed 27 Mar. 2026.%

\bibitem[67]{Elegoo}\hypertarget{Elegoo}{}Elegoo. "Elegoo UNO R3." Elegoo, 2026, www.elegoo.com\slash{}products\slash{}elegoo-uno-r3. Accessed 27 Mar. 2026.%

\bibitem[68]{octave}\hypertarget{octave}{}GNU Octave. GNU, 2026. https:\slash{}\slash{}www.gnu.org\slash{}software\slash{}octave\slash{}.%

\bibitem[69]{sheets}\hypertarget{sheets}{}Google Sheets. Google, 2026. https:\slash{}\slash{}www.google.com\slash{}sheets\slash{}about\slash{}.%

\end{referencelist}
\end{references-chapter-numberless}
\end{document}
